% ============================================================
%  R. Nielsen, "Grundideernes Logik. I." [The Logic of the Fundamental Ideas, I]
%  Kjøbenhavn: Gyldendal (F. Hegel); Trykt hos J. H. Schultz, 1864.
%  TRANSLATION (English).
%  Source of truth: transcription.tex (Danish), fully collated against the
%  scan 2026-09-22 (see its ACCURACY header). Mirrors it 1:1: every printed
%  page carries \opage{N} (or \apage{N}) at the corresponding point in the
%  English; the front matter carries roman \opage{VII}..\opage{XXXIV}.
%  Translated 22 Sep 2026 in 27 batches. Conventions: TRANSLATION-PLAYBOOK.md.
%  Latin, Greek, French and German quotations kept as printed; Danish work
%  titles in citations kept in Danish. Printer's errors in the Danish are not
%  imitated in the English; the book's own errata are flagged by % Errata:.
% ============================================================
\documentclass[12pt]{book}
\usepackage[utf8]{inputenc}
\usepackage[T1]{fontenc}
\usepackage[english]{babel}
\usepackage[protrusion=true,expansion=true]{microtype}
\usepackage[margin=1.4in]{geometry}
\usepackage{setspace}
\usepackage{amsmath}
\usepackage{libertinus}
\usepackage{libertinust1math}
\usepackage{textalpha}
\usepackage{graphicx}
\DeclareUnicodeCharacter{0254}{\reflectbox{c}}
\usepackage{enumitem}
\usepackage{fancyhdr}
\usepackage[colorlinks=true,linkcolor=black,urlcolor=black,
  pdftitle={The Logic of the Fundamental Ideas I},
  pdfauthor={Rasmus Nielsen}]{hyperref}
\setlength{\headheight}{15pt}
\pagestyle{fancy}
\fancyhf{}
\fancyhead[LE,RO]{\thepage}
\fancyhead[RE]{\textit{The Logic of the Fundamental Ideas}}
\fancyhead[LO]{\textit{\leftmark}}
\setlength{\parindent}{1.5em}
\setlength{\parskip}{0pt}
\setstretch{1.3}
\title{\textbf{The Logic of the Fundamental Ideas}\\[4pt]\large I.\\[8pt]
  \normalsize (\textit{Grundideernes Logik}, translated from the Danish)}
\author{R.\ Nielsen}
\date{Copenhagen: Published by the Gyldendal Bookshop (F.\ Hegel).\\[2pt]
  \small Printed by J.\ H.\ Schultz. 1864.}
\usepackage{marginnote}
\newcommand{\opage}[1]{\marginnote{\footnotesize\textit{[#1]}}}
\newcommand{\apage}[1]{\marginnote{\footnotesize\textit{[#1]}}}
\begin{document}
\maketitle
\tableofcontents

\chapter*{Preface}
\addcontentsline{toc}{chapter}{Preface}

\noindent Although this work is only a small part of a larger whole --- scarcely two-thirds of the first ninth of the entire Logic ---, it is nevertheless not a fragment, but a whole in itself. If the plan and method set out in it have any scientific significance, then the principle of the work as a whole must be clearly recognizable from the present part. But should it prove, upon thorough examination, that the plan and procedure of the present work are either of no significance whatever for a further development of philosophy in the present day, or at any rate suffer from considerable fundamental defects, then the author will at least not attempt to excuse himself on the ground that criticism has passed judgment on the single part without knowing the whole.

\bigskip
\noindent Copenhagen, 18 August 1864.
\nopagebreak
\par\hfill R.\ Nielsen.

\chapter*{Introduction}
\addcontentsline{toc}{chapter}{Introduction}
\opage{VII}

\noindent ``What thinking is can be learned only by thinking; what philosophy is can be learned only by philosophizing; if thinking ceases, thought must vanish; if philosophizing ceases, philosophy must vanish. Just as thinking is not only the beginning of thought but also its continuation and completion, so philosophizing is itself the beginning, continuation and completion of philosophy. If a philosophy is untrue, we must seek the ground of this in an untrue philosophizing; if a manner of philosophizing is untrue, we must see the proof in the untrue philosophy. But when, in order to prove the truth of a philosophy, we must appeal to true philosophizing, and in order to establish the truth of philosophizing must in turn appeal to true philosophy, then we are surely moving in a circle, and the question becomes how we are to determine the fixed point of departure where we are then to find the decisive criterion of truth.''\footnote{Forelæsninger over ``Philosophisk Propædeutik'' fra Universitetsaaret 1860--61, p.~1.}

In order to answer this introductory question, we consider philosophy partly from a propaedeutic, partly from a systematic standpoint. What is to be elucidated from a propaedeutic standpoint can then on the whole be summed up in an introduction \opage{VIII} to philosophy in general; as regards what is systematic in the stricter sense, the preliminary elucidations can be given in a special introduction to logic.

\section*{General Introduction.}

There are chiefly three main points on which attention must here be fixed, namely the determination of philosophy's \emph{a priori foundation}, of its \emph{empirical foundation}, and of the \emph{goal} toward which it strives.

\emph{The a priori foundation of philosophy.} The old γνῶθι σεαυτόν, which recalls the Seven Sages, also recalls Socrates, and contains a summons to philosophize, insofar as it contains a summons to become familiar with one's own life of consciousness. Against the proposition that everyone must surely know himself best stands the proposition that nothing is more difficult than to know oneself; self-knowledge is thus the easiest thing of all and the most difficult thing of all: here philosophy begins at once.

Every human being who is reasonably in full possession of his faculties is involuntarily aware of what goes on in his circle of representations, and since each individual must naturally have his own interests, his own plans and thoughts, the task of self-knowledge, seen from this side, certainly seems to fall outside philosophy.

But philosophy itself is one thing, its presuppositions and preliminaries another. Philosophy in the narrower sense is one thing, philosophy in the wider and widest sense another. In the widest sense, every turning of attention to the spiritual, every attempt to think about one's own thinking, every start toward becoming conscious of one's consciousness, may with full right be regarded as a preparation for philosophy, \opage{IX} indeed even as a kind of preliminary philosophizing. In this sense, then, there is not a single rational human being who does not to a certain degree philosophize and have his philosophy. According as we take the word philosophy in a wider or a narrower meaning, we can thus say that there are many who actually philosophize without knowing what philosophy is, and that there are in turn many who imagine they know what philosophy is without ever having philosophized.

``In order to comprise all particular considerations under one comprehensive standpoint, propaedeutics begins with the conceptual determination that philosophy is a scientific doctrine of consciousness. The doctrine of consciousness shows what it means to become conscious of something, and must to that extent distinguish between consciousness itself and what comes to consciousness. This separation in turn demands the union of the separated sides. An empty consciousness, a consciousness that is conscious of nothing, is a contradiction: consciousness must have a content.''\footnote{Philosophisk Propædeutik i Grundtræk. Kbhvn.\ 1857. p.~4.}

Whoever turns his attention inward with only some reflection and perseverance will soon make the discovery that a thoroughgoing self-knowledge is something more than what is called, in a practical sense, a private affair. The life of consciousness has, like everything in the world, its laws: however much that is peculiar the individual's nature may have, each individual self-consciousness nevertheless stands under the general laws valid for all human knowledge and consciousness. But if the cognition of the peculiar cannot be separated from the cognition of the universal, then conversely the cognition of the universal cannot be a cognition in spirit and in truth, a cognition of life for life, if by one-sided abstraction it is kept separate from \opage{X} the cognition of what is individual and peculiar to the particular personalities, the particular peoples and the particular times. The peculiar and the common, the special and the general, the contingent and the necessary enter into so many combinations in the life of consciousness that no one can know his own inner state from the ground up without also knowing the spiritual life of his contemporaries in its fundamental features and conditions. Only when self-consciousness is cognized in its relation to contemporary consciousness does the doctrine of consciousness become scientific; and only the scientific doctrine of consciousness is philosophy in the proper sense.

If one finds it comic when a man is so distrait that every other moment he comes as a stranger to his own door and gets lost in his own house, then one will surely also find it comic when a human being who otherwise lays claim to be counted among the rational and cultivated runs astray every moment in his own circle of representations, is disoriented in his own consciousness, and in a quarter of an hour speaks out of seventeen views of life. The first and most essential demand that true cultivation must make of every human being is surely this, that he must be reasonably at home in his own consciousness; for he who is not at home in his own consciousness, in which regions of existence should he properly be at home?

But to be at home in one's own consciousness presupposes a high degree of spiritual self-activity. Indeed, if the life of consciousness could stand still at a certain point, and its content could once and for all be grasped in fixed forms and distributed under standing rubrics, then it would not even be so insurmountable to get one's bearings; for with what stands still one can find one's way tolerably well. But the life of consciousness cannot possibly stand still; it has its being in becoming, its subsistence in an unceasing stream of changes: \opage{XI} what alternation of thoughts and representations, of moods and states, is not the consciousness of individual human beings in every age, and the consciousness of every age in turn, subjected to! The more developed, the more alive a human being is, the more currents, the more movement there is in such a person's consciousness; the more developed, the more alive an age is, the more streams, the stronger the movement there is in the consciousness of the age: the increasing movement is proportional to the growing life and its increasing development.

The life of consciousness must, however, under all its changes also have something constant: fixed centers, unshakable principles, determinate purposes belong in the world of spirit to the first necessities, the daily bread. But now it happens in the course of the advancing movement of culture that the factors of constancy, the constants of consciousness themselves, become variable and are swept along into the stream. Then the foundations of existence tremble; the old supports fall, and the new ones do not stand firm; everything runs round: now it is difficult to get one's bearings.

``The present state is properly no state, it is a transition'': so people thought toward the end of the last century. ``It is a transition'': so people thought at the beginning of this century. ``It is a transition'': so people think even to this day. And it really is ``a transition.'' The transition takes place not only in \emph{the outer} world, in the political upheavals; the transition occurs with still far greater emphasis in \emph{the inner} world, in the realm of ideas, in what concerns religion, poetry and art. ``For it is of no use for us to wish to conceal or embellish the truth: we must confess to ourselves that religion in our time is for the most part only the affair of the uncultivated, while for the cultivated world it belongs to what is past, what is left behind. \opage{XII} This sad result had already been attained before the first French Revolution, which only brought it nearer to consciousness, and brought it about that a great mass of the uncultivated passed over into the class of the cultivated, for the cultivated are in recent times those who politicize, and through the Revolution almost all became politicians'' \dots\ ``What they'' (those who politicize) ``must in all circumstances see is that the political character of the age itself constitutes its crisis, and consequently is a contradiction; for after having cast away religion, and thus the cognition of the infinite, it has only finite determinations left, and with these it nonetheless wishes to build up an eternal state and an eternal constitution. It is therefore also in the feeling of this contradiction that the better sort have returned to the rejected religion, but to no avail, for they themselves cannot possibly acquire the faith of past times, but at most do as certain merchants who, although they are near bankruptcy, nevertheless do not dare to confess their poverty to themselves; and with their followers things look still more doubtful, for to these religion has now become a matter of fashion, but a matter of fashion is never a serious matter, and moreover it has become so only for the fashionable, not for the cultivated'' \dots\ ``If it goes thus with religion, then things will not stand much better with art and poetry, those spheres in which man, just as in religion, should have the immediate certainty of the divine and eternal, and enjoy as present what one otherwise, under the name of the other life, awaits as something to come. It is only on the presupposition of that certainty that art and poetry have a true existence, and that their effects, even in the most profane jest, can be explained. But when the cultivated world has abolished religion, one cannot expect that its sisters should count for much. Art and poetry can at most be only an agreeable \opage{XIII} luxury; and even if, following the example of the Romans, one in certain countries puts plays in the same class as bread, this has only the meaning that luxury nowadays belongs to life's first necessities.''\footnote{J.\ L.\ Heiberg. Om Philosophiens Betydning for den nærværende Tid. Pros.\ Skrft.\ First vol. p.~396 ff.}

It is now roughly a generation since that judgment was pronounced. Much has changed in that time; even the judgment ``on the significance of philosophy'' has in many respects changed. But in all these changes there is nevertheless something here that has remained unchanged; it is the necessity of spiritual self-activity, it is a demand for a priori certainty, it is the need to have the essence of consciousness elucidated really thoroughly and many-sidedly.

\emph{The empirical foundation of philosophy.} As a scientific doctrine of consciousness, philosophy certainly refers us to our own interior, to what is original in our subjective nature; but this is not to be understood as though the interest in fathoming the world that is within us should dull our eye for the world that is outside us. The one-sidedness of wishing to solve the riddle of existence by immersing oneself in one's pure I and constructing the objects of experience a priori is something philosophy has long since got beyond. Hegelianism has, to be sure, by no means been able to appropriate what is actual in actuality itself; but it has nevertheless contributed mightily to carrying through the principle that the concept of existence must be determined by existence, since existence is originally determined by the concept, that the subjective can only be comprehended from the objective, since the objective in concreto is the concept of the subjective. ``If the individual,'' says I.\ L.\ Heiberg,\footnote{Indlednings-Foredrag til det i November 1834 begyndte logiske Cursus paa den kongelige militære Høiskole. Pros.\ Skr.\ vol.~1, p.~490.} ``is to confine himself merely to learning to know \opage{XIV} himself and the repetition of himself in other individualities; if both nature and state, both science and art, both God and the world are to be separated from his knowledge, then it is not to be expected, nay not even to be hoped, that any but a few persons, whom one is accustomed to call philosophers by profession, will find themselves satisfied by the cloister life in which philosophical contemplation consists, and whose insurmountable wall is to separate us from all that which is the goal of man's desire or the object of his longing. But that is not how matters stand. Even if we will grant that philosophical contemplation is a kind of cloister life, we must at the same time remember that in recent times many cloisters have been changed into educational institutions, which, even if for a time they seclude from the world, nevertheless do not make this seclusion their lasting aim, but on the contrary only a means whereby the pupil is to be enabled more easily to understand life and actuality when he leaves the school and exchanges it for them. This service philosophy is to render the young man who strives to attain the crown of cultivation: it is to lead him to the quiet regions where, undisturbed by the movements of the outer world, he can turn his gaze upon himself; but when it has thus cleared his vision and heightened his energy, it is to send him back to the activity that has claims upon him; then he will understand himself and his surroundings, \dots\ feel his calling, and recognize what and how he is to act. For just as he who has a cultivated musical ear knows how to find, even in the most confusing tumult of instruments and voices, the unity that binds them; or just as he who knows the laws of gravitation knows how to find order and design in the countless multitude of heavenly bodies which a blind chance seems to have strewn about in space, so the philosophically \opage{XV} cultivated man can hear the harmony, inaudible to the uninitiated, that runs through the confusing tumult of the world, and see the thread, invisible to the uninitiated, that shows us the way through life's innumerable shapes, cast about as if by chance.''

The significance of a philosophical propaedeutic could hardly be presented in a truer, clearer and more beautiful way than has been done in these words quoted from Heiberg's ``Indlednings-Foredrag''; but when the Hegelian conception of the relation between the empirical and the a priori is to be judged from the standpoint of objectivity and actuality, and the investigations go into particulars, then it becomes apparent how one-sidedly this philosophy is still entangled in the abstractly speculative. ``Religion,'' it says,\footnote{Heiberg. Om Philosophiens Betydning. Pros.\ Skr.\ vol.~1, p.~407.} ``and likewise art and poetry, must, because they are unable to \emph{enter into} the endeavors of the age, necessarily \emph{set themselves above them}; but with this the age would not be content, for what it fights for it will not let itself be robbed of, and in this it is right, since the infinite must now, as before, be won through the very development of finite endeavors, not through their sudden cessation. Only philosophy is able to go into the many details of our finite, especially political, purposes; only it is able to see their direction toward the infinite, and by this cognition to clear up what is obscure in them. Thereby it alone becomes able to sublate them without annihilating them; on the contrary, in their sublation into the infinite it precisely confirms their validity. In this way our finite striving is grafted with the infinite, the human with the divine, and limitation has vanished; our sciences become philosophy, and our states regain the ordering form.''
\opage{XVI}
This fundamental thought is correctly Hegelian and shows in broad outline the purpose toward which the philosophy of ``the concept'' as a whole strove; but there is one expression in this statement which involuntarily betrays the one-sidedness in which ``the system'' with ``the immanent method'' is entangled. ``Our sciences,'' it says here, ``become philosophy''; this expression says more than one would believe at first glance. --- In his philosophical propaedeutic Hegel concludes his survey of the encyclopedia of the sciences with the declaration that the encyclopedia is properly ``a presentation of the general content of philosophy; for what in the sciences is grounded in reason depends on philosophy, whereas what is in them of that which rests on arbitrary and external determinations, or, as it is called, is positive and statutory, lies outside it.'' As insignificant as such an expression would be if it were merely tossed off as a passing remark, so characteristic does it become when, as in the present case, it is to count as a fundamental thought running through an entire system. The misjudgment of the independence of the particular sciences which is thereby expressed is equally striking whether we look to the history of these sciences or to their special methods and peculiar purposes.\footnote{Cf. Philosophisk Propædeutik i Grundtræk. Kjøbhvn.\ 1857, p.~46--67. p.~99--117. p.~158--178. Mathematik og Dialektik. Kbhvn.\ 1859, p.~3. Forelæsninger over ``Philosophisk Propædeutik'' fra Universitetsaaret 1860--61, p.~30--39. Forelæsninger over ``Phil.\ Propæd.'' fra Universitetsaaret 1861--62. p.~43--75.} There is then also no danger that the special sciences, at their present stage of development, should let themselves be troubled by philosophy's pretended superiority and condescend once more to walk in its leading-strings. But if the misunderstanding cited does no harm to the special sciences, it is all the more calamitous \opage{XVII} for philosophy itself, which precisely thereby robs itself of its empirical foundation.

If a philosophy of ``the concept'' is to be able to secure the reality of the concept, it must surely above all attend to how a real concept is formed; but to what degree Hegelian philosophy, precisely as regards this essential point, has miscarried in the formation of concepts is conspicuous when one casts a glance at the philosophy of nature. ``That Hegel did not form his first, elementary concepts of nature under the guidance of any mathematical physics, but in a purely logical-speculative manner, is evident at every point. This is seen not merely from the fact that he expressly combats the designations and representations of mathematical physics; but it is seen chiefly from the fact that he did not even find it necessary to secure an authentic apprehension of the categories and ways of viewing which he wishes to combat in the mathematicians \dots\ For if philosophy were in earnest with propositions such as these: that essence is to be known from the phenomenon, the ground from existence, the thing from its properties, force from its manifestations, the cause from its effects, then the thinker would surely above all have to be concerned to secure a really reliable, exact and complete apprehension of natural phenomena, of the modes of manifestation of forces, of physical causes with their conditions and laws \dots\ The shortcomings pointed out here do not concern errors in this or that particular; such errors are always more or less inessential in a philosophical presentation of essential substance; --- to make an essential matter of the inessential must (cf. Begrebet: Nøiagtighed p.~3--11) naturally be left to pedants and pettifoggers ---; no, the shortcomings pointed out here concern the method itself, and betray a weakness which is, so to speak, congenital to the system. It is the manner in which the categories of nature are from the outset apprehended and treated that \opage{XVIII} has here miscarried. The concepts of nature are concepts of actuality; but in order for concepts of actuality to be formed, the various special sciences, each with its observations, analyses and calculations, must go ahead and show the way. There is in the whole realm of nature not a single conceptual sphere of which it holds that the philosophical thinker, without special negotiation with the special science concerned, should himself be able to show the origin of the categories and find their roots in the soil of reality.''\footnote{Forelæsn.\ over ``Phil.\ Propæd.'' 1861--62. p.~58--72.}

It has, as the writings just cited show, for a number of years been my endeavor to have the relation between philosophy and the particular sciences elucidated and demonstrated as many-sidedly as possible. On a thorough apprehension and exact determination of this relation the future of philosophy will essentially depend. The study of philosophy, the philosophical method, philosophy's scientific independence and its influence on the way of thinking of the cultivated, all this will on the whole depend on the position in which ``the science of the sciences'' will come to stand in relation to the sciences.

\emph{The purpose of philosophy.} The pettiness with which in earlier days one was accustomed to immerse oneself in introductory investigations concerning the utility and use of each individual science has gradually had to yield and give place to the more essential and more rational view that every true science must first and foremost have its worth in itself, and be cultivated in earnest for its own sake, before it can yield any considerable return, whether theoretical or practical. Now if even the most useful sciences feel themselves too exalted to recommend themselves by their utility, then philosophy must so much the more consider itself entitled, if such a question \opage{XIX} is addressed to it, to answer the question: of what use? with the candid declaration that it is properly of no earthly use at all. Should one then on this occasion loudly testify one's contempt for it, it can only answer that, as regards utility, it finds itself in the same case as art, poetry, ethics and religion. To commend religion for utility's sake is to mock the holy; to love art for utility's sake and to do one's duty for utility's sake is to give up the spirit for utility's sake.

With religion, poetry and art philosophy has the purpose, common from the standpoint of humanity, of procuring recognition for the idea and fighting for the right of spirit; but while each of the ideal powers fights with its own weapons, it is a distinguishing mark of philosophy that, as a science of spirit, it can use only scientific weapons. Philosophy has often been reproached for its negative nature: that it reasoned about everything, raised doubt about everything, shook every authority and undermined faith. This reproach strikes philosophizing at certain transitional stages, but by no means philosophy itself. Philosophy must certainly acknowledge the Cartesian proposition: de omnibus dubitandum est, insofar as from the standpoint of the scientific cognition of the idea it is to subject everything to a testing. But such a testing, which aims at purging out what is spurious, what is untenable in itself, aims by the same token at preserving what is genuine, what is valid in itself, and procuring recognition for it. Without a comprehensive testing of the axioms, principles and presuppositions of human knowledge, philosophy as the general theory of science could not possibly solve its task. ``It is thus not presumptuous arbitrariness when philosophy brings the various doctrinal edifices before the tribunal of the theory of science, and demands that with regard to each of them in particular an account be given both of the foundation of the building and of the manner of its \opage{XX} construction; it is not self-assumed authority when, after a carefully conducted testing, it permits itself to assign to those theories which in foundation and manner of construction bear more resemblance to aesthetic castles in the air than to scientific doctrinal edifices `a place of their own' outside science.''\footnote{Phil.\ Propæd. p.~72.}

Among the doctrinal edifices about whose scientific foundation philosophy in recent times has expressed the strongest doubts belongs, as is well known, theology. Does theology then perhaps want no theory of science? ``If one could at one and the same time hold to Anselmus and yet profit from Cartesius, could get credo, ut intelligam to say: de omnibus dubitandum est, and conversely get de omnibus dubitandum est to proclaim itself the true credo, ut intelligam, then theology would surely gladly unite with philosophy in commending a thorough and sound theory of science; but on these conditions it is an impossibility for philosophy to unite with theology, since the thing itself is a contradiction. Although philology, criticism and history, which in and for themselves belong to science, constitute components of theology, they by no means constitute what is theological in theology. Insofar as the disciplines named are treated scientifically, they are treated untheologically; insofar, on the other hand, as they are treated theologically, they are treated unscientifically. For what is theological in theology is recognized by this: that the cause from which all effects in the realm of nature as in that of spirit are, partly immediately, partly mediately, derived and explained is sought not in the nature of things or the essence of reason, but in the absolute will, in the Almighty who creates out of nothing. But even the purest, most acute and most consistent application of such a principle (occasionalism; Malebranche) makes science unscientific. When \opage{XXI} the theory of science is dealing with a principle, the matter cannot be decided by a more or less, as if someone, by following `the middle way,' should think that one could still get through scientifically by assuming a little of the supernatural, so long as one did not make too much of it. The question of a principle is a question of an unconditioned, all-conditioning presupposition; such a question must therefore always be answered with yes or no. Either theology must then once and for all acknowledge supernatural causality, and has thereby at the same time cast off the objectively rational foundation, or else once and for all acknowledge this foundation, and has thereby at the same time renounced every footing in the supernatural. In the first case theology will thus cease to be science, in the latter case to be theology.''\footnote{Phil.\ Propæd. p.~73--74.}

It is not philosophy in a single direction or a single philosophical school from which the objections to theology arise; they present themselves from all sides. It is not ``unbelief,'' not ``spiritual pride''; it is the impossibility of admitting the supernatural into science which in this matter leads thinkers of otherwise highly different spiritual orientation to one and the same negative result. That theology's doctrinal edifice, in spite of the billows of criticism and the storms of philosophy, has remained standing and stands still, as if nothing had happened, is then a singular phenomenon that can be explained in two different ways. For it may indeed be that theology remains standing like the strong house upon the rock, supernaturally strong, because it is built supernaturally and yet scientifically upon a supernatural and yet scientific foundation; but it may also be that theology remains standing, not like that strong house upon the rock, but like the \opage{XXII} old house in the fairy tale: ``it did not know to which side it should fall, and therefore it remained standing.''

Let one judge philosophy's attack on theology for the time being as one will; there is nevertheless one thing here on which all are agreed, and that is that theology is not the same as religion; an attack on theology is then also not the same as an attack on religion. Philosophy can thus attack theology and still honor religion; indeed, it can attack theology out of reverence for religion. In the 15th, 16th and the beginning of the 17th century, when minds truly began to break the chains of authority and emancipate themselves from medieval bondage, the philosophers in their dispute with the theologians introduced the well-known separation between two kinds of truth (duæ veritates), a theological and a philosophical truth. This separation was so far from bringing reconciliation that on the contrary it made the dispute more involved. Scientific and religious cognition were at that time so ecclesiastically grown together that they could not possibly be separated. When the philosophers then, in order to preserve scientific character, declared that in all matters of faith they were willing to submit to the authority of the Church, the Church was certainly entitled to call in question the sincerity of such a concession: how should anyone be able to hold fast to a truth of reason and yet acknowledge a theological truth, when they stand in decisive conflict with one another, and yet both are to count as scientific? But if the philosophers with their two truths had no good conscience, then the theologians, who by taking reason captive under the obedience of faith held to the theological unity of truth, certainly had no good conscience either. He who declared the truth of faith comprehensible was a heretic; he who declared \opage{XXIII} the truth of faith incomprehensible was also a heretic: on these terms a Pomponatius was threatened, and a Vanini burned.

But the assumption of duæ veritates ceases to signify a conflict between reason and revelation or between philosophy and religion when it is seen that the two heterogeneous truths belong to two spheres so mutually different that what is decided in the one must be left undecided in the other. ``In open conflict with everything that scientific insight can in any way acknowledge as objective, faith is so far from obtaining any support from science that precisely through its collision with objective knowledge it must be thrust back into itself, and confess that its only ground is a need of subjectivity, an inner, irrefutable demand \dots\ Life is subjective, the eternal true life absolutely subjective; the will is subjective, the eternal, unconditioned will absolutely subjective; the objective is thus, as objective, not the true: `Subjectivity is the truth' (S. Kierkegaard). --- In that the conflict between faith and knowledge is thus driven to its extreme, intellectual reconciliation is by the same token initiated: the absolutely heterogeneous principles cannot, in an essential sense, do each other harm. The faith that hates knowledge because it must strengthen itself by ignorance is a false faith. Actuality will come along; it is of no use to shut one's eyes: if the earth really goes round, then it goes round with the believers just as well as with the non-believers. In that faith declares that subjectivity is the truth, it certainly loses all significance for science; but in that science declares that objectivity is the truth, it by the same token renounces all influence on faith. Faith and knowledge can then without contradiction be united in one consciousness; their mutual conflict does not stem from what is a matter of principle. The conflict is psychological; it is human frailty, which has neither resignation enough \opage{XXIV} to reconcile itself with the objective of science, nor inwardness enough to hold fast to faith.''\footnote{Philos.\ Propæd. p.~76--77.}

Since the reconciliation between faith and knowledge, seen from this standpoint, might perhaps nevertheless contain a good deal that would awaken misgivings in a serious and independently thinking reader, we shall add some further elucidating remarks. Although science can neither build on miraculous presuppositions nor prove the possibility of the miracle, it nevertheless cannot straightforwardly and apodictically prove the impossibility of the miracle either. That it may seem to the physicist as though the miracle were a sheer impossibility is natural, since the causes of physics are causes of things, and its laws laws of necessity; but the highest cause of philosophy is an idea-cause, and the idea-cause a teleological self-cause. In that the absolute self-cause makes the laws of necessity the ground of the law of freedom, there is seen from the height of the cognition of the idea a depth of possibilities which no human knowledge is able to measure. Instead of fantasizing about the indeterminable and brooding over the unfathomable, philosophy here holds fast to the essential task: to secure the right of science by stating the limit of knowledge. ``By virtue of the inner, original connection of empirical research with the a priori there is thus an essential knowledge; but in consequence of the heterogeneity of thinking and sense-perception, real cognition is subject to finite conditions, and all our knowledge is limited. The limit of knowledge is a twofold one: a \emph{relative} one, insofar as there is always something here that we do not yet know, but which we can nevertheless come to know in time; an \emph{absolute} one, insofar as a finite separation of the content of consciousness and of actuality is grounded \opage{XXV} in our organization: here, then, is something that we in this existence never come to know.''\footnote{Phil.\ Propæd. p.~67.}

The relation between scientific truth and the truth of faith can be determined still more closely by the relation between the idea of truth and the idea of the good. The true and the good cannot, to be sure, fall so far apart that truth should not be good, and the good not true; but in spite of this there is nevertheless an essential opposition between the two ideas. When science, which stands under the idea of truth, answers me the question: what can I know? it takes in this its answer no account whatsoever of my individual subjectivity. ``In science subjective moods: inclination, disinclination, fear, hope and the like, have no significance. Insofar as the question here is of the actual, we do not want to know how this actual, in our opinion, should preferably be, but how it is in itself and necessarily must be; insofar as the question here is of the possible, the question is not whether this possible best satisfies our hopes, but whether it accords with the conditions of existence.''\footnote{Phil.\ Propæd. p.~74.} Religion is grounded in the idea of the good; in answering me the questions: what shall I believe? what may I hope? it aims unconditionally, not at subjectivity in general, but at my existing subjectivity. The task is to know and cognize the innermost, deepest need of subjectivity, of existing subjectivity; to find the means and conditions that are necessary for the satisfaction of this need; to demonstrate the highest purpose of subjectivity and to behold its eternal exemplar: for the solution of such tasks much \opage{XXVI} reflection, much deliberation, much subjective knowledge may be required; but this knowledge, insofar as it is religious, cannot possibly be scientific. The religious truths can certainly also become the object of a scientific treatment; but science must then judge their validity not from the standpoint of the idea of truth, as if it were able to justify them by universally valid proofs, but under the point of view of the idea of the good, in that it apprehends the ideals of faith as ideals of life and finds in the purpose the standard for a judgment of the means. ``But,'' one will ask, ``is the philosophy of religion then not a science? is ethics not a science? If there cannot be a scientific philosophy of religion, what then becomes of the philosophy of spirit? But if there can be a scientific philosophy of religion, why then should there not also be a scientific doctrine of faith, a scientific dogmatics?'' These and similar questions have long since found their answer. ``The separation of the scientific and the unscientific is made difficult by the fact that what according to its nature must be rejected by science nevertheless becomes the object of scientific treatment, and thus in a way does go along into science. Madness, e.g., is incompatible with science; nevertheless the symptoms of madness and the treatment of the mad become an important object for psychology and medical science \dots\ Scientific in the wider sense is everything that becomes the object of science's treatment: the scientific lies in the treatment; scientific in the narrower sense, on the other hand, is only that whose validity science sees, whose truth it can judge, test and acknowledge.''\footnote{Phil.\ Propæd. p.~78.}
\opage{XXVII}
The wide gulf between philosophy and religion is filled by ethics. In ethics the thoughts of life meet the fundamental scientific thoughts; rational ethics stands under the influence of science, religious ethics under that of faith: the difference is determined by the accent that is laid on subjectivity. ``It is just as impossible for scientific criticism to acknowledge the gospel of the Risen One as it is impossible for science to acknowledge the myth of Herkules.'' What then can prevent science from declaring the gospel of the Risen One a myth? The superiority of religious ethics over rational ethics! If there is, on the rational standpoint, no contradiction in making a distinction between the impersonal truths of physics and the personal truths of ethics, then there is, on the standpoint of personality, likewise no contradiction in holding fast to the scientific truths of philosophy and the supra-scientific truths of religion in their mutual opposition, decisive for the subjectivity existing in temporality.

\section*{Introduction to Logic.}

In the case of philosophy too it is necessary to distinguish between phenomenon and essence. The popular objections to the essence of philosophy are, in a superficial sense, applicable to its phenomenon: one thinks one is attacking philosophy, and attacks in an external manner what is striking in the phenomenon of the many mutually conflicting philosophies. But it is especially the system, the philosophical system, that gives offense. ``A little philosophical improvisation, a little explicative philosophy of fermentation can still be tolerated, partly on account of its wit, partly on account of its curiosity; but a fundamental thought strictly carried through in form, a philosophical system, that is a tower of Babel; it is horrible, it is almost like a blasphemy. In philosophical improvisations thoughts \opage{XXVIII} and ideas have free play; but in the system freedom is lost: a philosophical system is a penitentiary of thoughts.''

That the popular polemic against the system likewise takes the semblance for the essence, what is one-sided, contingent and transient in the matter for the matter itself, can easily be seen. One confuses the system with the schema; one attacks the system, but at bottom means schematism. Every thought is essentially systematic in the same degree as it is philosophical: the philosophical thinker is always a systematic thinker. But from the fact that every thorough thought is systematic in its essence it does not yet follow that it is always developed in systematic form; there can thus be a philosophy that is at bottom systematic but has not yet assumed the form of a system. Since, however, every content demands its form, it goes without saying that in philosophy at every stage of development there must be an original drive to stamp the content in an adequate form; thus there forms a philosophical system in the proper sense. The formal expression of the system is the type, the formal expression of the type the schema. In the schema the system exhausts its possibilities and is thereby made finite. The finite physiognomy is now the object of finite consideration, finite understanding. Those who adopt and pay homage to the system believe that they have therein grasped eternal truth in a completed form, and do not notice that the schematic completion of the form is precisely its being made finite. The system is like a temple built by human hands; but the idea itself, the true deity, does not dwell in temples built by human hands. To that extent criticism is right when it attacks schematism.

But schematism is one thing, systematics another. The systematic rests on the reciprocal relation between the finite schema and the infinite content of the idea; the systematic form is a living schema, at once finite and yet infinite. That \opage{XXIX} the form is finite is shown by its schematic determinateness; that in its finitude it is infinite is shown by systematic mobility. If determinateness is lost while mobility maintains itself, the system dissolves into a fermentation; if mobility is lost while determinateness remains, it crystallizes into spiritless schematism.

In the fundamental features of the system the typical is seen; the more the typical is reduced to the first universal, the more possibilities it contains, and the greater constancy it has. The old division of philosophy into \emph{logic, physics and ethics} is so typically pervasive that it may be said to hold even to this day; for the concept of ethics can be taken in so extended a meaning that it furnishes the foundation for the whole philosophy of spirit, just as the concept of physics does for the philosophy of nature. However great the transformations the philosophical system may otherwise undergo, it will nevertheless always be possible to refer its particular divisions to the three comprehensive main parts: \emph{logic, philosophy of nature and philosophy of spirit}. The more, on the other hand, the articulation is pursued into its particulars, the more the schema varies, the more what is changeable in systematics shows itself. ``The whole of science,'' says Hegel, ``divides into three main parts: 1) \emph{logic}, 2) \emph{natural science} and 3) \emph{science of spirit}. Logic is namely the science of the pure concept and the abstract idea. Nature and spirit constitute the reality of the idea, the former as external existence, the latter as the self-knowing essence (or: the logical is the eternal simple essence in itself; nature is this essence become external; spirit its return into itself from its externality). The sciences of nature and of spirit can then be regarded as the \emph{real} or \emph{particular} sciences, as distinct from pure science or logic, since they are the system of pure science in the shape of nature and of spirit.''
\opage{XXX}
In this survey there is undeniably something typical, and in this typical something abiding; but when the closer development of what is indicated in the survey begins, the doubtful at once announces itself, and herein lies the germ of manifold deviations and changes.

As the transition from logic to natural science Hegel sets up mathematics with regard to space and time. Natural science itself falls into: a) \emph{mechanics} with regard to matter and motion. b) The \emph{physics of the inorganic} with regard to light, magnetism, electricity, physical bodies and the chemical process. c) The \emph{physics of the organic}, under which in turn belong \emph{geology}, which treats of the formation of the earth; \emph{botany and plant physiology}, whose object is vegetable nature; \emph{physiology}, \emph{comparative anatomy}, \emph{zoology}, which treat of healthy animal nature, just as \emph{medicine} treats of diseased nature. The science of spirit too falls into three sections: a) spirit in its concept, \emph{anthropology, psychology} (language --- philology.) b) Practical spirit, which manifests itself in deed and action; under this \emph{jurisprudence, political science, history}. c) The completion of spirit in art, religion and science, to which correspond \emph{aesthetics, doctrine of religion and philosophy}.

If one considers this systematic layout more closely, features show themselves which are in such a degree peculiar to Hegelian philosophy that they must unavoidably change when this philosophy is displaced by another philosophy. To these belongs, among other things, the last part of the third section, on ``the completion of spirit.'' What is doubtful in this systematics shines forth, moreover, from the unclear relation in which philosophy and the particular sciences are here placed to one another. If one looks at the trilogy belonging under ``the completion of spirit'': aesthetics, doctrine of religion and philosophy, then doctrine of religion and aesthetics must necessarily be regarded \opage{XXXI} as special sciences falling outside philosophy; the same must hold of history, political science, jurisprudence, etc. But how then are the subdivisions brought into agreement with the main parts? As the first main part of the whole of science \emph{logic} is given; as the last part of the last part of the last section \emph{philosophy} is given: a strange systematics!

So much, however, can be taken as settled, that logic, the ideal foundation of all the sciences, is and must be the first general main part of philosophy. But now the question arises again: can one begin with logic straightaway, or is a special preparation not required here? In this respect it is a matter of determining the relation between the propaedeutic and the systematic beginning. The propaedeutic beginning can in turn have either a more popular or a strictly scientific form. As a propaedeutic in strictly scientific form we can regard Hegel's ``Phänomenologie des Geistes.'' Hegel expresses himself on this in the introduction to the Logic\footnote{Wissenschaft der Logik. 1ster Theil. Werke. 3ter Band. Berlin 1833. p.~33--35.} as follows: ``In der \emph{Phänomenologie des Geistes} habe ich das Bewußtseyn in seiner Fortbewegung von dem ersten unmittelbaren Gegensatz seiner und des Gegenstandes bis zum absoluten Wissen dargestellt. Dieser Weg geht durch alle Formen des \emph{Verhältnisses des Bewußtseyns zum Objekte} durch, und hat den \emph{Begriff der Wissenschaft} zu seinem Resultate.'' Further: ``Der Begriff der reinen Wissenschaft und seine Deduktion wird in gegenwärtiger Abhandlung also insofern vorausgesetzt, als die Phänomenologie des Geistes nichts anderes als die Deduktion desselben ist. Das absolute Wissen ist die \emph{Wahrheit} aller Weisen des Bewußtseyns, weil, wie jener \opage{XXXII} Gang desselben es hervorbrachte, nur in dem absoluten Wissen, die Trennung des \emph{Gegenstandes} von \emph{der Gewißheit seiner selbst} vollkommen sich aufgelöst hat, und die Wahrheit, dieser Gewißheit, so wie diese Gewißheit, der Wahrheit gleich geworden ist.''\footnote{Anf.\ Skr.\ p.~35.}

The manner in which the phenomenology of spirit has introduced logic is in systematic agreement with the manner in which logic is divided and developed. Without insight into the ``absolute knowledge'' with which the phenomenology of spirit concludes, it will be impossible to understand the ``self-movement'' of the concept through the three parts of logic:
\begin{quote}
``I Die Logik des Seyns,\\
II Die Logik des Wesens, und\\
III Die Logik des Begriffs''.
\end{quote}
But from this it follows in turn that this whole systematics must become doubtful when the doctrine of the ``self-movement'' of the concept begins to become doubtful. It is not the profound fundamental thought of the immanent method or the dialectic of the categories that will be the object of attacking criticism in the present work; nevertheless the deviation from the Hegelian fundamental view will show itself to be so pervasive that the altered way of viewing must necessarily entail an alteration in the systematic type. The objections that are here directed from \emph{Grundideernes Logik} against the Hegelian view and the Hegelian method turn chiefly on the following main points:

1) The opposition between the subjective and the objective is not carried through in Hegel's logic; therefore the self-comprehending concept is a mystical personification; therefore there is talk, in all ambiguity, subjectively of the objective, objectively of the subjective.
\opage{XXXIII}
2) The opposition between concept and existence is mystified in a panlogistic manner. Absolute knowledge has indeed a content, but no object; for the logical devours the antilogical, and the idea that was to be all actuality shrinks, in spite of all mediations and ideal concretions, into an empty abstraction of the absolute.

3) The movement of the concept in reflection and the movement of the thing in existence are identified in so unclear and ambiguous a manner that the most important problems of all (the relation between idea-movement and real movement, the ontological problem of objectification, etc.) are pushed back and veiled.

That a pervasive alteration in the fundamental logical view must entail a corresponding alteration in the logical type goes without saying. It has through many preliminary studies been my task to have the relation between the idea and actuality determined in a different way than has hitherto been done. ``The idea is the truth, and truth the identity of essence and actuality. He who in the actual sees something essential only insofar as he sees something universal in it sees only a half truth; he who in actuality looks only at the individual, the special, the merely factual, likewise sees only a half truth. When a half-truth is posited as the whole truth, untruth arises. As long as idealism as the doctrine of essence and empiricism as the doctrine of actuality keep to their finite opposition, no cognition of truth, no true cognition of the idea, will be possible.''\footnote{Phil.\ Propæd. p.~158.}

If, now, a true cognition of the idea requires that the opposition between knowing subjectivity and objective actuality be developed, in order that the unity with its rich content may properly \opage{XXXIV} be cognized, then it is seen from this that a logic of the idea must transform itself into a \emph{Grundideernes Logik}. In that the idea of actuality assumes on the one side the form of subjectivity, on the other side that of objectivity, it divides into two original ideas: \emph{the Idea of Knowledge} for subjective, \emph{the Idea of Power} for objective actuality. Only when the extremes of concept and existence, of idea-movement and real movement, are sufficiently developed through this opposition of the ideas does the concrete unity of the opposed ideas receive its full explanation in \emph{the Idea of Truth}. The doubling indicated in the type will typically repeat itself in every division, and since the ideal at all stages of concretion reflects itself in the real, the logical in the antilogical, and conversely, the possibility of cognizing the idea in actuality, actuality in its idea, and both in their truth, will gradually be realized. Each of the fundamental ideas of logic mirrors the others in its own way; to the mirrorings in the whole there correspond in turn mirrorings in the individual members. Further, the logical type may be said to reflect the type of the whole system in such a way that, while the Idea of Knowledge in its circle gives us the logic of logic, the Idea of Power in its circle gives the logic of nature, and the Idea of Truth in its circle the logic of spirit. So much for the plan. Whether the schema indicated here will otherwise lead to true systematics or to barren schematism, the execution must show. To the three fundamental ideas correspond the three parts of logic.
\begin{quote}
First Part: \emph{The Idea of Knowledge.}\\
Second Part: \emph{The Idea of Power.}\\
Third Part: \emph{The Idea of Truth.}
\end{quote}

% Errata (printed p. XXXV, "Trykfeil og Rettelser"; not reproduced as body text; the corrections apply to the body):
% Errata: p. 87 note, line 5 from bottom:  "Functionsbegreb" -> "Functionsbegrebet"
% Errata: p. 89, line 9 from top:          "reflecterende" -> "reflecterede"
% Errata: p. 202, line 13 from top:        "nævne" -> "nævne:"
% Errata: p. 290 note, line 12 from bottom: "sind außer," -> "sind, außer"
% Errata: p. 308, line 3 from bottom:      "+x5" -> "x + 7"
% Errata: p. 335, line 14 from top:        "det i Mechaniske" -> "i det Mechaniske,"
% Errata: p. 337, line 2 from bottom:      "Inertie, ved Stød og Bevægelse," -> "Inertie og Bevægelse ved Stød,"
% Errata: p. 363, line 9 from bottom:      "cos.t," -> "cost,"
\opage{1} \begin{center}
{\large\bfseries Part One.}\\[3pt]
{\Huge\bfseries The Idea of Knowledge.}
\end{center}
\begin{center}---\end{center}
\begin{center}
{\large\bfseries Section One.}\\[3pt]
{\LARGE\bfseries Subjective Knowledge.}
\end{center}
\begin{center}---\end{center}
\begin{center}
{\large\bfseries Chapter One.}\\[3pt]
{\Large\bfseries Knowledge and the Knower: The Logic of Subjectivity.}
\end{center}
\addcontentsline{toc}{part}{Part One. The Idea of Knowledge}
\addcontentsline{toc}{section}{Chapter One. Knowledge and the Knower: The Logic of Subjectivity}
\markboth{The Idea of Knowledge}{The Logic of Subjectivity}

\noindent With what shall philosophy begin? This fundamental question has in recent times been conceived and answered in highly different ways. Insofar as the talk is of a systematic beginning, one must certainly concede that philosophy is to begin both with and without presupposition, that it is to begin at once with everything and with nothing; but it now depends essentially, to be sure, on how this demand is more precisely understood. To begin with everything and yet with nothing is, superficially seen, undeniably a contradiction; but if one wished to reject the demand on account of the contradiction, one would thereby reject all genuine beginning, for the contradiction lies in the beginning itself. Philosophy must, if its beginning is to have methodical validity, necessarily begin with a presupposition that in a hidden way contains everything, namely everything that science is able to develop in scientific form; for in strict \opage{2} scientific form nothing can, in the course of the forward-striving development, be taken up from outside. But with everything as its presupposition, the beginning of the development nevertheless has nothing that it may lay down as something immediately given. What, after all, should such a given really be? An assumption with which the development began, but which was not itself developed! Thus the development would have to begin by passing over a development, that is, begin by skipping its own first member, the genuine beginning; but to begin in this way by passing over the genuine beginning is surely also a contradiction. What is contradictory in the beginning arises from the fact that the beginning itself is conceived as something that is, as a concept closed in itself, whereas this concept is nevertheless, in infinite becoming, infinitely relative. One begins by presupposing everything, in that one begins by cancelling all presupposition; for one begins with the presupposition that everything is to be developed, and that nothing is yet developed.

The concept of beginning, the concept with which the development begins, thus cannot be the concept of the empty beginning, a concept that would arise from beginning with a beginning in order to determine what beginning is; no, the concept of beginning must, by the nature of the case, be the first adequate concept of development, the concept which, as far as philosophy is concerned, satisfies the demand of the beginning in the most decisive way. Which, now, is this concept? Hegel answers: ``\emph{it is pure being}''; we answer: ``\emph{it is pure knowledge}''. The beginning thus takes place in that the concept of beginning itself is put to the test.

Of pure being it undeniably holds that it is not merely ``the uttermost abstraction from all that is'', but also the most general and most comprehensive determination under which all that is can be subsumed; pure being can then, with respect to what is, be said to be the richest in content and yet the poorest of all concepts. Being embraces everything, and yet contains, considered by itself, nothing; to that extent \opage{3} pure being can certainly be conceived as the concept of beginning; but the beginning itself is, from this standpoint, conceived one-sidedly. As the uttermost abstraction from all that is, pure being must surely presuppose abstraction, namely the abstraction by which it itself arises; but to presuppose abstraction is to presuppose thought: pure being presupposes pure thinking. Just as little as the beginning can take place with a thinking without being, just so little can it take place with a being without thinking; but when the beginning itself presupposes unity of thinking and being, then the beginning here is surely made, not with pure being, but with pure knowledge.

By beginning with being, Hegel wants to purge away everything finite, the contingent, the merely psychological and subjective, which otherwise clings to human thought. But in that our thought in this way begins by thinking universal objective being, it itself becomes abstractly objective, makes itself without further ado one with being, and forgets its own subjective essence to such a degree that all determinations of thought appear as pure determinations of being, all movements of thought as movements of being. This is a logical mystification.

From the standpoint of real being, and not without ground, Hegel has been reproached that the being with which his philosophy begins is not an objective, essential and actual being, but an abstractly subjective one, a pure being-of-thought, a general thought, so that thought, under the pretense of developing nothing but determinations of being, moves from first to last in the circle of its own forms and thus, instead of comprehending what is, spins itself into its own web. The mystification thus rests, to that extent, on a misapprehension of the objective.

From the standpoint of ideal knowledge, however, one can with just as much right reproach Hegel that he fundamentally misapprehends the subjective essence of thought, of objective thought; for he treats the matter as though thinking \opage{4} subjectivity could by ``the immanent method'' be reduced to an idle spectator who in an objective logic might witness how objective being is able to think itself, objective essence to reflect itself, the objective concept to comprehend itself. The mystification thus rests, in this respect, on a misapprehension of the subjective.

If the mystification is to be removed, philosophy must become clearly conscious of what it means that it expressly begins, not with pure being, but with pure knowledge. By placing the beginning in knowledge, the subjective essence of thought and the objective essence of being are equally asserted. In that the beginning proceeds from knowledge, it proceeds from a thinking and knowing subjectivity, and has in subjectivity its ideal independence over against objective being, at the determination of which it essentially aims. To be sure, objective, real being cannot, according to this beginning, acquire significance except by becoming an object of knowledge; but the objective element has by no means been pushed back by this. On the contrary! real being is essentially asserted; for while it is by no means self-evident that being must everywhere presuppose knowledge, it is on the other hand self-evident that knowledge must everywhere presuppose being. From the Cartesian cogito, ergo sum there follows a cogitat, ergo est, but not conversely an est, ergo cogitat.

That knowledge as concept of beginning presupposes everything and contains nothing is then also easy to show. For in knowledge everything that is in any way to become an object of philosophical treatment must necessarily be taken up. Were there now a kind of being that fell wholly outside knowledge, then this being would surely also have to fall outside philosophy. Such a being would then have to be --- for the question here is not of my and your imperfect knowledge, but of the eternal essence of knowledge --- conceived by the knower as that of which nothing, absolutely nothing, could be known. But to know of something that one knows nothing \opage{5} of it is a sheer contradiction: to know of a thing that nothing is known of it is, after all, to know something of it. Our pure knowledge is in its beginning nonetheless just as empty as Hegel's pure being. This follows directly from the fact that here abstraction is made from all content. Here it is known, indeed, that here there is a knowledge; here it is known in all generality what knowledge is, namely: a unity-in-difference of thinking and being; but the unity itself is floating and indeterminate, because thinking in this unity is still just as indeterminate and floating as being. If pure being is a nothing because it vanishes in an indeterminate thinking, and the indeterminate thinking in turn a nothing because it vanishes in an indeterminate being: then knowledge, as the still indeterminate unity of thinking and being, must in this its beginning surely be absolutely a knowledge of negative content, a knowledge of its own indeterminacy, of its own emptiness, a knowledge of nothing.

As concept of beginning, knowledge must also be a concept of movement; this demand is then fulfilled in that what is still indeterminate is determined, in that the unity-in-difference of thinking and being is developed. Being by itself cannot be a concept of movement; pure being, when abstraction is made from the abstracting thought, cannot possibly move itself. Movement is, according to Hegel, a transition from being to nothing, and from nothing to being; such a transition, however, is a movement of being only insofar as it is a movement of thought; but a movement of being determined by thought is expressly a movement of knowledge. Could being without thinking make itself into nothing, and the nothing thus arising in turn turn about and transform itself into being: then the first logical movement would undeniably be not merely objective, in objective unity with being, but so objective that it fell outside knowledge; in order to develop the idea of ``absolute knowledge'', logic would thus have to begin with a movement that fell outside knowledge, and so make a transition from non-knowledge to knowledge; but to make a \opage{6} transition from an absolute non-knowledge to an absolute knowledge is an absolute impossibility.

It is this impossibility that is impressed, in an energetic if somewhat sophistical manner, by the ancients, when it is said: ``The knower can, insofar as he is knowing, come to know nothing; for as knower he surely already has knowledge, and what one already has, one cannot get''. If one objects ``that one can surely very well know something without knowing everything, and, precisely by already knowing something, be prepared to come to know something he does not know'', then the following answer is given to this: ``No one can be called a knower except in relation to what he knows; in relation to what he does not know he is an ignorant one. It was thus not the knower but the ignorant one who came to know something. But neither can the ignorant one, insofar as he is ignorant, come to know anything. In order to come to know something, the absolutely ignorant one would surely have to begin absolutely with a knowledge of his ignorance; in order to begin with a knowledge of his ignorance, he would have to begin with a knowledge of a non-knowledge; but to begin with a knowledge of a non-knowledge is to begin, not as one ignorant, but as one knowing.''

What is true in this is that knowledge presupposes itself to infinity, moves from itself through all else to itself, so that its whole development is essentially a self-development. That knowledge --- not my or your imperfect knowledge, but knowledge in its originality --- presupposes itself and begins with itself, means then, in other words, that it has its principle in itself, that its essence is selfhood. If knowledge were not itself principle, it could not possibly be the genuine and first beginning. How closely \emph{principle} is related to \emph{beginning} can already be seen from the fact that words such as ἀρχή and principium in free usage also mean beginning. If knowledge cannot be its own beginning unless it is \opage{7} its own principle, then neither can it be its own principle unless it is its own beginning.

But although there is here unity of principle and beginning, the principle itself is nevertheless not simply the same as the beginning: the unity is a unity-in-difference. If the difference is emphasized, then the principle relates to the beginning as the first mover to the movement. The beginning is not the first mover but the first movement; in that the movement begins, the beginning does not fall outside the movement but coincides with it. It is this point that the Eleatics overlooked; and only because they overlooked this point did they consistently come to deny movement. Before Achilles can reach the tortoise, he must already have covered half the way; but before he can have covered half of the way, he must have covered the half of the half of the half, to infinity. ``In order that a movement can be completed, it must of course have begun; if now the beginning movement is posited as identical with the completed one, then we have, as the cardinal point of the contradictions, the admittedly absurd proposition: `in order that a movement can be completed, it must already beforehand be completed'. One demands, in other words, an absolute beginning; but to demand an absolute beginning is a self-contradiction, for no beginning is or can be absolute.''\footnote{Forelæsn. ov. ``Phil.\ Propæd.'' 1860--61. p.~70.}

What is absolute in the beginning is the principle, and what is relative in the principle itself is the beginning. Since now the beginning is movement, the principle falls at once outside and inside the movement. That the all-moving principle falls outside the movement, Aristotle expressly emphasized, though to be sure in an altogether one-sided way, in conceiving the first mover as that which is in itself unmoved (τὸ πρῶτον κινοῦν ἀκίνητον ὄν). To conceive the relation \opage{8} of the moved and the unmoved in such a way that they immediately exclude each other is to conceive it without idea; movement presupposes movement to infinity; only from what is moved in itself can a movement proceed: the all-moving principle must itself be in movement. ``In imperfect energy, rest and movement fall apart from each other; what rests does not move; what moves itself does not rest; what rests must be set in movement, what is moved must come to rest. In perfect energy, on the other hand, movement coincides with rest: here there is rest in movement, movement in rest.''\footnote{Phil.\ Propæd. p.~149.}

If the principle is to be determined as what is unmoved in its movement, then the movement of the principle must necessarily be conceived as self-movement, and the principle itself as original selfhood, as knowing unity. Only as knowing unity, as the knowledge that knows itself in everything, is knowledge its own principle; but the knowledge that knows itself in everything is, in pure generality, knowing subjectivity. Just as little as knowledge can be determined and developed without determining itself, by means of its other, as a knowledge of something, just so little can it be determined and developed as a knowledge of its other without being determined as an original knowledge of itself. If the other, the real, which is to constitute the content of knowledge, is taken away, there arises empty knowledge, a knowledge of nothing that is not at the same time a knowledge of something; but a knowledge of nothing without a knowledge of something is an absolute non-knowledge, an annihilation of all knowledge. The real other, which in actuality is to constitute the content of knowledge, must thus lie in knowledge's own principle; but the principle of knowledge is knowing subjectivity; the objective content belongs originally within subjectivity.

That the principle of everything subjective is also the principle of everything objective, that objectivity with its infinite \opage{9} content of being must to that extent be just as original as subjectivity with its infinite knowledge, is incontestable; but it is nonetheless a misunderstanding, and a gross misunderstanding at that, when one wants, with Schelling, to seek a principle lying behind both\footnote{``If one immerses oneself in the thought that subject and object mutually condition and limit each other (cf.\ Schelling: Vom Ich als Princip der Philosophie, oder über das Unbedingte im menschlichen Wissen), then they both vanish as moments in an unconditioned unity, an absolute identity, a subject-object, in which the difference is originally grounded. But by leading the ideal-real development back to a principle that is just as abstractly dualistic as monistic, just as immediately subjective as objective, one is compelled to personify the dark real ground, to insinuate the understanding's flash of light into the selfless dreaming reason, to anticipate freedom without owning to the anticipation, and thus to transfer the predicates of subjectivity to the presupposition, which as yet can have no subjectivity.'' Den propædeutiske Logik. Kjøbhvn 1845. p.~206.

In other places, notably in ``System des transcendentalen Idealismus'', Schelling has undeniably laid emphasis on the I as the principle of knowledge; one sees in this the after-effect of Fichte. ``Der Begriff des Ich, d.\ h.\ der Act, wodurch das Denken überhaupt sich zum Objecte wird, und das Ich selbst (das Object) sind absolut Eins; außer diesem Act ist Ich nichts.'' \dots\ ``Ich als reiner Act, als reines Thun, ist im Wissen selbst nicht objectiv, deswegen, weil es Princip alles Wissens ist. Soll es Object des Wissens werden, muß dieß durch eine vom gemeinen Wissen ganz verschiedene Art zu wissen geschehen.'' \dots\ ``Ich ist nichts Anderes als ein sich selbst zum Object werdendes Produciren. Die Wissenschaft kann von nichts Objectivem ausgehen; sie muß vom `Nicht-Objectivem, das sich selbst zum Objecte wird', ausgehen.'' --- But to any clear and methodical development of how subjectivity originally relates to its objective content, and of what the principle of knowledge really means, the Schellingian philosophy, for all its attempts, has not attained.}, whether one then otherwise wishes to call this \opage{10} principle ``the indifference'' or ``the absolute'' or ``ein unvordenkliches Sein''. To make the indifference the principle is to presuppose an original unity in which the opposition between the subjective and the objective has not yet come to be, a fundamental unity from which this opposition itself was supposed to be derivable. Now at first glance it certainly looks as though by such a procedure one impartially wished to assert for the subjective an equal right with the objective, namely by deducing both from a third, higher thing. But the impartiality is only apparent; by beginning with the indifference as the third, higher thing, one has in actuality begun by taking sides with the objective. The third higher, or rather: the first deeper, which is to lie at the ground of both the objective and the subjective, is precisely, when we look more closely, nothing other than the first, still indeterminate objective itself, the objective from which everything is to be deduced, in opposition to the derived objective, coordinated with the derived subjective. To deduce the subjective is, in the principle itself, just as impossible as to deduce the objective. Instead of giving the objective, by abstract shadow-sketches of an indifferent background, a one-sided advantage over the subjective, we must on the contrary, acknowledging the originality of both, concede to the subjective the ideal priority, insofar as subjectivity, as knowing selfhood, nevertheless always is and remains that which overreaches.\footnote{``The exclusive criterion of subjectivity is comprehending selfhood, which consists in reflection; that of objectivity, on the other hand, a selfless objecthood, which lacks reflection. Without objectivity the subject's self-comprehending would vanish into emptiness, for all representations and intuitions refer to objectivity. Without subjectivity the concept of objectivity would be unthinkable, for the object is, properly speaking, neither for itself nor for another object, since all reflection is lacking. If we say: the one object refers to the other, then we employ a determination of reflection, and keep the one in mente when the other is represented. There is indeed an essential connection in the objective world; but since the individual object neither has itself in its power nor the others in mente, the objects cannot themselves objectify themselves for one another, cannot master their own totality, but must on the contrary be mastered and objectified by a comprehending subjectivity.'' Propæd. Logik. p.~204--205.}

\opage{11} Since now all the presuppositions of the beginning, the objective together with the subjective, lie hidden in the principle of knowledge, knowing subjectivity, the development of the original relation between knowledge and the knower must necessarily coincide with the systematic development of these presuppositions. According as the presuppositions then appear in abstract diversity, in mutual transparency and in totalized unity, the principle will develop for us as \emph{immediate}, as \emph{reflecting} and as \emph{comprehending subjectivity}. The distinguishing mark of \emph{immediate subjectivity} is on the whole its psychological standpoint. From this standpoint it works its way gradually deeper and deeper into universal determinations of essence and ontological problems. The distinguishing mark of \emph{reflecting subjectivity} is its conscious apprehension of its own ontological essentiality: reflecting subjectivity is essentially ontological. With reflection carried through, subjectivity is indeed led back to its first psychological point of departure; but in such a way that it now recognizes both its eternal ground and its temporal origin. That subjectivity, ontologically seen, is eternal and divine, but, psychologically seen, temporal and human, awakens the consciousness of a dualism in its own essence. This dualism must then find its explanation at the standpoint of \emph{comprehending subjectivity}. Only the concept of a perfect comprehending can lead thinking subjectivity to see its double face in the mirror of the idea, and thereby to see what it essentially is on which all insight, comprehending and knowledge rest.

\opage{12} \begin{center}
{\bfseries A.}\\[2pt]
{\large\bfseries Immediate Subjectivity.}
\end{center}
\addcontentsline{toc}{subsection}{A. Immediate Subjectivity}

Knowledge is unity-in-difference of thinking and being. With respect to pure being, knowledge is immediate intuiting; with respect to pure thinking, it is negation of all immediate intuiting. Intuiting determined by being is intuition; intuiting cancelled by thought is negation. Intuiting and thinking relate as being and nothing; immediate being cannot be thought; pure nothing cannot be intuited. To intuit nothing is the same as not to intuit; but to think nothing is by no means the same as not to think. By means of intuition, thought becomes determined by being; by means of negation, being becomes determined by thought: all determinations of intuition are positive, all determinations of thought negative. That being objectively converts into nothing and nothing into being, rests, then, originally on the fact that intuiting subjectively converts into thinking and thinking into intuiting. Already from this one sees what it means that subjectivity is that which overreaches. If objective being could not in any way be determined, it could not possibly be; but by what, now, is being to be determined? By its negation, by nothing, thus by thought! But thought is thought only by being thought, and it is thought only by there being a thinking subjectivity; insofar as being is determined by its negation, it is thus determined objectively by being determined subjectively. Or is being perhaps to be determined, not by negation, but by position, in such a way that being itself objectively distinguishes itself from being, as the one being from the other, as something from other? But to get something separated from other without negation is surely altogether impossible! The separation of something and other rests precisely on the fact that something is what other is not, and conversely. If being were now to determine itself objectively by itself, then it would also have to be able objectively to negate \opage{13} itself; but objective negation is a subjective determining, an act of thinking subjectivity. Or is it perhaps a logical blunder to let being, pure general being, determine itself into a different being, and thereby into separation of something and other: then let us for once turn the relation around and assume that there is originally a swarm of this-beings, absolutely existent things, from which all determinations of being and of knowledge are to be derived. Instead of beginning with a subjective unity, we thus begin with an objective multiplicity --- not of abstract and to that extent non-independent determinations of being, but of independent beings. In this multiplicity of independent beings each member then has a determinateness which it has neither given itself nor received from other beings, an original determinateness which can just as little be traced back to a determining as the existing of the member in question can be traced back to a coming-to-be; we then have in this abstractly objective presupposition a foundation of multiplicity that recalls at once Democritus's atoms and Herbart's reals.

But if a multiplicity of determinations which have not arisen, and cannot arise, through any objectively determining influence whatsoever is nonetheless to be determinable, then this can happen only in that the determinations themselves show themselves determining for a hovering apprehension, for a knowledge peculiarly determined by them, that is, for an intuition. That the absolute determinatenesses, indeterminable in themselves, are determining only for a corresponding intuition will readily become evident on closer consideration. The absolutely determinate atom, the immediately existent real, can surely in its absolute determinateness be determined neither by itself nor by the other, equally absolutely determinate, atoms. That the atom cannot be determined by any self-determining follows directly from the fact that its determinateness coincides immediately with its existing; that it then cannot be determined by other determining \opage{14} atoms either is self-evident, if the determinateness is to be one with the determinate existing, and this existing is to be absolute. So far is it from being the case that the absolutely determinate atoms could exert any influence whatsoever on one another and thereby in any way determine one another, that on the contrary they withdraw from every possible objective relation. Were they actually to come into relation to one another, they would naturally have to become dependent on one another; but by becoming dependent on one another they would surely have to become relative; but how should they become relative, when they are originally absolute?

The fixed idea of absolute atoms, immediately existent reals and the like, stems not from an actual objective being, but from a one-sidedly subjective thinking, a thinking that is beguiled by intuition.

That intuition, without knowing it, can beguile thought comes about, however, in that thought, without thinking of it, distorts intuition. Thought, abstract thought, never comes into immediate or direct contact with what is itself; every transaction between thinking and being is mediated by intuition. In its pure originality intuition is naive; it is determined by what is as it is, apprehends the given just as it happens to show itself, and sees nothing else. That no intuition can be determined by what is without in turn determining it, of this intuition itself in its innocence knows nothing; it is a secret that first dawns upon thought. But thought is negative: the negative is determined by the positive only insofar as it determines the positive and thereby itself. Thought is determined by intuition in that it negates intuition in order to be determined by what is; but since it is expressly through a negation, thus through its own negative self-determining, that it can be determined by what is, the determining thought has surely already thereby got the immediate determination of being into its power, insofar as it has got intuition into its power.

\opage{15} If thought, with this its formal power resting on negation, is nonetheless itself unclear about negation and the negative, then it easily happens that in negative naivety it immediately posits something positive and fixes it in pure immediacy; it thereby at the same time fixes itself, since the positive thus fixed is, after all, admittedly of negative origin.

It is not immediate intuition, naive in its immediacy, that has discovered the atoms and established the reals; on the contrary! the absolute atoms are thought intuitions, and one need only cancel the last remnant of something immediately intuitable in order to transform the absolute atoms into absolute positions, thus into pure reals. It is thought which, in the fancy that it has got immediate hold of what is itself, uses intuition and forces the representations, the individual intuitions, into its own web of fancy.

It is the contradiction hidden in change whose resolution is here a problem for thought. At the standpoint of naive knowledge there is no difficulty with change; in accordance with the daily experience that things constantly change, intuition and thought agree with each other that something must naturally become other, since it surely never stands alone, never rests on itself alone, but is everywhere in connection with other, leans on other, and, as this something determined in its peculiar appearance, always has an other hidden within it; intuition and thinking agree in a natural way that in existence there is indeed an opposition between something and other, but that the opposition is relative, and that it is precisely this relativity of the opposition that expresses itself in change.

It is otherwise when thought, with its abstract negation, its Not and Not-Not, begins to separate itself from intuition, while it is nonetheless still naively grown together with it. At this standpoint of negative naivety the problem of \opage{16} change, though under the pretense that it is a pure problem of being, becomes a one-sided problem of the understanding, a fixing task for self-fixing thought. --- Something, a this-being, is then, in its immediate opposition to other, which is surely also a this-being, given for intuition and thereby determined as object for thought; what does thought now do with every such this-being? It takes it positively, determines it negatively, and then, without heeding what is negative in the determination, passes off what is negatively determined as something positive in itself, as so absolute a position that there is to be in it not even a shadow of negation. This is naturally an enormous contradiction; and yet it is precisely this contradiction that abstractly critical thought lays at the ground, the foundation of self-contradiction on which it originally takes its stand in order to reject all the contradictions of existence, and therefore also those of change.

When A as a something, B as an other, each by itself, are intuited as a this-being, what can pure thought then teach us about each of these members? That A is A and not B, that B is B and not A! One sees at once that thought is altogether negative, that of itself it knows not the very least about the positive. Thought does indeed say that A is A; but when it is then asked: what, then, is A? it merely repeats what it has from intuition. What, then, does it have from intuition? That A, a this-being, is something positive; but by no means that it is something absolutely positive. The absolutely positive, a positive something in which there is not even a disposition to negation or otherness, is an invention of abstract thought, but an invention which thought nonetheless ascribes to intuition, in order to get it to count as what is. For by what means does it come about that the positive given by intuition becomes, although it is only relative, something absolutely positive?\footnote{It will be made clear in a later context that this, as it seems, absolutely positive for thought is, in its objective relativity, an element of power by which knowledge is determined, a real negation of the merely logical, an antilogical factor which is decisive for the relation between concept and existence. The Kantian ``Ding an sich'' arises precisely from the antilogical factor being specially apprehended and held fast by itself. The Herbartian mania with the absolute reals, and the endeavor to twist and turn the logical so that the determinations of thought take on the appearance of being pure determinations of existence, evidently stem from a reflection which is dazzled by the antilogical reflex that existence casts, and which fancies that this antilogical is precisely the true logical, and which then, in order to reach the true logical, labors to drive thoughts outside thinking and to reflect itself out beyond all reflecting. Let \emph{change} be an example. As transition from something to other, change is a categorical determination of everything existent. A, something existent, must thus, according to its concept, be able to run through the states: $A_0$, $A_1$, $A_2$ etc. The contradiction that A during the change is the same and yet not the same finds its logical resolution in the very concept of change. But if now the reflex of the antilogical is noticed, what change must then not take place in the apprehension of the concept: change! In existential immediacy $A = A$, which expressly conflicts with A's being able to become either $A_0$ or $A_1$ or $A_2$; further, $A_0 = A_0$, which conflicts with its being's passing over into its non-being, so that it could become $A_1$ etc. In this way thought, which has not yet seen through its relation to existence, can torment itself with wanting to understand itself as an existence, thus as non-thought, by wanting to understand itself as thought, thus as non-existence. The whole atomistic unnaturalness of thinking can moreover be seen from the ``reals''. The real is not a material but a logical atom. The material atom is, as object for thought, as thought's real negation, that about which one thinks but which does not itself allow itself to be thought, an antilogical element, and has as such its external existence in space. In thought's apprehension of the atom, the logical is thus united with the antilogical. The real, on the other hand, the logical atom, is burdened with the intolerable contradiction that its logical essence is to be absolutely antilogical.} It comes about, \opage{17} remarkably enough, through nothing but negations! Since A, a this-being, is according to intuition A, thought sees that when \opage{18} A is what A is, then A is not what A is not. With the same clarity thought sees that when B, a this-being, is according to intuition B, then this B can be what it is only by not being what it is not. The immediate difference between A and B in intuition, a difference which in all its immediacy is nonetheless relative, is thus in this way transformed into an absolute opposition. That A is not what A is not means that A is not B and in all eternity can never so much as approach becoming B; that B is not what B is not proves with the same emphasis that B in all eternity can never so much as approach becoming A. The absolutely positive: A and B, this something and this other, have then by the repeated negations, just as by logical conjurations, become separated by so yawning an abyss that every transition from something to other, thus every possible change, has in the matter itself become altogether impossible.

All difficulties with such pure concepts as: becoming, change, movement and the like stem essentially from a confused apprehension of the relation between intuiting and thinking. That pure thinking should, as Hegel assumes, be able to move itself without help from intuition and representation is just as untrue as that pure being should be able to set either itself or thought in movement. When Hegel nonetheless gets pure thought, as it seems, to move without any intuiting, this evidently happens only because the intuiting cancelled by thought is nonetheless preserved in thought, so that the intuition excluded from the movement is after all, as a logical confirmation of the old saying: \textit{naturam expellas furca}, included in the movement. How unthinkable, on the other hand, movement is when thought wants sharply to separate itself from intuition can be made clear from many sides, both by Herbart's and by the Eleatics' procedure. In every movement of thought, however pure, there is in actuality always an accompanying intuition, however infinitely fleeting, of space and time. If,
\opage{19} e.g.\ the transition is made from being to nothing, then one must of course begin by holding the two concepts being and nothing apart from one another, or, as it is also expressed, outside one another; but such determinations as \emph{apart from, outside}, evidently rest on an accessory representation of space; since, moreover, the transition itself: being, nothing, being, signifies that being alternates with nothing, just as nothing in turn with being, and follows from nothing, just as nothing in turn from being, this alternating and following must necessarily presuppose not merely an accessory representation of space, but also, and just as originally, an accessory representation of time: all actual thinking presupposes space and takes place in time.

Accessory representations of time and space come forward still more strongly in the determination of something and other, since something and other are more concrete expressions than being and nothing. If something and other are to be held apart from one another and beside one another as two ``reals,'' then the intuition of space forces itself upon us of its own accord. Now since the two reals, however decidedly they are set against one another, are nonetheless to be held fast, not as external things, but as absolutely positive objects of thought, indeed as absolutely positive thoughts, they must necessarily be held apart from one another, and that not in an ordinary space with natural intervals, but in a pure space of thought, a space in which thought can freely lay down its ``between,'' without anything ``between'' properly having to be thought of here on that account. In this way Herbart, for the sake of the reals, has enriched metaphysics with an intelligible space\footnote{``The one real, being for itself, lies wholly outside, and not inside or partly inside, the other. The many reals, or, what is already sufficient, two reals, can be together, but also not together. If we thus have A and B, we can connect with A the image of B, or conversely. Hereby we have obtained four concepts, namely the two actual ones A and B and their images. After this operation has been performed, one can again connect the actual A with an already connected pair, namely B and the image of A, so that these two are now together, while the image of B is not together with them. We want no interval between A and B; they are not together, but there is nonetheless nothing between them. Complaints that one cannot represent this to oneself are of no help here. The concept is to remain pure; and we desire no images, as though one could intuit them; but we demand concepts and their linkage.'' Against those who would introduce ordinary representations of space, Herbart expresses himself in the following way: ``When we demanded that one should think A and B as \emph{not being together}, those people, after their accustomed fashion, moved A and B apart from one another to an arbitrary distance, just as though there were already space enough from which one could insert an arbitrary magnitude between A and B. This insertion, no doubt unavoidable, we nonetheless forbid. Everything that could in any way be \emph{between} A and B, when it is not those empty images which we ourselves produced, is to vanish, and must be annihilated again in thought.'' The history of philosophy will not easily offer anything more arbitrary, far-fetched and perverse than what Herbart, otherwise acute and in other respects deserving of science's gratitude, has here (Metaphys. II. 119 ff.) put together. Cf.\ P.\ S.\ V.\ Heegaard, den herbartske Philosophie p.~25--30; p.~111--121.}.
\opage{20} But what, then, is such a space properly supposed to mean? In sharp separation from the space of the senses, intelligible space is of course a thoughtless thought, the intuition of an annihilated intuiting. That the space of the senses itself has an intelligible side must be conceded; but when intelligible space is taken as the thinkable side of the space of the senses itself, then it precisely reconciles intuiting with thought and furnishes an essential medium for the transition of the intuitable into the thinkable, of the positive into the negative, of being into its nothing, of something into its other.

While Hegel places movement within pure thinking, and Herbart strives to reconcile thought with a represented \opage{21} change by contriving artifices with the reals and intelligible space, the Eleatics have --- which, when discord is at all to be set between thinking and intuiting, is indeed also the only consistent course --- most strictly excluded the categories of movement, becoming and change from pure thinking, and thrust all such determinations of sensibility under mere opinion (δόξα). What can be intuited, then, cannot, according to this view, be thought. Here, e.g., a certain distance, a length of road, is intuited; now what, in sharp opposition to what is thus intuited, is the task of thought? To apply its negation, to set limit! Upon the given length of road countless limits are then set, and yet the length itself cannot be composed of limits; in this discord between intuiting and thinking we meet our ``Achilles'' again. The merely limit-setting thought is as though without movement; only with the help of intuition can thought move, only the self-moving thought can think movement.

``In the limit the limited is reflected. To think the limit without understanding by it something limited thereby is just as impossible as to grasp the meaning of a negation without thinking of that which is thereby negated. The immediate zero-limit, too, is posited directly under the presupposition that in intuition there is something given, whose limit it is. That such limits as points, lines and surfaces, although in and for themselves they are pure negations, are nevertheless \emph{intuited} in positive form: the point as the smallest part of the line, the line as the edge of the surface, the surface as a side of the body, is a striking proof that the limit goes into one with the limited. But when the limit nonetheless distinguishes itself from the limited as the negation from the positive, when this distinction is equally valid for thought and for intuition: then this contains the task of following that which is to be limited to the limit. Now how does one get from limit to limit? By a movement! Here there is \opage{22} originally no other way out, no other means. If one attempts in thought to mark off, e.g., on a line two limit-points, A and B, at a certain distance, then one must, even if unconsciously, perform a movement, in that, in order to intuit the arbitrarily chosen distance, one lets the gaze run through the way from A to B. The gaze may fix itself where it will, it meets on its way only limit, for every point in the line is in and for itself a limit-point \dots\ To set limit immediately beside limit is impossible, for as abstraction is made from the limited, the two points coincide \dots\ Between two limits following upon one another there is here, then, something limited. Now if this limited coincided immediately with the limits, so that there were here nothing, nothing at all, in between, then the limits themselves would have to coincide and could not follow upon one another; if the limited fell, as something immediately existent, between the limits distinct from it, then there would be here a finite distance, hence countless limits between two limits immediately following upon one another, which is a contradiction: in the dialectical unity of the limit and the limited all movement is grounded.''\footnote{Forelæsninger over ``Phil.\ Propæd.'' 1860--61. p.~56--58.}

To be sure, all thinking is in itself a movement, but the merely limit-setting thought moves in just such a way that it cannot itself notice its own movement. In order that the movement may be noticeable, there must be in the limit something more than limit. For insight into the movement of thought, the doctrine of the differential, the mathematical element, is for just this reason of decisive logical significance; for the differential, the intellectual moment of intuition, is limit and yet more than limit. The differential is neither zero nor finite magnitude, neither an abstract negative nor an abstract positive, neither a product of thinking by itself without intuiting nor of intuiting by itself without thinking: the mathematical element is \opage{23} a point of movement in punctual movement, the restless unity of intuition and thought, of the finite and the infinite, of the positive and the negative.\footnote{``It no doubt makes an essential difference whether the transition is intuited in space or merely represented as an arithmetical transition from the finite to $0$, from the finite to $\infty$; but this stands firm: without reflexes of space and time no movement, without movement no quantifying, without quantifying no quantum, without quantum no mathematics. To exclude the category of movement from pure mathematics and still retain the categories: ``limit,'' ``transition'', ``approximation'', is a contradiction. How the categories are chosen, stressed and applied in mathematics, insofar as it is technically a matter of the most expedient execution of exact operations, must of course be decided within mathematics itself; but the logical connection and dialectical transitions of the operational concepts are investigated in philosophy. If it then happens --- for it can indeed happen --- that a specialist, a routinier, whose mechanics of concepts has long since gone into one with his mechanics of operations, becomes loud-voiced on such an occasion and, by the loud-voiced application of the Saracen dilemma, expressly demands that philosophy shall give up dialectic and unconditionally pay homage to the mechanics of concepts: then such an occurrence cannot further interrupt, but rather hasten, the dialectical development.'' Forelæsninger over ``Phil.\ Propæd.'' 1860--61 p.~59 ff. Cf.\ Phil.\ og Math. p.~23--24; Begrebet Nøiagtighed p.~45--47. Prof.\ Steen's ``Afregning'', reviewed by R.\ N., p.~7, appendix p.~29.}

If it now stands firm that the true relation, both between thinking and being and between thinking and becoming, thinking and change, thinking and movement, can only be mediated by intuition, then it is seen from this that the relation of subjectivity to the actual world, and through it to itself, must then also, from the standpoint of knowledge, be determined by the principal relation between intuition and thought. But intuition, in its first concrete immediacy, is a sense-perception, namely a sense-perception of what is immediately objective, a first objectification. By giving itself over to sense-\opage{24}perception and sense-impression, knowing subjectivity will then reduce thought to an empty shadow of the sensible and misapprehend its own originality: this one-sidedness is fully expressed in thoroughgoing \emph{sensualism}. By unconditionally holding fast to its own infinite self-thinking, knowing subjectivity will not only reduce sense-perception and the sensible to vanishing accessory determinations in pure thinking, but will also set about transforming all objective actuality into a merely negative condition for its own self-development: this one-sidedness has found its consistent expression in \emph{subjective idealism}. As immediate subjectivity is thus seen to throw itself from the one extreme over into the other, it will become evident in the course of the movement what must then necessarily be assumed to lie in its original nature and to constitute its ontological characteristic. Thus:

\begin{center}
{\bfseries a)\quad Sensualist Subjectivity.}
\end{center}
\addcontentsline{toc}{subsubsection}{a) Sensualist Subjectivity}

By sensualist subjectivity one cannot here, in this connection, think of a consciousness lost in sense-intuitions and resting in the sensible; such a subjectivity, saturated with sensible content, has only slight significance for an essential determination of knowledge and the knower. The philosophy of sensibility, the sensualism with which we are dealing here, always has as its presupposition a definite, finite separation of sense-perception and thought, a separation which already recalls that the thought which abstracts from the sensible has become empty and faint and unfruitful. The principal task of sensualism is: \emph{to derive thought from sense-perception and sense-perception from the sensible}; in other words: \emph{to deduce subjectivity from the sensibly objective}.

How impossible it is to solve such a task is seen clearly enough from the attempts directed at it. When I.\ Locke, in order to explain the origin of ``ideas,'' begins with the empty soul as with ``a clean paper,'' a ``tabula rasa,'' then \opage{25} he nonetheless begins with a priori subjective presuppositions, insofar, namely, as he proceeds from a given separation between external and internal sense-perception, between ``sensation'' and ``reflection.''

Instead of beginning with an empty soul, that is to say: with a soul without dispositions, Locke on the contrary begins, though he will not properly acknowledge it, with a soul full of dispositions, indeed with a whole system of psychical facts. For he indeed finds sense-perception present as a fact and lets it count as a fact; he finds reflection present as a fact, and lets it count as a fact. And how many psychological facts are not then contained here, when we look more closely, in these two that are stated! So far is Locke from having been able to deduce the subjective from the non-subjective that he rather begins, not merely with a subjectivity, but with a subjectivity that is in original self-activity.

Locke's nearest predecessors, Baco, Hobbes, Gassendi, proceeded on the whole from the presupposition that our subjective concepts could be corrected and enriched by the facts of experience, without investigating how the matter really stood with the psychological origin of concepts. To this point Locke became attentive, and by this turn of the matter his Essay concerning human understanding may well be said to mark a turning point in the history of modern philosophy.

But so one-sidedly is Locke occupied in his attempt with the refutation of the Cartesian doctrine of ``innate ideas'' that, in his zeal to show how concepts must have been gradually formed under impressions that are owed to experience and consequently are a posteriori, he quite forgets to lay weight on the a priori and, if one will, ``innate'' in the activity of subjectivity, without which the formation of concepts, in spite of all sense-impressions and empirical facts, would be \opage{26} impossible. So far, then, is Locke from having been able to derive sense-perception from the sensible, that in order to determine what the sensible is he has expressly had to presuppose sense-perception; he is so far from having derived thought from sense-perception that on the contrary he has had to use thought in order to explain sense-perception; so far from having brought something active out of the merely passive that, without knowing it, he has given us proofs precisely that the passive subjectivity reduced to a tabula rasa originally is, and must be, active.

The immediate independence which Locke had conceded to ``reflection'' alongside ``sensation'' was, however, so conspicuous and so contrary to the essence of sensualism that ingenious attempts to explain it away could not fail to appear. It is in this respect a real advance when Condillac leads reflection back to sensation and thus makes earnest of deducing thought from sense-perception. Sense-perception, it is said, is in its first immediacy and indeterminate generality sensation. Now if a single sensation forces itself upon us with peculiar strength, attention arises; but the attention thus awakened is not a faculty of the soul distinct from sensation; on the contrary! it is sensation itself that has transformed itself into attention. Sensations alternate; the faculty of sensation, sensation, divides itself by means of attention between the sensation we have had and the sensation we now have. The sensation we have through the immediate impression is feeling; the sensation we have had, but whose real impression has vanished, is memory; thus memory, too, is not a faculty of the soul distinct from sensation: memory is a transformed sensation. Through the opposition between feeling and memory, attention in turn transforms itself, for here a double attention now forms, one which is determined by immediate sense-\opage{27}perception, and one which is determined by memory. But when there are here two kinds of attention, then there is indeed also comparison here; for to be attentive to two ideas and to compare them is one and the same. To compare, however, is impossible without noticing likeness and unlikeness; but to notice mutually equal relations is to judge etc. Through this transformation of sensations (transformation des sensations) there thus arise in essential sequence: first attention, from this comparison, and from this in turn judgment; in this way, according to Condillac\footnote{Traité des sensations,}, thought is supposed to let itself be deduced from sense-perception.

But this whole deduction evidently rests on a misunderstanding. It is not sensation that, without any other condition whatsoever in the soul, transforms itself into attention; but it is the attention slumbering in subjectivity, the disposition to attention itself, that is awakened by the sense-impression; it is not sense-perception that of itself gradually transforms itself into thought; but it is the thought hidden in sense-perception that little by little develops itself out of sense-perception. If sense-perception is in principle its own presupposition, then thought is no less so. As little as an animal sense-consciousness can, under favorable conditions, transform itself into human self-consciousness, just as little can the sensible thought of the animal be transformed into a free, rational, human thought. Only the thought which, according to the peculiar nature of the soul, is originally laid down in, with and under sense-perception can, through the self-development of subjectivity, develop itself out of sense-perception.

Suppose even that in a certain superficial sense one let the alleged deduction stand: would the task of sensualism then be solved by this? By no means! However much one thinks one can derive from immediate sense-perception, so long as one \opage{28} in no way finds oneself able to derive sense-perception itself from the sensible, the originality of subjectivity, the a priori validity of selfhood and of the self, is tacitly presupposed. But how, indeed, should it be possible to derive sense-perception itself from the sensible? To derive sense-perception from ``sensibility'' and to determine sensibility as a property of a ``ganglion tissue'' is a psychological evasion which in this connection betrays thoughtlessness. For ``the first thing peculiar to the soul as soul, the mark by which the soul-principle distinguishes itself precisely from the general principle of life, is sensation. But what is sensation? Immediate self-reflection, the self's unconscious self-discovery, its impression of itself under the impression of another, a state in which it finds itself posited.''\footnote{Forelæsninger over ``Phil.\ Propæd.'' 1860--61 p.~335.}

Had sensualism really been able to solve its task: to deduce sense-perception from the sensible, then materialism would undeniably have been its true and higher consequence; now, on the other hand, since it has been shown that the task could not be solved, because in order to explain the sensible one necessarily had to presuppose sense-perception, it is not materialism but spiritualism that in the stricter sense states the consequence; for spiritualism indeed begins precisely by fixing attention on the senses and sense-perception in order to understand sensible things. ``Sensible things,'' says Berkeley, ``are such as are immediately (immediatly) perceived by the senses; whereas the cause of things, which cannot be perceived by sense, must expressly be thought and, in being thought, be in consciousness. If I see that the one part of the sky is red, the other blue, I conclude that there must be a cause of this difference of color; but the cause cannot be said to be a sensible thing. Sensible things are aggregates of sensible qualities (combination of sensible qualities); \opage{29} these sensible qualities are not outside the soul, but in the soul.'' That things outside us are transformed into ``ideas in our own soul'' is of course not to be understood as though sensible representations rested on our arbitrary choice; they must rather be referred to a cause independent of us. But since cause and effect must be of like kind, and since a thing cannot possibly communicate to us what it does not itself have, it is just as unreasonable to suppose that a thing which does not itself have sensation should be able to awaken sensation as that an impression on the brain should be able to become a representation. There exist, then, no bodies, but only spirits; it is not a material cause, but an almighty spirit, that through all sense-images, representations, sensations and impressions lets finite spirits communicate themselves to one another.

If sensualism is in its essence spiritualism, and spiritualism in its original presupposition sensualism, while in their immediacy they nevertheless appear as two mutually exclusive extremes, then it goes without saying that the subjectivity which cannot discover the unity in the sharply expressed opposition must, by swinging between the extremes, be seized with dizziness and end in skepticism. The contradiction in the cause, that as a real concept it seems to spring from sense-perception when one would derive it from thought, and yet can as little descend from pure sense-perception as from pure thought, the contradiction which on the whole furnishes the foundation of Hume's critique of the concept of cause, is essentially one with the contradiction that spiritualism is at bottom sensualism and sensualism spiritualism.

\begin{center}
{\bfseries b)\quad Idealist Subjectivity.}
\end{center}
\addcontentsline{toc}{subsubsection}{b) Idealist Subjectivity}

If skepticism now rests essentially on this, that a subjectivity which has fallen into discord with itself must give up all assurance and with it all true knowledge, because it has given up assurance of its own a priori validity, then \opage{30} the first condition for regaining the lost knowledge must evidently be this: that self-certainty returns, that the disheartened subjectivity takes courage again. Subjectivity must by losing itself win itself; in order to win itself it must assure itself of the necessity of holding fast to its own forms and its own functions, in that it at the same time assures itself that all possible contradictions nevertheless in the last instance have their origin in a hidden self-contradiction.

``If one were to ask,'' says Kant, ``the cool-blooded David Hume, created properly for the equilibrium of judgment, what then moved him, through scruples which only laborious brooding could discover, to undermine the persuasion, so comforting and useful to men, that their rational insight suffices when the concept of the highest being is to be determined, then he would probably answer: nothing but the intention of bringing reason further in self-knowledge etc.'' It was, as is well known, the skeptic, ``der Zuchtmeister des dogmatischen Vernünftlers,'' who led Kant to a priori self-reflection and thereby to a sound critique of understanding and reason; and it was in turn, as is well known, the Critique of Pure Reason that led Fichte to the I as I and thereby to an idealist development of the infinite self-validity of subjectivity.\footnote{However great the significance of the Kantian critique of reason for the \emph{theory of cognition}, it can nevertheless in the present connection only be named as a transitional stage from Hume's skepticism to Fichte's subjective idealism. When one has considered what the Critique of Pure Reason proves, that there are: a priori sense-intuitions, a priori categories of the understanding and a priori ideas of reason, and what it means that a ``Ding an sich,'' as the uttermost reduction of an objectivity valid in itself, was placed in sharp opposition both to phenomena and to subjectivity's a priori forms of cognition, then there is at bottom nothing further here to dwell on; for the critique of reason has indeed not engaged in any attempt to deduce psychological subjectivity either from something non-subjective or from a divine presupposition.}
\opage{31} ``In every experience,'' thus the Fichtean train of thought begins, ``two factors are found set against one another, the I and the thing, intelligence and its object. If one abstracts from the I, one obtains \emph{a thing in} and \emph{for itself}, and must then take the representation as a product of the object; if one abstracts from the object, one obtains an \emph{I in and for itself}, and must then keep to the self-active consciousness which produces its content. The former is \emph{dogmatism}, the latter \emph{idealism}. These two are mutually incompatible; a third there is not: here one must choose. In order to judge between dogmatism and idealism, let one note the following: the I occurs in consciousness, the thing in itself is a fiction of consciousness. When dogmatism is to explain the origin of representation from the thing in itself, it performs the dubious feat of going behind consciousness and invoking being; but being effects only being, not consciousness; hence only that standpoint can be true which is chosen, not in being, but in consciousness, not in the object, but in the I itself.''

If sensualism, as a kind of real dogmatism, had begun by placing the unity of thinking and being in the immediate unity of thought with sense-perception and the sensible (nihil est in intellectu, quod non antea fuerit in sensu), subjective idealism now sets about dissolving everything sensible into sense-perception, all sense-perception into thought, and all thinkable being into the self-thinking of the I. ``If one concedes,'' thus subjective idealism begins, ``that the proposition A = A is evident of itself, then one concedes at the same time the capacity of consciousness to posit something. One does not posit that A is, but only that, if A is, then A is. If now, instead of looking at what is posited, we will look at the one positing, or rather: at the identity of the one positing and what is posited, then the proposition: A = A transforms itself into: I = I. While, instead of A = A, we could not \opage{32} say that A is, we can now, instead of: I = I, say: I am, and in this which is grounded upon itself see the ground of all action in the human spirit, the original distinguishing mark of the activity of consciousness in and for itself. Just as the proposition: A = A signifies a fundamental action, so the proposition: Non-A not = A likewise signifies a fundamental action, a primal action unconditioned as to its form. Yet, since the proposition Non-A not = A relates to the proposition: A = A as antithesis to thesis, it is already by this to be taken as conditioned as to its content, its matter; for antithesis is necessarily conditioned by thesis. In the first fundamental proposition: A = A, the form of the primal action was a \emph{positing}, in the second it is a \emph{counter-positing}. What Non-A is, is not yet known; only this is known, that Non-A is the opposite of A; but A is posited by the I; here only an \emph{I} is presupposed: what is opposed to the I is therefore \emph{Not-I}. From the two fundamental propositions a third follows with necessity. The task for the action which is expressed by the third fundamental proposition is given; but not the solution of the task. On the one side the I is completely canceled by the Not-I: the I cannot be posited insofar as the Not-I is posited. On the other side the Not-I, posited by the I, is itself an I, posited in consciousness. The I is therefore not canceled by the Not-I; the I is canceled, and yet not canceled: this is the contradiction. To solve this contradiction we must then try to find an X in which the mutually exclusive fundamental propositions can be united. But how can being and non-being, reality and negation, be united with one another without annihilating one another? By mutual limitation! The X in question is therefore a limit, and self-limitation the primal action of the I.''

From these first basic features the logical character of subjective idealism will on the whole be recognizable. The principle is the I, the pure I, and the beginning takes place in this, that the I at one and the same time posits, presupposes and deduces itself. The re\opage{33}lation between principle and beginning is thus expressed with quite another emphasis in Fichte's pure I than in Hegel's pure being. To make the pure I the principle is essentially the same as to make subjectivity the principle; to begin with the pure I is, in other words, to begin with pure knowledge. Whence comes it, then, that the Fichtean beginning has not been able to satisfy the demand of systematics? In what, then, does Fichte's original mistake consist?

From three sides the defect of ``the Fichtean I'' is conspicuous. In the first place there is here an ambiguous relation between the I itself and its existence, or, to use the current expressions, between the pure and the empirical I. The empirical I, my and your actual subjectivity, has its existence in its infinite limitation: in this limited existence my I falls outside your I, my thinking outside your thinking, my knowledge outside your knowledge. In existence itself there is here, then, an existing multiplicity, a multiplicity of I-existences; but where is the existing unity? In I-hood itself, the pure I! No, I-hood itself, the pure I, does not exist. The pure I is an abstraction which each existing I must perform for itself. The one abstracting is necessarily prior to his abstraction; the empirical I must therefore in actuality be prior to the pure I. But, according to the principle, it was indeed precisely the pure I that was to be prior to the empirical: here is a striking conflict between principle and actuality. In the second place, the original self-positing of the I, the primal action by which the principle of I-hood posits itself as an actual I, is an insoluble riddle. To solve the riddle one would then first of all have to obtain an answer to the question whether the original self-positing of the I is from eternity or an action in time. To let the primal action be from eternity is, in other words, to ascribe to the I a mythical-Platonic pre-existence; to let it take place in time is to let the unconscious furnish \opage{34} a presupposition for the conscious, is to let being go before knowledge. If the primal action is placed in eternity, then the actual I can only cognize its principle insofar as it can remember a pre-existence of its own; but such a remembrance of pre-existence is in actuality impossible. If the primal action is drawn down into time, then in every individual in whom a consciousness is to arise, a state of being indeed goes before consciousness; but ``being effects only being, not consciousness''; by letting a state of consciousness develop out of a state of being, one has at bottom canceled the principle. In the third place, the primal action of consciousness, whether one would otherwise place it in time or in eternity or in the boundary between time and eternity, is burdened with the fundamental contradiction of being an unconscious act of consciousness, an absolute act of knowledge without knowledge.\footnote{Fichte himself has a feeling of this, but helps himself by letting the ``Doctrine of Science'' make realism a concession. ``Die Wissenschaftslehre ist demnach \emph{realistisch}. Sie zeigt, daß das Bewußtseyn endlicher Naturen sich schlechterdings nicht erklären lasse, wenn man nicht eine unabhängig von denselben vorhandene, ihnen völlig entgegengesetzte Kraft annimmt, von der dieselben ihrem empirischen Daseyn nach selbst abhängig sind \dots\ Ohnerachtet ihres Realismus aber ist diese Wissenschaft nicht transcendent, sondern bleibt in ihren innersten Tiefen \emph{transcendental}. Sie erklärt allerdings alles Bewußtseyn aus einem unabhängig von allem Bewußtseyn vorhandenen, aber sie vergißt nicht, daß sie auch in dieser Erklärung sich nach ihren eignen Gesetzen richte, und so weit sie hierauf reflectirt, wird jenes Unabhängige abermals ein Produkt ihrer eignen Denkkraft, mithin etwas vom Ich abhängiges, insofern es für das Ich (im Begriff davon) da seyn soll. Aber für die Möglichkeit dieser neuen Erklärung jener ersten Erklärung wird ja abermals schon das wirkliche Bewußtseyn, und für dessen Möglichkeit abermals jenes Etwas, von welchem das Ich abhängt, vorausgesetzt: und wenn jetzt gleich dasjenige, was fürs erste als ein Unabhängiges gesetzt wurde, vom Denken des Ich abhängig geworden, so ist doch dadurch das Unabhängige nicht gehoben, sondern nur weiter hinausgesetzt. --- Alles ist seiner Idealität nach abhängig vom Ich, in Ansehung der Realität aber ist das Ich selbst abhängig; aber es ist nichts real für das Ich ohne auch ideal zu seyn; mithin ist in ihm Ideal- und Realgrund eines und eben dasselbe, und jene Wechselwirkung zwischen dem Ich und Nicht-Ich ist zugleich eine Wechselwirkung des Ich mit sich selbst.'' Grundlage der Wissenschaftslehre, Leipzig 1794, p.~261, 272, 273} The \opage{35} contradiction hidden in the primal action makes itself known clearly enough, too, in this, that with the self-limitation of the I there forms a self-contradictory selfhood, a subjectivity which can neither do without nor endure the objective. Insofar as the act of consciousness takes place unconsciously, a world of things is here produced in obscure necessity; and to immediate consciousness the things appear as externally existing objects, as actual objects. But every time the I becomes conscious of its own unconscious action, the objects vanish; it then turns out that the things are not things, but appearances of a Not-I varying to infinity, limits in the self-limiting I. The unceasing conflict between the I and the limits is from this point of view an unceasing conflict between the subjective and the objective, an original conflict in subjectivity itself.

\begin{center}
{\bfseries c)\quad The Problem of Subjectivity.}
\end{center}
\addcontentsline{toc}{subsubsection}{c) The Problem of Subjectivity}

The opposition between a sensualist and an idealist subjectivity is qualitative, and to that extent immediate; nonetheless the two so opposed views have this in common, that they both seek subjectivity in human self-consciousness, and up to a certain point agree in determining it psychologically. But psychological subjectivity in turn has a qualitative and to that extent immediate opposite in the abstractly logical. From the psychological standpoint subjectivity is essentially sensing, thinking, knowing; from the 
\opage{36} abstract logical standpoint, by contrast, subjectivity is only a general unity of concepts, namely the unity of concepts in which the subject has all its predicates in such a way that it is itself the sum of these predicates. The abstract logical subjectivity is thus, with its abstract form, a subjectivity in determinations of sense without sense-perception, in determinations of thought without thoughts, a selfless selfhood, a subject that exists only as object.

How impossible it is to solve the problem of subjectivity when one one-sidedly holds fast to the psychological starting point has already been shown in the assessment of sensualism and subjective idealism; but how impossible it is, on the other hand, to solve the problem when one one-sidedly proceeds from the universal-logical standpoint is seen, if possible, even more clearly from the misdirections to which objective idealism, and in particular the Hegelian panlogism, has already led. To begin with the universal-logical is to begin with an abstraction from everything psychological; it is to begin with universal thinking in objective, universal being. Objective, universal being then develops into objective, universal essence, which in turn develops into objective, universal substance, which then in turn develops into the objective, universal concept, which thereupon at once begins to develop into objective, universal subjectivity. Now, that such an objective subjectivity, a selfless selfhood, which---after having gone through ``the mediation of the syllogisms''---must seek its higher concretion in objects such as gold and iron, should of itself move further forward, first through the plant and then through the animal, until in man it at last reached itself in the form of self-consciousness, is an unreasonable assertion. It is only the fruitless attempt discussed above---to want to derive the subjective completely from the objective, consciousness from the unconscious, knowledge from non-knowledge---that we here meet again in another turn.

\opage{37} ``The determination of the subject as object,'' says J.~L. Heiberg,\footnote{Forelæsninger over Philosophiens Phil. eller den speculative Logik 1831--32 p.~102--103. (Prosaiske Skrifter, 1ste Bind 1861, p.~322--325.)} ``is the most important point in the whole of logic. It is the center around which all philosophy, from the oldest to the most recent times, has turned; it is the problem that all philosophers, each in his own way, have attempted to solve. In the present exposition it is the syllogism by which the subject is determined as object. Seen from the merely psychological standpoint, this determination has no other significance than that actuality must correspond to the judgment that comes about through a syllogism, since the syllogism is a means to the cognition of truth, while truth is placed in the agreement of the concept with actuality, so that, when the latter deviates from the concept, truth is missed, which can arise solely from the syllogism's having been incorrect. But this way of viewing it is highly one-sided and limited. Not to speak of the fact that truth is the concept itself, and that actuality, only insofar as it likewise is itself the concept, is true actuality, or more than phenomenon, so that there can be no question of the concept's agreement with what is only itself, this determination presupposes objectivity as something given, which the subject is to take possession of, and in this opposition the subject is then determined as human cognition. But seen, not from the psychological, but from the logical standpoint, the determination of the subject as object acquires a far deeper and more comprehensive significance. For logic does not presuppose the object, but lets the subject determine itself to it, just as it also does not determine the subject as a human, cognizing subject, but rather as the divine, creating one \dots\ The way in which philosophers have conceived the relation of the subject to the object characterizes the different systems. All idealistic philosophers, older as well as \opage{38} more recent, have regarded the subject as identical with the object, but have, from Plato to Schelling, presupposed this unity, whereas it is instead to develop itself, namely as the Idea. But among the philosophers who, instead of taking this unity as the basis, have regarded either the subject or the object as the absolute, Fichte and Leibnitz are the ones who have presented the opposition most sharply, the former by determining the subject as the absolute, and the object only as a limit for the subject; the latter, on the other hand, by determining the object as the absolute, namely as monad, i.e. substance itself, regarded as that which is in objective respect one and indivisible, but including subjectivity in the form of \emph{the different degrees of the faculty of representation}. Since the monads, like the atoms, constitute a plurality, and are independent, Leibnitz sublates this determination through the new determinations of a central monad or \emph{monad of monads}, and of \emph{the pre-established harmony}. Both Fichte and Leibnitz are idealists, but the former is a subjective, the latter an objective idealist. In the present exposition the idealism is both. Here, where the subject is determined as object, it is subjective; later, when the object returns to the subject (in the Idea), it is objective. It is the first of these determinations that we have to consider here. Now although this may well be said to be subjective, it is so in an entirely different way from Fichte's exclusively subjective idealism; for the opposite determination already lies in the fact that it is the same concept that is both subject and object, so that the concept is already seen as their unity, i.e. as the Idea. Thus the object is not seen here as a mere limitation of the subject's activity, a limitation about which it is doubtful whether one should rather regard it as a product of the subject, or (as on the psychological standpoint) as a presupposition which the subject cannot overcome and take possession of, since in both cases the object is determined only as the indifferent, \opage{39} as limit, and to that extent as the subject's quality; on the contrary, in the present exposition the object is the being which the subject, after having previously taken it up, now develops out of itself, and thus at once both the subject's presupposition and its product. Further, the object is not the subject's limit or quality; in this sense it is still determined as an other for the subject, and is only predicate in the immediate judgment; but inasmuch as the immediacy of the judgment is mediated by the syllogism, and this mediation in turn is sublated by the system of syllogisms and returned to immediacy, the predicate has become object, and the object is in this relation the whole objective world, which has the same reality as the subject, since it is not an other for the latter, but stands in the relation of immediacy to it. The determination of the subject as object thus has here the significance that all being proceeds from thought, and is nothing but visible thought.''

For a characterization of the problem of subjectivity, seen from a Hegelian standpoint, the above statement is one of the clearest and most unreserved that one could wish for in this connection; we have included it here in its full extent because it presents precisely such main points as are necessary if it is to become really clear what it is that is at issue here. The difference between the psychological and the logical standpoint is thus placed in this: that the former requires that the judgment which comes about through a syllogism shall correspond to actuality, whereas the latter lets the subject itself determine itself to object. How a judgment comes about at the psychological stage of consciousness is a well-known matter; but how does it come about in ``the speculative logic''? The judgment surely cannot possibly either judge itself to stand in logic or judge logic to procure for the subject the predicate it needs in order to determine itself to object. If there can be no talk here of ``the concept's \emph{agreement with} \opage{40} what is only itself,'' how then can there be talk here of a \emph{transition} of the concept into what is only itself? That the subject, seen idealistically, is identical with the object must be granted. But identity, after all, has difference; the difference between subject and object is more than difference; it is an express opposition, indeed an opposition totalized in the concept. But how then does it come about that subject and object can, in ``the speculative logic,'' by a little mediation, be exchanged in such a way that what at one time is subject, wholly and entirely subject, is at another time object, wholly and entirely object? It is expressly said that the object, which has the same reality as the subject, ``is not an other for the latter, but stands in the relation of immediacy to it.'' How then can the subject for its part stand in a relation of immediacy to that which in reality ``is not an other for it,'' is nothing other than itself? What kind of subject is it that is spoken of here?

It is said that logic ``does not determine the subject as a human cognizing subject, but rather as a divine creating one.'' How many subjects do we then actually have before us here? Here is a human cognizing subject, which is excluded by logic; a divine subject, which is acknowledged by logic; and then an excluding and acknowledging subject, the subject that determines ``the subject'' not as human but as divine; this subject determining ``the subject'' is logic itself; so here there are three subjects in all. But how, now, are these three subjects actually related to one another? Which of them is the independent one, and which are the dependent ones? At first glance one would suppose that ``the divinely creating subject'' must be the independent one; for that a subject is able of itself to determine itself to object undeniably points to plenitude of power and great independence. But on closer inspection, the divinely creating subject can after all determine itself to object only \opage{41} insofar as logic ``lets it determine itself to it''; but if logic really is the subject that determines ``the divinely creating subject,'' then logic must surely be the mighty, all-capable subject, hence the independent one! Yet this supposition too will not really satisfy. Certainly one can, by an abbreviation in the manner of expression, speak of what logic does, in contrast to what mathematics, physics, and the other sciences, each from its own point of view, must and ought to do; one can, in other words, quite well conceive the different sciences as thinking, willing, and acting subjects, without any misunderstanding being thereby occasioned, provided only that the abbreviation is not applied in such a way that the figurative expression is taken literally, and that precisely in a connection where the point is to get the concept ``subject'' determined with logical exactness. The subject in logic certainly cannot stir from the spot without logic's most gracious permission, cannot determine itself either to predicate or to object without logic's standing by it with counsel and deed, promoting it from one appointment to another, and, by steady advancement, determining it to decide upon every consistent and essential determination.

But is logic, ``the speculative logic,'' this subject which determines all logical subjects, predicates, objects, and subject-objects, then also really a logical self, an independent being, a thinking, judging, knowing subjectivity? If we assumed this, then we would surely also have to assume that ``the speculative logic'' could develop itself; but if it really could develop itself, then it would surely also have to be able to present itself; but if it really could present itself, then it would surely also have to be able to lecture itself, and then presumably also to write itself, to see itself through the press, etc. But with this ``immanent development'' of logic's logical subjectivity \opage{42} we have after all already gone beyond logic as subject to a subject hovering above logic with its subjects, namely to the subject that is in a way overlooked when it is ``determined as human cognition.''

Instead of divesting itself of human subjectivity along with human cognition, thinking consciousness must precisely set about finding the solution of the riddle in its own subjective being. If this consciousness can recognize its finitude and psychological limitation, then this recognition is precisely a proof that it is by that very fact raised above finitude and merely psychological limitation; if it can, in all its human imperfection, recognize the imperfection, then by this recognition it is already in itself raised above itself, and can through the imperfect discover the perfect, can see the divine in the human.

Between psychological and universal-logical subjectivity it must here absolutely be possible, if the problem is to be solved, to demonstrate a middle term of principle, a fundamental subjectivity from which the subjective in all relative subjectivities is to be derived, a principle-knowledge to which all psychological consciousness and finite knowledge can be traced back. The thought then lies close at hand of making the fundamental subjectivity one with the Spirit of God and deriving human knowledge from divine knowledge, man's self-consciousness from God's self-consciousness: man is, after all, created in the image of God. Nonetheless, such an immediate setting-up of the opposition between God's and man's subjectivity would only contribute to increasing the confusion. For by introducing God himself as a scientific presupposition we introduce a supernatural factor into science, a dogmatic principle that makes all scientific grounding impossible.

To be sure, from a ``theocentric standpoint'' one can, with a certain appearance of thoroughness, reason as follows: \opage{43} ``We human beings certainly know immediately what it is to know something, that it is expressly to perceive, to feel, in general to sense something, to intuit something, to represent something, and above all to think something; but we know at the same time that all our knowledge is limited and has a beginning in time. Now since the conscious cannot possibly be derived from the unconscious, knowledge cannot possibly be derived from non-knowledge, our limited knowledge must necessarily presuppose an unlimited and perfect knowledge; our spirit, awakening in time and under temporal conditions, cannot be a product of the blind elements of matter; it must necessarily have an eternal spirit, exalted above time and temporality, hence God himself, as its author.'' With such reasoning the psychologically limited subjectivity certainly seems at first glance to want to raise itself above its own limitation and, as it were, to go behind itself in order to seek its source; but the whole movement is only apparent. The reasoning subjectivity will, when we look more closely, in reality remain standing within finite bounds and, in the midst of finitude, continue to measure the infinite with its own psychological yardstick.

How, then, does the reasoning subjectivity come to insight into divine knowledge? Via negationis: by taking away all bounds; via eminentiæ: by thinking the imperfect forms of our knowledge raised to absolute perfection. By the via negationis we are then referred to an abstract qualitative opposition of finite and infinite, limited and unlimited; the unlimited must thus become recognizable at the limit, as the infinite at the finite. By negating the finite in consciousness one quite rightly arrives at a consciousness that one must call infinite; but a consciousness that is infinite in the sense that it has no finitely determinate content is empty. By sublating the limits of knowledge one undeniably arrives at a knowledge without limits; but a knowledge without any limits, hence without determinateness, \opage{44} is an indeterminable knowledge of something in itself indeterminable: τὸ ἄπειρον is in and for itself an ἄλογον.

In order to remedy the limitless emptiness to which the via negationis must consistently lead, the reasoning subjectivity next turns onto the via eminentiæ; the analogy between human and divine knowledge is held fast: divine knowledge corresponds to human knowledge as the perfect to the imperfect. But now human knowledge rests, after all, not merely on thinking but also on sense-perception. Without sight and hearing our cognition of the things in the world would be partly impossible, partly incomplete. For a perfect knowledge, then, a perfect sight and a perfect hearing will be required. From a theocentric standpoint, accordingly, one has not hesitated to ascribe to God both sight and hearing. But how, now, do matters stand with divine sight and divine hearing?

That religion in its peculiar figurative language speaks of God's eye, God's ear, etc., causes us no trouble in this connection, since religion is not and does not want to be science. But theology does want to be science; theology must surely, then, find itself able to give us a scientific concept of God's omniscience: what contributions, then, has theology made to a deeper understanding of the doctrine of God's eye and God's ear and the perfections of knowledge that correspond to them? None, none at all! To declare the figurative expressions anthropomorphisms and let the emphasis fall on pure omniscience is an easy matter; but the question is: what becomes of pure omniscience when one gets serious about removing all so-called anthropomorphisms?

Let us consider omniscience more closely. Divine omniscience must naturally be not merely a knowledge about all things but also a knowledge of the true essence of things: as things are known by the omniscient one, so \opage{45} and only so are they in themselves. From this it follows that, if there really exists a world of sensible things, there must in omniscience itself be a sensible knowledge; for only through a sensible knowledge can the sensible be known. If the sensible things are visible, tangible, audible, then the corresponding sensible knowledge must be conditioned by a seeing, feeling, hearing, etc.

But what a difference must there not now be here between human and divine seeing, between human and divine hearing, in short: between human and divine sense-perception! For human sight there is an unavoidable opposition between light and darkness, between white and black, between red and yellow, between transparent and opaque bodies; this opposition is decisive both for the subjective in sight and for the objective in visible things. But the opposition at the same time shows an essential limitation, an imperfection in our seeing. For it is, after all, an imperfection of our sight that we do not see what is in the distance as clearly as what is near us;\footnote{In order that this imperfection may be remedied, Dante's sight in Paradise must undergo manifold transformations. Thus in the 31st canto, where he sees Beatrice on the throne \textit{nel terzo giro dal sommo grado}:
\begin{verse}
Da quella region, che più su tuona\\
Occhio mortale alcun tanto non dista\\
Qualunque in mare più giù si abbandona,\\
Quanto li da Beatrice la mia vista;\\
Ma nulla mi facea; chè sua effige\\
Non discendeva a me per mezzo mista.
\end{verse}} it is an imperfection that we do not see things as clearly at night as by day; an imperfection that we cannot see through thick walls as easily as through windowpanes. \opage{46} Such imperfections must therefore, via eminentiæ, disappear for the eye of the omniscient one.

But as the imperfections disappear, all the corresponding oppositions, objective as well as subjective, disappear along with them; the eye of the omniscient one must therefore see all visible things in equal nearness and equal distance, equally clear, equally transparent, hence: without shadow, without color, without outline, without shape; for the all-seeing eye the visible things then, to put it plainly, run together into one: the all-seeing eye closes; the visible world disappears. We arrive at similar results by examining any of the other sense-perceptions and determinations of sense.

To be sure, we ought not in this connection to forget either that there is a speculative dogmatics which, with the assistance of the Holy Spirit---``inasmuch as we ascribe to the testimonium spiritus sancti not merely practical but also theoretical significance''---searches the depths of God, at least ``to a certain degree'' (for if the Holy Spirit really wished to concern itself with science, there would surely be no doubt that it would be in dogmatics that we should come to seek the real solution of the problem of subjectivity); but the half-true remarks which dogmatics can adduce on this occasion as signs of the many testimonies the Spirit is supposed to have given here, in ``Christian science'' as well as in ``the Christian world-view'' and in ``Christian art,'' all rest on a confusion of spheres, inasmuch as Christianity itself is confused with the realm of ideas in which it has found a reflection, and for whose development it has, as a world-historical power, furnished essential conditions.

Now if the dogmatics that ``searches the depths of God'' really had so much religious inwardness that it was able to pose and treat the question of eternal salvation as a decisive question of life, then the principle of subjectivity which it would in that case have to employ would be so existentially subjective that for that reason alone it could not be appropriated by \opage{47} science. This difficulty, however, falls away of itself. For the objectively moved dogmatics does not even have so much ethical inwardness that, even if it would condescend to make use of philosophy's leavings and to deck itself out with an imposing apparatus of learning, it would be able, from its objective standpoint, to treat the question of the immortality of the soul with any particular emphasis.\footnote{``The question of immortality is, then, a learned question? Honor to learning, honor to him who can treat learnedly the learned question of immortality; but the question of immortality is essentially not a learned question, it is a question of inwardness, which the subject, by becoming subjective, must put to himself. Objectively the question cannot be answered at all, because objectively one cannot ask about immortality, since immortality is precisely the potentiation and highest development of developed subjectivity. Only by truly willing to become subjective can the question properly come forth; how then should it be able to be answered objectively? Socially the question cannot be answered at all, because socially it cannot be posed, since only the subject that wills to become subjective can grasp the question and ask rightly: do \emph{I} become, or am \emph{I}, immortal? Look, several people can very well club together for various things; thus several families can club together for a box at the theater, and three single gentlemen can club together for a saddle horse, so that each rides every third day. But it is not so with immortality; the consciousness of my immortality belongs to me entirely alone; precisely in the moment I am conscious of my immortality I am absolutely subjective, and I cannot become immortal in company with two other single gentlemen taking turns. Subscription collectors who procure a numerous subscription of men and wives who feel a general need to become immortal likewise obtain no advantages for their trouble, for immortality is a good that cannot be extorted by a numerous subscription. Systematically, immortality cannot be proved either. The fault lies not in the proofs, but in one's not being willing to understand that, seen systematically, the whole question is nonsense, so that instead of seeking further proofs one must rather see about becoming a little subjective. Immortality is subjectivity's most passionate interest; precisely in the interest lies the proof; when one, quite consistently systematically, systematically abstracts objectively from it: God knows what immortality then is, or even what it means to want to prove it, or even what sort of fixed idea it is to trouble oneself further about it. If one could systematically manage to hang up an immortality, like Geßler's hat, before which we all took off our hats in passing, that is not to be immortal, or to be conscious of immortality. The system's incredible pains to prove immortality are wasted trouble and a ludicrous contradiction: to want to answer systematically the question that has the peculiarity that it cannot be posed systematically. It is like wanting to paint Mars in the armor that makes him invisible. The point lies in the invisibility, and with immortality the point lies in subjectivity and in subjectivity's subjective development.'' Afsluttende uvidenskabelig Efterskrift til de philosophiske Smuler af Joh. Climacus. Udgiven af S. Kierkegaard, Kjøbenhavn 1846. p.~126--127.}

\opage{48} On the other hand, the dogmatics that commends itself with the testimonium of the Holy Spirit has this questionable feature: that, despite its scientific cut, it everywhere on the preserves of science conducts itself as if it stood in an innate relation of superiority to science itself; it seems on the whole to regard its scientific investigations as a kind of diplomatic negotiations, about which it knows that they can never put it, as a spiritual great power, in any serious embarrassment; should it really here and there get into some embarrassment, it can always, in its capacity as a great power, produce its testimonium and break off with a peremptory decree.

How matters stand in scientific respect with the alleged testimonium of dogmatics can be seen here, among other things, from its treatment of God's omniscience. Should the testimonium really enable it to search the depths of divine knowledge, then it must surely be able to give a new and quite distinctive contribution toward \opage{49} determining the concept of knowledge and illuminating for us the relation of knowledge to the knower. This distinctive contribution of its would then have to be of such a nature that every honest thinker could grant that the new concept of knowledge really stood in the closest connection with the presuppositions: revelation and faith. Further, it would have to be possible to make it evident that the new knowledge built on faith in revelation was in its kind higher and more glorious than that with which science, without faith and revelation, must otherwise content itself. Only when this had become truly evident could there be talk here of the opposition between a believing and a non-believing standpoint within science itself. Those who rejected the higher concept of knowledge, conditioned by revelation and faith, would then, as surely as they wished to proceed scientifically, have to admit that the concept of knowledge, considered for itself, really was the higher; but that they rejected it solely because in their unbelief they could not hold fast the presuppositions on which its reality came to rest.

Now there is, after all, for every dogmatics, and especially for the speculative one, a quite peculiar call to unfold such a higher concept of knowledge when the discussion concerns God's omniscience. What, then, does the dogmatics ``that searches the depths of God'' have to teach us about the omniscient God? It is not without a certain suspense that we listen to a communication such as this: ``The omniscient God is the God manifest to himself, the God to whom all is manifest.'' One expects here to see the concept of revelation shed a higher light over the concept of omniscience. But when this higher light now consists solely in the dogmatic enlightenment that one has ``from of old represented God under the image of an eye,'' and that with reason, because ``God does not have an eye'' but is eye, because ``his essence is knowledge,'' then one becomes perplexed, and has difficulty in finding out wherein the enlightenment is really supposed to consist. If it is now the dogma: God is eye, that is to \opage{50} explain to us the dogma: God is knowledge, then one may well see in this a proof that there really must be ascribed to the testimonium spiritus sancti not merely practical but theoretical significance; for however often the holy Scripture speaks of God's eye, indeed of God's eyes, one will hardly find any place where either the Old or the New Testament has risen to the speculative knowledge that God does not have an eye but is eye.

But in what, now, does a proposition such as this: ``God does not \emph{have} an eye, but he \emph{is} eye'' find its scientific grounding? Is it perhaps self-evident to all, or at least to believers? Is it evident to all believers or only to those who do not lay so much weight on what the Bible teaches about God's eye, hence about the eye God \emph{has}, as on the circumstance that one has from of old represented God under the image of an eye, hence---speculatively!---of the eye God \emph{does not have}? Is it perhaps this last practical circumstance in which the theoretical significance of the testimonium spiritus sancti is especially to be sought?

``But,'' a fair-minded reader may object, ``why cling precisely to so pregnant and ingenious an expression as: `God is eye'? The meaning is obviously that his essence is knowledge!'' To this we have the following answer: as soon as it is merely granted that the whole ``dogmatic enlightenment'' comes down to an ingenious paraphrase of what is already posited with omniscience itself as presupposition, we at once give up the pregnant expression. The omniscient God is thus, according to the higher concept of knowledge of dogmatics, the God ``whose essence is knowledge,'' the God ``whose being is transfigured in eternal thinking,'' ``the God manifest to himself, to whom all is manifest''; the theoretical testimonium of the Holy Spirit thus comes down, through several pregnant paraphrases, to a dogmatic justification of the well-known principle of formal \opage{51} logic: everything is what it is, \emph{the omniscient one is the omniscient one}!

Yet---for this too we may not overlook---``the dogmatic enlightenment'' need not necessarily consist in nominal definitions of omniscience; the dogmatics ``that searches the depths of God'' shall be permitted to employ as many pregnant, figurative, half-figurative, prefigurative, non-figurative paraphrases of the concept of omniscience as it pleases, if only it will then, by virtue of the ``theoretical significance'' that is to be ascribed to the testimonium spiritus sancti, help us to a solution of the problems stated above. How is it possible that finite differences and sensible oppositions, such as the opposition of day and night, of light and darkness, can subsist for infinite knowledge, ``the all-penetrating seeing''?---As things are known by the omniscient one, so---we repeat it---so are they in themselves. That the omniscient one must know both light and darkness, both the sensible and the non-sensible, if these oppositions have true validity, and if there really is an omniscient one, is not the problem; for that is something that follows of itself and can be seen without a ``theoretical'' testimonium from the Holy Spirit. No, the problem is this: how is it possible that there can be an omniscient God, if there is in truth a finite, qualitative, and essential opposition between light and darkness, between transparent and opaque bodies? or: if there is in truth a finite, qualitative, and essential opposition between transparent and opaque bodies, how then is it possible that an omniscience can be thought to be? That problems of this nature are not contrived arbitrarily but force themselves upon us with necessity can be seen from any example whatever. What is light? Object and medium for sight, hence for seeing knowledge. And darkness? A bound for sight, object for a non-seeing knowledge, hence for a knowledge that ``an \opage{52}all-penetrating seeing'' is not possible! How then can an ``all-penetrating seeing'' at one and the same time both see in the darkness and yet see the darkness? How can omniscience, to cite a palpable example, subsist with a boulder? Seen in the light of omniscience, the stone that is opaque for us must surely become perfectly transparent, and all its material properties be transformed into determinations of knowledge, reflexes in pure knowledge; but in this way the stone disappears as stone. If the stone is to remain and keep its material properties: opacity, hardness, weight, etc., then it will present itself as an external finite existence, as an object for knowledge; but such an object of knowledge is at the same time a bound for knowledge: the object then limits the knowledge whose object it is; but omniscience, after all, tolerates no limitation. If one now wishes to appeal to the claim that omniscience will surely be able to master the boulder, because the omniscient one is at the same time the almighty one, then such an immediate dogmatic appeal to the unity of omnipotence with omniscience has, as it seems, not so much resemblance to a theoretical testimonium spiritus sancti as rather to a practical humbug, so long as no arrangements are made to shed light on anything concerning the relation between omniscience and omnipotence.

What contributions, then, has the objective dogmatics, objectively searching God's subjective depths, made to science toward a closer understanding of omniscience ``in relation to creation''? None, none whatever! For declarations such as these: ``Inasmuch as God cognizes all in the eternal unity, he cognizes it just as well in its inner oppositions and differences. God sets division between light and darkness; he knows essence as essence, semblance as semblance'' etc., are not, after all, scientifically grounded cognitions; on the contrary! they are purely groundless assertions, commonplaces that can be set up, accepted, and repeated by anyone, without the one who sets them up therefore in any way \opage{53} needing to have the testimonium spiritus ``theoretically''; unless, theoretically, it were to be a testimonium paupertatis.

Deterred by the failed attempt of dogmatics with the testimonium, one might almost feel tempted to put an end at one stroke to all quackery with subjectivity and its knowledge by setting up and holding fast the presupposition that the absolute spirit is in its omniscience as exalted above all human thinking as above all human sense-perception. But, not to speak of what is unscientific in assuming a divine knowledge that is to fall so entirely outside our human knowledge that we, strictly speaking, could have no knowledge of such a thing, the presupposition thus set up surely places so yawning an abyss, not merely between human and divine knowledge, but between human knowledge and truth, that on such terms it must become impossible for science to form a true concept of what truth is. Now if on this standpoint one cannot know what truth is, then neither can one know what belongs to a true knowledge. But a subjectivity that in this way gets the true outside itself and to such a degree goes astray from all knowledge of the truth that it finally does not know what true knowledge is, goes astray in itself, and as a subjectivity astray in itself cannot possibly escape the self-contradiction from which all other contradiction springs. If the self-contradiction is to be sublated and the problem of subjectivity solved, then it must here be possible to demonstrate in the principle a unity-of-opposition of original and derived subjectivity, a unity-in-difference of divine and human knowledge. But the principle of knowledge, the fundamental subjectivity in which divine knowledge can become human and human knowledge divine, cannot develop at the stage of immediate oppositions; here higher forms are required.
\opage{54} \begin{center}
{\bfseries B.}\\[2pt]
{\large\bfseries Reflecting Subjectivity.}
\end{center}
\addcontentsline{toc}{subsection}{B. Reflecting Subjectivity}

As knowledge is, so is the knower; as the development of subjectivity is, so is subjectivity. If for a characterization of reflecting subjectivity one demands only a consciousness that reflects, then we already have in sensualism, and still more in subjective idealism, telling evidence of a reflection that is inseparable from the essence of subjectivity; but when the expression ``reflecting subjectivity'' designates a distinct stage of subjectivity's logical self-development, it is not enough that subjectivity reflects; in order to become transparent to itself, it must expressly reflect its reflecting. And not even that is enough; for subjective idealism has sufficiently demonstrated that a subjectivity can reflect its own reflecting in so abstract, formal and one-sided a way that, by incessantly consuming its content, it feeds upon itself and consumes itself. The subjectivity that is clear in self-reflection must know how to develop the forms of reflection in indissoluble unity with a content reflected in the forms; with the content so reflected it must not merely mirror itself in its other, but see through the opposition of itself and other so completely that otherness, from its side, can conversely be seen to reflect its independence in being reflected in selfhood, and the objective is confirmed by the subjective as the object of knowledge is by the knower. That precisely this is the task becomes evident as soon as one properly begins to reflect upon reflection. What, then, is reflection?

According to ordinary usage, thinking and reflecting pretty much coincide: to reflect on something is to think about something; the difference is at any rate only that by reflection one usually designates a particularly conscious thinking, a deliberation, an express use of afterthought. Further, according to ordinary usage, one distinguishes between \opage{55} thinking of something and reflecting on something: ``reflecting hereupon,'' etc. From this point of view, then, reflection with its various meanings is a subjective act, just as thinking is. In an entirely different sense the word reflection occurs in optics, where the talk is of the reflection of light, in the sense of the ``throwing back'' of the rays of light; this reflection is wholly objective. Reflection in the objective and reflection in the subjective sense are, according to the usual way of representing things, so wholly heterogeneous that they scarcely even seem to concern one another. It is otherwise in the absolutely objective, immanent development of the concept; here one has gone to the opposite extreme of lumping subjective reflection together with objective. The latter is just as one-sided as the former. From the standpoint of logic, therefore, it becomes an essential task to develop the difference between reflection and thought and at the same time to demonstrate the unity, namely the unity-in-difference, of logical and optical, of subjective and objective reflection.

If we consider the relation between thought and reflection more closely, the question at once arises whether reflection is to be determined as a modification of general thought, or thought as a particular moment, a side, of reflection itself. In Hegel, reflection is expressly determined as a particular kind of thought, an objective intermediate form between immediacy and speculation. Thus all the categories of being correspond to immediate thought, all categories of essence to reflection, and all the categories of totality belonging under the concept to speculation. What holds of the trichotomy on the large scale holds also on the small, and repeats itself in each single triad. The place of reflection in the Hegelian system is so typical that every deviation from the type betrays itself as an inconsistency. ``In Hegel'' --- says Heiberg\footnote{Perseus No.~2, p.~44. (Prosaiske Skrifter, 2nd vol., 1861, p.~165)} --- ``absolute being is presented in the categories: \opage{56} a) being, b) nothing, c) becoming. --- But this order must be regarded as a small oversight, for it conflicts with the whole remaining architecture of the system. Becoming is indeed the unity of being and nothing, but only the dialectical unity of the two, their --- as Hegel himself calls it --- \emph{restless} unity, and thus not the resulting or speculative unity. Becoming is in the Hegelian system evidently the fundamental category for everything finite; but precisely for that reason it must, according to the entire method of the system itself, be kept in the place of finitude.'' But the place of finitude is, in this context, precisely the place of reflection. That reflection has in this way become an intermediate form of thought, a transitional form, designated by ``the dialectical,'' from categorial immediacy to the resulting ``speculative'' unity of developed difference, would not have been touched on here at all, had not this subordination of reflection to thought, as of the finite to the infinite, as of understanding to reason, had its ground in an unmistakable overestimation of the abstract essence of ``pure thought.''

In sharp opposition to intuiting, sensing and representing, pure thinking by itself is one-sided, just as pure negation is one-sided over against immediate position. Thought must, according to Hegel, be the negative in which all sense-perception and representation vanish: ``he who properly begins to think loses sight, hearing and all the senses''; but if the negation is to be nothing but an abstraction, if one will really be in earnest about casting away all the heterogeneous elements of sense-perception, representation and intuition, so that thought itself, as the universal, may move from the universal through the universal to the universal, then one may spend as many negations and negations of negations on this caput mortuum of abstraction as one likes; it is dead and stays dead. If, on the other hand, the liberation of thought from sense-perception and representation is to be understood --- as Hegel indeed also insists --- in such a way that the elements of representation which thought by \opage{57} its negation sublates are preserved in concrete thought, then thought is thereby at once transformed into reflection, and there is no longer any ground for giving thought, as abstractly opposed to representation, any unconditional precedence over representation. Through reflection the positive becomes a reflex of the negative, but in such a way that the negative thereby at once becomes a reflex of the positive; through reflection intuition becomes a reflex of thought, but in such a way that thought thereby at once becomes a reflex of intuition; through reflection being becomes a reflex of nothing, but in such a way that nothing thereby at once becomes a reflex of being. Instead of determining reflection as a subordinate form of pure thinking and taking the side of abstraction, we must on the contrary recognize that it is precisely reflection which is the developed thought, integrated through its oppositions, the thought which alone is fitted to pass over into concretion.

If it is already necessary from the subjective side to recall that all thinking in concreto is a reflecting, it is, from the objective side, with express regard to the Hegelian misdirection, still more necessary to recall that all reflecting, logically understood, is a thinking. For however strictly Hegel everywhere holds to ``das Denken und das Denken und das allgemeine Denken,'' it nevertheless happens to him that in his objective logic he comes upon thoughts of which he does not quite know whether they are properly thought or not thought. ``Wenn man sagt,'' --- it says in the ``Vorbegriff'' to Part 1 of the Encyclopaedia\footnote{Hegels Werke, 6ter Band. Berlin 1840, p.~45.} --- ``der Gedanke, als objectiver Gedanke, sey das Innere der Welt, so kann es scheinen, als solle damit den natürlichen Dingen Bewußtseyn zugeschrieben werden. Wir fühlen ein Widerstreben dagegen, die innere Thätigkeit der Dinge als Denken aufzufassen, da wir sagen: \opage{58} der Mensch unterscheide sich durch das Denken vom Natürlichen. Wir müßten demnach von der Natur als dem Systeme des bewußtlosen Gedankens reden, als von einer Intelligenz, die, wie Schelling sagt, eine versteinerte sey. Statt den Ausdruck \emph{Gedanken} zu gebrauchen, ist er daher, um Mißverständniß zu vermeiden, besser Denkbestimmung zu sagen.'' It is a great problem that Hegel touches on in these few words, a fundamental problem that already hovered before Plato: ``If the ideas,'' it says in Plato's Parmenides, ``are only subjective concepts (νοήματα), and thought must necessarily have its object, then one must assume either that each object consists of thoughts and that everything thinks, or else that there are indeed thoughts but no thinking (ἀνάγκη, εἰ τἆλλα φῂς τῶν εἰδῶν μετέχειν, ἢ δοκεῖν σοι ἐκ νοημάτων ἕκαστον εἶναι καὶ πάντα νοεῖν, ἢ νοήματα ὄντα ἀνόητα εἶναι).'' That things themselves should be able to think, or that there should be able to be thoughts in things which were not thought, Plato found meaningless to such a degree that, in order to be rid of so unreasonable an assumption, he did not hesitate to put the eternally existing ideas in place of the unthought thoughts. What, now, has Hegel, with his incessant insistence on ``das Denken und das Denken und das allgemeine Denken,'' actually done to solve that ancient Platonic problem? Nothing! Instead of sharpening the problem, he has on the contrary blurred it with a reasoning remark such as this: ``Statt den Ausdruck \emph{Gedanken} zu gebrauchen, ist es daher, um Mißverständniß zu vermeiden, besser Denkbestimmung zu sagen''. What in this respect could blind a thinker otherwise so energetic, vigilant and persevering as Hegel is not difficult to show; it is evidently the unclear relation in which he from the first places reflection to thought proper, and the use --- ambiguous because of the unclarity of the relation --- that he then makes of objective reflection. A striking proof of this is the following. Of the category of existence it is said \opage{59} in the Logic: ``Die Existenz ist die unmittelbare Einheit der Reflexion-in-sich und der Reflexion-in-Anderes''; the thing, that which exists, is now, in consequence of this, determined thus: ``Die Reflexion-in-Anderes des Existirenden ist aber ungetrennt von der Reflexion-in-sich; der Grund ist ihre Einheit, aus der die Existenz hervorgegangen ist. Das Existirende enthält daher die Relativität und seinen mannigfachen Zusammenhang mit andern Existirenden an ihm selbst, und ist in sich als Grund \emph{reflectirt}. So ist das Existirende Ding.''

It is incontestably clear from the Logic that Hegel not merely lets the existence of things rest on objective reflection, but even makes things themselves outright reflecting, insofar namely as every thing reflects itself into itself by reflecting itself in its properties, etc. But when Hegel does not hesitate to make things into reflecting beings, while he does hesitate to make them into thinking beings, it is plainly to be seen from this that he tacitly makes an essential separation between reflection and thought, and indeed a separation according to which, despite all appeal to ``das Denken und das Denken und das allgemeine Denken,'' it is after all not thought\footnote{That according to ordinary usage there may be reason to use the word \emph{thought} as the more general expression for the logical activity of spirit may readily be granted; only one must remember that thinking proper, determinate thinking, is a reflecting.}, but reflection, namely optically objective reflection, that from first to last is taken as the universal.

For a closer determination of the relation between thinking and reflecting, and thereby also of the relation between subjective and objective reflection, it will therefore be necessary to let reflection clarify itself from the subjective side as \emph{reflection of the understanding}, from the objective side as \emph{essence-reflection}, and, through energetic reciprocal determination of the subjective and the \opage{60} objective, as \emph{actuality-reflection}. The unity-in-opposition of subjectivity and the objective then develops at once in such a way that reflection of the understanding is seen to be essentially a self-reflection conditioned by otherness, essence-reflection a reflection of the other grounded in selfhood, and actuality-reflection a reflection of independence, determined for otherness by the overreaching selfhood.

\begin{center}
{\bfseries d)\quad Reflection of the Understanding.}
\end{center}
\addcontentsline{toc}{subsubsection}{d) Reflection of the Understanding}

All understanding of another is essentially self-understanding; for understanding rests precisely on this, that the understanding subjectivity stands, so to speak, before its object as before an other, something foreign, and yet in this other recognizes its own, in this foreign thing discovers itself: in understanding, otherness is reflected in selfhood precisely in such a way that subjectivity expressly knows itself as reflecting. In understanding, the understanding at the same time reflects itself. To be sure, the reflecting understanding can be developed only through the development of an objective content which is at the same time understood; but the development of the object is nevertheless, in reflection, different from the understanding that develops itself thereby. Here, within understanding itself, an intelligent distinction must be made between an understanding of the object and an understanding of the understanding by which the object is understood. The sharper the distinction, the clearer in turn is the reflection of what is distinguished. If it were an impossibility for the understanding to reflect its eternal determinations of form in reflected separation from its real content, there would then be so infinite a gulf between the human and the divine understanding that the understanding would thereby necessarily become absolutely a riddle to itself. But who does not see that the determinations of form which prove absolutely necessary for the human understanding must also have validity for the divine, since they belong to the essence of the understanding? It is thoughtless if anyone should say: ``our understanding is indeed, in its human limitedness, \opage{61} compelled to distinguish between being and becoming, rest and motion, essence and seeming, ground and consequence, whole and parts, inner and outer, cause and effect, means and end, and so forth; but it is not thereby said that such a distinction should be necessary for the divine understanding as well. We, with our imperfect understanding, must of course grant that the whole, mathematically regarded, is greater than a part and equal to the sum of its parts; but what proof have we that the divine understanding, incomprehensible to us, does not perhaps reverse the relation and make the part greater than the whole and the sum of the parts smaller than each single part? We, with our human understanding, are indeed unable to think an essence without phenomenon, a thing in itself (ein Ding an sich) without properties, an absolute inner without a corresponding outer; but a divine understanding, an understanding that is one with a divine power, might for that reason well find itself able to cast away the phenomena and keep the essence, annihilate the properties and conserve the thing, abolish the outer and settle down in the inner?'' --- That such a separation of divine and human understanding is thoughtless and unintelligent needs no special proof; the proof lies in the matter itself. The endeavor to make the understanding superhuman and divine is, through the projected separation, transformed into an endeavor to make the divine understanding into a monstrosity of the understanding, a supreme unintelligence. ``Yes, to our limited power of comprehension,'' it is said, ``the divine understanding must necessarily look like unintelligence; for otherwise it would not absolutely surpass our understanding.'' But the ambiguity in such turns of phrase ceases to dazzle when we bethink ourselves of the decisive difference between: grasping the infinite reflection of the understanding in its infinite content, and grasping reflection of the understanding itself in its formal intelligibility. The first task designates the opposition, \opage{62} the last the unity, of the divine and the human understanding.

Insofar, then, as the whole system of real determinations of difference which particularly concern the content falls away or, more precisely, loses itself and vanishes in reflexes of the understanding's ontological self-reflection, we have in the ontological selfhood so determined the principle of the understanding as a principle of subjectivity. The self-development of the principle, therefore, can be comprehensively grasped only in the same degree as reflection of the understanding itself is comprehensively developed.

\begin{center}
{\bfseries α)\quad Positing and Presupposing Reflection.}
\end{center}
\addcontentsline{toc}{paragraph}{α) Positing and Presupposing Reflection}

Were anyone to ask: ``how old, then, is the understanding of existence? when did it come to be? how long can it last?'' then of course anyone of understanding would at once notice the absurdity of such questions. The understanding itself, the understanding of existence, can surely have no finite coming-to-be: it is and must be eternal. But what, now, are we to say of an understanding that exists for an entire eternity without understanding either itself or anything else, thus: without the least understanding either of itself or of anything else, thus: an eternal understanding without any understanding? That view looks somewhat unintelligent. In a temporal, finite way there may indeed be much that is intelligent in things themselves which your understanding and mine does not grasp; but here the talk is not of an understanding like yours and mine, but of the understanding itself, the eternal understanding. What holds of thought, intuition and representation holds in this respect also of the understanding. Without thinking no thought, without intuiting no intuition, without representing no representation, without understanding (as act and as achievement) no understanding (as faculty).

From the standpoint of psychology it holds that the understanding subject is one thing and the subject's understanding another; but what is finite in this separation belongs only to finitude. Seen in the principle, that is, from the standpoint of ontology, \opage{63} the understanding that understands itself is subjectivity itself. It is the divine understanding --- the νοῦς of the ancients --- that eternally thinks itself, and this thinking is, as Aristotle expresses it, the thinking of thinking: αὐτὸν ἄρα νοεῖ, εἴπερ ἐστὶ τὸ κράτιστον, καὶ ἔστιν ἡ νόησις νοήσεως νόησις.

But if the understanding itself is original subjectivity, how does it come about that one does not see the eternity of subjectivity as immediately as that of the understanding? It evidently comes from this, that here one has to contend with a twofold misunderstanding, a psychological and an ontological one. The psychological misunderstanding is that all thinking subjectivity, all knowledge and consciousness, must as a matter of course be analogous to human subjectivity, conditioned as it is by the activity of nerves and brain, and to its limited knowledge and consciousness. The ontological misunderstanding is that the eternal being of the understanding should be able to be immediate, like the being of an eternal space, whereas one cannot conceal from oneself that the being of subjectivity is and must be infinitely reflected. To attribute to the understanding an original immediate being, be it in a given order of things or in matter and its forces, is unintelligent in the very highest degree; for the infinite understanding is only in an infinite reflection. Psychology says that the understanding is a faculty of reflection; ontology sees that the understanding itself is reflection. The problem of the originality of subjectivity is precisely the problem of the originality of the understanding, the problem of the eternal νοῦς as ἡ νοήσεως νόησις.

If such categorial expressions as ``thought thinking itself,'' ``the understanding understanding itself,'' ``the concept comprehending itself'' are to be more than figures of speech by which one seeks to avoid the difficulties with the subjective in order to give oneself undisturbed to the consideration of the objective, then it is in truth not enough that one now and then remembers that there is a thinking subjectivity, and \opage{64} then, in all psychological-ontological ambiguity, says with Hegel: ``Das Denken als Subject vorgestellt ist Denkendes, und der einfache Ausdruck des existirenden Subjects als Denkenden ist Ich\footnote{Hegels Werke, 6ter Band. Berlin 1840, p.~34.}.'' Already Fichte's beginning of a construction of the pure I shows clearly enough what it means that the understanding which understands itself is originally a subjectivity. The first two propositions, $A = A$, and Non-A not $= A$, are, considered by themselves, the well-known principles of formal logic, namely the principles of identity and of contradiction, and thus pure propositions of the understanding, principles to which the reflecting understanding would lead all reflections back in order properly to understand itself. But whereas formal logic only sets up the principles in such a way that it holds fast, in a separate presupposition falling outside the logical positing, to the psychologically thinking subject, a subject that uses its understanding without seeing its essential unity with the understanding itself, it is on the contrary Fichte's merit to have shown from the ground up that the thinker does not stand outside the thought, that it is precisely the self-understanding subjectivity that, in the very first intelligent reciprocal determinations of positing, presupposing and opposing, is essentially concerned with itself.

According to formal logic, positing something and presupposing something must be regarded as two things that fall outside one another. To posit something is to give it existence; what I myself posit I have in my power: my thought can hold it fast or let it go as it sees fit. To presuppose something, on the other hand, is to let something be given to one, to take it as a given. The reality of the given shows itself precisely in this, that what is presupposed in thought governs the presupposing thought and makes it dependent on itself. In its positing, in other words, the understanding shows itself as active, in its presupposing as passive. As long as positing and presupposing are \opage{65} set against each other in this way, it is certainly an impossibility to grasp the original self-positing of subjectivity, the primal act by which ontological I-hood is from eternity I $=$ I.

In order in truth to be able to be a self, subjectivity must posit its own selfhood; but in order that a subjectivity should actually be able to perform such a self-positing, it must of necessity already be beforehand; a subjectivity that is not cannot possibly posit either itself or anything else. Actual subjectivity must, in order to be able to posit itself, begin by presupposing itself, begin by finding itself as something already posited. Thus, thesis: subjectivity is subjective only in virtue of an original self-positing; its positing goes infinitely before its presupposing. Antithesis: the original self-positing presupposes a positing self; the presupposing of subjectivity goes infinitely before its positing. This antinomy would certainly be insoluble, and the problem of the ``primal act'' of the I a hyperspeculative problem, if the reflecting understanding could not naturally recognize in that seemingly mystical primal act its own well-known operations, its own positing and presupposing.

For the pure reflection of positing and presupposing it is indifferent what subjectivity may otherwise be able to posit, whether a whole solar system or only the letter A. The point is that every positing rests on a presupposing. One cannot, to recall Fichte again, posit $A = A$ without presupposing I $=$ I; but neither, conversely, can one posit I $=$ I without presupposing that, if A is, then $A = A$; in that $A = A$ rests on positing, I $=$ I rests on presupposing, and conversely; positing and presupposing are thus in reflection not two acts, but one: presupposing a beginning positing, and positing a self-completing presupposing. By reflecting itself in its presupposing, every positing does indeed reflect itself in its opposite; but since its \opage{66} opposite is precisely a presupposing, and presupposing, as a beginning positing, likewise reflects itself into itself by reflecting itself in its opposite, every positing which thus continually reflects itself in an opposition reflecting its opposite must assert for itself an infinite self-duplication, and in this self-duplication show itself as original self-positing. On original self-positing, the unity of positing and presupposing transparent in reflection, there then rests selfhood's fixing of itself and its other, and on this fixing there rests in turn all essential determining, of the objective as well as of the subjective.

\begin{center}
{\bfseries β)\quad Opposing and Unity-Positing Reflection.}
\end{center}
\addcontentsline{toc}{paragraph}{β) Opposing and Unity-Positing Reflection}

If original self-positing is held fast in the form of a presupposition, the corresponding positing of the other must be held fast as a consequent proposition. All positing of the other rests on self-positing; but whereas self-positing, considered by itself, keeps to the uniform, the positing of the other on the contrary designates something multiform and diverse; for the other is reflected not merely as the opposite of the self, but in otherness itself there is reflected at the same time a self-repeating opposition of other and other (τὸ ἕτερον) to infinity: positing reflection is hereby determined as opposing. What in existence itself is the originally first: unity or duality, the finite or the infinite, indifference or the different, the imperfect or the perfect? This many-sided and ambiguous question, which has so often repeated itself in the history of philosophy and will continue to repeat itself, cannot possibly receive any satisfactory answer so long as one has not yet properly accounted for the unity-in-opposition of principle and beginning; but in order to get clear about this unity-in-opposition, one must above all get clear about the opposition of opposing and unity-positing reflection.

\opage{67} If one proceeds without further ado from the presupposition that unity, pure absolute unity, must absolutely be posited as the properly first, then the subsidiary representation easily creeps in that pure, absolute unity should be, if not exactly an entire eternity, then at least some moments within eternity itself, prior to the duality derived from it and to the multiplicity derived in turn from duality. This view is untenable. It is evident, when one reflects, that the unity which is to be able to develop duality, multiplicity and opposition out of itself must already have opposition, multiplicity and duality in itself. As little as the understanding can begin by positing an original duality without presupposing the opposite of duality, an original unity, just as little can it begin by positing an original unity without presupposing the opposite of unity, an original duality.

The fundamental contradiction, that duality always shows itself as the first when one would begin with unity, and unity as the first when one would begin with duality, finds its resolution in this, that duality can begin only insofar as unity is in the principle, since unity can begin only insofar as duality is in the principle. The primal act, the original self-positing in which self-reflecting subjectivity has eternal subsistence, must on no account be confused with a positing of I $=$ I corresponding to $A = A$; such a first general positing is only an initial proposition, not a primal act but an act of abstraction, a one-sided moment in the many-sided, infinitely concrete primal act itself. Nor can one say that the primal act was from the beginning completed in the sense that it could be set up as a first member in an infinite series of acts; one could with equal right maintain that the primal act was first completed with the last member of the series; indeed, one could perhaps with still greater right maintain that the member which designated the completion of the primal act must properly constitute the middle of the series, infinite \opage{68} in both directions. The truth is that the primal act is precisely the act of subjectivity repeating itself in all acts of the understanding, the unity-positing reflected in all the oppositions of reflection.

The opposition of opposing and unity-positing reflection, then, cannot be concluded either by immediately positing one of these reflections as the negation of the other; for in reflection, as shown above, the abstractly negative is just as much a reflex as the abstractly positive: immediate negation is only a limit for the immediate, not for what is reflected into itself. By letting the I itself posit the Not-I as its original limit, Fichte certainly believes he has fixed the qualitative opposition once and for all; but that is a misunderstanding. As long as the limit of the I can obtain no other independence than the abstractly negative one of being Not-I, and indeed a Not-I which, posited by the I, falls within the I itself, so long will the consequence also press itself upon us that the I is just as much a limit for the Not-I as conversely; that the I, consequently, can just as well be said to fall within the Not-I, and to be posited by the Not-I, as conversely. But when the Not-I in this way becomes both Not-I and I, while the I for its part is to be both I and Not-I, then the opposition of I and Not-I must thereby at once cancel out\footnote{``Ist Ich = Ich, so ist alles gesetzt, was im Ich gesetzt ist. Nun soll der zweite Grundsatz im Ich gesetzt seyn, und auch nicht im Ich gesetzt seyn. Mithin ist Ich nicht = Ich, sondern Ich = Nicht-Ich, und Nicht-Ich = Ich.'' J.~G. Fichtes Wissenschaftslehre. Leipzig 1794, p.~25.}.

In order to hold fast the nominal opposition between I and Not-I, Fichte appeals immediately to the fact that both members of the opposition are in consciousness, and that their unity must consequently also be posited in \opage{69} consciousness\footnote{``Die Gegensätze, welche vereinigt werden sollen, sind im Ich als Bewußtseyn. Demnach muß auch X im Bewußtseyn seyn.'' Anfr.\ Skr.\ p.~26.}. But consciousness as a psychological fact is far too finite and concrete a magnitude for reflection to be able without further ado to set it in unity with that first abstract opposition. If natural understanding is to judge between the I and the Not-I, it will surely not fail to award the Not-I a world of objects. By appealing to the immediate consciousness of the opposition, Fichte has from the beginning precisely diverted attention from the insufficiency in the opposition itself; herein evidently lies the real reason why, with all his exertion to develop the whole of actuality out of the I, he has not been able to carry it further than a series of categorial paraphrases of that first, continually vanishing and continually returning strife between I and Not-I. It is a doubling of reflection that is demanded here; for it is not enough that an abstract opposition of I and Not-I be shown here; within the Not-I there must also be shown an express opposition between the one Not-I and the other. Only when the Not-I is in this way doubled just as the I is will it be seen from the two-sided doubling that the opposition in which the one Not-I stands to the other by means of the I is of an entirely different kind from the opposition in which the I, by means of the Not-I, originally stands to itself. In the self-opposition of the I by means of its Not-I, reflection is \emph{unity-positing}, for the opposition itself is a reflex in the posited unity: the posited unity is selfhood reflected in its other, subjectivity itself. In the opposition of Not-I to Not-I by means of the I, reflection is on the contrary \emph{opposing}, for the unity of Not-I and Not-I is negative, a reflex in self-continuing otherness. The self-positing that holds exclusively \opage{70}for the I now stands, then --- but also only now --- in decisive opposition to the positing of other against other, and so forth, and thus to the positing of otherness by which the Not-I is essentially characterized. As the opposition is carried through and otherness is stressed, the emphasis falls no longer on a uniform repetition of Not-I and Not-I, but on an express opposition of the one Not-I against the other. With the express opposition, the negative expression Not-I then also finally falls away; instead of negatively setting the one Not-I against the other Not-I, reflection now positively sets the one against the other, the one positive against the other positive; thus --- in opposition to self-positing subjectivity --- the one objective against the other objective, the one object against the other.

\begin{center}{\bfseries γ)\quad Indwelling and Hovering Reflection.}\end{center}
\addcontentsline{toc}{paragraph}{γ) Indwelling and Hovering Reflection}

Insofar as subjectivity itself is essentially reflection, reflection is \emph{indwelling} in subjectivity; for the object, on the other hand, which cannot thus reflect itself, but must be determined by the subjectivity reflecting upon its opposite, the determining reflection is \emph{hovering above}. If the opposition of indwelling and hovering reflection could be held fast in full immediacy, we should have to acknowledge unconditionally the Kantian separation between reflection of the understanding and the thing in itself; we should have consistently to deny the validity of an objective reflection. But the opposition of hovering and indwelling reflection is in turn so reflected that the indwelling itself shows itself as hovering and the hovering as indwelling.

If indwelling reflection were not at the same time hovering, subjectivity could have no knowledge, neither of itself nor of anything else. There could be no knowledge at all; for since knowledge, as shown above, presupposes itself to infinity, it must everywhere rest on this, that the apprehension of something already apprehended is apprehended again, that the representation of \opage{71}a representation is again represented, etc. Understanding of the finite always begins with a representation of how one representation is determined by another representation, and so it goes on to infinity. Just as representing subjectivity at every moment, even when it would expressly represent itself, goes beyond what is represented, so thinking subjectivity at every moment, even when it expressly thinks itself, goes beyond the thought thought. Since, now, reflecting subjectivity has representation in the reflex of thought and thought in the reflex of representation, subjectivity must necessarily have its indwelling reflections in such a way that it incessantly hovers above them, and thereby at the same time incessantly hovers above itself.

As regards the reflection of the objective, the reflection of objects, the critique of the understanding has hitherto placed itself in a highly characteristic opposition to ``speculation.'' Whereas ``speculation'' does not hesitate to attribute to things themselves, to substances, objects, etc., so immanent or indwelling a reflection that they are thereby enabled to reflect themselves in one another and with one another and each by itself into itself, without any subjective or hovering reflection whatsoever being needed for this, the critique of the understanding for its part has carried matters to the opposite extreme, and declared all objective reflection, indwelling in things themselves, to be a flat absurdity. Hegel's and Herbart's diametrically opposed treatment of the principles of identity and of contradiction shows clearly enough what the extremes here mean.

The proposition $A = A$ signifies, according to Herbart, not that $A$ reflects itself into itself, but only that our thought doubles itself in an empty repetition, a tautology. Posito A, nulla ponitur relatio. Posito A = A, res non duplicatur, sed cogitatio. Itaque res non refertur ipsa ad sese (quasi Fichtianum Ego), sed duplex vel multiplex cogitandi actus, quem iterare licet vel hodie vel cras vel minimo \opage{72}quoque temporis intervallo interjecto, semper in idem punctum reverti dicitur sine ullo discrimine, cui relationis aliquid affingi debeat vel possit\footnote{Herbarts De principio logico exclusi medii inter contradictoria non negligendi. Kl.\ phil.\ Schr.\ II.\ P.~724. Cf.\ P.~S.~V.\ Heegaard, den Herbartske Philosophie, Kbhvn.\ 1861, p.~4 ff.}.

This view set up by Herbart, sound common sense will of course at first glance find just as irrefutable as it will at first glance find Hegel's objective reflection unreasonable in the very highest degree. Pure $A$ surely cannot possibly reflect itself, any more than the letter $A$ can think itself, or the magnitude $A$ divide itself by $B$. But once what sound common sense must thus have granted at first glance is granted, the problem is seen at the next glance; and as soon as the problem is seen, it will also be seen that the critique of the understanding is just as one-sided as speculation. For what sort of standpoint is it from which Herbart would reject the reflection indwelling in the object $A$? The empirical-psychological one! Reflecting subjectivity here places itself, in all finitude, over against the existing objects; from this standpoint, as we have seen above, positing reflection cannot see through its own presupposing: in that psychological subjectivity presupposes A as an actual object, reflection forgets itself to such a degree that it without further ado regards this its presupposition as something that is, whose existence has nevertheless not been posited beforehand. The real, that which truly is, is indeed precisely such a thing presupposed, which must not have been posited beforehand. The fundamental contradiction, that the reals are presupposed not to have been posited beforehand, repeats itself in the reflection of their mutual opposition. If the reals were not in absolute mutual opposition, there could not be a plurality of absolute reals. But if the reals actually are in opposition, they must surely be opposed to one another; but opposed they can
\opage{73} yet be impossible without being posited. Opposition is, moreover, impossible without relation, and relation impossible in an absolute opposition, since an absolute opposition, critically understood, is itself impossible. The contradiction is conspicuous from every side: insofar as the proper essence of the reals is assumed to fall outside the relations in which, according to our ``zufällige Ansichten'', they may happen to meet, each real is transformed into a ``Ding an sich'' and must vanish as a reflex in a Not-I; insofar, on the other hand, as the determinate being of the reals is to find its expression in determinate relations, they must necessarily be permeated by reflection, for all thoroughly carried-out relation is, essentially regarded, immanent reflection.

For clarifying the relationship between relation and reflection the mathematical analyses are especially suited. Should the representation of indifferent relations with corresponding ``zufällige Ansichten'' find application to any kind of terms of a relation whatsoever, it would surely have to be to the mathematical magnitudes, when, be it noted, each magnitude with its immediate limitation is considered in and for itself. Between the parts of matter there is, after all, an interaction and to that extent an essential relation; but between the pure mathematical units, the pure mathematical points and the pure mathematical lines there is, insofar as each of these elements is considered in and for itself, no essential relation whatsoever: in this respect they may pass for genuine, reliable reals. But who does not now see that an unconditional holding-fast, each for itself, of the individual terms as pure elements in their pure singleness is a pure idée fixe? What the equation $\frac{y}{x} = c$ already sufficiently shows, that the magnitudes $y$ and $x$ count for themselves only for what they count in the relation, becomes, if possible, still clearer when we reflect that a differentiation of the equation $y = f(x) = cx$ likewise gives us the relation, namely: $\frac{dy}{dx} = \frac{0}{0} = c$. Even if one might, \opage{74} beguiled by the immediate representation, be tempted for a moment to hold fast the finite magnitudes $x$ and $y$ as terms for themselves, which were to have existence and significance outside the relation as well, every attempt to hold the terms of the relation fast outside the relation must nevertheless prove altogether absurd when we see the relation, the relation itself, remain standing, although the pure elements, the immediate terms of the relation, vanish into the infinite and lose themselves in zero. What holds of the magnitudes in abstracto holds of things in concreto: they are for themselves what they are in and through their relations; an absolute real, a ``Ding an sich'' without properties, is just as pure a fiction as an absolute unit, a magnitude ``an sich'' without relation to other magnitudes.

The reciprocal relationship between magnitude and relation now shows itself at the same time as a reciprocal relationship between relation and reflection. That the empirical world is a world of relations everyone will surely admit; yet even acute thinkers have not always been attentive to the truth hidden herein: that it is precisely the relations themselves, the objective relations, by which an objective world is to be known. Had Fichte, instead of restricting the relation $A = A$ to I $=$ I, seriously undertaken to investigate the relations between $A = A$, $B = B$, $C = C$, $D = D$ etc., then an equation such as $\frac{A}{B} = \frac{C}{D}$ would already have sufficed to prove to the great thinker that in the Not-I itself there is a system of relations valid in and for themselves, that, in other words, there is an objective world. To posit $a = b$ and $b = c$ may well be said to depend on the I; $a$, $b$ and $c$ are to that extent without objective independence. But when reflecting subjectivity has in this way made itself master over certain presuppositions, there at once present themselves consequences over which it is not master. From the presuppositions $a = b$, $b = c$ there follows, after all, the consequent proposition $a = c$ with so \opage{75} irrefutable a necessity that subjectivity itself, in such consequent propositions, must bow to the objective relations as to an objective power.

The objective power is nevertheless no alien power, no blind barrier against which the I is in turn to be broken; on the contrary! it is precisely in the objective power thus determined that subjectivity cognizes its own subjective power. The relations are immanent in the system of consequences and become objective only through their express opposition to reflecting subjectivity; the reflections take place in subjectivity and become subjective and hovering-over only through their express opposition to the relations immanent in the consequences.

Understanding begins with the objective relations and the subjective reflections corresponding to one another point by point. But that is only a beginning. Were reflecting subjectivity definitively to remain standing at a finite paralleling of our reflections with the relations of things, and to confine itself to referring the agreement, in pious admiration, either immediately to a divine truthfulness that ``will not deceive'', or to a harmonia præstabilita that cannot deceive, then the understanding would on these terms remain standing at a striking misunderstanding of its own essence. As the terms of the relation are transformed into reflexes in the relation itself, the relation is at once transformed into objective reflection; the relation of the terms is valid for the terms themselves only insofar as it is valid for a subjectivity reflecting over the terms, and conversely. Here there is not a series of relations to set up alongside a series of reflections; for just as no relation can be without reflection, so too no reflection can be without relation. To assume a system of objective relations without wishing to see therein a system of objective reflections immanent in the reflexes of the terms is just as absurd as to assume a system of objective reflections without seeing therein a system of subjective relations \opage{76} hovering over the reflex-determinations. The objective reflection, immanent in the terms themselves, is the reflected unity of the relations and the hovering-over reflections; the subjective reflection, hovering over the relations and their terms, is the reflected unity of the relations and the immanent reflections.

\begin{center}{\bfseries b)\quad Essence-Reflection.}\end{center}
\addcontentsline{toc}{subsubsection}{b) Essence-Reflection}

The reflection of the understanding is everywhere accompanied by the subsidiary representation that selfhood and otherness, the subjective and the objective, stand in a finite separation from one another. It is the finite in the separation on whose sublation the self-understanding understanding incessantly reflects. The very separation of positing and presupposing reflection already shows that if otherness, the non-subjective, has only a negative, vanishing being, then the same must hold of selfhood, the non-objective: if otherness is semblance, then selfhood is semblance too, and conversely; if selfhood is essence, then otherness is essence too, and conversely. Essence is, as Hegel has aptly shown, being that affirms itself through the negation of its own negation. Being-reflection is, on the whole, masterfully carried through in Hegel as far as the formal side is concerned. One can only wonder that analyses which in subtlety and acuteness still stand unequaled have not been able to convince the thinker that essence cannot reflect itself in objective-abstract universality; that what is original in essence itself is precisely selfhood, subjectivity. What it means that otherness, the objective, is posited as the derived essence, while selfhood, self-understanding subjectivity, expressly posits itself as the original, the difference between opposing and unity-positing reflection can already show. From the thoroughly reflected opposition, doubled in reflection, it becomes clear, namely, that the reflection of otherness, the optical-objective reflection whereby the one term in the infinite series of otherness mirrors its determinations in determinations of the other, is essentially different \opage{77} from self-reflection, the essence-reflection whereby subjectivity, presupposing everything, asserts in the multiplicity of oppositions its unity-positing, its infinite and concrete self-positing. As every finite separation of opposing and unity-positing reflection is in turn sublated in reflected unity, it is expressly seen that the opposition of the reflection of selfhood and of otherness is a unity of opposition. Thus there finally emerges a finite separation of the understanding between reflections which, according to the subsidiary representation, are merely subjective, and relations which hold for the objective. The finite in this separation of the understanding is not removed by the difference between subjective and objective reflection fading into an ambiguous unity; the finite is precisely dissolved and yet preserved by the opposition between reflection as something merely subjective and relation as something merely objective being transformed into an opposition falling within reflection itself, namely the opposition between hovering-over and immanent reflection. Instead of falling outside reflection as a being indifferent to knowledge, the system of objective relations is on the contrary transformed into a system of intermediate determinations in which immanent and hovering-over reflection develop their reflected unity, thus into a system of determinations that belong to reflection itself. Through the all-sided sublation of the finite separations according to their immediacy, the reflection of the understanding is gradually clarified into essence-reflection.

With particular regard to the relationship between essence and phenomenon, phenomenon and ground, ground and existence, we conceive essence-reflection in this connection as \emph{phenomenal}, as \emph{grounding}, and as \emph{founding} reflection.

\begin{center}{\bfseries α)\quad Phenomenal Reflection.}\end{center}
\addcontentsline{toc}{paragraph}{α) Phenomenal Reflection}

That thought in its abstract negative form falls outside intuition has significance only insofar as the heterogeneous moments of reflection separate themselves for reflection itself; \opage{78}but such a separation is then just as infinitely vanishing as it is infinitely returning: the moments are the reflected unity's own reflexes.

Phenomenal reflection naturally presupposes an immediate separation of sensible and intellectual intuition. If, then, a reflecting understanding, which makes itself finite by holding that over which it reflects outside the reflection, is allowed to fix the separation, then the critical philosophy is certainly right when it rejects intellectual intuition as a separate cognoscendi genus; but the finite separation of sensible and intellectual intuition must, through the reflection of appearing, be incessantly converted into unity. What, indeed, is intellectual intuition other than a unity of thinking and intuiting reflected in what appears? If intuition were not in itself essentially intellectual, it could not possibly be determined by anything intellectual whatsoever; if the intellectual were not in itself essentially something intuitable, it could not possibly be determined by the addition of anything intuitable whatsoever. The one of two heterogeneous things cannot immediately determine the other; color-determinations, e.g., cannot be applied to tones, since colors and tones are heterogeneous. Were the case with the separation of the pure intuitions and the pure forms of the understanding as Kant assumes, then it would be impossible, as Kant nevertheless demands, to determine intuitions by concepts and concepts by intuitions.

If the unity-in-difference of intuition and thought has first become transparent on the subjective side in phenomenal reflection, then on the objective side the reflective unity of the phenomenon itself and the phenomenon's substrate follows with irrefutable consequence. That in every phenomenon, everything that appears, a distinction must be made between the appearing for itself and that which appears in it is self-evident\footnote{``Dies war das Resultat der ganzen transcendentalen Aesthetik, und es folgt auch natürlicher Weise aus dem Begriffe einer Erscheinung überhaupt, daß ihr etwas entsprechen müsse, was an sich nicht Erscheinung ist, weil Erscheinung nichts für sich selbst und außer unserer Vorstellungsart sein kann, mithin, wo nicht ein beständiger Zirkel herauskommen soll, das Wort Erscheinigung schon eine Beziehung auf Etwas anzeigt, dessen unmittelbare Vorstellung zwar sinnlich ist, was aber an sich selbst, auch ohne diese Beschaffenheit unserer Sinnlichkeit (worauf sich die Form unserer Anschauung gründet) Etwas d.\ i.\ ein von der Sinnlichkeit unabhängiger Gegenstand sein muß''. Kritik der reinen Vernunft, Leipzig 1838, p.~247--48 note.}; but \opage{79} the question is how matters stand with this doubling in essence-reflection itself. It is a unity-in-difference of essence and semblance on which the emphasis falls. If there could be determinations of essence so absolutely essential that they had nothing at all to do with semblance, not even with a semblance of semblance, then one would undeniably have to grant Kant that an essence in itself could be thought which kept outside every phenomenon. But what becomes of pure essence when every semblance vanishes? It itself becomes a semblance. Likewise semblance, when essence has become semblance, will double itself into a semblance of semblance, a semblance subsisting in fleeting semblance, and thus show itself at once as both essence and semblance. As the reflected negation of essence, semblance is the form of essence; as the reflected position of semblance, essence is the content of semblance: semblance filled by essence is, in reflected immediacy, a phenomenon.

Here, then, there cannot possibly be two series of determinations: pure determinations of essence that were not to concern the phenomenon, and pure phenomenal determinations that were not to concern essence itself. If it now holds in the objective sense that all determinations of essence are essentially phenomenal determinations, then it holds at the same time in the subjective sense that all essential phenomenal determinations must be determinations of essence.

An essence in itself whose objective content was altogether heterogeneous with our subjective forms of intuition and of the understanding would \opage{80}not in any way be able to fill the empty subjective forms.

The absurdity of the view that behind the phenomena there should hide an inaccessible substrate, a noumenon, which was so absolutely positive that it could only be determined negatively, Fichte has excellently demonstrated; but in transforming the positive thing in itself into the negative Not-I, and drawing the objective factor, in the form of an obscure necessity, in under the I itself, he too missed phenomenal reflection, only from an opposite standpoint.

For to see through phenomenal reflection it is not enough to distinguish between the reflected and the immediate, between reflection itself and that over which there is reflection; it must be expressly brought to consciousness that the immediate, in which reflection has its opposition, is posited and presupposed by reflection itself, that the concretely immediate, as something reflected in itself, must precisely for reflection's sake obtain relative independence.

For the Kantian critique, what is immediate in the object remained standing over against the reflection reflecting upon it; likewise, as regards the subjective, the one of the main sides of reflection was set as something immediately given over against the other, the intuition of space over against the intuition of time, the pure sense-intuitions over against the categories of the understanding: it is a rational cutting-up and a rational piecing-together that characterizes ``the Critique of Pure Reason''; it does not get beyond patchwork. The advance marked by Fichte\footnote{``Kant geht in der Kritik d.\ r.\ Vst.\ von dem Reflexionspunkte aus, auf welchem Zeit, Raum und ein Mannigfaltiges der Anschauung gegeben, in dem Ich, und für das Ich schon vorhanden sind. Wir haben dieselben jetzt a priori deducirt, und nun sind sie in dem Ich vorhanden.'' Wissenschaftslehre 2te Lieferung p.~108. Schluß-Anmerkung.} is \opage{81} a sublation of the patchwork, a deduction of the immediately given; but as reflection labors from every side to assimilate the immediate, it consumes immediacy itself and so entangles itself in the contradiction that it can neither digest what it has thus consumed nor do without it. It is essence-reflection in subjective form, infinite self-reflecting, that is one and all in Fichte's I; but this reflecting still stands in abstract opposition to the multiplicity of phenomena. In order to reflect upon itself the I must naturally at the same time reflect upon the phenomena; but one-sided self-reflecting, despite all its exertions with ``a priori syntheses'' of $X$ and $Y$, of $A + B$ by $a$ and $b$ etc., cannot after all achieve an affirmative determination of what the phenomenon is. It goes with the phenomenon in Fichte's Not-I as with the elf-maiden in the fairy tale: it is hollow in the back. If it was impossible for the critique of reason to see through the phenomenon because it wanted the thing itself in absolute immediacy, it is now impossible for subjective idealism to grant the phenomenon relative independence because it does not see that self-reflection presupposes a reflection of other, a phenomenal reflection, in which the reflected itself is immediate and the immediate reflected.

That the immediate is something reflected in itself, and what is thus reflected something immediate, is certainly one of the hardest contradictions that can be offered to the reflecting understanding. So long as the contradiction has not been seen through, the understanding will accordingly be constantly tempted to assume one of two things: either that the contradiction lies in reflection, or that it lies in immediacy. The critique of reason assumes the former; subjective idealism the latter. From both standpoints the phenomenon is misjudged; and yet, so far is the contradiction from lying in the phenomenon that the phenomenon is precisely the resolution of the contradiction.

\opage{82} As reflection of other in other, the essence of the phenomenon is a reflected immediacy. The contradiction of the critique of reason is now this: that the phenomenon is to be an appearing of a thing which nevertheless does not appear. By fixing the thing in itself, the reflecting understanding has driven reflection's own opposition to the utmost point of petrified immediacy: it has secured for itself a decisive reality to the same extent as it has secured for itself an indissoluble otherness. But as reflection in this way seeks to secure reality for itself in the thing, the thing in itself withdraws into the unknown and leaves only the phenomenon standing, as an after-image distorted in reflection. The understanding has here quite rightly seen that essence in and for itself is one thing, and the semblance in which it reflects itself another; but it has not yet seen that the difference immediate in reflection is reflected in an equally immediate unity, that the difference is apparent because the unity is apparent, that here on every side there is reflection of other in other, and that this reflection is precisely the phenomenon in which essence itself immediately appears.

Seen from the standpoint of subjective idealism, the phenomenon is, of obscure necessity, an appearing in which nothing appears. What is deceptive in the phenomenon, then, does not consist in its being like a mask under which the thing itself preserved its incognito; the deception is rather that reflection, dazzled by immediacy, will constantly seek a thing in the phenomenon, and lo, there is after all no thing in the phenomenon. The I that reflects itself in everything has rightly seen that there can be no immediacy so firm that it does not let itself be permeated by reflection; but, entangled in one-sided self-reflecting, it has not yet seen that there is a reflection of other in other, a reflection that affirms its own opposition in granting the reflected immediate existing and cognizing, in the reflected immediacy of what exists, the determinate phenomenon. It \opage{83} is not otherness, nor is it immediacy, that is here simply rejected; on the contrary! as there is constant talk here of barrier and many barriers, of reality and negation, of impulse and check, counter-impulse etc., it is noticed clearly enough that immediate otherness cannot in any way be rejected. But why, then, are the phenomena so hollow and essenceless? It evidently comes from this, that immediacy is indeed dissolved by reflection, but not appropriated by reflection, not cast again in the forms of reflection, not stamped into peculiarly determinate phenomena. The immediate, whose immediate form is constantly disturbed without its thereby gaining independence and shape in a new and higher form, has only an elemental existence, can only pass for a stuff. It is precisely in this stuff-like immediacy that the categories barriers, check, counter-impulse and the like are used in subjective idealism; precisely because reflection of other is empty, self-reflection is empty too.

The reciprocal determination of self-reflection and reflection of other first receives the positive content of its development in phenomenal reflection. Only from phenomenal reflection is it seen that there are outwardly existing essences, essences which reflect themselves in one another objectively without, however, being able to reflect themselves for themselves subjectively. Every such essence looks at first glance like something immediately existing; its immediate existing is, however, only an appearing: it is in itself a negation of the immediate, something reflected, an essence. As the understanding now grasps that the first immediate cannot be the true objective, it may certainly very easily come to declare what immediately exists to be mere semblance. But when it then, following Herbart's ``wieviel Schein soviel Hindeutung auf Seyn'', also imagines that it can catch the truly existent behind the semblance, that is a deception: essence is in the semblance itself, in reflection; that which truly is appears. If, finally, ``speculative \opage{84} reason'' supposes that, since immediacy is reflected on every side in the appearing, phenomenal reflection must necessarily be so all-sided that, without hovering-over reflecting, it can reflect itself in an immanent manner, and thus, in other words, transform itself into a self-reflection immanent in the phenomena, then this too is a deception. If the immediacy of the phenomenon in opposition to reflection is an appearing; if the reflection of the objective essence in its phenomenon is an appearing; if the reflection of the peculiarly determinate phenomenon in its difference from other phenomena is an appearing: then the self-reflection of the self-reflecting phenomenon is in truth also an appearing. What is it, then, that uninterruptedly prevents the essence reflecting itself as phenomenon from attaining self-reflection? Answer: it is otherness, it is immediacy. Instead of a self-determining selfhood the phenomenon has an apparent selfhood; this apparent selfhood of it is precisely its immediate independence, its existence. The phenomenal essence is an existence, and what exists a thing. That things reflect themselves in their properties and, by means of the properties, in one another mutually, then means, in other words, that things appear in their properties, and that through the properties they appear to one another. But appearing is on every side reflection in other; things cannot possibly appear to one another without thereby appearing to a contemplating reflection: the objective in appearing presupposes the subjective, and conversely.

If the appearing of things to one another mutually cannot be abstractly objective, then neither can the individual thing's appearing in its properties be so: only by appearing in a reflecting that is not itself apparent can the thing appear in its properties. An appearing that is merely subjective is false; it is a hallucination, a semblance, a deception. But an appearing that was supposed to \opage{85} be merely objective is also false; phenomenal reflection is objective only insofar as it is subjective; it is immanent only insofar as it is hovering-over: things reflect themselves in their relations by reflecting themselves for reflecting subjectivity.

\begin{center}{\bfseries β)\quad Grounding Reflection.}\end{center}
\addcontentsline{toc}{paragraph}{β) Grounding Reflection}

Between self-reflection and reflection of other, or, more intuitively, between subjectivity itself and the objective world of phenomena, there has thus developed an essential and deep-going opposition; if this opposition is not to turn into a dualism that makes all understanding impossible, then a thoroughgoing unity of opposition must now show itself here as well. It is again reflected immediacy around which the question turns. Phenomenal reflection has certainly already illuminated this immediacy for us; but the light still falls in from one side only. Phenomenal reflection in a certain manner took the side of the phenomena and their reality as existing things; it gave us everywhere the immediate reflected in immediate and positive form. But reflection is positive only insofar as it is essentially negative, positive through its moment of intuition, negative through its moment of thought; reflected immediacy must consequently necessarily also have its negative form. The negative form of reflected immediacy develops in \emph{grounding} reflection; for grounding reflection dissolves the content of immediacy into self-reflection, and thus forms a significant opposition to the phenomenal. All grounding rests essentially on a reflected unity of self-absorption and absorption in the content of the matter, i.e.\ of immediacy: grounding reflection seeks at one and the same time the essence of subjectivity in things and that of things in subjectivity.

\opage{86} Grounding reflection seeks the ground; but what is the ground? From whatever side we take the concept of ground, it constantly points us to an immediacy that has been, an immediacy that reflects itself in a negative form. Thus, e.g., ``the ground of things'' is certainly in one respect the same as ``the essence of things'', and yet there is here an essential difference between essence and ground. In essence the immediate is negated by being transformed into semblance; in semblance as semblance immediacy has vanished: the negation of the immediate is in its semblance immediate only as negation, but precisely for that reason not itself a negative, existing immediacy. What holds here of semblance holds also of essence: if one would immediately grasp it in semblance, one grasps something negative and is referred to what is behind the semblance; if one would immediately grasp what is behind the semblance, one again grasps something negative and is referred to what is seen in the semblance. Essence is in itself only what it is in its semblance: it shines in semblance, a semblance of semblance; it is on every side negatively determined, and consequently without immediacy. Otherwise with the ground. Insofar as essence is conceived as ground, it has, in the negation of the immediate, an immediacy of its own. This is seen clearly enough from the change that takes place in reflection when the unity-in-difference of essence and semblance is transformed into a unity of opposition of ground and consequence. How decisive immediacy is in the determination of the relationship between ground and consequence becomes clear at once from the fact that the ground can be conceived by itself as presupposition, and the consequence as the conditioned consequence of the presupposition. From this standpoint subjectivity places the foundation of its knowledge in a system of fixed, universally valid presuppositions, the so-called principles [Grundsætninger]. In the system of principles all determinations of essence are reflected as immediate determinations of being; the foundation, the existing ground, is precisely a reflection of the immediacy of the ground. A corresponding im\opage{87}mediacy is seen in the system of consequences. Certainly nothing can emerge in the consequence other than what has already lain in the presupposition; the consequence is to that extent only a presupposition converted in reflection, repeated in formal paraphrase. Nonetheless there is here an immediate separation of presupposition and consequence whose significance must not be overlooked. The same presupposition can, namely, have different consequences, and the same consequence proceed from different presuppositions. For between the presupposition and its consequences there fall, immediately regarded, a multiplicity of varying conditions; these conditions can in turn immediately be partly separated from, partly coincide with, the determinate presupposition, etc.

The immediacy here pointed out is not fetched from outside, as from impressions of a merely empirical multiplicity; it lies in the ground as ground and constitutes a main side of its logical characterization. The exact and physical sciences essentially remain standing at this immediacy, and many philosophers, e.g.\ Herbart, at bottom mean nothing other than a universal presupposition in a sum of conditions when they speak of the ground. But just as little as essence, in order to be ground, can do without the immediate, just so little can immediacy, in order to be ground, do without negation. The mere presupposition is not yet ground, and the conditioned consequence only an abstract consequence\footnote{``Since the concept of function precisely shows us the reflected unity of the presupposition and the consequence, of the independent and the dependent, the difference between the mere presupposition and the ground proper can be elucidated in this connection.
% Errata: the errata list (p. 87 note, line 5 from bottom) corrects ``Functionsbegreb'' to ``Functionsbegrebet'' (``concept of function'' to ``the concept of function''); the translation reads the corrected form.
The presupposition is something given, something assumed, which is apprehended in representation, and falls under reflection only insofar as it is seen that, in order to be a presupposition, it must have its consequences. One can thus assume a presupposition, hold it fast in representation as something known, and yet be greatly surprised by its consequences. The surprise arises from the essence of the presupposition not having been seen through from the beginning; the presupposition was indeed made distinct in a representation, and to that extent apprehended from a certain side, but not clarified from all sides, not apprehended according to its concept. Only when we see through the presupposition and discover the consequences hidden in it do we cognize the essence of the presupposition, posit it as ground and the consequence as the grounded consequence.'' Mathematik og Dialektik p.~7--8.}. It is precisely by clarifying immediacy in negative reflection that philosophy goes further than the special sciences when it is a matter of an investigation of ``the grounds of things''.

\opage{88} In order to be conceived as ground, the presupposition and its conditions must be converted as moments in negative reflection; without such a conversion the relationship between ground and consequence is unthinkable. ``The consequence'' --- says Herbart --- ``is to lie in the ground. But it is also to follow from it, that is to say: it is to separate itself therefrom. Now if it really lies in it, then it belongs to the ground, and he who arbitrarily separates it therefrom has not extended his knowledge; he has rather merely repeated what he already knew, since he knew the ground. If, on the other hand, the consequence teaches something new, then this new is not the old, and did not lie in the ground; it is then wrongly called a consequence of the same.'' But the way out that Herbart then chooses, namely to conceive the consequence as a part of the totality of the ground, that is, identical with a part of the same, so that it thereby becomes ``a repetition in a separated thought'', is only an empty evasion; for if the determinate consequence is simply a part of the ground, then one of two things indisputably holds: either we have, in the sum of consequences corresponding to a certain ground, only a sum of divided grounds, and thus at bottom no consequences at all, but a sum of grounds; or else, if the parts of the ground are not themselves grounds but consequences, we have in the ground itself no ground, but mere consequences, consequences without any ground. Should there remain here, in the sum of such consequence-less grounds and ground\opage{89}less consequences, even a mere semblance of ground, then the ground would necessarily have to be sought in the mutual connection of the consequences and the unity determined by the connection. But who does not now see that such a unity of connection, reflected in the totality of the consequences, is precisely the totality in which the presupposition with its conditions is posited with negative immediacy?

As the one and the same totality is expressed in the two different forms of reflecting immediacy,
% Errata: the errata list (p. 89, line 9 from top) corrects ``reflecterende'' to ``reflecterede'' (``reflecting'' to ``reflected''); translated as printed.
the negative form of the ground and the positive form of the consequence, it must naturally look like two totalities, for the essential difference in form is at the same time a difference in content. To see one square equal in size to two others is to see something immediate reflected in positive form; for though the squares themselves are immediate, their equality in size is only for reflection. To see the ground of the two squares' equality in size with a third, on the other hand, is to see the immediate reflected in negative form. From this standpoint attention can no longer dwell on the immediacy of the squares or on what is immediate in the fact that the equality in size in question actually takes place; every immediate determinateness transforms itself into a determination of relation. Thus each square's magnitude is seen in relation to its side; the sides of the three squares are seen in their mutual relation insofar as they are seen as sides of a triangle; the sides of the triangle are furthermore seen in relation to the angles, it becomes, after all, a right-angled triangle etc. Since relation is now a categorial intermediate determination between reflection and immediacy, it is at the same time evident from this how easily the ground itself can be confused with the conditions and the immediate presuppositions.

That the ground must be sought in the proof is indisputable; but it depends on how one sees the proof, and what one sees in it. There can be many proofs of the same proposition, but in all of them there can be only one totalized ground. Because one can prove a mathematical theorem,
\opage{90} it does not yet follow from this that one at once sees through the relation between proof and ground. Just as the consequence is the unity of the consequents, expressly posited in the positive form of reflected immediacy, so the ground is the unity of the presuppositions and conditions, expressly posited in the negative form of reflected immediacy; if the consequence is simple, then so is the ground; if the consequence as consequence is the totality in the outer, then the ground as ground is the totality in the inner: the outer is seen, the inner is seen into; the consequence is seen, the ground is seen into.

When the totalized consequence is posited as existence, the ground is at the same time presupposed as the real ground. That the real ground does not keep itself behind existence, but is absorbed into it, comes into existence in and with what exists, is expressly seen from the fact that what exists, in its becoming and change, is incessantly in the act of arising and passing away. But to arise is surely to proceed from the ground, and to pass away is to go to ground, and so to go back into the ground. Existence is an intermediate form, and what exists is the real middle term through which the ground, in an infinite circuit, goes out from itself and returns to itself. In opposition to the immediate and positive reality of what exists, the ground is essentially ideality; but this ideality is the negative reality of the phenomenon.

As negative reality the ground is akin to the Schellingian indifference; nevertheless the ground is not simply an indifference, and the indifference is not ground: in the dawn of indifference difference is already dawning; both are originally seen in the same twilight. An indifference in itself without difference is a mystical ``un-ground'', a hyperspeculative ``primal ground'', which can no more bear the light of reflection than ``the subterraneans'' can bear the daylight. If one conceives pure indifference as a first indeterminate ground, a groundless ground, from which difference with the determinate grounds has proceeded, what then is one doing? One is seeking a ground for the ground, while in an \opage{91}illogical manner lumping the categories ``ground'', ``principle'' and ``beginning'' together into one. The real ground is, considered by itself, neither principle nor beginning, but a middle term between the two, reflected in negative immediacy. It is true that the word ἀρχή occurs in Aristotle also in the meaning of ground; but from this one can no more infer that the ground is simply principle than that principle is beginning.

That the ground, although it has originality and the original content in common with the principle, is nonetheless not the principle itself; that, in other words, the principle, conceived as ground, is conceived only one-sidedly and incompletely, can be discerned already from the fact that the content of the ground rests just as essentially on a resulting as on an anticipating. Nothing can proceed from the ground but what is contained in the ground; but nothing can be contained in the ground but what has previously gone to ground: the rich content of the ground derives from the countless existences that incessantly go to ground; the outflowing (emanation) of things from the ground is conditioned by a corresponding backflowing. Whereas the principle, the unconditioned, always acts in its system of conditions in such a way that it masters these conditions, transforms them into moments, and in the midst of all becoming and change reflects itself as the absolute and unchangeable, the ground, the real ground, is on the contrary so sunk in becoming, change, finitude and relativity, that it must precisely be said to reflect changeability itself in the unchangeability of the principle.

Reflected in becoming and change, the ground is akin to the beginning; whatever is in truth to begin must begin from the ground up; the ground of things is the real beginning of things: the ground is the reality of the beginning. But what is said here of the beginning holds just as emphatically of the end. The thing that goes to ground ends in the ground; the thing passes away; but this passing away is not absolute annihilation: annihilation has its reality in the ground.

\opage{92} Negatively determined, the ground can now indeed be conceived as an indifference, but this indifference is no ``un-ground''. It is the negative form of reflected immediacy that gives us indifference: in this form the ground is indifferent to the difference between principle and beginning, since as the essential intermediate determination of both it has equally much and equally little of both; in this form the ground is indifferent to being and becoming, to the unchangeable and the changeable, to beginning and end, to selfhood and otherness, to immanent and transcendent reflection, etc.

Insofar as the grounding reflection with its negative form dominates the phenomenal, subjectivity is one-sidedly absorbed in itself, and reflection, in this one-sidedness of self-absorption, is a brooding. Insofar, on the other hand, as immediacy in the negative form at every point presses itself immediately forward with existential elements that demand shape without yet being able to take on determinate shape, otherness gains the upper hand, and the grounding subjectivity, absorbed in the essence and becoming of things, finds itself in a ferment.

\begin{center}{\bfseries γ)\quad The Founding Reflection.}\end{center}
\addcontentsline{toc}{paragraph}{γ) The Founding Reflection}

During subjectivity's unclear brooding and ferment, unity, as the Ginnungagap of the un-ground, will incessantly swallow up multiplicity, and multiplicity, with its ground-points and germs, will in turn fill, overfill and as it were overwhelm unity; at this point the ground is mystical: the ground as ground is the one and all of mysticism. The clear obscurity which mystical speculation so readily spreads over the questions belonging here has had such a deterrent effect on the critical understanding that many an acute thinker still regards inquiries into the ground-relations of things as idle speculations. But instead of letting itself be deterred by the emptiness of the ``un-ground'' and the superfluity of indifference, the grounding reflection must precisely set about clarifying itself; for in \opage{93} clarifying itself it at the same time works its way out of the unclarity of the ground. The task is to determine the ground by what is founded in such a way that the emphasis falls on the foundation.

The founding reflection can in this connection rely on the following three propositions: 1) \emph{All is in the ground one}; 2) \emph{All is in the ground different}; 3) \emph{All must have its sufficient ground}. But it lies in the nature of the case that such propositions, set up singly, may not be taken for more than indicative points of departure; what is essential is not how the propositions are formulated, but how they are developed.

1) \emph{All is in the ground one.} The proposition, set up in this way, must not be confused with the Eleatic proposition that only the One (τὸ ἕν) is; for with the Eleatics, who let the One coincide immediately with what is (τὸ ὄν), there can properly be no question either of any being in the ground or of being of the ground. It is true that by what is the Eleatics meant what in truth is (τὸ ὄντως ὄν), and so what is in the ground, for otherwise they could not possibly have set it against what is not in propositions as sharp as these: \emph{only what is, is; what is not, is not}. ``But what then,'' Plato already asks, ``is one to judge of what is not? If what is not is not at all, then it can surely be neither thought nor uttered; if what is not is a seeming, then this seeming must surely after all be. This is precisely the difficulty in thinking what is not, that both he who denies it and he who defends it are compelled to contradict themselves. If one grants that there is a false opinion (ὄντος δέ γε ψεύδους ἔστιν ἀπάτη), then one at the same time presupposes the representation of what is not, for only that opinion can be called false which either declares what is not to be something that is, or what is to be something that is not \opage{94} (κινδυνεύει τοιαύτην τινὰ πεπλέχθαι συμπλοκὴν τὸ μὴ ὄν τῷ ὄντι).''

What is characteristic in this connection is that the Eleatics use foundation to annihilate all ground. If the world of phenomena were to be founded, then the One would surely have to be the ground of the many; but the world of phenomena is, according to the Eleatics' ``opinion'' (δόξα), groundless, the being of the many is a non-being, and what is not, is not. These propositions, offensive to common sense, have nevertheless to be proved, and so founded; in founding them, the Eleatics thus come, in the most preposterous manner, to found the groundless. Why can the One not be ground? why can the many not be founded? The Eleatics answer correctly: ``all ground and foundation presupposes becoming.'' This answer does not, it is true, lie outright in their words, but it lies in their train of thought: it lies in their attack on ``becoming''. ``Becoming (arising and passing away) is unthinkable; nothing becomes or passes away. Nothing can arise: it would then have to arise either from something that was like it, or from something that was unlike it. It cannot arise from the like, for the like is related to the like in a like manner\footnote{Poul Møller remarks on this: ``I understand this as follows: the like cannot arise from something else like it; for then they would be unlike in this, that the one was that from which something arose, and the other that which itself arose: consequently the two alleged likes would have unlike predicates and would thus not be alike.'' Udkast til Forelæsning over d.\ gl.\ Phil.\ Xenophanes. (Efterladte Skrifter, vol.~4, p.~59).}. Nor can it arise from the unlike, for then something would have to come from nothing. The weaker cannot come from the stronger, nor the stronger from the weaker: for then something would have to come from a thing that it itself was not, and so something from nothing.'' In this manner, then, \opage{95} the Eleatics seek to found that there is neither ground nor foundation.

But the proposition, All is in the ground one, cannot be understood either as if there were a Gnostic pleroma in whose depth the ground-unity rested. We must indeed grant the Gnostics that the original ground is a unity rich in content, and unity that which founds all; but what is immediate in the ground-unity is reflected, and the reflection negative. What is profound in the Gnostics', and especially in the Valentinians', intuitive and naive presentation of the primal ground and of the ideal-real oppositions that proceed from it and stir within it is, moreover, not to be misjudged. ``In the depth of the primal ground (βυθός, προπάτωρ)'' --- so it runs --- ``were self-consciousness (ἔννοια) and silence (σιγή). The Godhead hidden in the depth, the primal ground itself, revealed itself in 3 series of aeons: \emph{Reason} and \emph{Truth} (νοῦς and ἀλήθεια), \emph{the Word} and \emph{Life} (λόγος and ζωή), Man and the Church (ἄνθρωπος and ἐκκλησία); they emanated in pairs (σύζυγοι) as man and woman, and together they present the `fullness' (πλήρωμα) of ideal-reality over against the essenceless chaos, the infinite emptiness (κένωμα), etc.'' --- Such a dramatic-intuitive presentation, in grand-style naivety, surely makes a better impression than those phrases about ``the lack of love'' in ``the superabundance of love'', about love as ``the ground of creation'' and the ``final end'' of the ground in creat sibi, creat nobis, creat omnibus etc., by which a speculative dogmatics, which must secure its orthodoxy in petty sensibleness before it dares to let itself loose in Gnostic ingenuity, goes fully to the bottom in the ``fullness'' of the grounds and reveals its hollow scientificality.

The decisive opposition between Gnostic-dogmatic speculation on the one side and the founding reflection on the other is, however, conspicuous enough. Instead of a foundation, Gnostic speculation gives us images of fancy which represent a foundation; but a \opage{96} foundation cannot be represented; the represented foundation suffers from the contradiction that the existent that was to be founded is already found as existent in the ground. The primal ground, the ground-unity, is by this external conception transformed into a mere setting for the many, a deep basin from which the different flows out; but the difference in this different is not founded.

2) \emph{All is in the ground different.} In order to understand the ground-unity it is necessary to give an account to oneself of what lies in difference. Two extremes are to be noted here: either one gives up unity entirely and lets a multiplicity of existent real points furnish the presupposition for founding the difference of phenomena; or, so that nothing shall be unfounded, one lets the many and the different come about through pure negations and negations of negations in essence itself, the universal ground-essence. The first extreme characterizes realism, the latter idealism.

If we begin with idealism, it is undeniably a very correct thought that the different can be founded only by its proceeding from a unity in which it has, indeed, its ground, but not its existence: this ground-unity is identity. But what now is difference by itself as difference? ``A negation of identity!'' and identity by itself as identity? ``A negation of difference!'' Now if identity by itself were to be the ground of difference, then difference by itself would surely also have to be the ground of identity. In this way, then, difference could not be founded before identity, nor identity before difference. In order for a foundation to become possible, the difference between identity and difference must itself be founded in the unity of both; the founding unity then, as unity of difference, itself becomes an identity; and thus we understand what Hegel means when he calls the ground an identity of identity and difference. \opage{97} But the defect of this foundation is that what is founded lacks reality.

To lack reality is to lack the immediately positive. The Anaxagorean homoeomeries, the Democritean atoms, the Herbartian reals therefore draw attention anew in this connection. ``The Hellenes,'' says Anaxagoras, ``wrongly assume becoming and passing away, for no thing becomes or passes away, but things are compounded from existing things, and so they would more correctly call becoming composition and passing away separation''\footnote{Lectures on ``Phil.\ Propæd.'' from the academic year 1860--61, pp.~122--167.}. Instead of pure grounds we thus get pure realities, that is, existent ground-parts: to the reality of the ground-parts, then, corresponds the reality of difference. Now that the ground-parts also have their reality shall not be denied; only it must be remembered that this reality may not be groundless. There is in the world of things a real difference which cannot be explained by logical negations and condensed reflexes. Or how, pray, should one conceive the possibility that a difference so expressly determined as the difference between the atoms of oxygen, of hydrogen, of sulfur, etc., could come about, if there were not in the ideal ground-unity of all existence something real, something determinate in itself, which originally lay at the ground of such differences. The atomists' error is by no means that they presuppose plurality and difference in the ground; but the fundamental error is that they ascribe to the individual real members, the individual atoms, an absolute, in itself unfounded, independence; the fundamental error is that the physicists, just as much as the Gnostics, let what was to be founded in and by the all-founding unity of connection have an immediate, and so groundless, existence in the ground.

On the original relation between reality and ideality, change and transformation, all \opage{98} foundation essentially rests. The criticism of the understanding in general lets negation fall outside reality; that is a misunderstanding: if reality and negation kept outside each other, then the one reality could surely not itself be the negation of the other. But surely no one will deny that oxygen and hydrogen maintain themselves in real difference precisely by the fact that oxygen is a negation, namely a real negation, of hydrogen, and hydrogen likewise of oxygen. The fundamental difference, the thorough difference, rests no more on immediate realities than on pure negations: the thorough difference rests on real negations. As reflexes in the founding reflection, the real negations are at the same time moments in the founding unity: thus, and only thus, does the thorough difference itself become founded.

That difference becomes founded is expressly seen from the fact that the one difference goes to ground so that the other may proceed from the ground. Thus the real difference between oxygen and hydrogen goes to ground in the chemical union that founds the formation of water; but the properties, the real differences, that rest exclusively on the formation of water vanish again and go to ground when the water is decomposed, and the real difference between oxygen and hydrogen, with all the differences flowing from it, is thereby founded again, and so forth. The unity of difference of reality and negation is only a more determinate expression for reflected immediacy; the opposition of negation and reality is, in foundation, transformed into a unity of opposition of ideality and reality.

The unity of opposition of ideality and reality in turn founds a unity of opposition of transformation and change, and thereby at the same time a unity of opposition of the grounding subjectivity and the objective content of the foundation. That one-sided idealism, whose negations and determinations of difference always lack reality, misjudges otherness and the change corresponding to otherness is evident from many sides. It everywhere volatilizes the reality of change and keeps \opage{99} endlessly to transformation; therefore, for all its grounding and founding, it can never get anything founded in actuality; for all its attention to the objective in the universal and the universal in the objective, it never gets beyond the universal determinations of reflection of the grounding subjectivity.

With one-sided realism the case is the reverse; from this standpoint ``transformation'' is only a mystical figure of speech: all ground and foundation here rests on nothing but ``changes''. By exclusively holding fast to change, however, realism, without itself knowing it, transforms all real consequences into formal consequents, all real grounds into immediate presuppositions and finite conditions. Thus everything again comes to rest on cutting up and piecing together. In order to secure itself against mystical transformation, the understanding, unclear about its own reflections, has in this respect permitted itself the most impermissible of all transformations: the transformation by which the objective, in dead mechanical externality, falls outside the subjective; the transformation by which the light of ideality, that is, of reason, is extinguished; the transformation by which all ground and foundation is essentially abolished\footnote{Cf.\ Lectures on ``Phil.\ Propæd.'' 1860--61, pp.~160--167, where examples in concreto are found.}.

3) \emph{All must have its sufficient ground.} The proposition of sufficient ground (principium rationis sufficientis) was, as is well known, set up and defended by Leibnitz; but its logical development has in several respects been hampered and bungled. In the first place it is a bungling of the ground-relation when the proposition of sufficient ground is set in an external manner over against the proposition of identity, as a principle for truths of experience over against a principle for truths of reason. This forms a chasm between the \opage{100} a priori and the empirical, the abstract and the concrete, which makes all true foundation impossible. As long as formal identity, $A$ is $A$, is to count for more than a pure initial proposition, an abstraction destined to be integrated by an: $A$ is not $A$, but $B$, there can be no question here at all of seeing into what ground and foundation are. The abstract proposition, $A$ is what $A$ is, is just as groundless an expression of a formal agreement as the proposition, $A$ is what $A$ is not, is a groundless expression of a formal contradiction. Only with the help of the proposition of contradiction, $A$ is not what $A$ is not, does a shadow of foundation appear. It is easy to see where the shadow comes from; it comes from the contradiction contradicted in the proposition of contradiction. It is in a high degree characteristic that the so-called principle of contradiction, --- which naturally is not a principle, but a proposition of abstraction, --- must take contradiction twice, and so consist in a doubled contradiction, in order that identity, pure, sheer identity, may fall outside all contradiction. That propositions such as: $A$ is not what $A$ is; $A$ is what $A$ is not, are, considered by themselves, both contradictory and groundless, is surely self-evident. But how now do we sublate the contradiction? In a superficial manner by the proposition of contradiction; the superficiality is seen in the fact that thought will simply repel contradiction, without considering that it is precisely contradiction, contradiction alone and nothing else, that it has to thank for having nonetheless --- namely by contradicting contradiction --- come into a kind of formal movement. In a thorough manner, contradiction can naturally be sublated only by the ground itself and the corresponding foundation. $A$ is not --- in the ground, what $A$ is --- in existence. If a thing were in its ground what it is in its existence, then there would be neither ground nor existence; and yet the ground of the thing is identical with its existence, for the ground goes along into the foundation, and the existence is founded.

\opage{101} That formal logic by its so-called principles thinks of nothing but axioms is clear enough; but it is in a high degree characteristic that precisely these axioms everywhere place themselves in a thoughtless relation to ground and foundation. They fit the abstract in its pure, formal abstraction, for the abstract has not yet received its foundation; they fit the concrete in its immediate concretion, for the immediately concrete is something existing, and what exists has already received its foundation: foundation thus falls in both cases outside reflection. If a foundation is to take place, the formal axioms must be sublated. Thus principium exclusi medii undeniably holds for what exists; but for the ground of what exists holds principium inclusi medii. Against the disjunctive judgment of existence: $A$ is either $B$ or $C$, the ground sets its conjunctive judgment: $A$ is --- in the ground --- both $B$ and $C$. The human being is in the ground both woman and man, for woman is just as essentially a human being as man; but in existence the human being is either man or woman, for the existence of the androgyne is mystical.

The leap from the pure truths of reason to the truths of experience, by which ground and foundation are leapt over, coincides immediately with a leap from ground to cause, by which the founding intermediate determinations are likewise leapt over. It is true that, as we shall see later, all causal relations are in themselves at the same time ground-relations; but not all ground-relations are for that reason simply to be conceived as causal relations. It is against Malebranche's one-sided denial of natural intermediate causes that Leibnitz asserted the proposition of sufficient ground. Now if the fundamental thought hidden in the proposition had been fully developed from the standpoint of ground and foundation, then the development would undeniably have had to be fruitful in two respects. In the first place, the doctrine of sufficient \opage{102} ground would have shown us what it means that things themselves stand in essential, objective relations to one another; for the essential relations surely consist precisely in this, that things are founded by one another mutually, that one thing arises only by other things going to ground. Foundation would then, from this side, have had to show itself as an objective movement, immanent in the phenomena. But a monadology that would abolish the interaction of things and put its artfully devised harmonia præstabilita in its place cannot possibly attain to such an insight. In the second place, the doctrine of sufficient ground, had it been sufficiently developed, would have opened to us a deep insight into the reciprocal determination of the multiplicity of grounds with the ground-unity, of the separately conditioned, derived grounds of things with the original ground of all existence. For it is not the case, as Hegel has remarked in fleeting reflections, that every ground must simply count as a sufficient ground\footnote{``Daß der Grund \emph{zureichend} sey, ist eigentlich sehr überflüssig hinzusetzen''. Werke. Vol.~3, Berlin 1833. p.~76. Cf.\ p.~103.}. For one can bungle a foundation in widely different ways: either by keeping to the abstractly universal or by confining oneself to the all too particular. What, e.g., is the ground of $A$'s having become a drunkard? If one answers: the ground is that we are all sinners, then one keeps too close to the abstractly universal; if one answers: the ground is the disproportion $A$ sets every day between the quantity he consumes and the quantity he tolerates, then one would be keeping too one-sidedly to the abstractly particular. From the particular standpoint every peculiar thing must necessarily have its peculiar ground; for the existing phenomenon cannot possibly contain anything other than what lies in its ground.

\opage{103} That carbonate of lime crystallizes at one time as calcite, at another time as aragonite, has its special ground in the conditions that vary with the temperature. This special ground is naturally a specified utterance of the universal ground to which the crystallization of substances as a whole must be referred. But the universal ground of crystallization, with the system of relations and conditions under which its phenomena appear, is in turn to be determined as a separate ground in relation to other grounds, likewise each universal in its own sphere and separate through the opposition of the spheres, e.g., the ground of the change of the forms of state through heat, the ground of the formation and dissolution of chemical compounds through affinity\footnote{Such factors as affinity and heat are, namely, considered by themselves, not immediately grounds, but causes; nor can the ground be said to lie in affinity by itself or in heat by itself; for the ground is properly determined by the relation, the existential essential relation between heat and the forms of state, between affinity and the conditions that hold for the compounds, etc.} and so forth. Only as the founding reflection finds itself able to articulate the system of the particular grounds in each particular sphere can the sufficient ground come to light in the founded connection of the sphere.

In the foundation of the single phenomenon the foundation of its position in the whole system of phenomena must be reflected, and conversely. This Leibnitz saw insofar as he saw that each single natural phenomenon is conditioned, that the one conditioned points to the other to infinity, that the endless series of the conditioned can thus be determined only by an all-determining ground of grounds in which the conditions are absorbed as moments; thus he came to posit a final end of foundations, a ground of grounds, a ratio finalis. But once a leap had been made from grounds to causes, it is not surprising that similar leaps were made from causality \opage{104} to teleology and, with the help of teleology, from natural to supernatural causes, from physics to theology.

If, however, we disregard these leaps and keep to what lies in the problem of foundation, we are convinced that the doctrine of ``sufficient ground'' coincides with the doctrine of completed foundation. In completed foundation the objective first fully obtains its reality by means of the subjective. Could the foundation of things in and by their mutual connection be concluded by itself, then the founding reflection would lead subjectivity not merely to an absorption, but to a complete self-loss in the grounds of things; there would then have to be a reason factual in things, a reason so objective that it would indeed be able to found all things, but not able to fathom its own essence. It is this fundamental contradiction in reason itself that comes to light in the world of things through the fact that the conditioned is to be founded by the conditioned in endless series. According to sufficient ground, on the other hand, objective founding goes together as one with a subjective determining; for subjective ground-determining is a motivating. The system of the objective grounds of things is, from the original standpoint of subjectivity, of ontological subjectivity, a system of motives. The ground of existence and the motive are in foundation one and yet fundamentally different. The ground of existence, the objective, conditioned ground, corresponds to the finite reality of things; the motive, the subjective ground, corresponds to the significance of things for the system of relations. Only in unity with the real ground can the motive become sufficient; separated from objective founding, the motive by itself would on the contrary be abstract-subjective and arbitrary. Only in unity with the motive can the real ground become sufficient; for, separated from the motive, the real ground by itself would, precisely in its immediate, factually given presuppositions and conditions, have something groundless about it. That the motive is ground proves that the self-absorption of subjectivity \opage{105} is an absorption in the objective essence of things; that the ground is motive proves that the foundation of the essence of things must be grounded in the infinite self-absorption of subjectivity.

\begin{center}{\bfseries c)\quad Actuality-Reflection.}\end{center}
\addcontentsline{toc}{subsubsection}{c) Actuality-Reflection}

Insofar as the phenomena of existence are sufficiently founded, existence is actual; the founding reflection, which determines the phenomena, thus determines actuality. Now since an actuality cannot be without activity, nor an activity without acting, the question arises here: how do we determine the relation between reflecting and acting? ``All acting is in the ground an energetic reflecting and consequently one with an objective thinking. Thoroughgoing reflection will always demand that the members reflected in each other not only be transformed into non-independent moments, into pure reflexes, but in the transformation at the same time show themselves in their difference and preserve a certain relative independence. Now the greater the independence with which the members appear, the stronger must also be the reflection that is to transform them into moments and master them as reflexes; we have here a measure of the energy of reflection. In merely intellectual reflection the reflected members, as images, representations and concepts, are only fleeting seemings, non-independent reflexes; hence is explained the impotence of pure, that is, abstract, thinking to produce anything whatsoever that exists. In the reflection of natural force the reflected members have a prominent independence and, precisely by this independence of theirs, become striking expressions of the power that marks their impotence and holds them in non-independence. Reflection is, in its energetic and substantial essentiality, power. By power the moments reflected in the ground are posited as actual existences, as separate things: from this stems material discretion, from this the swarm of atoms, from this the existence of aggregates and masses in space. By power the actual existences, \opage{106} the separate things, are posited as moments and reflexes in reflected interaction: from this all union and dissolution, from this all arising and passing away, from this the multiplicity of changes with all transitions, from the fleeting alteration to thoroughgoing transformation! The power that posits and the power that sublates the independence of things is indeed different in the form of existence, but yet in the ground a power identical with itself''\footnote{Lectures on ``Phil.\ Propæd.'' from the academic year 1860--61. pp.~167--177.}, i.e., an infinite self-power.

In the thoroughgoing unity of reflecting and acting we then have a unity of difference of knowledge and power. This unity of knowledge and power can, as far as reflection itself is concerned, be developed from the standpoints of a \emph{substantial}, a \emph{causal} and a \emph{totalizing} reflection.

\begin{center}{\bfseries α)\quad The Substantial Reflection.}\end{center}
\addcontentsline{toc}{paragraph}{α) The Substantial Reflection}

The substantial reflection is in general a reflection of independence, a reflection of what subsists in itself. Just as the grounding reflection sought to determine the ground, so the substantial reflection now seeks to determine substance. In order to determine substance, we again consider what exists, the sufficiently founded phenomenon, yet from a new side. What is independent in existence, what immediately subsists, is then the thing itself, and the qualities of this which subsists by itself are the thing's properties. Now is the thing as thing to be called substance in the logical sense? No, for that its independence is still too immediate and, in its immediacy, again too fleeting. If the thing were in truth to count as substance, it would have to have an essential subsistence in itself; but the thing subsists in itself only by subsisting in its properties. The thing in its unity for itself is one thing, and the properties in their multiplicity are another; by subsisting in its \opage{107} properties, the thing thus subsists in its other, and is, like the properties, subject to changes.

That the properties change is founded in the fact that the one property must always rest on the other, the property of the one thing on the property of the other. Thus the thing could not, e.g., have the property: red, if it did not also have the property: extension; nor could it have the property: red, if light did not have a corresponding property. To a superficial consideration it does indeed look as though the properties could very well change and the thing still remain the same: the leaf that was green before is now yellow; but the yellow leaf is still also a leaf\footnote{``How long can one go on calling an object `one'? That is to say: how long can one regard it as the same object, in spite of the incessant changes to which it is subject? Do the changes not in the end become so many and so great that the object thereby becomes another `one'? --- From `a forest' one can fell many trees, and it still remains the same forest; `a house' can partly be enlarged, partly fall into decay, and it still remains the same house. But that there is a limit to all these changes, beyond which the object becomes another, is evident, and is also expressed in the well-known popular irony about the knife which, after having got a new blade a few times and a new handle other times, still always counted as the same knife.'' Heiberg, Det logiske System. Første Afhandl. Perseus. Nr.\ 2. p.~38. (Prosaiske Skrifter, vol.~2, 1861, p.~158.)}. This separation between the changeability of the properties and of the things holds, however, only partly in the finite. Here, to be sure, there is with regard to the thing itself a difference between essential, less essential and inessential properties; but in the logical sense it holds that a change in any one of the properties whatever must eo ipso be a change in the thing itself, since the thing in its existential immediacy
\opage{108} is just as immediately one with as it is different from its properties.

As the changeable is determined, so must its opposite, the unchangeable, also be determined. If something is apprehended in abstract qualitative opposition to another, then the ceaseless transition from something to another is a categorial expression for change itself; but the unchangeable then also falls, through this transition, qualitatively outside the changeable. In essence-reflection the transition from something into another is a transition of semblance into semblance, a transition which, to use an expression of Hegel's, is no transition; the changeable, which in this way has been transformed into sheer semblance, has thereby obtained so perfect a transparency that the unchangeability reflecting itself in all changes of semblance, pure selfhood, can shine through at every point.

Now if semblance in its transparency did not at the same time reflect its immediate opposite, the immediate itself, selfhood would be just as essenceless as otherness, and the pure clarity of self-reflection would lose itself in emptiness. But with what emphasis immediacy forces itself upon reflection, the reciprocal determinations in phenomenal, grounding and grounded reflection have sufficiently shown. The fleeting semblances, the empty reflexes, were thus gradually displaced by real phenomena, by existing things; but as each member in the series of finitude obtained heightened independence, the opposition between the independent and the dependent also became all the more pressing.

In existence, grounded being-there, the difference between independent and dependent is still only relative; one thing is independent, another dependent; the one is more, the other less independent: the transition is seen in the changes. But when this relativity is now to be thoroughly reflected through, we come upon a contradiction which is the very hardest of actuality, \opage{109} namely, that the independent itself is a dependent, the dependent an independent.

Changing is acting, and all acting a reflecting. In the active unity of reflecting and acting, the conversion of independent into dependent is like an infinite vibrating; each member, each part, each atom vibrates in the activity itself between independence and dependence. It is this vibrating, this trembling and quivering in punctual reflecting, that in the present context particularly characterizes actuality and the active. It is no longer pure shining, as in the reflection of semblance and essence, that we have before us here; nor can the difference between the dark and dawning in the ground on the one side and the half-clear, clear and intuitable in phenomena on the other any longer supply distinguishing marks; for ground-determinations and phenomenal determinations are sides of the same actual thing: activity itself is at once ground-determining and existence-determining. Accordingly, the essence active in actual changes does not present itself to intuition as a ground fermenting in the depths; the unity of becoming and acting reflected in activity itself is rather to be compared with fire, the ``ever-living fire'' (πῦρ ἀεὶ ζῷον), of which Heraclitus says that it is extinguished in measures (ἀποσβεννύμενον μέτρα) and kindled in measures (ἁπτόμενον μέτρα).

In all actual changes the activity is essentially one with the active itself, and the active itself is the unchangeable. The active itself is what is one in itself; for otherwise it could not possibly act in all actual changes and in them all assert its unchangeability. The active itself has an absolute subsistence; for it acts in all and upon all, and its self-acting is its subsistence. But the active itself does not have an absolute subsistence in finite opposition to the multiplicity of relatively subsisting things; what relatively subsists is only a moment of actuality, a point of passage for the \opage{110} infinite acting. The active itself, which in everything works out its own independence, is a subsisting selfhood; for subsistence is selfhood, and actual selfhood, self-activity, is one with reflected independence.

If, now, what is unchangeable in actuality, what subsists in itself, is substance, then absolute self-subsistence is absolute substantiality, and absolute substantiality at this stage of development identical with the absolutely actual subjectivity active in all things. It is not, as Hegel assumes, substantiality that is the original and subjectivity that is the derived; but it is the reverse. It is not substantiality that has become subjective; on the contrary, it is subjectivity that in the development of the concept has become what it eternally is, substantial. To \emph{derive} subjectivity \emph{from} substantiality is just as impossible as to derive knowledge from non-knowledge, selfhood from the selfless, the subjective from the objective. On the other hand, it is in good accordance with the nature of the matter to \emph{develop} subjectivity \emph{into} substantiality, for that is precisely to develop out of subjectivity itself what lies eternally in its essence.

The reflected unity of the independent with its dependent is power; the transparency of the moments for the reflecting self is knowledge: in the reflected unity of power and knowledge substantial subjectivity has its absolute subsistence.

All subsisting and perishing rests on the reflecting of independent in dependent. What subsists in actuality must subsist by a power; which is also why one says of such a subsisting thing that it remains in force, meaning thereby that it counts as something actual, that it actually holds good. What perishes in actuality perishes only by giving way to a power; it sinks into impotence, it succumbs in the struggle and disappears. At the stage of actuality power steps into the place of the ground; instead of grounds and groundings we meet here everywhere acts of power, that is, words of power, \opage{111} in which a commanding necessity declares itself. This must not, however, be misunderstood as though the ground were abolished and power groundless; for power itself is precisely the ground completed in all moments of presuppositions and conditions, a ground so actual that the grounding and the grounded must coincide in the act of power itself. Power is in infinite reflecting; and yet it looks as though it rested in eternal being; power is in its reflected being exalted above all becoming: its acts are its revelations. What this means becomes clear when we examine how power itself stands toward the powerful and toward the powerless.

As identity is, so must difference also be; if the identity of power and the powerful is substantial, then there must be a corresponding difference here. This Spinoza overlooked. In him the infinite substance is an absolute unity of power and the powerful: the powerful itself, substance (id quod in se est et per se concipitur), does not have power; for substance (cujus essentia involvit existentiam) is nothing other than power itself, and just for that reason infinitely powerful. If the powerful, power itself, is the sole being, then it naturally follows that existing things in their finitude, diversity and multiplicity must in themselves be absolutely powerless. The immediate opposition of power and the powerless rests precisely on the immediate unity of power and the powerful. But the immediate opposition of power and the powerless is a substantial contradiction. How can things in their multiplicity have even so much as a shadow of existence, when in their powerlessness they stand in direct opposition to power? From Spinoza's standpoint the answer to this must be: that the powerlessness of things does not fall outside, but on the contrary coincides with, the sole power of substance; that for precisely this reason things must be determined as modes and modifications (modi, modificationes) of substance itself. The thing, the single mode, is thus at once both \opage{112}one with and different from substance: one with substance insofar as its power, that is, the power in it, is substance's own power, its form substance's form, its being substance's being; different from substance, because it is finite, substance infinite, because it is powerless, substance all-powerful.

The substantial contradiction can now be more precisely determined thus: the one and the many, the infinite and the finite are two sides of the same substance; the infinite is the all-powerful, the finite the powerless: the powerless is thus a side-determination of the powerful, and what is powerful in power itself a side-determination of the powerless. And yet it lies in substantial reflection that the two sides must at the same time show themselves as two independents. ``The substantial relation consists in power; the absolute independence of substance rests on the infinite self-power with which it posits and sublates the accidental totality. But pure self-power is an empty abstraction. Power does not have itself, and for that reason cannot be hypostatized; it is the attribute of independence, which presupposes the independent. Substance, which is sui causa, is at the same time sui compos, and since it is itself a determinate something, it must at the same time have power over something else. Substantial sole power thus rests on the power that can be conceded to accidental multiplicity. It is precisely the hard contradiction of actuality and of power that it imparts to each of the members of the relation an adamantine independence, while at the same time making the one a reflex of the other. Let $A$, e.g., have complete power over $B$, and the contradiction will be evident. For if $A$ has the whole absolute power, then $B$ is of course only a shadow; but now $A$'s infinite self-power consists precisely in the power with which it masters its $B$; if $B$ is only an empty shadow, then $A$ has power only over a shadow, and its whole self-power is consequently a shadow-power. If $B$ is on the other hand more than a shadow, then \opage{113}it has at the same time an actual independence, and thus a power in itself; $A$, which now masters $B$, can then indeed be said to have an actual power, since it masters something actual. But since this actual thing itself has a power, and to that extent is not mastered by $A$, $A$ thus masters its $B$ insofar as in actuality it does not master it.''\footnote{Den propæd.\ Logik by R.\ Nielsen, Copenhagen 1845, pp.~166--67.}

This contradiction resolves itself into a substantial opposition between power in the form of reflection and power in the form of immediacy. In the form of reflection power is ideal, in substantial immediacy it is real. The immediately real always has its other, its finite opposite, outside itself; the ideal has its opposite within itself. In real immediacy power can never be absolute; the real powers are all relative. The ideal power, on the other hand, can be absolute, for what is absolute in power is precisely its ideality: absolute power is original self-power. The real powers are substances; the ideal power is not substance but substantiality: substantiality itself is subjectivity.

The unity of self-power and the relative powers, of substantiality itself and the existing substances, is a reflected unity of knowledge and power. Knowledge is the ideality of power; power is the reality of knowledge: knowledge is power. If absolute substantiality were not, there would be no substances; if knowledge were not substantial, there would be no substantiality. Only self-power, the subjective unity of power, knowing power, powerful knowledge, is able to bear the substances, the relative powers, the objective multiplicity of powers, in which the actual moment is an independent member and the independent member a moment.

\opage{114} \begin{center}{\bfseries β)\quad Causal Reflection.}\end{center}
\addcontentsline{toc}{paragraph}{β) Causal Reflection}

Aristotle has, in one of his aporias,\footnote{Metaphys.\ III.\ 5.} treated in a peculiar way the question whether numbers, bodies, planes and points are substances (οὐσίαι τινές) or not. ``In comparison with the things to which one would most of all ascribe substance, such as water, earth and air, of which the composite bodies consist, heat and cold and similar qualities (πάθη, affections) are of course not substances; only the body which has these affections remains as something that is, a substance (ὡς ὄν τι καὶ οὐσία τις οὔσα). And yet the body seems to have less substance than the surface, this less than the line, and the line less than the unit and the point. For by these the body is bounded, and it seems as though these could well subsist without the body, but the body on the other hand could not possibly subsist without them. If, however, one grants that lines and points have more substance than bodies, one falls into difficulty; for these boundary-determinations evidently rest on divisions of the body (διαιρέσεις ὄντα τοῦ σώματος), partly in breadth, partly in depth, partly in length. Moreover, either every form must be contained in the body or none: if, then, Hermes is not in the stone, and the half of the cube is not found marked off in the cube itself, then neither will the surface be found there; for if any surface were found in the cube, that surface too would have to be found by which the cube is halved, etc.'' The difficulties that arise in this respect, when it is asked what is in truth the original: the point or the line, the line or the plane, the plane or the body, the body itself as an existing thing or the qualities which make up the thing's essential properties, recur in new turns and with renewed emphasis when it is expressly asked: what in the absolute sense is the original, independence or the in\opage{115}dependent, actuality or the actual, substantiality or substance, causality itself or the cause? To treat such problems one must keep a sharp eye on the shifting points of view and be very exact with the categorial determinations.

Seen from the point of view of finitude, the immediately concrete is to be apprehended as the original, as what is independent in itself, in opposition to the form-determinations abstracted from it. In this respect it holds without contradiction that independence must rest on the independent, not the reverse. The Aristotelian critique of Plato's Ideas shows us from many sides how questionable it is to make pure essentiality independent, and, e.g., to let tableness (ἡ τραπεζότης ἐν τῇ τραπέζῃ) be the original, thus what is substantial in itself, and the table, the actual table, the derived. But however much Aristotle has otherwise insisted that the existing individual beings are the proper substances, that substance (ἡ πρώτη οὐσία) must always be a determinate this (τόδε τι), an individual, he himself has nevertheless not been clear either about the relation between the concrete and the abstract. Thus, e.g., in his well-known construction of the four elements he makes the abstract the original, the concrete the derived. ``The elementary properties: hot, cold, dry and moist, are namely brought forward in such a way that each of them is posited as an essence by itself, a substantial form which determines matter; thus a simultaneous union of dryness and heat gives the element of fire, of heat and moisture the element of air, of moisture and cold the element of water, of cold and dryness the element of earth.''\footnote{Cf.\ Forelæsn.\ over ``Phil.\ Propæd.'' 1860--61, p.~84 ff.} But if pure dryness can in this way dry the thing, then pure tableness can of course also manufacture the table.

\opage{116} From the standpoint of infinity, that is, of essentiality, the matter looks otherwise. Here it is now the opposition, not between the abstract and the concrete, but between two forms of the concrete, that of reflection and that of immediacy, which is at issue. It does indeed seem, when one takes the expression reflection in the vulgar sense, as though the concrete absolutely had to exist in an immediate way: ``a stone,'' one says, ``is something concrete; a concept, an idea, on the other hand, is something abstract.'' But as regards this misunderstanding, Hegel has sufficiently shown that the ideal can be concrete fully as well as the immediately real. If we look more closely, it is so far from the case that the form of reflection should be the form of the abstract, that thoroughgoing reflection, in opposition to immediacy, is precisely the form, the only form, of the absolutely concrete. That immediacy cannot be the completed form of concretion follows directly from the fact that the one immediately concrete thing must always find itself in a finite opposition to the other. Finite opposition means fragmentation, and fragmentation is the opposite of concretion. If the stone is to represent the concrete, then in the single stone we have only a fragment. And how one-sided, then, is concretion in such a fragment? Quartz, e.g., is a different concrete from feldspar, which again is a different concrete from hornblende, etc. Only on one condition can a concretion of all, an infinite and absolute concretion, take place; what it requires is namely that the determinations of existence which scatter themselves in external existence be apprehended in inner connection. The more manifold the outward is in its real fragmentation, the richer then is the inward in its ideal unity. The ideal unity, enriched in reflection of the multiplicity of realities, is inwardness; inwardness is essential concretion: the deeper the inwardness, the richer the concretion. Seen from this point of view absolute substantiality is not merely concrete but all-sidedly concrete, the only \opage{117}thing absolutely concrete in itself; every substance existing in immediate form is, in opposition to this, only a one-sided concrete, a single one of the countless members into which the external concrete is fragmented.

One more misunderstanding must be taken into account here, when the question is of the original and the derived. It holds as a law at the stage of actuality that all being is acting. The existing substance cannot exist without acting upon another existing substance; here, then, there must in actuality be at least two substances, of which the one acts, the other is acted upon: what is active in the agent is determined as cause, the change in what is acted upon as effect. Naturally regarded, then, the cause ought to be first in time, and the effect follow after. The determination of time announces itself in this relation with such insistence that an actual succession can easily seem to be what is characteristic of the relation itself. From this point of view Kant too, in his ``Critique of Pure Reason,'' set up his ``principle of temporal sequence according to the law of causality.''\footnote{``Der Begriff, der eine Nothwendigkeit der synthetischen Einheit bei sich führt, kann nur ein reiner Verstandesbegriff sein, der nicht in Wahrnehmung liegt, und das ist hier der Begriff des Verhältnisses der Ursache und Wirkung, wovon die erstere die letztere in der Zeit, als die Folge, und nicht als etwas, was bloß in der Einbildung vorhergehen (oder gar überall nicht wahrgenommen sein) könnte, bestimmt.'' Kritik der reinen Vernft.\ Leipzig 1838, pp.~196--97.} Now if it really were the case that the cause as cause had to be some moments ahead of its effect in time, then the categories cause and effect would have to cease to be correlative concepts. But a cause can no more be thought without an effect than a ground without a consequence: in the same now in which the cause is, the effect also is; in the same now in which the effect ceases, the cause also ceases. What the ground is \opage{118} in essence, the cause is in actuality.\footnote{``That the thing must have its ground means that its own inner presupposition cannot itself be a thing, but must be a unity reflected into itself, a system of active relations, a system of essence, in which the thing-causes are reduced to causal moments. Thus, e.g., oxygen cannot be the ground of water, any more than hydrogen. Oxygen and hydrogen, as the combination is entered into, are contributing causes, and their conditioned reciprocal action thus the proximate cause of water; in this respect it holds that every thing must have its cause. But the cause is one thing, the ground another. The ground of the formation of water, namely the real ground, is the transformation that takes place in the atoms of oxygen and hydrogen at the moment of union. Without such a transformation the coming-to-be of water --- if, be it noted, it is to be actual water, not merely water apparent to `our imperfect senses' --- is impossible.'' Forelæsninger over ``Phil.\ Propæd.'' year 1860--61, pp.~165--66.} The causal relation therefore cannot be set up by itself alongside the relation of ground; for the causal relation is in itself nothing other than the relation of ground's own actualization.\footnote{On the categorial relation between \emph{ground} and \emph{cause} Heiberg remarks in his Logic (Anfr.\ Skr.\ p.~47) with his usual sharpness and clarity: ``When one gives \emph{grounds} for something, one has not yet stated \emph{the ground}. In this the ground differs from \emph{the cause}. The ground can have becoming, and so it is undetermined whether it is to obtain existence or not. (One can have ground for undertaking an action and yet leave it undone, for it is arbitrary whether and where one will break off the infinite series --- of grounds pro et contra). But the cause is itself existence, namely of necessity, and must therefore necessarily have effect. A ground for acting thus need not be active; it first becomes so by being taken up into a will, which gives it existence, or perhaps even makes it a cause.'' --- Further on (Anfr.\ Skr.\ p.~75) he remarks concerning ``reciprocal action'': ``When one regards the cause as merely active, and the effect as merely passive with respect to the cause, though indeed as active with respect to a following effect, then one remains in the infinite series of causes and effects, thus in finite causality, in the sublation of which the truth of causality consists. For in this sublation the cause is also passive, in that it is sublated in the effect, and thereby the latter itself becomes active, or reacts as cause against its own cause.'' Further: ``By reaction one must not here think of that which takes place between two substances, each of which acts as cause (e.g., when a ball, shot against a hard object, springs back, or when the heat of the sun, by means of the reaction of terrestrial substances, becomes the cause of the emergence of plants), for here there is as yet no question of one substance, and still less of two or more. The reaction spoken of here is that which follows from a single original cause, in that its effect, merely as its own effect, not as a new cause, reacts.'' --- Against this last, however, several misgivings might arise here. For finite causality it must of course necessarily be required that the cause not merely differ formally from its effect, but that the cause have existence in one existence and its effect in another. Now whether one, with Hegel, will have this existence expressly determined beforehand as substance, or for the time being be content to apprehend it as a thing, may well be indifferent; but when the effect in ``finite causality'' has its existence for itself, just as well as the cause, is it not then an unavoidable consequence of this that the existence of the effect must at the same time function as cause and that of the cause as effect? The reaction in the finite must, after all, be conditioned precisely by two existences, ``each of which acts as cause.''} Now that, as regards the relation of \opage{119}ground, there can be no finite temporal separation between the ground and its consequence is evident. But whence comes it, then, that the representation of a temporal sequence forces itself forward so strongly in the causal relation? It comes from this, that the semblance of succession which forms in every movement of thought and reflection obtains an all the more essential significance the greater the independence the moments attain, the more thoroughgoing reflecting passes over into acting. To a superficial observer of actuality it looks as though the cause itself were a substance, or, to use a \opage{120} more current expression, a thing. But to make the cause a thing is just as unreasonable as to make force a matter. The thing-cause actualizes the concept of cause only in a wholly external, finite and one-sided way. The thing is indeed a cause insofar as it is the material unity of forces, the existential center from which the effect proceeds; but the cause itself is nonetheless not a thing. The thing perishes, but the force, the causing, does not perish; it only acts in another connection and in another way, identical with itself as another force, when the thing has gone to ruin.

That the cause cannot possibly have existence by itself in the one thing and the effect in the other can be seen, among other things, from impact. The acting is acted upon, and what is acted upon acts; impact is reciprocal action: in reciprocal action the cause is effect and the effect cause. But although the cause is one with its effect in impact, two colliding balls do not for that reason become one ball. Rather, the cause in the one does not only become one with the effect in the other; but in each of them the cause, by means of the reaction, becomes its own effect, and the effect conversely its own cause.\footnote{In the collision of $A$ and $B$ the cause in $A$ corresponds to the effect in $B$, and conversely, the cause in $B$ to the effect in $A$. This must, if causal reflection is to be carried through, be expressed again thus: the cause in $A$ is related --- by means of the reaction from $B$ --- to the effect in $A$, just as the cause in $B$ --- by means of the reaction from $A$ --- is related to the effect in $B$.}

The difference between the thing and the thing-cause is further recognizable from this, that one and the same thing, e.g.\ sulfur, can in several different respects become a cause, a cause of friction, of sound, of heat, etc.; there can in this way be many and different causes in one thing. Conversely, one and the same thing-cause can repeat itself in \opage{121}countless things; if a pressure, e.g., is to be the cause of a motion, then a pound of water can of course exert the same pressure, and thus supply the same cause, as a pound of lead, a pound of iron, etc.

As existing substances active in existence, actual things are thus at once causal moments and existing causes. The thing-cause is in this light both identical with the thing and different from the thing. If, now, the difference is brought out, the existing causes transform themselves into causal moments, substantial points of passage for the infinitely concrete self-cause. To this corresponds, then, the Kantian thesis that besides the natural series of conditioned causes presupposing one another ad infinitum we must necessarily presuppose a cause unconditioned in itself and determining itself, ``an absolute spontaneity of causes,'' a ``transcendental freedom.'' If, on the other hand, one lays the whole emphasis on the unity of thing and thing-cause, then the self-cause disappears; we are then led to accept the antithesis set up against the thesis, namely that ``there is no freedom, but everything in the world happens solely according to the laws of nature.''\footnote{Cf.\ Forelæsninger over ``Phil.\ Propæd.'' 1861--62, pp.~139--43.}

It is an entirely correct thought that Kant has expressed in the proof of his thesis, namely that there would be no true reality in the relative causes if there were not a first, original cause; but this thought, true and sound in and for itself, has in the proof of the antithesis been bungled and made unrecognizable by the logical monstrosity of a transcendental freedom which is not freedom but abstract arbitrariness. If $A$ is to be an effect of $B$, $B$ of $C$, $C$ of $D$, and so forth, then we get, quite rightly, what Kant in his way maintains, an effect of an effect of an effect --- without cause. But the cause fleeing away into infinity returns at once, of course, if we merely let the effects form a cycle.

\opage{122} So that the external multiplicity shall not have a disturbing effect, it will be sufficient in this context to let the whole cycle be closed with three effects. Let $A$, then, be an effect of $B$, which again is an effect of $C$, which again is an effect --- of what? Of $A$! $A$ is thus in the cycle its own cause and its own effect, and to that extent an independent member; but $A$'s whole acting is in the infinite cycle at the same time something effected ad infinitum. That $A$ has its cause in itself will then properly mean that it has its cause outside itself, namely in $B$ and $C$; as something thus effected, $A$ is only a moment in the cycle. What holds of the single member holds of them all: they are independent and yet moments.

With the reflective unity of the moment and the independent member, reciprocal action must necessarily double itself. Here in the cycle there is a double reciprocal action, a finite or real one and an infinite or ideal one. Finite reciprocal action determines the mutual dependence of the independent and yet conditioned members; it shows us the endless series of thing-causes, it presents causal reflection made finite in the substantial immediacy of the members. Ideal reciprocal action, on the other hand, is seen in infinite reflecting itself, which bears the cycle and gives it its omnipresent center. For this reflecting the substantial members are only points of passage; here we have what is original in the conditioned series of causes, not as a single existing member, but as self-cause, as absolute causality.

Absolute causality is thus not, as Kant assumes, an unconditioned by itself in finite opposition to the finite conditions; on the contrary! the unconditioned acts in an infinite system of conditions and is to that extent itself conditioned. But conditioned and conditioned mean something entirely different for the finite and the infinite. The finite member is conditioned in such a way that its conditionedness expresses precisely its \opage{123} dependence, since it is always conditioned by its other. The infinite is indeed also conditioned; but in such a way that through the cycle of conditions it relates itself to itself, and determines the conditions by which it is conditioned. The conditionedness of the unconditioned is actuality's expression for its ideal concretion, its originality and absolute independence. The self-cause is, in other words, the ideality of the finite causes, and these again are the reality of the self-cause. Causal ideality is a powerful knowledge; therefore all actual causal relations are intellectual, all causal determinations original determinations of knowledge. Causal reality reflects a knowing power; therefore the intellectual causal moments are existing substances, existing things in an actual order of things.

\begin{center}{\bfseries γ)\quad Totalizing Reflection.}\end{center}
\addcontentsline{toc}{paragraph}{γ) Totalizing Reflection}

``There is in our immediate knowledge much that must of course fall outside thinking insofar as it belongs to sense-perception, and, as regards sense-perception, much that falls outside sight insofar as it, e.g., falls under hearing, etc.; thus'' --- someone might say --- ``how much must there not exist in the actual world that falls entirely outside all reflection and thereby outside all knowledge!'' We answer: there are naturally manifold things that fall outside your and my reflection and thereby outside your and my knowledge; but as regards ontological reflection, there is nothing, nothing in the world, nothing either in the heavens or on the earth or in the abysses, that either falls or can fall outside reflection in itself, the substantial reflection that eternally acts as the infinite self-reflecting of causality. It is with ontological reflection in its relation to actuality as with our psychological reflection in relation to the objects we become conscious of: every outside goes into one with its inside, and yet an out\opage{124}side and an inside must constantly be distinguished. The object of sight actually falls outside sight insofar as it falls outside the one who sees; and yet the object in itself is only an object of sight insofar as it falls within sight. The image in the mirror is different from the mirrored object, and yet the object in its mirroring coincides with the image: we cannot become conscious of the image as image without becoming conscious of an object corresponding to the image as object. Reflection everywhere has the infinite suppleness that in its doublings it takes up into itself what it posits outside itself, and appropriates what is set against it. If immediacy is set against reflection, then it is not the immediate that determines itself without reflection, but it is conversely reflection that determines the immediate by determining itself. If a higher, that is, richer immediacy is to be recognized by the fact that the reflection from which it emerged is again sublated, then the sublated reflection nevertheless continues to preserve the reflex of something sublated; and so in all other respects.\footnote{Heiberg aptly remarks in his Logic (Anfr.\ Skr.\ p.~36) that ``in the following categories, of which the two that make up each pair stand in a relation of reflection to one another, such as the thing and the properties, the whole and the parts, matter and form, cause and effect, substance and accident, the first member \emph{has} the second, and conversely. E.g., a) The cause has effect. Thus it is the resulting ground, in which the being of the effect is sublated, i.e., the cause has itself been effect, or an earlier cause has had it. The effect that it is itself to produce is still immanent in it, whereby it wholly coincides with the already sublated effect. But b) the effect has cause; for when the effect has obtained being, the cause is sublated in it. And finally c) reciprocal action is the unity in which the cause has effect, and the effect has cause. --- That the ground has being, and being has ground, can also be expressed by saying that the ground \emph{has been}. It has namely been a) as immediate being, before this was sublated in the ground, and b) as ground, before this was sublated in immediate being. Both belong to the past, and are therefore expressed by \emph{having been}. The word `have' designates, as we have seen, sublated being, and this is the past. Therefore it is a correct consequence of certain languages, e.g.\ Danish, to express past time with the verb `have,' while others, e.g.\ German, do it with the verb `be.' (`It \emph{has} been'; `es ist gewesen').''

But when one now properly reflects on what has been said here, it must surely at last become evident that all objective logical reflection must necessarily take place in a reflecting subjectivity. Objects in their existential immediacy cannot possibly either have or preserve either ``sublated being'' or ``sublated reflection.'' The thing does indeed have an essence, thus ``a sublated being''; but it has this essence only in and by means of substantial reflection, which is eternally borne by subjectivity.}
\opage{125} If it is now evident that all determinations of actuality, as reflexes and determinations of reflection, must be absorbed into reflection, then it is by the same token evident that the whole infinitely rich content of actuality must fall within reflecting subjectivity. The subjective content is in reflection itself objective: reflection is \emph{totalizing}.

It is the substantial unity of the countless relative substances, the causal unity of all relative causes in their many-sided acting and interacting, with which the totalizing is completed. Here, then, there are in actuality --- according to what we have seen --- two actual totalities, an actually inward one and an actually outward one; and yet the outward is actual only by means of the inward, and conversely: the two totalities are in actuality one.

In the unity of the totalities the principle is seen. The principle itself, the absolute principle, is not a principle of subjectivity by itself, any more than it is the one-sided principle of the abstractly objective; for subjectivity in its reflecting has both the subjective and the objective, and subjectivity is the principle. But so infinite is totalizing reflection that subjectivity, \opage{126}although it is the principle itself, can at the same time hover above the principle, in that it can hover above its own two-sided development. The principle, whose unconditioned acting is absorbed into the development and the developed, then separates itself, with particular regard to the different modes of development of the two totalities, into two relatively different principles, the principle for the subjective or ideal and the principle for the objective or real totality.

In the inner, subjective and ideal totality all determinations of existence are reflected as pure determinations of knowledge; this totality is, in the form of unity, identical with the originally concrete selfhood, and must therefore be called a totality of selfhood. In the outer, objective and real totality the determinations of knowledge are reflected as real determinations of power; this totality is, in the form of multiplicity, a system of other against other, and must be called a totality of otherness. The reciprocal determination in principle of the two developments of totality is more than an empty mirroring; it is an acting articulated in all its moments and members, an interaction whose distinctive character in principle can be precisely determined.

The principle for the development of selfhood is determined by a movement of the same through other back to the same, a movement in which reflection continually sublates succession. Thus the sight of this is indeed different from the sight of that; but the sight of this is only this sight, the sight of that only that sight; this and that sight are determinations of difference within the same seeing, in the reflected unity of the sights, in the seeing selfhood. In a corresponding way the thought of this and the thought of that, that is to say: this and that thought, are two thoughts in one thinking; the representation of this and the representation of that, that is, this and that representation, two representations in one representing. The outward other reflects itself at every point as an inward one, and the inward as a determination of difference in the self-articulating unity. To be sure, this must be determined by a seeing when that is determined \opage{127} by a thinking, this rest on a representing when that rests on a sensing; but the one who sees thinks, the one who thinks represents, the one who represents senses: sight and thought, representation and sensation are side-determinations of the same knowledge, reflexes in the same reflecting. --- The principle for the development of otherness goes in the opposite direction. Here the original unity is seen only by a mirroring in the immediately appearing difference; the inward other is as a semblance over against the outward, for the emphasis of otherness falls on the outward. When the existing members of otherness are posited as the expression of the real power, the multiplicity of power is reflected in the outer totality, and the outer totality in turn separated into a multiplicity of relative totalities, differing in kind through their mutual opposition.

In the most outward of all outward forms, namely the mechanical, the existing member has so immediate an independence that the independence, seen from the side of finitude, really looks as if it were absolute. Atoms and masses assume in this light the semblance of an eternal existence resting on itself. What is selfless in this immediate independence has its categorical expression in the inertia of matter. ``Matter,'' says physics, ``is will-less,'' that is, it cannot determine itself to any change. By its dead independence, its indifference toward the state of rest or motion in which it is placed, the material member is on all sides subject to the relation of otherness: it cannot move itself, but it can move another and be moved by another; it does not attract itself, but it attracts another, and is attracted by another, etc. The revelation of power is in this outward form so finite, the parceling-out so predominant, that every whole seems to be collected and composed of preexisting parts.

As is well known, the distinction is drawn between a mechanical and an organic whole, that the parts in the former are always \opage{128}the first, whereas the parts in the latter are on the contrary developed from the rudiment of the whole, so that the wholeness in the rudiment itself is the first. This distinction is, as far as the logical is concerned, misleading. There is in the logical sense no totality, and hence no mechanical totality either, of which it can with reason be said that the parts are in actuality the first and that the whole emerges only as a result of the composition of the parts. The examples one is accustomed to adduce on this point betray, one and all, an evident misapprehension partly of the way in which a whole is formed, partly of the very concept of wholeness. ``A clock,'' one says, ``is a mechanical whole''; and this must be granted. But on the other hand it must not be granted that any of the clock's parts whatsoever should be formed in complete independence of the whole; on the contrary, it is well known that the form and constitution each single part receives is determined solely by the way in which it is to serve its whole. If one would fetch examples from stones, crystals and earths, such examples serve only to show how superficially and thoughtlessly the concept of totality is grasped. To be sure, a stone, a sandstone e.g., can be said to be a result of the composition of previously existing parts; but such an aggregate is after all only a fragment, not a whole. The concept of totality\footnote{Here the real totality is expressly meant, not the abstract formal development of the category. From the formal side the concept of totality can emerge, e.g., through a development of thing and properties, ground and conditions, etc. Every logical concretion can, in relation to the moments it takes up into itself, be determined from the side of form as totality. --- ``The unity, that the thing is only (as sublated) in the properties, and that the properties only are, (as sublated) in the thing, is \emph{reality} or \emph{unity of essence} (corresponding to unity of being'', --- namely One). --- ``When totality, which is the unity of phenomenal being, just as reality is the unity of essential being, develops its moments, content and extension, into a new and higher cycle, then the extension is determined as \emph{the whole}. The content, on the other hand, is determined as \emph{the parts}; and in this reflection the whole has parts, and the parts have (constitute) the whole.'' Heiberg. Anfr.\ Skr.\ p.~49; p.~63--64.

``Unity of essence'' can, in opposition to ``unity of being,'' be called totality with the same right as ``the unity of phenomenal being'' in opposition to ``the unity of essence.''

How abstractly and formally the concept of wholeness is conceived here is seen, among other things, from the following explanation. ``The category of the parts determines the contingent in actuality in such a way that the contingent is the \emph{parts that fall outside the whole}. In the following example

{\par\centering
\begin{tabular}{ccccccc}
\multicolumn{7}{c}{$C$}\\ \hline
\multicolumn{3}{c}{$A$} && \multicolumn{3}{c}{$B$}\\ \cline{1-3}\cline{5-7}
1.&2.&3.&4.&5.&6.&7.
\end{tabular}
\par}

\noindent the thing $A$ is divided into the parts 1, 2, 3, and the thing $B$ into the parts 5, 6, 7; but the part 4 belongs neither to $A$ nor to $B$, but falls outside them both. The part 4 is thus not seen as a part at all, except insofar as it is a part of the whole series from 1 to 7. If it is thus to be seen as a part, then it must be of a higher whole $C$, which includes both $A$ and $B$ and the part excluded by both. The part 4 is thus contingent for the two wholes $A$ and $B$, for it falls outside both, but it is not contingent in the higher whole $C$, of which it is a part.'' Anfr.\ Skr.\ p.~64. (Prosaiske Skrifter, vol.~1, p.~241.)} expressly demands that there must \opage{129} be, not merely a quantitative, but also a qualitative difference between the whole and its parts. In this respect not even the crystal can be said to be a true whole; for every part into which the crystal is divided is itself a crystal, hence in quality the whole; from which it follows in turn that the whole, in quality, is a part. But when the part is as the whole and the whole as the part, then there is here, in the essential sense, neither part nor whole. Everywhere, on the other hand, where the concept of wholeness is satisfied, it turns out that every development of totality, the mechanical one included, is a development of principle.

But the principle is surely essentially selfhood; how then is it possible that a selfhood can constitute itself in what is actually selfless? It is reality and the immediately real \opage{130} on which the emphasis falls. Considered by themselves, the individual members of reality in their outward multiplicity do indeed form a, so to speak, blind opposition to the self-unity acting in them all; but in the totality itself, the mechanical system of motions and moving causes, the unity of the opposition is motivated. Mechanical selfhood, reflected in the form of power, proves precisely what totalizing reflection is capable of. For so almighty is the infinite reflecting that the self-cause, the unconditioned itself, comes to subsist in a multiplicity of power of conditions and relative causes; and so reflected, in turn, is the infinite power that the relative causes with their conditions are transformed into points of transit for the totality of mechanical determinations in which the self-power is objectified.

In the mechanical totality the opposition of selfhood and otherness is abstract; the unity of the opposed determinations is therefore likewise abstract. We ascribe to the earth, e.g., a self-motion, insofar as during its motion in its orbit it at the same time moves about its axis; but the selfhood in this motion is immediately one with otherness; the earth's determination by itself is on all sides a determination by other; its self-motion is only a proper motion. What is abstract, and to that extent rigid, in the mechanical form must then be sublated by the development of the intermediate determinations still held back. From this point of view the organic totality forms, in the objective series of totalities, a distinctive opposition to the mechanical.

As regards the terms of the relation in their finitude, the difference between mechanical compositions and chemical combinations is already characteristic. In mechanical composition the members retain their outward finite independence; however closely they are united, the union is still only a unity of coherence: the one member falls, in the coherence, outside the other. Only in chemical combination is it \opage{131}expressed in a real way that the existing terms of the relation are actually moments for one another reciprocally. The multiplicity of transitions from existing members (material, product, educt) to moments and from moments to existing members, which is effected by the chemical processes, therefore also furnishes an intermediate stage of elementary formations, in which the development of totality is carried through with greater inwardness and richer articulation than in the abstractly mechanical circular motions. If we take the terrestrial whole by itself as a geological totality, then real selfhood, the principle of the earth's development, has surely already begun here to work upon the elementary forms; in this working there is so thoroughgoing an interaction of wholeness and parts, of general and particular totalities, that the terrestrial body may well be called an elementary organism.

The elementary organism is, however, afflicted with the contradiction of being an organic inorganic; the imperfection of the organization is seen in the fact that the parts still have the character of fragments and outwardly separated masses. The heterogeneous formations take place everywhere in such a way that the producings and the produced fall outside one another; the selfhood in this whole is only elementary. If the contradiction is to be resolved and the elementary articulated, the principle must by new breakthroughs lay out the development in such a way that the whole can reflect itself ideally in each single part, and each single part be a moment for the whole. This happens with the thoroughly formed organism.

But the organism in its immediacy is nevertheless afflicted with the contradiction: of being at one and the same time subject in equal degree to otherness and to selfhood. The union of the members is from the outward side a unity of coherence, and the existing of the whole is conditioned by thoroughgoing, incessant changes; seen from the inward side, on the other hand, the union of the members is a self-unity, an articulating acting that masters the changes and determines the outer to a mirroring of the inner. According to \opage{132} the outer, the individual members and parts with their distinctive character are in finite opposition both to one another and to the whole; for the outer whole, which consists in the coherence of the parts, is only partial, that is, partly in each particular part: the totality is a totality of otherness. But according to the inner, the individual members are reflected in the active unity in such a way that the wholeness is undivided in each part: the totality is a totality of selfhood.

As the totalized unity of two dissimilar totalities, the organism posits a decisive tension between the inward and the outward, selfhood and otherness. According to the self-unity and in its inward aspect the organism is a living being; according to the unity of coherence and in its outward aspect the living being is an organism. In the abstractly living or merely vegetative organism, however, the opposition of the inward and the outward is only intimated by two sides of the same living being; but with the ensouled organism the two sides become two totalities, and the tension sets in. The determinateness in which the two mutually opposed unities, the unity of coherence and the self-unity, here meet is a sensation. In sensation the selfhood is an immediate self, and this self is reflected in an immediate unity with, and an equally immediate difference from, its other. Sensation is the first distinctive utterance of the soul's selfhood; here is the principle of psychological subjectivity.

Submerged in the stream of immediate sensations, psychological subjectivity is a dawning consciousness; this dawning consciousness is an animal sense-consciousness. But animal sense-consciousness is afflicted with the contradiction of being an unconscious consciousness, a consciousness of a self which is nevertheless not conscious of itself. This contradiction is resolved by a new breakthrough of the principle; for the principle of consciousness is, at the stage of human life, a principle of self-consciousness. Only at this stage is reflecting subjectivity able \opage{133} again to posit otherness as a transparent moment, and to develop systems of determinations of knowledge that correspond to systems of determinations of power.

But when psychological subjectivity develops into an all-round totalizing reflecting, the objective determinations of power fall outside the subjective determinations of knowledge; it becomes evident that reflection with the abstract forms of ideality must become formal and powerless over against the system of existences in which the multiplicity of power has its reality. The unity in principle of knowledge and power, on which the originality of ontological subjectivity rests, is then, on the psychological standpoint, transformed into so finite a separation that a system of determinations of power without knowledge must reflect itself for the reflecting self in a system of determinations of knowledge without power. Reflection is carried through; but the dualism that has thereby arisen now brings with it problems that can only gradually be illuminated from a higher point of view.

\begin{center}
{\bfseries C.}\\[2pt]
{\large\bfseries Comprehending Subjectivity.}
\end{center}
\addcontentsline{toc}{subsection}{C. Comprehending Subjectivity}

Already at the stage of immediacy, thinking, psychologically reflecting subjectivity strives in two directions: partly to understand its surrounding world with its facts, partly to understand itself. But both endeavors must fail as long as the standpoint is one-sidedly psychological, that is, as long as it is still immediate in the ontological sense. The sensualist consciousness, which with all its thinking directs itself only by sense-impressions, is a selfhood lost in its other, and can for that reason comprehend neither sense-perception nor the sensible. The idealist consciousness absorbed in its thinking does indeed make the discovery that everything sensibly given can be dissolved by reflection; but as a selfhood \opage{134}consuming its other, psychological subjectivity is incessantly in the act of developing the ontological, yet without having any inkling that with this ontological it will again meet itself in a higher subjectivity. Objective idealism can indeed in a way be said to have developed with formal consistency out of the subjective; but in reality it has made a leap, in that in the logic itself it has leapt off from its subjective foundation. If we set Hegel's ``Logic as Science'' against Fichte's ``Science of Knowledge,'' one can indeed say that both works at bottom aim at the same thing, insofar as they both aim at an analysis and immanent deduction of a priori categories; and yet what a difference! In the ``Science of Knowledge'' it is the I first and the I last around which everything turns; the development of the categories here goes in so subjective a direction that the development itself is to be the proof that an objective cannot be given. In ``Logic as Science'' the I, thinking subjectivity, has become superfluous to such a degree that the objective categories here develop one another in an objective manner, and perform their logical play for one another, indifferent as to whether there is any spectator or not.

The ambiguity of the objective logic is seen, among other things, from the fact that it can pass for a general philosophy of essence, a metaphysics in the stricter sense of the word, and for a mere philosophy of words, a kind of dialectical synonymics, according to how one takes it. To illustrate this ambiguity we will in the present context refer to its treatment of three categories developed in the preceding section: \emph{the ground}, \emph{substance} and \emph{causality}.

\emph{The ground.} If the logical development is laid out in the following way: ``The ground can be taken 1) abstractly or separated from its content; so taken it is \emph{essence}; 2) concretely, or as being in its content; so taken it is \emph{phenomenon}; 3) as the unity in which essence has phenomenon, and phenomenon \opage{135}has essence, but neither of them is except in and through the other; so taken it is \emph{necessity}''\footnote{J.\ L.\ Heiberg, Ledetraad ved Forelæsningerne over Philosophiens Philosophie eller den speculative Logik 1831--32, p.~28 ff. (Prosaiske Skrifter, vol.~1, p.~183.)}, --- then one is equally entitled to expect an explanation of essence as an explanation of words. Both can indeed to a certain degree be given together, and that with a clarity, sharpness and evenness that makes the presentation logically plastic, without its therefore coming to a development of what is properly the metaphysical knot of the problem. If we look at the meaning of the words and at the formal side, we must indeed grant that several grounds are the mark that one has not yet found the single sufficient one, and that ``when one gives grounds for something, one has not yet stated the ground''; but if we lay particular weight on the explanation of essence, hence on what is objective and ontological in the matter itself, what a strange leap is there not then from a reasoning with grounds that takes place in psychologically subjective reflection to a grounding of the objective itself, which is designated categorially by the word \emph{existence} and by what exists in existence: \emph{thing}! It must indeed be granted that ``the existing ground conditions or is the condition for existence''; but when it is to be investigated in full earnest how the ground then properly obtains existence, or how the existent is grounded, then it is truly not enough to let ``the existing ground'' emerge as the result of a reasoning hovering between $+$ and $-$; not enough that by a formal paraphrase one transforms the ground thus ``come into being'' into a ``condition,'' and then goes on with an explanation of the words \emph{thing} and \emph{property}. If one will really enter upon the question of the ontological fundamental relation of existence, i.e.\ upon the relation between the multiplicity of existing phenomena and the \opage{136}original unity of the ground, then the fundamental problem becomes, as has been shown above, not merely an objective problem of existence, but an essential problem of subjectivity.

\emph{Substance.} The dialectical synonymics can, with reason, adduce concerning substantiality and substance e.g.\ the following: ``Substantiality is necessity's \emph{unity of being}, hence the unity that is first designated by One, but then by Many, because One by its repulsion or exclusive independence conditions the Many. In the same way the one substance conditions \emph{the plurality of substances}; every cause is substance, and there are as many substances as causes \dots\ Since the many substances, each by itself, are nothing other, with respect to substantiality, than what the one substance is, or, since the effect which the one substance has on the other becomes, by virtue of the interaction, its effect on itself, the plurality of substances flows together in the determination of identity, not of the one, but of the sole substance. The sole substance can, because it is the sole one, act on no other than itself. Its effects thus do not fall outside it, but are immanent. In this determination the effects are its \emph{accidents} \dots\ Further, the accidents are nothing other than the many finite substances, which constitute the determination of difference over against the determination of identity of the sole substance, in the same way as in reality the many matters became forms in opposition to the sole matter, and as in totality the multitude of parts became wholes outside the sole whole.''\footnote{J.\ L.\ Heiberg. Anfr.\ Skr.\ p.~77--78. (Prosaiske Skrifter, vol.~1, p.~271--273.)} In this explanation regard has undeniably been had here not merely to the logical meaning of the words, but to the ontological fundamental meaning of the categories. But when we now really apply existence's own substantial measure in judging this explanation, how is it then possible that \opage{137}``the accidents'' can be ``finite substances,'' when there is here in actuality only ``a sole substance''? Is it only some accidents that can become finite substances, which ones might these be, then? Or is it perhaps all accidents? But then in finitude all substances are surely transformed into sheer accidents, so that in actuality there remains no difference between substances and accidents. Actuality, then, should still actually be able to subsist after the relation of substantiality is dissolved? Or should the relation of substantiality perhaps not yet be dissolved, although the substantial opposition between the sole substance and the many substances, between the substances and their accidents, has vanished? No, if one will really give oneself an account of the reciprocal determination of substantial unity and substantial multiplicity, instead of dissolving the relation of substantiality, exchanging the word substance for a new word, and going on in the text, then one will, according to what we have said above, surely be compelled to grant that the problem of substantiality is essentially a problem of subjectivity.

\emph{Causality.} As a particular contribution to a characterization of the ontological explanations of words of the logic, or of dialectical synonymics, the following is still to be adduced in this connection, with express regard to Hegel himself. In his ``objective logic'' Hegel, as is well known, introduced the doctrine of actuality with what he calls ``an exposition of the absolute.''\footnote{Wissenschaft der Logik. 2te Abtheilung. Berlin 183.\ p.~185 ff.} Over against the absolute now stands the relative; this opposition is repeated in each of ``the absolute relations,'' both in the relation of substantiality and in that of causality. One should thus, as regards causality in particular, be entitled to expect a thoroughgoing development of the relation between the multiplicity of relative causes and the absolute \opage{138} unity of cause. For such a development there was here all the greater call, as Kant, in the antinomy treated above --- and Hegel otherwise readily treats the Kantian antinomies where occasion offers --- has posed the problem in a sharp and definite manner. The Kantian problem can, as regards specifically the unity of opposition of the absolute and the relative, be expressed in all brevity thus: ``Causality cannot possibly be relative without at the same time being absolute; since now an absolute causality contradicts itself, the concept of causality is without objective validity.'' But without entering further into this problem, Hegel concludes with a sketched presentation of general interaction, and then thinks to arrive, in the category of interaction, at a thorough solution of the knots of causality. The whole solution is summarily contained in the following explanation: ``Die Kausalität ist bedingt und bedingend; das Bedingende ist das Passive, aber ebenso sehr ist das Bedingte passiv. Dieß Bedingen oder die Passivität ist die Negation der Ursache durch sich selbst, indem sie sich wesentlich zur Wirkung macht, und ebenso dadurch Ursache ist. Die Wechselwirkung ist daher nur die Kausalität selbst; die Ursache \emph{hat} nicht nur eine Wirkung, sondern in der Wirkung steht sie als \emph{Ursache} mit sich selbst in Beziehung. Hierdurch ist die Kausalität zu ihrem \emph{absoluten Begriffe} zurückgekehrt, und zugleich zum Begriffe selbst gekommen.''\footnote{As further proof that it is not from carelessness or forgetfulness that the Kantian problem is passed over here, but that its solution is really to be sought in general interaction, we refer to Hegel's philos.\ Propædeutik, Berlin 1840 p.~144.} If one may really make the demand of Logic as science that it shall expressly orient us with respect to the difference between the validity of the same categories for the absolute and for the relative, then that solution of the problem of causality is undeniably very superficial. That the \opage{139} finite causes must be moments in the infinite interaction is understandable; but how then does it come about that these causes, although they are moments, nevertheless in actuality continue to subsist as existing members, and thereby prove that the relative will not let itself be logically sublated? When the relative causes in all their relativity continue to exist, what then becomes of absolute causality? Must it not also have a form of actuality in which it to all eternity continues to master the relative causes? These questions cannot possibly be answered without the problem of causality, as we have seen in its place, becoming a problem of subjectivity.

If, on the other hand, Logic as science is to be philosophy of essence only insofar as it can, in general, be so in the capacity of a philosophy of words, dialectical synonymics, then the Hegelian solution of the problem of causality must on the whole be considered satisfactory. Once it has been clearly seen what logical meaning must be connected with the words \emph{cause}, \emph{effect} and \emph{interaction}, it is surely quite in order that one passes over to new words, in order to develop their logical meaning in connection with what has already been developed.

``But'' --- someone might say --- ``the problem whose solution is here demanded from so many sides, namely the problem of subjectivity, has after all, according to Hegel's method, found its solution in its place, namely under the first section in the doctrine of the \emph{concept}. The transition from \emph{actuality} to \emph{concept}, which is expressly designated as a transition from \emph{causality} to \emph{subjectivity}, surely shows clearly enough that Hegel has subsumed, not merely the problem of substantiality and causality\footnote{As proof that Hegel has actually recognized the unity of the problem of causality and the problem of subjectivity, and in this unity the significance of the Kantian antinomy of freedom and necessity, one could refer to his Wiss.\ d.\ Logik 2ter Theil p.~214. But a glance at the bottom of p.~215 will suffice to show how little is gained thereby.}, but in \opage{140} the progressive stream of the development all problems of being and of essence under the problem of subjectivity. ``The concept,'' he says, ``is the truth of the substantial relation, in which being and essence have attained their full independence and determination through one another reciprocally,'' but the concept surely begins precisely as subjectivity. --- That Hegel must be understood thus, and has actually also been so understood by competent judges, can easily be proved. ``As the transition of substantiality'' --- says Heiberg\footnote{Anfr.\ Skr.\ p.~80. (Prosaiske Skrifter, vol.~1, p.~276--277.)} --- ``the concept is \emph{subjectivity}, as that of necessity it is \emph{freedom}, and as that of the whole ground or reflection (since these are the becoming of the concept) it is the \emph{existence} of abstract thought. With the concept, subjectivity is immediately given. The first immediate standpoint of the concept is therefore the \emph{subject}. And herein lies likewise \emph{freedom}; for since substance is cause of itself, it lacks nothing for freedom but subjectivity. By this, substance is raised from necessity to freedom; for the subject is that which has developed from necessity, as its presupposition, hence the free. Subjectivity is expressed by \emph{I}, but this is not yet consciousness; it is only the subjectivity common to the shapes of nature and of spirit (hence logical), and thus the root from which consciousness develops, though not yet this itself.''

These and similar elucidations prove precisely the opposite of what they intend. If there is any section in the whole Hegelian philosophy which shows to what degree the concept of subjectivity has here been mishandled and botched, it is precisely the section of the logic just cited on subjectivity. How forced and contorted the whole conception is appears immediately from utterances such as this: ``Die Gestalt des \emph{unmittelbaren} Begriffes macht den Standpunkt aus, aus welchem der Begriff ein subjectives Denken, eine der Sache äußerliche Reflexion ist. \opage{141}Diese Stufe macht daher die \emph{Subjectivität} oder den \emph{formellen Begriff} aus.'' The subjective concept, logical subjectivity, thus emerges, according to this explanation, by the objective movement of thought being suddenly laid into the psychologically-subjective medium. The immanent method, which at the stages of being and essence has been able to do without everything subjective, even the ontologically subjective, suddenly acquires at the summit of the concept such a need for the assistance of a subjectivity that, as in a tight corner, it even makes do with psychological subjectivity.

Ashamed at having to need such emergency aid, the immanent method naturally seeks as best it can to give the subjectivity of the concept a more objective tinge. The subjectivity is then not properly supposed to derive from the fact that its ``shape'' precisely ``constitutes the standpoint from which the concept is a subjective thinking, a reflection external to the matter''; in other words, it is not the standpoint in the subjective, but the shape in the objective, on which the emphasis is to fall. The subjective concept is thus in the strictest sense not subjective, for it is objective; but whence comes this? Naturally from the fact that there is, bluntly put, no essential difference between the subjective and the objective. ``Both in the unconscious shapes of nature and in the self-conscious shapes of spirit'' --- it is said\footnote{Heiberg. Anfr.\ Skr.\ p.~80. (Prosaiske Skrifter, vol.~1, p.~277.)} --- ``causality has determined itself as substantiality and spontaneity; herein lies subjectivity, and every substance is an I, whether this I has come to consciousness of itself or not. The determination of I lies already in reaction; for that every substance reacts happens precisely by virtue of its subjectivity. (A table, a chair, a cupboard, etc.\ are all of wood; but in those contingent forms which a joiner has given the wood there lies no subjectivity or any I; for it is not as chair, table, etc.\ that they \opage{142} react, but as wood, hence by their substance; and the I of this substance shows itself in the natural, organic form of the wood, just as the I of the inorganic world already shows itself in crystallization).''

By ascribing to finite substance (the wood, the crystal) in its immediate and material existence a subjectivity, an I, the objective logic has botched from the root the opposition between selfhood and otherness, between the subjective and the objective; for if there is anything that in actuality must be said to exist in the form of otherness, to express the selfless and thereby what is in itself objective, then it must surely be material substances such as the crystals and the wood in chair and table. To demonstrate ``the subjectivity common to the shapes of nature and of spirit (hence logical)'' surely cannot possibly be the same as to efface the logical difference between subjectivity and the objective both in ``the shapes of nature and in those of spirit.'' But that the latter is precisely what happens in the Hegelian logic can be seen in the following way. What is a crystal: subject or object? From Hegel's standpoint one must answer: ``in the first section of the concept the crystal is subject, in the second section of the concept, on the other hand, it is object.'' But now there is surely no Hegelian who will maintain that the crystal as subject must first go through the course of subjective logic on judgments and syllogisms that is treated in the first section, before it can in the second section be incorporated into the series of actual natural objects and pass for a proper object. Here there is, as far as the crystal itself is concerned, not a shadow of conceptual movement from subject to object; whereas in our conception of what a substance such as the crystal is there is quite rightly a conceptual movement, and that a true conceptual movement: from conceiving the crystal from the point of view of a logical-grammatical subject to conceiving it from the point of view of an actual object occurring in nature. But when the conceptual movement, however one turns and twists it, always falls, and must fall --- not in \opage{143}the object as object, --- but in thinking subjectivity, then it must surely necessarily become the task of logic to show how the ontological and the psychological in thinking subjectivity determine each other reciprocally; in other words: to bring subjectivity to logically all-round self-understanding by presenting it as \emph{comprehending subjectivity}.

However decisive the logically objective content now surely is for the distinctive development of the forms of the concept, we may nevertheless, with the problem of subjectivity in view, at no point separate the concept from the comprehending activity; the points of view characteristic of the development we therefore designate by: \emph{Comprehending and Concept}, \emph{The Phases of Comprehension}, \emph{Comprehending and Knowledge}.

\begin{center}
{\bfseries a)\quad Comprehending and Concept: Logical Movement of the Idea.}
\end{center}
\addcontentsline{toc}{subsubsection}{a) Comprehending and Concept: Logical Movement of the Idea}

There is no thought that has not arisen through a thinking, no reflection that has not arisen through a reflecting, no concept that has not arisen through a comprehending; if thinking stops, the thought must vanish; if comprehension stops, the concept must vanish: the concept is as inseparable from comprehension as the thought from thinking. What cannot think itself does not have thought in itself either; what cannot comprehend itself does not have the concept in itself either. To be sure, thoughts, reflections and concepts can all be preserved, e.g.\ in writings, without the writings being read; but the thoughts in a writing exist only according to possibility, not according to actuality; this existence according to possibility is expressly designated by the fact that the writing, as the expression of the thoughts, is a logical intermediary between the author and the reader.

How thinking in general relates to reflection has already been shown; but how now does reflection relate to comprehending proper? A certain reflecting is indeed possible without comprehending, but no comprehending is possible \opage{144}without a reflecting; for reflection is the becoming of the concept, and, as the becoming of the concept, itself a comprehending in becoming. Since now comprehension must at every stage direct itself by that which is to be comprehended, it is seen from this that to every stage that reflecting subjectivity runs through there must correspond a particular, if one-sided, development of the concept. As reflection becomes more totalizing, comprehension also becomes more many-sided; only with completed reflecting does comprehending proper begin. And yet comprehending proper is essentially different from reflection, however many-sidedly the latter be carried through. Reflection culminates in the transparency of the moments; even the most fully carried-through reflection never reaches beyond the system of pure reflexes. For pure reflection everything becomes logical, the content as well as the form. Wherever it is merely a matter of a formal comprehending, it is enough to carry reflection through; but when it concerns reality, then the opposition between mere reflecting and comprehending proper shows itself. In order to separate the sphere of reflection from that of the concept, the Hegelian logic draws attention, and rightly so, to the difference between the categories, the one-sided concepts, and the concept itself, the totality of the categories. Each of the categories of reflection by itself is only one side of the concept, but in the concept's own sphere each of the concept's sides is itself the whole concept. But with pure logical totalities logic can never reach beyond reflection; and when a logic cannot reach beyond reflection, and thereby beyond the merely logical, then the whole corresponding philosophy cannot reach beyond the merely logical either, but becomes a panlogism. The Hegelian philosophy is panlogism, not because it holds to ``the concept,'' but because it has not yet arrived at comprehending proper.

For in comprehending proper there is a something that forms a decisive opposition to reflection and thereby to the merely logical; what sort of something, then, is this? The \opage{145}illogical! No, the illogical cannot possibly form an opposition to reflection and the logical, since it must simply be rejected by logical reflection. The immediate! To be sure, the immediate, as we have seen in the foregoing, forms an opposition to reflection; but this opposition is from a logical point of view so mobile that it falls precisely to reflection, since it is surely reflection itself that both posits and sublates it. The real! But when it is to be seen what the real is, reflection comes again and transforms this real for us into a logical reflex in the category of reality; and then we are no further. In order, then, to make the opposition to reflection and the logical so strong that no reflecting shall be able to dissolve it for us and transform it into a merely logical reflex, although the opposition itself naturally must continue to fall under reflection and the logical, we designate it negatively by the category \emph{the antilogical}. Only when logical reflection reflects the antilogical without being able to transform it into a merely logical reflex does comprehending proper set in, and only with comprehending proper comes the concept.

For completed reflection there is always an opposition between the form of the totality and its content; but since the opposition is logical, so is the identity; if the form is logical, so is the content, and conversely. That the mechanical form is determined as logical means, then, on this standpoint, that the content, the mechanical whole, is itself determined as logical; the same holds of the semi-organic and organic totalities: they can one and all be reflected through in such a way that the content's mirroring of the form and the form's of the content makes them both logical reflexes. With comprehending proper, on the other hand, the turn sets in that the content of the totality is antilogical when the form is logical, and the content logical when the form is antilogical. To comprehend the logical by the logical is formal; to comprehend the antilogical \opage{146}by the antilogical is impossible; to comprehend the logical by the antilogical, and conversely, is real. When the logical is comprehended by the antilogical, the concept is obtained in logical form, and all the determinations are determinations of negation; when the antilogical is comprehended by the logical, the concept is obtained in antilogical form, and all the determinations are determinations of position. Only in the Idea of the concept, the concrete unity of the logical and antilogical determinations, can comprehension be completed.

\begin{center}{\bfseries α)\quad Determinations of Negation: The Concept in Logical Form.}\end{center}
\addcontentsline{toc}{paragraph}{α) Determinations of Negation: The Concept in Logical Form}

Pure thinking is negative, and every negation a subjective act. That, independently of subjectivity, there should be objective negations, that an existing thing, an actual object, should actually be able to negate either itself or something else, is, when one properly bethinks oneself, just as absurd a fancy as that a door should be capable of denying either itself or that the inhabitants of the house are at home. ``Negations,'' says Herbart, ``are only in the representation of him who seeks in the object something he does not find; the lack itself is nothing.'' Indeed, so subjective is negation that it presupposes not merely a consciousness, but a thinking consciousness that can use language: negation has existence only in the word. Even if the subjective logic was wrong against Hegel in refusing to see that its principle of contradiction, like every other formal axiom, needed a grounding, it has nevertheless, in a certain sense, been right to hold fast to what is subjective in all thinking, reflecting and comprehending. How the subjective logic, however, from its side also misunderstands what is logical in subjectivity, insofar as in all narrow-mindedness it clings to its psychological standpoint, has already been sufficiently discussed; here it need only be added that its blindness in this respect would rather be increased than diminished by the polemic against the mystical objectivity of the immanent method.
\opage{147} The Hegelian doctrine of negation and of the negative is, in the history of philosophy, a remarkable, perhaps one of the most remarkable examples of how strangely a great scientific discovery can coincide with a great scientific misunderstanding. It is Hegel's immortal merit to have discovered ``the negation of negation,'' ``immanent negation,'' ``infinite negativity''\footnote{``One not seldom hears the opinion expressed that negative knowledge does not satisfy, that philosophy ought to be positive. Now it depends on how one understands the positive. If the positive is taken so immediately that it has the negative outside itself, it is taken in finite form, and is then, in this its finitude, forfeit to contradiction, and thereby to the negative. A philosophy that is dogmatically positive by keeping the negative out is without idea, and positivism without idea is, precisely by the Idea, destined for dissolution in the negative. Therefore dogmatism is always attacked by skepticism. Dogmatism is without idea because it lacks the negative. Skepticism is without idea because it lacks the positive. The positive is in the Idea infinitely negative, and this negative is for that very reason the true positive. The objection that a philosophy is not positive enough is then, in the concept, to be understood as meaning that it is not negative enough. Skepticism is, in the same degree as it has the negative, nearer to reaching the true positive than dogmatism.'' Phil.\ Propæd.\ i Grundtræk.\ Kjøbenhavn 1857, pp.~152--153.}, and yet immanent negation has occasioned much logical confusion.

The point is that ``negation is immanent.'' If negation is not taken as ``immanent,'' then ``the negation of negation'' is an empty play on words. Thus the principle of contradiction in formal logic, $A$ is not not-$A$, must necessarily become comical when it is to be applied to determinate cases or, in other words, to serve for developing something determinate. What is a mirror? ``A negation of a table, for a mirror is not a table, hence a not-table.'' And what, then, is an elephant? ``A negation of a mirror; hence an elephant-negation of a mirror-negation of a not-mirror.'' Just as \opage{148}in a kaleidoscope one can produce from the same bits of glass the most diverse combinations, so by means of the negations set up one can carry on the most arbitrary play with the negation of negation. To avoid arbitrariness, formal logic then clings so fast to the double negation of the principle of contradiction that the progress announced by the negation of negation falls into a ridiculous conflict with the tautology: $A$ is not what $A$ is not; the whole movement of thought is feigned.

``But,'' Hegel concluded, ``if negation is to be immanent, then it must necessarily be \emph{objective}; the immanent is thus the objective, and this objective independent of subjectivity.'' Here is the misunderstanding, an enormous misunderstanding! The misunderstanding arises from the ambiguity in the objective. It goes with the object and the objective as with concept and thought: the object in actuality exists only by being objectified: if one thinks away all objectifying, then there are no objects and nothing objective. But does the converse now also hold, that if one thinks away the objects, then there is no objectifying? Yes, so long as one reflects expressly on actuality, actual existing; but so overreaching is the power of subjectivity that it can objectify to itself not only the existing but also the non-existing, and moreover even objectify to itself this objectifying. If there were not a subjectivity objectifying objects, and that a subjectivity which could at the same time objectify its own objectifying, negation and the negative would never have come into being. Negation is not merely of subjective origin, but the purest, most abstract expression of subjectivity's liberation from everything else, indeed from its own immediacy. For there is nothing that is so positive, no truth so striking, that the hovering subjectivity cannot, by applying its negation, prepare for it a formal annihilation. But precisely as subjective freedom reaches its extreme point in \opage{149}arbitrary negation, subjectivity, wounded by self-contradiction, must react against the arbitrariness of negation by a new negation; thus negation then becomes immanent. Immanent negation is objective; but here is precisely the turning point where the ambiguity of the objective must strike the eye. The actual object is objectified by the objective's determining the subjective and being determined by the subjective. If one now reflects on its being the objective that determines the subjective, then the being of the object is antilogical; if on the other hand one reflects on its being the subjective that determines the objective, then the being of the object is logical. The objective is thus both logical and antilogical; but only the antilogical can be set independently over against subjectivity: the logical, however objective it may be, belongs to subjectivity, wholly and entirely to thinking, comprehending subjectivity. Now when Hegel lays weight on immanent negation's being objective, we gladly grant him that he is right; but when he then acts as if it went without saying that the logically objective must be just as indifferent toward thinking subjectivity as actual objects toward a chance spectator, we must see in this a misunderstanding extremely dangerous for logic and for all sound philosophy.

How immanent negation stands related to the reciprocal determination of the logical and the antilogical can be seen by a glance at the beginning-concept of the Hegelian logic: being. We have in the proposition ``being is'' an immediate unity of the logical and the antilogical. The antilogical is understood, insofar as one here thinks of a being of something existing in actuality, a this-being. The existing in its singularity and punctuality defies thinking subjectivity and resists the logical; but this antilogical again becomes logically determined by subjectivity's apprehending it and holding it fast in a determinate category. In the category this-being is found in turn the category a being, and in the category \opage{150} a being, again the category being. Now when being, thus separated off by abstraction and freed from all antilogical accessory representations, is considered by itself, then this pure being, seen in an antilogical light, is quite rightly the same as a not-being, a nothing. In the much-disputed proposition ``being is nothing'' the Hegelian logic thus gives us the first sample of immanent negation. In what sense this negation is objective emerges from the following consideration. If we reflected on the logical alone, the category being would necessarily continue to be what it is, a determination by the category that which is, hence only in a very improper sense a nothing. If we reflected on the antilogical alone, the determination being would have to be absorbed immediately in its this-being, and here not even the shadow of negation would appear. Only as one looks here, on the one side, at the extreme abstraction of a logical being and, on the other, at the determinate claim of antilogical being, is the opposition between being and nothing driven to its extreme. With regard to the antilogical, nothing is an absolute not-being; nothing has nothing whatever to do with what is; a nothing that is is as meaningless as a green tone; from this standpoint it holds: from nothing comes nothing. With regard to the logical, nothing in its objective not-being is a negative expression for being itself; the same subjective thinking that bears being, the extreme abstraction in positive form, by means of pure intuiting, bears also nothing, the extreme abstraction in negative form, by means of sublated intuiting; the negative side of abstraction thus has just as objective a being as the positive. In this sense it holds that nothing is being, since being is nothing.

It is then incontestable that the becoming of which the Hegelian logic speaks, when it lets the continued alternation of being and nothing sublate itself in its third member: 1) being $=$ $A$, 2) nothing $=$ not-$A$, 3) being $=$ not-not-$A$, \opage{151}is a movement which indeed has objective validity but nonetheless takes place in subjectivity; it holds as a logical determination for the antilogical, but is not therefore itself antilogical.

In close connection with Hegel's misunderstanding of the objectivity of immanent negation stands his thoroughgoing misunderstanding of the objective relation between the \emph{universal} and the \emph{singular}. What is the universal in its essence, subjective or objective? From Hegel's Logic it is impossible to get any definite answer to this question. ``The pure concept,'' it says\footnote{Hegel, Wissenschaft der Logik, 2ter Theil, Berlin 1834, p.~36 ff.}, ``is the absolutely infinite, unconditioned and free. \emph{Essence} has its becoming from \emph{being}, the \emph{concept} from essence, hence from being. But this becoming of the concept has the significance of a counter-thrust in itself, so that what has emerged through becoming is precisely the unconditioned and original. Being has, in its transition to essence, become a semblance, its becoming or transition into other has become a positing; conversely, this positing or the reflection of essence has in turn sublated itself and returned to an original being. The concept is the interpenetration of these moments in such a way that the qualitative and originally being is only a reflection into itself, and this reflection into itself the becoming of other which expresses determinateness: this self-determining determinateness is infinite. The concept is thus in the first place the \emph{absolute identity with itself}, but in such a way that this identity is expressly the negation of negation, the infinite unity of negativity with itself. This bending back of the concept into itself and resting upon itself by means of negativity is the \emph{universality} of the concept. It is precisely the nature of the \emph{universal} to be such a simplex, whose simple being, by virtue of absolute negativity, contains the highest difference and determinateness. The simplicity of being is its immediacy; one can point to it by an opinion, but cannot say determinately \emph{what} \opage{152} it is; it is therefore immediately one with its other, with not-being. Just this is its concept: to be a simplex that immediately vanishes into its opposite: it is \emph{becoming}. The \emph{universal}, on the other hand, is the \emph{simple} which is nonetheless just as much what is in itself richest, because it is the concept.''

In order to understand the Hegelian movement from being through essence to concept, it suffices to understand the movement in the very first trilogy: \emph{being, nothing, determinate being} or \emph{existence}. In this trilogy immediacy occurs twice, immanent negation twice, the universality visible in negativity, or the developed concept, once. As the objective is in this trilogy, so is it in the whole Logic; and the same holds of the subjective. As the universal is in this trilogy, so is it in the whole Logic; and the same holds of the singular. If we now reflect on immediacy, it is evident that the objectivity of the universal in the third member is just as far from the antilogical as in the first member. From this it is seen that the immanent movement from more abstract to more concrete concepts in logical form shows no approximation to determinations of the concept in antilogical form. Now if what really exists expresses its conceptual content in antilogical form, then it is surely proved hereby that the concept in the Hegelian logic is, at its last and highest stage of development, just as far removed from what really exists as in its first abstract beginning. If we next reflect on the negative, it is evident that, if the first immanent negation, not-being $=$ nothing, is, as we have seen, a pure and immanent determination of subjectivity, then the same must hold of that negation of negation whereby the third member reproduces the first in a higher and richer immediacy. Now if negation itself together with the negation of negation in the first trilogy is an immanent determination of subjectivity, then the same holds of all other negations and negations \opage{153}of negations at the following stages right up to the conclusion of the Logic.

It is pure imagination when one supposes by this route to come nearer to what is objective in objectivity by means of concepts such as ground and phenomenon, thing and properties, cause and effect. The negation whereby the cause is converted into its acting is, in logically objective essentiality, just as subjective as the negation whereby being is exchanged for nothing; the negation of negation whereby the cause, through its effect, is restored to unity with itself as reciprocal action stands at just as infinite a distance from the objectively actual as determinate being, determined by the negation of negation, from what objectively has determinate being. It is all subjective in the concept and yet all meant objectively. Here, in the objective logic, is a naive collision between concept and opinion.

``\emph{The singular},'' says Hegel, ``is the universal.'' This proposition is, as we shall later make clear from another standpoint, entirely false if the singular is to be taken in immediate existence; for in immediate existence the singular is antilogical, while the universal in its negations always is and remains logical. If, on the other hand, we do not require the immediate singular but keep to its logical reflex, then the proposition, the singular is the universal, can with full right hold as an expression of the concept. However, in order not to botch such an expression of the concept, we dare not see in the copula ``\emph{is}'' an all too immediate ``return to original being.'' For here there still holds not only, as regards the judgment as a whole, the opposition of reflection, that the identity of the singular in the subject and the universal in the predicate must have its essential difference, but there is here also, as regards the \emph{singular} in particular, an opposition between the antilogical and the logical characteristic of the concept itself, an opposition whose unity must be grasped precisely in comprehending proper. The \opage{154} totality of singular determinations in the existing subject will from the logical side reflect itself in the totality of the universal determinations of the complete predicate, while the singular determinations themselves will from the antilogical side resist reflection: \emph{thus is an antilogical content grasped and comprehended in logical form.}

If, on the other hand, the antilogical is to vanish entirely and the singular be wholly absorbed in the universal, then we must not, with Hegel, think of the grammatical subjects in which the concept has ``immediate determinate being as a something or a thing''; a higher form must then be chosen here. For there really is a form in which the singular is identical with the universal, and that is infinite negativity, subjectivity's own form. In accordance with infinite negativity the universal already appears as the concept of the first trilogy. For what, indeed, is the universal here? The universal is neither being by itself nor nothing by itself nor determinate being by itself. By being a single one of these members by itself, the universal would surely exclude the other members, and by this exclusion be, not the universal, but the particular. That the universal also cannot restrict itself to negating a single member and affirming the others is likewise evident: infinite negativity shows itself in this, that the universal equally negates and, by the negation of negation, affirms all single members, in that in them all it equally negates and affirms itself. In infinite negativity the universal is quite rightly a singular, for the negation that transforms all members into determinations by the universal itself is, as subjectivity's most subjective acting, the categorical expression of pure selfhood, logical singularity. Likewise the logically singular is the universal; for the negatively active selfhood does not maintain itself in determinate being as an immediate individual that is sated by devouring the single members, made particular by their mutual diversity; on the contrary! the negative selfhood negates precisely its own immediacy in negating the \opage{155} immediate members, and in turn affirms the immediate members as expressions of determination in which it affirms itself. In these bendings-back, then, all particular determinacies become universal expressions for one and the same singular, which in them all determines itself as the universal.

From this standpoint one would thus expect that the Hegelian logic, with its essential judgment \emph{the singular is the universal}, would open to us a rich insight into the infinite self-deepening of comprehending subjectivity, and thereby derive the logical determination of the opposition between the subjective and the objective; but how disappointed we are! Only at one or another single place is there something like a run-up to an apprehension of the universal in its subjectivity. ``Das Allgemeine,'' it says\footnote{Hegel.\ Anfr.\ Skr.\ p.~39.}, ``ist die \emph{freie} Macht; es ist es selbst, und greift über sein Anderes über; aber nicht als ein \emph{Gewaltsames}, sondern das vielmehr in demselben ruhig und bei \emph{sich selbst ist}. Wie es die freie Macht genannt worden, so könnte es auch die \emph{freie} Liebe und schrankenlose Seligkeit genannt werden, denn es ist ein Verhalten seiner zu dem \emph{Unterschiedenen} nur als zu \emph{sich selbst}, in demselben ist es zu sich selbst zurückgekehrt.'' That this, which sounds so expressly of the subjective, is nonetheless meant of the objective, that here in the development of the relation of the singular to the universal there is a confusion of the subjective with the objective, of the logical with the antilogical, will be further cleared up when in what immediately follows we consider the matter from the objective standpoint.

\begin{center}{\bfseries β)\quad Determinations of Position: the Concept in Antilogical Form.}\end{center}
\addcontentsline{toc}{paragraph}{β) Determinations of Position: the Concept in Antilogical Form}

The fundamental principle of formal logic, $A$ is $A$, points at once in a logical and an antilogical direction. If the proposition is understood \opage{156} of the general, it becomes abstractly logical: \emph{the general is the general}; if it is understood of the singular, it becomes abstractly antilogical: \emph{the singular is the singular}. In both cases it is without concept; for however formally one may take the concept, formal comprehending will still always require that the general develop itself into a totality of determinations of difference. In such a development the general must then show itself both as one with and as different from itself. But the conceptlessness is most conspicuous from the standpoint of the singular; for a proposition such as \emph{this is this} points so directly to what is objective in objectivity, the antilogical, that the singular in its singularity even becomes downright unnameable.

From this it is explicable that abstract thinking, by holding strictly to the tautology $A$ is $A$, can become so zealous for the logical that, without itself comprehending how, it finds itself captive to the antilogical. In this respect too the logic of the Eleatics and the Megarians presents remarkable examples. Of Stilpo it is related that he regarded only the general, not the individual, as the true. As then, out of zeal for pure identity, \emph{the universal is the universal}, he meant to hold firmly to the logical, he came, as is well known, to put forward the assertion that one could say nothing but identical propositions, e.g.\ ``man is man, horse is horse''; but not ``the horse runs,'' since ``horse'' and ``run'' are not identical; --- he was thus in this way thrown from the logical over into the antilogical, in that he came to deny the validity of the logical judgment, which posits the predicate as one with and yet different from the subject. The same holds of the Eleatics with regard to the ``\emph{One}.'' The reality of the One is built, as we have seen above, on the two tautological propositions: ``what is, is,'' ``what is not, is not''; but tautological propositions easily bring about a sudden reversal of the logical into the antilogical. The \opage{157} manifold is here to be the unthinkable, hence the illogical, the One on the other hand the thinkable, the only thinkable, hence the only logical; and lo, then it turns out that this only logical is precisely antilogical. What the antilogical, what is objective in objectivity, properly means, Plato indeed did not see; but he has nonetheless, with dialectical art and admirable acuteness, especially in his ``Parmenides,'' sufficiently shown that the Eleatic One is on all sides in conflict with every possible form of the logical, that what is by hypothesis absolutely positive on all sides dissolves into the purely negative. ``If it is assumed,'' --- it says --- ``that the One is (εἰ ἔν ἔστιν) in sharp separation from all possible plurality, then this One cannot in the first place be a whole; for if it were a whole, it would have to have parts, and would thereby come to consist in the plurality of the parts. If the One has no parts, then neither has it either beginning (ἀρχή) or end (τελευτή) or middle (μέσον), for such moments would then surely be its parts. But beginning and end are surely the limits (πέρας) of every thing; if the One has neither beginning nor end, it must be unlimited (ἄπειρον), and without shape (ἄνευ σχήματος). It can then partake neither of the round nor of the straight, can be surrounded neither by another nor by itself, can be neither in motion nor at rest, can in its infinite self-contradiction not tolerate any predicate whatsoever at all, can, in a word, not be at all.''\footnote{Phil.\ Propæd.\ p.~131 ff.}

What holds of the absolutely One, insofar as this One is to be --- by express exclusion of all negations --- the absolutely positive, holds also for a plurality of equally positive ones, holds, in other words, for the absolute atoms and the Herbartian reals; only that the proposition \emph{the universal is the universal} is preferably exchanged for the proposition \emph{the singular is the singular}. Every thing that is is here a this-\opage{158}being, this singular, this positive. The one absolutely positive can consequently have nothing whatever to do with the other; whether there are in actuality countless ones here or only a single one is indifferent to the single positive: the difference between one and many concerns only an external consideration, a hovering reflection. The absolute contradiction that is in the single absolute atom, the single absolute real, does not vanish because an infinite multiplicity of such absolutely positive ones is presupposed here; it proves on the contrary that here there is an infinite multiplicity of absolute contradictions. By setting negativity in motion, by requiring connection and therewith unity-in-difference and relativity, the logical everywhere reacts against the abstractly antilogical. And so strong is the reaction that even a thinker as acute and as self-willed as Herbart must make concessions to actuality, must abate somewhat of consistency and go in for the presupposition that the reals, with all their positivity, nonetheless stand in a certain relation to the negative, that with all their peculiarity they nonetheless take up reflexes of otherness, that in their absolute absoluteness they are nonetheless relative, all of which is to say: that the reals, despite what is antilogical in their being, are nonetheless logical.

But just as decisively as the logical reacts when one would one-sidedly hold fast to the abstractly antilogical, just as decisively does the antilogical react in turn when one would dissolve its peculiar determinateness in the merely logical. Hegel's treatment of sensing consciousness in his ``Phänomenologie des Geistes'' is in this respect characteristic\footnote{Cf.\ Philos.\ Propæd.\ p.~37.}. The movement of thought is as follows: ``sensing consciousness is immediate certainty of an \emph{external object}. The expression for the generality of such an object is that it \emph{is}, and is \emph{this}, \emph{now} according to time, \emph{here} according to space. The object is thereby wholly different from all other \opage{159} objects, and sufficiently determined by itself. But this \emph{now} is no more, another has stepped into its place, which however is also a now: the now is in its vanishing something abiding, namely the universal. This \emph{here}, which I mean and point out, has a right, a left, an above, a below, a behind, a before, to infinity. This \emph{here} is thus not the singular, the immediately determined, but: now here, now here, now here! a sum of many, and in these many nonetheless the one, hence the universal.'' But however excellently Hegelian thinking knows how to discover the universal in the immediately individual and to find the logical where it seems to be lost in sensuous representations and contingent opinions, just so one-sided, again, is Hegel's endeavor to dissolve everything, even the antilogical, into logical reflexes. The panlogism which Hegel already laid the foundation of in the first section of his Phenomenology of Spirit, and which he carried through in the Logic and in the other parts of philosophy with such astonishing comprehensiveness, consistency and perseverance, was bound at last to call forth a reaction. This reaction, which had its most talented representative in L. Feuerbach, now went to the opposite extreme, let the universal vanish in the immediate singular, and ended by appealing to ``the unsayable.'' ``Das Sein'' --- says Feuerbach\footnote{Grundsätze der Philosophie der Zukunft \S~26 ff.} --- ``ist an und für sich selbst Nichts. Aber eben deswegen ist auch das Sein, wie es die speculative Philosophie in ihr Gebiet hereinzieht und dem Begriffe vindicirt, ein pures Gespenst, das absolut im Wiederspruch steht mit dem wirklichen Sein und dem, was der Mensch unter Sein versteht.'' Further (\S~28): ``Welch ein gewaltiger Unterschied ist zwischen dem Diesen, wie es Object des abstracten Denkens, und eben demselben, wie es Object der Wirklichkeit ist! Dieses Weib z.\ B. ist mein Weib, \emph{dieses} Haus \emph{mein} Haus, obgleich jeder von seinem Weibe und seinem Hause wie ich sagt: dieses Haus, \opage{160} dieses Weib\ldots Das Sein, gegründet auf lauter solche Unsagbarkeiten, ist darum selbst etwas Unsagbares.'' --- Here we stand, then, at the antilogical again.

That thought, after so many systematic attempts, can continue to swing between two such decided extremes as the Hegelian and the Feuerbachian is a proof that the concept of an actual comprehending has not yet been brought to clear consciousness. If the two propositions \emph{the universal is the universal} and \emph{the singular is the singular} are each by itself without concept; if the former is the abstractly logical, the latter the abstractly antilogical, then it is seen from this that comprehension must rest on a reciprocal determination of the universal and the singular, hence on a reciprocal determination of the logical and the antilogical.

The proposition \emph{the singular is the universal} cannot express the reality either of the concept or of comprehension so long as the antilogical element hidden in ``the singular'' is overlooked. But this element gets its logical validity only by the opposite extreme being designated, not by the proposition \emph{the singular is the singular}, but by the proposition \emph{the universal is the singular}. Only when it becomes clear what properly lies in this proposition will it be able to become clear what the reciprocal determination of the logical and the antilogical has to signify.

Of the singular in its immediate existing it holds that it both is and yet is not the universal. What, e.g., is this singular? A table! As existence the table is the singular; but as category it is the universal. Now if it is to become really evident what the relation is between the singular and the universal, existence and category, then it is in truth not enough that we start from the singular and ask about the universal; here it is expressly required that we also reverse the relation, start from the universal and ask about the singular. What, then, is a table? Just such a thing as this existing one! The universal is this singular. To be sure, \opage{161} contradiction now breaks out from both sides; for the singular in its singularity will not be the universal: existence protests against its category, the antilogical against its logical. On the other side the universal in its generality will not be this singular either; for this singular surely falls outside that singular (this table has nothing to do with that table), and precisely for that reason the difference between the singular and the singular falls outside the universal: the category protests against its existence, the logical against the antilogical. But as contradiction thus breaks loose at once from two opposite sides, it dissolves itself and succumbs to the concept.

The logically positive judgment, \emph{the singular is the universal}, rests on a negative foundation; for it presupposes that the singular and the universal are two different conceptual expressions for the same infinite negativity, it presupposes that the singular is nonetheless not the singular, the universal not the universal. The logically negative judgment, the \emph{singular is not} the \emph{universal}, rests on a positive foundation; it presupposes that the singular determinations negate the universal determinations only because they themselves are not merely antilogical but logical, hence universal. From the positive as well as from the negative side the singular is comprehended by the universal; but the reality of the concept rests expressly on the antilogical's being comprehended by the logical.

But if the reality of the concept requires that the antilogical be comprehended by the logical, then the reality of the concept must surely also require that the logical be comprehended by the antilogical; thus we come then to the antilogical judgment: the \emph{universal is the singular}. That the antilogically positive judgment likewise rests on a negative foundation is easy to see. The universal is not the singular insofar as the positions in the universal are logical, in the singular antilogical. But if the antilogical now vanishes in the logical, this latter must itself, by means of the tautology, become antilogical. To this corresponds what Stilpo says of the cabbage: ``The cabbage that is offered for sale here is not; for cabbage already was \opage{162} 1000 years ago, and this cabbage is not that which was 1000 years ago.'' If, on the other hand, the logical is to vanish in the antilogical, this latter, despite its ``unsayability,'' must nonetheless, by means of the tautology, become abstractly logical. To this corresponds the Feuerbachian: \emph{``Dieses Weib ist mein Weib, dieses Haus mein Haus''}. The reality of the concept must then, as a consequence of this, always require that its form be logical when its content is antilogical, and conversely. Only in antilogical form does the content of the concept get objective existence; this Kant from his standpoint expressed in the well-known proposition that a hundred actual thalers do not contain the least bit more than a hundred possible ones. But, it goes without saying, the concept could surely not possibly be expressed in antilogical form if its logical determinations of content were not themselves determining for the existential expression. It is this side of the matter that Hegel so one-sidedly brings out when he again and again insists that the concept realizes itself through itself.

Kant and Hegel misunderstand, each in his own way, the original relation between the real concept and its existence. Kant is certainly right when he assumes that the logical form of the concept, considered by itself, must be subjective, and that all determinations by this subjective form can be given without anything being thereby yet decided concerning the existence of the concept. ``A concept does not lose a single one of its properties because its object does not exist, and does not gain a single property by its object's being found existing.'' This can in other words be expressed thus: the antilogical, by which each of the concept's logical determinations is to be integrated in order that the concept itself may obtain its complete reality, cannot be reached by the accretion of new determinations that are in turn fetched from the logical. But herein Kant then overlooks that a real concept cannot possibly rest content in a system of merely subjective determinations, that it can in no way be satisfied with an intellectual determinate being in subjectivity's pure understanding; for \opage{163} what is a real concept? Answer: a concept whose content is always antilogical when its form is logical, and conversely.

Now insofar as the content is presupposed in actual existence, the concept in its logical form is a concept of abstraction. In this sense the zoological concept \emph{horse} is a concept of abstraction. In finite, psychological subjectivity such concepts of abstraction may certainly be said to have a sort of abstract existence as pure logical forms. But the logical positions nonetheless have conceptual validity only in and through the antilogical ones; the abstract determination of form is a halved determination of the concept; such an abstract determination of form first gets reality as a determination of content in unity with its actual content. The concept of abstraction \emph{horse} is thus not, as Kant assumes, a complete concept, such that a hundred actual horses should contain no more than a hundred possible ones. The concept of abstraction is on the contrary a shadow-concept: so far is it from being the case that such a shadow-concept should be indifferent toward its existence, that its reality rests precisely on its existence; so far is it from being the case that a hundred possible horses should contain in the concept as much as a hundred actual ones, that the concept in one single actual horse contains, on the contrary, infinitely more than the pure concept, the concept of abstraction of all possible horses, since a real concept in its reality always contains infinitely more than its logical adumbration\footnote{Of this zoology as well as the other real sciences are excellently well aware.}.

Insofar as the content of the real concept is presupposed in a possible existence, the concept in its logical form is a concept of construction. In this sense, e.g., the concept \emph{house} is, for the subjectivity of the master builder, a concept of construction. Now that the concept of construction can just as little stop at the merely logical form as the concept of abstraction is evident. But while abstraction is always determined by the logical judgment \emph{the singular is the universal}, con\opage{164}struction must be determined by the antilogical judgment \emph{the universal is this singular}. When someone is asked what a triangle is and answers the question by drawing a triangle, he answers by a construction, that is, he expresses his judgment about the triangle not in a logical but in an antilogical way. By means of abstraction the antilogical positions are converted into logical ones; by means of construction the logical positions are converted into antilogical ones. Insofar as abstraction must let the content remain what it is in its existence and can only alter the form, the comprehending conditioned by abstraction will grasp the antilogical content in logical form. As regards construction, on the other hand, it surely cannot take place with full insight if, during the execution of the plan, something extraneous and disturbing to the concept intrudes; the logical determinations of position must then, from this standpoint, be taken as the concept's own determinations of content, the antilogical ones on the other hand as determinations of form. The comprehending conditioned by construction will therefore always grasp a logical content in antilogical form.

From the standpoint of abstraction as well as from that of construction, comprehension and the concept show us an express unity-in-opposition of the subjective and the objective. Here we come now again to Hegel's misunderstanding. For it is not at all the case, as Hegel thinks, that the reality of the subjective concept should in the Logic coincide with the vanishing of subjectivity itself in the objective. To be sure, subjectivity in Hegel vanishes in the objective in order to return in teleology\footnote{``Es ist'' --- remarks Hegel in connection with Kant's separation between ``objective propositions'' and ``subjective maxims'' with regard to the teleological (Logik.\ 2ter Thl.\ p.~215) --- ``auf diesem ganzen Standpunkte dasjenige nicht untersucht, was allein das philosophische Interesse fordert, nämlich welches von beiden Principien an und für sich Wahrheit habe; für diesen Gesichtspunkt macht es keinen Unterschied, ob die Principien als \emph{objektive}, das heißt: hier äußerlich existirende Bestimmungen der Natur, oder als bloße Maximen eines \emph{subjektiven} Erkennens betrachtet werden sollen'' \ldots\ But what is ``an und für sich Wahrheit'' when the subjective and the objective run out into one?}; but since this its return is in the Hegelian logic just as mystical and enigmatic as its vanishing, it is of no help to appeal to it.
\opage{165} How impossible it is, according to the Hegelian method, to arrive at any characterization whatever of what is objective in objectivity has already been shown in the foregoing, and can be seen still further from the fact that, with all its variations on the rich theme of immediacy, it does not know how to make the difference between logical and existential immediacy logically recognizable. ``Es sind, wie bereits erinnert worden, schon mehrere Formen der Unmittelbarkeit vorgekommen; aber in verschiedenen Bestimmungen. In der Sphäre des Seyns ist sie das Seyn selbst und das Daseyn; in der Sphäre des Wesens die Existenz und dann die Wirklichkeit und Substantialität, in der Sphäre des Begriffs außer der Unmittelbarkeit, als abstrakter Allgemeinheit, nunmehr die Objektivität \ldots Seyn ist überhaupt die \emph{erste} Unmittelbarkeit, und \emph{Daseyn} dieselbe mit der ersten Bestimmtheit. Die \emph{Existenz} mit dem Dinge ist die Unmittelbarkeit, welche aus dem \emph{Grunde} hervorgeht, --- aus der sich aufhebenden Vermittelung der einfachen Reflexion des Wesens. Die \emph{Wirklichkeit} aber und die \emph{Substantialität} ist die aus dem aufgehobenen Unterschiede der noch unwesentlichen Existenz als Erscheinung und ihrer Wesentlichkeit hervorgegangene Unmittelbarkeit. Die \emph{Objektivität} endlich ist die Unmittelbarkeit, zu der sich der Begriff durch Aufhebung seiner Abstraktion und Vermittelung bestimmt.'' By considering this scale of immediacies one will easily convince oneself that the immediate at every \opage{166} stage, from the lowest to the highest, is of purely logical origin. The formation of every such objective logical immediacy does not rest at all on anything existential in itself, but exclusively on the universal conceptual development of thinking subjectivity. The immediacies named have each, at its own stage, come about through a negation of negation, hence through a reflection sublated within reflection; they are nothing but affirmations without a single existential position. The immediacy which is here to represent objectivity is therefore in actuality further from being anything objective than the actor who in ``A Midsummer Night's Dream'' is to represent the Wall is from being a wall.

Existential immediacy is antilogical; the antilogical is what is objective in objectivity. Instead of abstracting from this objective element and developing an ``objective concept'' without objectivity, logic must aim precisely at clarifying the unity-in-opposition of the subjective and the objective by clarifying from all sides the unity-in-opposition of the logical and the antilogical.

\begin{center}{\bfseries γ)\quad The Reality of the Concept: Perfect Comprehending.}\end{center}
\addcontentsline{toc}{paragraph}{γ) The Reality of the Concept: Perfect Comprehending}

The unity-in-opposition of the subjective and the objective is determined in the concept itself by a unity-in-opposition of determinations of negation and determinations of position. All determinations of negation are moments of knowledge, all determinations of position moments of power. The real concept in its completion, that is, the concept itself according to its absolute reality, is a totalized unity of negative and positive determinations, of moments of knowledge and moments of power. Perfect comprehending is therefore not a merely human but a divine matter; here the opposition between psychological and ontological subjectivity is seen.

At the psychological standpoint power is separated from knowledge. The system of determinations of power yields the actual \opage{167} surrounding world; the system of determinations of knowledge is found in subjectivity; the former presents reality in a one-sided manner, the latter ideality; the former expresses the concept in antilogical, the latter in logical form. The divorce between the subjective and the objective is here so thoroughgoing that the concept of existence must seem at one and the same time both to demand existence and to oppose itself to existence. ``If the whole system of sensible intellectual determinations is placed in the intellectual form, then the sensible vanishes as a moment of reflection sublated in the thinkable; that is to say: what exists turns into a concept of existence, but existence goes out. If, conversely, the whole concept of intellectual sense-determinations is placed in the sensible form, corresponding to immediate sense-impressions, then the object steps into existence, but the concept goes out.''\footnote{Cf.\ Forelæsn.\ ov.\ ``Phil.\ Propæd.'' 1860--61, pp.~24--30.}

The concept, as concept, cannot possibly stand still: a stone can \emph{lie} still, a wall can \emph{stand} still; but a concept can no more enter into sensible rest than it can engage in the sensible difference between lying, walking and standing. And yet the concept is at one and the same time both at rest and in motion; its ideal motion is one with its ideal rest, for the motion is a comprehending. It is an obstinate prejudice, a prejudice to which ``the objective logic'' has on its side given just as ample nourishment as ``the subjective'' has on its own, that the concept, ``that which is in and for itself,'' should find itself objectively in an objective being, instead of resting on an actual comprehending.

To be sure, from the psychological point of view it looks as though a finite separation had to be made here between the becoming of concepts and their existing; concepts do become, after all, as we develop them logically, and as concepts of existence they have an actual existence in the objects from which we abstract them and in which we recognize them again. Thus, e.g., the concept feldspar has \opage{168} its existence in the corresponding stone, and in the stone's sensible existence a sensible rest. Further, the concept feldspar has a becoming in the mineralogist's consciousness when he either develops it himself or has it developed for him by another. What in this way holds of concepts of abstraction must then also hold of concepts of construction. In the master builder's consciousness the concept of the house may with reason be said to have an ideal existence, while during the building's erection it is in a real becoming; when the building is at last finished, the realized concept finds a sensible rest in the resting building. Here, then, as far as the concept itself is concerned, there is a finite separation between ideal existence, real becoming and real existence.

But how wavering and slippery such separations are will easily be seen on closer consideration. The existing object is not the same as the existing concept: for if the concept were without further ado the object, then by breaking off a piece of the object one would have to break off a piece of the concept; the red object would then be a red concept, the heavy one a heavy one, the wet one a wet concept, etc.; all of which is absurd. From this point of view it goes with the realized concept as with the realized possibility: it vanishes in its reality; its immediate existence is its destruction. To abstract the concept from existing objects is therefore, in other words, to call back to life the concept that has died out in the object's existence, to kindle anew the spark of ideality that is extinguished in reality; that is to say, to begin comprehension over again, for the concept is only in a comprehending. To realize the concept in a construction answering to its demands is to sublate the concept by the concept itself, and to put reality in the place of ideality.

In this way the real concept can never come to rest: when its form is logical, then, as a concept of abstraction, it has its content outside itself; for then the content is anti\opage{169}logical. When its content is logical, then, as a concept of construction, it has its form outside itself; for then the form is antilogical.

If this unrest is to cease, that cannot happen by stopping the movement, for by what should it indeed stop? By existence! But when existence enters, the concept goes out, since the form of the concept is logical. By the form of the concept! But when the concept puts on its logical form, it falls into discord with the form of existence, since the form of existence is antilogical. Instead of stopping the movement we must on the contrary infinitize it, and let it complete itself in a circuit.

The logical circuit: from concept through existence to concept, from existence through concept to existence, and so forth, can then be more precisely designated as a movement: from construction through abstraction to construction, from abstraction through construction to abstraction. What matters here is only to get the categories of abstraction and construction to resolve into fuller, more descriptive expressions. Insofar as abstraction in particular conditions that act of comprehension by which the antilogical form of the content is transformed into logical form, this category must be integrated by everything that belongs to the formation, appropriation and working-up of the ideal moments of the real concept, and be completed in what we shall here call \emph{logical recollecting}. By virtue of appropriation and logical recollecting, subjectivity thus possesses its real concepts in an ideal manner, and only on the presupposition of such an intellectual possession can concepts of construction form themselves. Seen in its relation to the construction demanded, the concept is a subjective or anticipated end; the constructing itself, with its conditions, is the carrying-out of the end by corresponding means; what is brought forth by the construction, the realized end, is finally the existence in which the concept dies, in order again to yield material for a new abstraction and to rise transfigured in logical recol\opage{170}lecting. The one-sided category of construction must here be integrated by everything that logically belongs to intellectual-real production.

How essential intellectual production is to formal comprehending we are reminded by the well-known proposition: \textit{scimus, quia facimus}; how necessarily an intellectual-real producing belongs to a real comprehending is characteristically indicated by Schelling's brilliant but foolhardy words: ``Ueber die Natur philosophiren heißt die Natur schaffen''.

At the psychological standpoint the movement of the idea corresponding to comprehension can indeed be recognized \textit{in abstracto}, but not carried through \textit{in concreto}. For psychological subjectivity, concepts of abstraction go in one direction, concepts of construction in another; here reality, the object of the concept, falls outside the concept itself; here power falls outside knowledge. In ontological subjectivity, on the other hand, power is identical with knowledge; and on this identity perfect comprehending rests. Since all comprehending rests essentially on recollecting and producing, real comprehending must necessarily be grounded in an ability, a having-power. The concept in its reality is totality; in totality the individual parts are at one and the same time members and moments. As far as the concept is concerned, this double-sidedness can no longer be explained by a pure separation of reflection between the immediate and its reflex: the existing members of the totality have, as expression of immediate and real power, an antilogical independence. With this immediate independence these members are now objects for perfect comprehending; but how do they become moments? They become transparent to infinite knowledge by means of power; in perfect transparency they are moments. But what then becomes of the antilogical independence? It stands by power, but is seen in the light of knowledge; for infinite power is itself transparent to knowledge. As \opage{171} finite objects of knowledge the existing members are elements in the real concept, in that they are moments for ideal power, infinite self-power: the ideality of power is the reality of knowledge. The movement of the idea is an unceasing conversion of the logical into the antilogical, of the ideal into the real, and conversely: \emph{an infinite reciprocal determination of recollecting and producing}! But this movement is in turn rest; the moment of knowledge is not to be disquieted in a finite manner by searching for its corresponding element of power, since it already has its integrating element in its power.

On the identity of the movement of the idea with rest depends the finite separation of a rest in which movement is lifted and a movement in which rest is lifted. The rest in which the movement of the idea is lifted is the distinguishing mark of outwardly existing things; here the concept is immediately lost in its existence; the concept is, in other words, not concept but object. Now if this outward, antilogical rest were actually absolute, an absolute negation of all movement of the idea, then the sensible understanding would be right when it asserts that things are not realized concepts, that there is in the world of things a system of forces and conditions that falls outside all knowledge. But how superficial this way of seeing is can be perceived when we reflect that things could not possibly be actual objects if they were not realized concepts, and that they could not possibly have subsistence as realized concepts if they were not borne by an infinite comprehending.

The existing thing is an actual object; the actual object rests on an actual objectification; but the objectification carried through in actuality rests in turn on an actual comprehending. However sure psychological subjectivity is of its separation of the subjective and the objective as long as it can keep to immediate impressions, it becomes just as unsure and irresolute when reflection is set in \opage{172} motion. --- ``Sound common sense demands that I should immediately become conscious that there is something outside my consciousness; but in becoming conscious of something, it is, after all, in my consciousness; insofar as it is outside, I cannot become conscious of it: thus I am to become conscious that what is in my consciousness is not there; what a contradiction! When the sun shines and the trees are mirrored in the stream, sound common sense thinks that the light is brightness whether there is anyone who sees it or not, and that the trees are actually mirrored in the stream whether this mirroring is reflected in a seeing eye or not. But the question is how one is to convince oneself of this. By sight? That is a contradiction, since what is at issue here is the phenomenon, not insofar as it is seen, but insofar as it is not seen. By thought? That is a contradiction; the phenomenon, as phenomenon, does not allow itself to be thought; for in that case the man born blind could think it too. By the testimony of others? An empty evasion, since here it must go with the one as with the other.''\footnote{Phil.\ Propæd.\ p.~41.}

In order to examine what lies in this, we shall in this connection reduce the general problem of objectification to three particular ones: the problem of motion, the problem of change and the problem of possibility. Thus:

\emph{The problem of motion.} Among the several proofs the Eleatics devised against the reality of motion we mention here ``the flying arrow.'' ``If a body were to move, it would follow that it must be at one and the same time in motion and at rest; for every body that moves must at each instant be in a space that is exactly equal in size to itself. But thus the body is at rest; for what finds itself, even if only for an instant, in a correspondingly equal space is, with respect to this space, at rest. If we imagine an arrow that with the greatest speed flies through a space, then \opage{173} it must at each instant be in another space and, insofar as it occupies this, be at rest at each instant. Such a flying arrow, which is always: here, here, now, now (ἐν τῷ νῦν κατὰ τὸ ἴσον), must then be thought to be at one and the same time both at rest and in motion, which is a contradiction.''\footnote{Cf.\ Mathm.\ og Dialekt, p.~137 ff.} One generally exposes what is sophistical in the proof by laying weight, after the example of Aristotle, on the reciprocal determination of continuity and discreteness; but by confining oneself to this one nevertheless basically overlooks the real point. When attention fastens on the flying arrow, it does not, after all, fasten on a formal concept, but on an actual object. It is not the concept of motion in its logical form, but what is objective in motion, the concept of motion in its antilogical form, actual motion, that according to the Eleatics' way of seeing is the unthinkable. From the logical side the continuity of a motion can very well be united with its corresponding discreteness, insofar as the past can always, by means of recollection, be united with the present; but what absolutely matters here is precisely the circumstance that ``the flying arrow'' is without recollection; it has in its present nothing past. What is objective in the object cannot possibly move without transition from a preceding to a succeeding time, from a preceding to a succeeding place; but the relation between a preceding and a succeeding is a reflection that falls outside the object itself; therefore the object can advance neither in space nor in time; its place is always here, here! its time is now, now! If it nevertheless holds that an actual object in its motion actually traverses a determinate space in a determinate time, the explanation is that what has been traversed, what is past, is preserved in subjective reflection, in the observer's recollecting. Actual motion, in other words, can \opage{174} only be objective by being objectified; but when it is objectified in actuality, it has objective validity.

\emph{The problem of change.} If the problem of motion can be solved only on the condition that the objective which determines objectification is determined by objectification itself, then the same must naturally hold of the problem of change, since every change must, after all, presuppose a motion. The most objective conception of what is objective in change relies on the presupposition that change rests on the displacement of parts, hence, when we go to the limit, of atoms. But whether this displacement otherwise consists in connected parts being separated, or in separated parts being united, in both cases change must presuppose a relation, an actual and active relation among the parts themselves. Here there is in change an attracting and repelling that must be referred to the parts themselves. But now it has appeared from many sides in our foregoing that objective relation is essentially a reflection, and reflection originally subjective. The contradiction that all objective relation rests on subjective reflection finds its complete solution in the ontology of objectification.

From the psychological standpoint it would certainly be an absurdity if someone were to assert that no changes could take place in actual things except when an observer was present who could objectify for himself the changes taking place; but, seen from an ontological point of view, the matter holds true. In actual change the reciprocal determination of the logical and the antilogical is conspicuous. $A$ changes by running through the states: $A_0$, $A_1$, $A_2$\ldots\ If we look here at the abstract logical identity $A = A$, then the difference between $A_0$ and $A_1$ is abstractly factual, hence antilogical; if the difference of $A_0$ and $A_1$ is determined abstractly logically, then the identity $A = A$ is abstractly factual, hence antilogical.

\opage{175} The difference between psychological and ontological subjectivity may be characterized in a peculiar way by the corresponding difference between a psychological and an ontological objectifying. Merely psychological objectifying presupposes a finite separation of knowledge and power, and therefore stands in a finite and contingent relation to objects; ontological objectifying, on the other hand, which presupposes an original unity of knowledge and power, stands in so essential, intimate and necessary a relation to actual existences that there can here be as little question of objects without objectifying as of consequences without grounding, effects without causality, concepts without comprehending.

\emph{The problem of possibility.} Ontological objectifying is immanent, that is, a determining inherent in the objects themselves; but as a determining that proceeds from subjectivity, immanent objectifying is at the same time a transcendent acting, a determining that hovers above actual objects. Only by this double-sidedness of its own can objectification master what is double-sided in objectivity itself, namely the necessary and the contingent. From this point of view the problem of possibility is essentially a problem of objectification. It is quite characteristic of the category of possibility that when it is fixed objectively it points toward the subjective, and when it is fixed subjectively it points toward the objective. According to the logic of the Megarians there could in actuality be no possibility. ``Nothing,'' says Diodorus Cronus, ``is true except what is actual or will at some time become so. Consequently nothing is possible either except what is actual or will at some time become so. What does not become actual is impossible, but what becomes actual was necessary.'' If something in nature is possible, then, from this point of view, it is also actual and necessary; if something is not actual, that is because its actuality is not possible (ὅταν ἐνεργῇ, μόνον δύνασθαι, ὅταν δὲ μὴ ἐνεργῇ, μὴ δύνασθαι). But what then does this prove? That \opage{176} the concept of possibility is of purely subjective origin. It begins, psychologically seen, as a concept of uncertainty, a concept by which subjectivity expressly makes known that it must let the objective itself decide. What is objective in objectivity is therefore not something possible; it is on the contrary something actual. The actual is because it is; but what thus is because it is, is from the one side something contingent, from the other something necessary. But if the contingent is to be objective, then the possible must be so too, for contingencies alternate, and the possible lies in the alternation of contingencies. If one now says: the possible is nevertheless subjective; it does not lie in the alternation of contingencies, for contingency follows upon contingency as one actual thing upon another, and the following takes place with necessity --- then it is, accordingly, what is necessary in actuality that is the objective. But if necessity, namely the hypothetical one: if a is, then b must be; if the contingency a has occurred, then the contingency b must occur, is actually to be the objective, then possibility itself must be objective; for possibility is then the expected but not yet occurred necessity. Conditioned necessity is in a constant becoming, and possibility is precisely this becoming of necessity. ``It is altogether a remarkable relation that obtains between real possibility and its conditions. Only with respect to the actuality demanded is possibility real; as long, therefore, as a condition is still lacking here, the possibility cannot become actuality and is to that extent impossible; but when all the conditions are present, the possibility is, after all, no longer possibility either, for then it is actuality,''\footnote{Cf.\ Forelæsn.\ ov.\ ``Phil.\ Propæd.'' 1860--61 pp.~98--100.} that is to say: the existence of a conditioned necessity.

As the becoming of necessity, possibility is objective and immanent, and yet possibility is at the same time subjective and hovering above, for where something is merely possible its opposite is \opage{177} of course also possible. That the actual rests in possibility and proceeds from it means: the objective is determined subjectively; that the possible rests in actuality and is conditioned by it means: the subjective is determined objectively. As the category that designates the transition from the subjective to the objective, possibility is the concrete category of objectification, which has taken up into itself motion and change. Le présent est gros de l'avenir, says Leibnitz; this pregnancy proves precisely the antilogical marriage of the objective with the subjective in infinite objectification; but this infinite objectifying is one with perfect comprehending.

\begin{center}
{\bfseries b)\quad Concept and Idea: The Phases of Comprehension.}
\end{center}
\addcontentsline{toc}{subsubsection}{b) Concept and Idea: The Phases of Comprehension}

``The Idea,'' it says in the Hegelian logic, ``is the unity of the subjective and the objective, the unity that expresses the return of subjectivity to itself after the opposition of the objective world has been overcome, etc.'' This determination of what ``the Idea'' is may very well be acknowledged as correct in a formal respect; but what does it help to refer to the unity of the subjective and the objective when these concepts themselves are empty forms? ``Die Idee,'' says Hegel, ``ist die \emph{Wahrheit}; denn die Wahrheit ist dieß, daß die Objectivität dem Begriffe entspricht''; but when objectivity is concept and the concept objectivity, then the unity of the two is, after all, merely formal, and the development of the idea a pure formalism. As the concept is determined, so too the idea; as the unity-in-opposition of the subjective and the objective is developed, so too the idea is developed. If one says: concept and idea are related as absolute form to its absolute content, this conception may likewise hold; but then a profound development of the relation between the absolute and the relative is presupposed here, and such a development must then coincide with the development of the idea itself. The objection that holds against the Hegelian concept, namely that with all its ideal concre\opage{178}tions it is nevertheless only formal, because the logical lacks its antilogical, the same objection holds also against the Hegelian Idea. In order to ward off formalism we shall in this connection determine the idea as the absolute unity of the logical and the antilogical; we have then at the same time determined it as the absolute unity of the absolute and the relative, of subjectivity and the objective, of the real concept and perfect comprehending.

The absolute unity of the logical and the antilogical, of subjectivity and its objective content, cannot be thought as an immediately existing one; its being is a movement of the idea; only insofar as we are able to present the movement of the idea are we able to make the concept clear through the idea, the idea through the concept. If the movement of the idea is to be presented, comprehension must move in a threefold circle, and in doing so exhibit the following three different phases. The antilogical becomes a moment for the logical by power being absorbed in knowledge: this circular movement shows us the originality of knowledge. The logical becomes a moment for the antilogical as knowledge is absorbed in power: this circular movement shows us the originality of power. The logical and the antilogical are articulated by all-sided reciprocal determination, as the last and highest circular movement of comprehension is determined by a unity-in-redoubling of the two preceding: subjectivity is seen in the idea itself.

\begin{center}{\bfseries α)\quad The Originality of Knowledge: universalia ante rem.}\end{center}
\addcontentsline{toc}{paragraph}{α) The Originality of Knowledge: universalia ante rem}

Since no knowledge is possible without a knower, it follows of itself that the originality of knowledge must presuppose an original knower, that here, therefore, in a logical sense, reference is made to the principle of knowledge, the original itself, knowing subjectivity. With knowing subjectivity the universal, what is logical in itself, is at the same time posited as the overarching; but whence, then, comes the antilogical? Naturally, according to our foregoing, from otherness, from the immediacy of the individual, from the objective! \opage{179} The movement in this circle is therefore to be determined as a transition from selfhood through otherness back to selfhood, from the universal through the individual back to the universal. That the universal, which in this movement is the first and the last, must be conceived as the eternal, what is independent in itself and absolutely valid, is indicated in the heading by the scholastic expression, characteristic of the phase: universalia ante rem.

The designation ante rem must not, however, be misunderstood here, as though universalia could actually be without res; on the contrary, a logical expectation of res is awakened when universalia are presupposed; so too, after all, logically seen, an expectation of the consequence is awakened when the ground is presupposed, of the effect when the cause is presupposed, etc.

The next misunderstanding that must be warded off here is that an expression such as universalia sunt ante rem should point to an immediate being of the universal, an eternal resting, without movement and without acting. This misunderstanding, of which the one-sided Realists of the Middle Ages were especially guilty, is removed by conceiving res, the existing individual, as the middle term by which the universal, the universal that determines itself in all things, posits and determines its unity-in-difference with itself.

One more disturbing misunderstanding we must see to removing, namely that the universal in its self-active ideality should be empty and uniform. To remove this misunderstanding, however, it is not enough to deliver logical eulogies over the universal, to declare it the eternally existing, the true, the divine, etc. The task is to show how it comes about through comprehension that the universal, although it must everywhere be the same, can be what is infinitely rich in content in itself, and, although in itself what is infinitely rich in content, can nevertheless continue to enrich itself infinitely. To posit the universal as the absolute is on the whole a distinguishing mark of idealism; --- for the scholastic realism of the Middle Ages corresponds precisely to the philosophical tendency which, in \opage{180} present-day usage, is called the idealist one; --- here, then, the solution of the task will require a characterization of the logical peculiarity of idealism.

With respect to the misunderstandings cited above, it is chiefly two forms of one-sided idealism that must attract attention in the present connection, namely the Platonic and the Hegelian. With all the difference that otherwise obtains between them, these two forms of idealism nevertheless have this in common, that they both conceive the universal as what is objective in itself. As regards the Hegelian misunderstanding on this point, the matter has already been elucidated in our foregoing. Here we shall merely cast a further glance at the Platonic misunderstanding. ``The certainty of the reality of the Ideas'' --- says Poul Møller\footnote{Phil.\ Hist.\ pp.~161--62.} in the section on Plato --- ``rests on the certainty of the reality of science, which aims at the thinking of the essence of things. There could be no science if one assumed with Protagoras and other Sophists that the vanishing sensible phenomena were the only true. The Ideas are the same as the essence of things. But there are several Ideas, because the Eleatics' doctrine of unity, which excludes all multiplicity, is incorrect. Yet there is unity, coherence in the plurality of Ideas; for rational thinking must combine concepts in higher, more general unities, and to this coherence between concepts there must also correspond a coherence in the real. The Ideas are not isolated unities that are for themselves without relation to one another. Their mutual relation consists in subordinate Ideas being comprehended in higher ones as their common unity. Thinking that seeks after unity then requires for its perfect satisfaction that all Ideas be united in one highest Idea.'' --- The system of the Ideas is thus, according to this way of seeing, one with the system of scientific concepts; the \opage{181} objectivity which must be ascribed to the scientific concept, as thinking's expression for the essence of things, is therefore with Plato, as with Hegel, to be ascribed to the universal. Now since what is objectified in knowledge is undeniably, with respect to what is universally valid, objective in quite another sense than what is immediately apprehended by sense-perception and determined by mere opining, it is from this alone explicable that the opposition between opining and knowing has been able to mislead thinkers such as Plato and Hegel into confusing the universal with the objective, although the universal itself is not merely the subjective, but the only thing that is absolutely subjective in itself.

How absolutely subjective the universal must be is seen already from what was elucidated in an earlier connection about its relation to the negative and its unity with subjectivity; for if there is anything that is absolutely subjective, it must surely be subjectivity itself. At the stage of reflection it was reflection, pure reflecting, that was the form of subjectivity; at the stage of the concept it is now the universal. Let us therefore reflect once more on the category of the universal; what, then, does the universal really mean? One could say: the universal is unity, the true, all-embracing unity of existence. The universal must therefore be the unity that has the manifold within it, but in such a way that it at the same time has multiplicity and diversity over against it. The universal must, in its categorial meaning, necessarily satisfy the following demands:

1) As the unity of all multiplicity, the universal must permeate the manifold in such a way that it is in all and over all one with itself; the universal must embrace and grasp the manifold in such a way that in the whole of multiplicity it seizes and grasps itself. But such a unity, grasping and comprehending all in itself and itself in all, is a \emph{self-unity}; the universal, as unity, is therefore self-unity, therefore selfhood, therefore subjectivity.

\opage{182} 2) As the infinite self-unity of existence, the universal must, in all diversities and finite oppositions, necessarily be the \emph{absolute identity}. It is only in a superficial sense that the universal can be said to be the common, the like in what is otherwise unlike, the general. Conceived as something merely common, the universal is only an empty schema, an abstractum; but an abstractum cannot be the universal; by its one-sidedness it is on the contrary only something particular.

3) The identity of all, the infinitely concrete identity, is therefore selfhood, subjectivity. But as the identity of all, the universal must, after all, be the same in all things, everywhere and in all ways, always the same; but this demand can be fulfilled only on the condition that the universal goes out from itself and through all else returns to itself, that it is, in other words, in \emph{infinite movement}, and that this movement is subjectivity's own movement of the idea.

But although Plato and Hegel, out of interest in knowledge and in what is objective in knowledge, both, so to speak, agree in overlooking the principle of knowledge, knowing subjectivity, their mutual divergence as regards the mystification of the universal, and thereby of selfhood, is nevertheless a necessary consequence of their mutual divergence in the determination of otherness. To determine how res depend on universalia is, in a logical sense, to determine how otherness originally relates to selfhood, for with otherness res come forth. What role, then, does otherness play in Plato?

``Otherness expresses that what is, in infinite self-separation, relates itself to itself by unceasingly falling away from itself. The other is then not something that is, but something infinitely changing, a mass unbounded in the multitude of its parts (ὁ ὄγκος αὐτῶν ἄπειρός ἐστι πλήθει). In this mass each smallest part may indeed seem by itself to be a one; but it is as in a dream-vision (ὥσπερ ὄναρ ἐν ὕπνῳ): suddenly there appears, instead of the apparent
\opage{183} One, in turn, the manifold, and instead of a smallest, something infinitely great in relation to the continued division. With the development of otherness in ``Parmenides'' the moment of sensibility is already intimated. This moment is particularly brought out in ``Philebus'', which distinguishes three kinds of being, namely: the limit (πέρας), the unlimited (ἄπειρον), and that which is compounded of both (τὸ ἐξ ἀμφοῖν τούτοιν τι συμμισγόμενον), to which is further added a fourth principle, namely the cause of the mixture (αἰτία); likewise in ``Timæus'', where a distinction is drawn between that which truly is, which does not become, and that which always becomes, which has no true being (τὸ ὂν ἀεί, γένεσιν δὲ οὐκ ἔχον, καὶ τὸ γιγνόμενον μὲν ἀεί, ὂν δὲ οὐδέποτε). That which always becomes is in Plato what we call matter, an invisible nurse of things (πάσης γενέσεως τιθήνη), which in an incomprehensible manner has come to partake of the rational (ἀνόρατον εἶδός τι καὶ ἄμορφον, πανδεχές, μεταλαμβάνον δὲ ἀπορώτατά πῃ τοῦ νοητοῦ), the mass (ἐκμαγεῖον) out of which all sense-phenomena are formed. From matter, the non-being, which can neither be sensed nor thought, but is apprehended only by a kind of spurious inference (μετ' ἀναισθησίας ἁπτὸν λογισμῷ τινι νόθῳ), things thus have their sensible constitution; from the ideas they have their fundamental form and reality; as imperfect copies of the ideas in the sensible, things are intermediate links between being and non-being.''\footnote{Phil.\ Propæd.\ p.~133--134.}

It is unmistakable from the above that Plato's doctrine of ideas has its antilogical element in otherness (ἕτερον δέ γέ ποί φαμεν τὸ ἕτερον εἶναι ἑτέρου, καὶ τὸ ἄλλο δὴ ἄλλο εἶναι ἄλλου); but it is at the same time evident that this antilogical element does not signify a co-determining power, but on the contrary a pure impotence. Certainly this impotence, according to the Platonic view, is not supposed to pertain to the ideas in \opage{184} their eternal being, but only to be characteristic of what is sensible, temporal, inconstant and perishable in things. A closer consideration, however, soon teaches us that the impotence which was supposed to fall exclusively upon things falls essentially upon the ideas themselves. Precisely because the ideas in themselves lack the element of otherness, the real negation that was to give them the impulse to activity and self-development, they are, as Aristotle has aptly remarked\footnote{Cf.\ Phil.\ Propæd.\ p.~137--138.}, without a principle of becoming and motion. The eternal ideas are in this respect not merely indifferent to temporal things, but in their ideal immediacy they are even indifferent to one another. In this bare immediacy each idea is an unconditioned for itself: likeness for itself, unlikeness for itself, motion for itself, rest for itself, etc.; thus there arises a multiplicity of substantial ones (ἑνάδες, μονάδες) which, although each of them is, with respect to the changing phenomena, a universal (τό ἐπὶ πᾶσι κοινόν) and has a being independent of sensible things (χωρὶς ἔστι), nevertheless mutually exclude one another, and cannot possibly be absorbed as moments into a single, absolute unity. Plato does indeed present the totality of the ideas in their highest unity as the Godhead, not as God; here, then, we should expect a higher illumination of the content of knowledge in the knower himself. But the naive manner in which the relation between God and the ideas is everywhere described shows clearly enough how the intuition of genius takes the place of comprehension, and what is really the reason why, wherever the talk is of the Idea of ideas, it always gains the upper hand over the philosophical, the aesthetic over the logical. It is said in the Republic: ``Just as the sun is the cause both of things being seen by its light and of their growing and thriving, so the Good (the Idea of ideas) is the cause for the soul of science and truth in everything that is an object of science. Just as the sun is not \opage{185} sight or the seen, but stands above both, so the Good is not truth or science, but stands above them both, since they are both akin to the Good.'' It is undeniably an apt image; but it is only an image.

Since the logical being of the ideas originally lacks that antilogical element of power which conditions energy, the categories ground and cause can be applied to them only in a figurative sense. Just as the element of otherness, the unlimited (ἄπειρον), is too impotent itself to grasp the ideas, so the ideas, which are to supply the limit, are for their part likewise too impotent to grasp the unlimited. The union of the limit and the unlimited therefore comes about in a rather external manner. For the union is regarded --- which expressly points to something external and contingent --- under the viewpoint of a mixture, and is referred to a fourth principle, a cause of the mixture, a cause, therefore, which arbitrarily hovers above its own effects. Although the ideas were supposed to constitute the essence of things and consequently to be one in essence with things, yet in their eternal being they have an essence different from things themselves; it is this loose connection of the ideas with things that Aristotle censures when he urges that the ideas of the Platonists, although they were supposed to be independent unities subsisting for themselves, nevertheless lack any content of their own. ``The ideal two and the empirical two have the same content. One speaks of the self-man, of the self-horse, of self-health (αὐτὸ ἄνθρωπόν φασιν εἶναι καὶ ἵππον καὶ ὑγίειαν); but by putting `self' in front of these nouns nothing is gained. The actual horse has essentially the same conceptual determinations as the self-horse; to define the actual man is the same as to define the self-man, etc.''

If we now turn to Hegelian idealism, the idea, absolute negativity in positive form, undeniably has enough of the negative; but this negative is, as we have already seen, \opage{186} the abstractly logical; real negation, the antilogical element of the concept, has entirely vanished. That the logical undermining of otherness, of singularity, of what is objective in objectivity, is not so easily noticed in the logic is explicable when one considers that the surrogate put in the place of antilogical otherness, namely immediacy, logical immediacy in its reciprocal determination by negation, can indeed in many ways present a deceptive imitation of the objective itself. Whence, pray, comes all matter? From immediacy! And the mechanical in the concept? From immediacy! And the opposition between nature and spirit? From immediacy! Immediacy has precisely the convenience that it can represent the sensible and yet be the thinkable, indicate the positive and yet transform itself into the negative, posit the finite, the external, and yet be sublated in the infinite and reflect itself as something internal.

With these presuppositions the Hegelian doctrine of ideas is naturally, in a logical respect, formed throughout quite differently from the Platonic. Instead of letting, with Plato, every concept, however abstract, come forward with unconditional independence as an idea for itself, Hegel makes the whole system of one-sided concepts into a system of moments for the concept itself; he does not assign to the individual logical ideal forms an eternal existence in rest and repose; on the contrary! he orders the categories in great circles from the most abstract of all to the most concrete of all, he awakens in each single category the slumbering logical impulse, he shows by his immanent method how the subsequent categories develop out of the preceding ones, as the higher out of the lower, the richer out of the poorer, and in this way transforms the whole logic into a reflection of the logical becoming of the idea. ``In the Platonic presentation the subjective prevails --- namely the psychologically subjective --- insofar as now this (εἰ ἓν ἔστι), now that (εἰ ἓν μὴ ἔστι) is arbitrarily assumed; so subjective a form cannot possibly be the adequate form of the idea. The \opage{187} Hegelian method, on the other hand, is objective --- within the abstractly logical --- and by its systematic unfolding of the moments of the content can furnish an objectively clear form of the idea. With his anatomizing dialectic of distinction Plato could not get the individual ideas to be absorbed as moments into the unity of the idea; with his immanent dialectic of identity Hegel can demonstrate the absolute unity of the idea in all particular ideas, since these are at once independent and yet moments.''\footnote{Phil.\ Propæd.\ p.~157.}

But with these formal advantages, great and unmistakable for the further development of philosophy, the Hegelian doctrine of ideas has in its essence a hollowness and inner emptiness which can only for a given time be veiled by a development of form that is in so many respects ingenious and grand in style. The fundamental damage lies in the logical, although it is least conspicuous in the logic itself, since, as has been said, it rests on an exclusion of the antilogical.

That Plato's ideas as universalia ante rem were not able to master the actual res is understandable from their logic and evident enough from Plato's physics. ``Natural things are a product of two principles: the totality of the ideas, God, who is the father of all things, and the passive principle, the mother or nurse of things, which receives shape from the formative principle. Nature is as it were a child (ἔκγονος) of these two.'' Of the ingenious, poetically beautiful and childlike manner in which the whole edifice of the world with all natural objects can, from this standpoint, be conceived as a copy of the realm of ideas --- how, in other words, idealistic physics can be transformed into a symbolism of nature --- Plato's Timæus furnishes a remarkable example. But if Platonic idealism is not able to impart to its universalia the capacity to master res, and that for the good reason that real negation, the antilogical itself, is recognized only as an image of impotence, then universalia in Hegelian idealism, however strongly the dialectic, \opage{188} if I may say so, whips the universal, can still less procure any reality for the actual res, since this idealism has cut away at the root the element on which the real itself in all reality rests, namely the antilogical. ``A philosopher who disdains natural science's express demonstration of the formation of springs and the activity of volcanoes, and with a remark such as this: `wie die Quellen die Lungen und Absonderungs-Gefäße für die Ausdünstung der Erde sind, so sind die Vulkane ihre \emph{Leber}, indem sie dieß Sich an ihnen selbst Erhitzen darstellen' thinks he hits upon the pure truth in the pure idea, has no eye for the causes of things and no sense for actuality; he who sees bitter earth in so speculative a light that he discovers in it `the subject of salt, an intermediate taste that has become a principle of fire', need pay just as little heed to what is written in chemistry as he who, in the metaphysical essence of granite, sees himself able to demonstrate `die irdische Dreieinigkeit, welche sich nun nach ihren verschiedenen Seiten entwickelt' ---, because it lies in the `granite principle' that this `Kern der Erde' must, with logical, dialectical necessity, consist of quartz, mica and feldspar --- needs to consult geognosy.''\footnote{Forelæsn.\ over ``Phil.\ Propæd.'' 1861--62 p.~71.}

Hegelian idealism, insofar as in the philosophy of nature it is compelled to engage with the sensibly concrete, naturally cannot escape the sense that here there is an antilogical element; but since Hegel has no place in the logic for its development, since the pure concept and the pure idea can become perfectly complete without dissolving this element, it is to that extent natural enough that this element must, in the realm of nature itself, be regarded as something impure, which can easily infect pure, speculative thinking. The disproportion between idea and nature is therefore described in Hegel's philosophy of nature in a manner which in a logical sense must arouse \opage{189} misgivings. ``Nature is in itself, in the idea, divine; but as it is, its being does not correspond to the concept; it is rather the \emph{unresolved contradiction}. Its peculiarity is posited being, the negative; as indeed the ancients on the whole regarded matter as a non-ens. Thus nature is also an expression of the idea's falling away from itself, since the idea, as this shape of externality, is in a state inadequate to and in conflict with its essence, etc.''

True idealism, conceived in all its aspects, must necessarily set out to show how the antilogical element, not as a moment of impotence but on the contrary as a moment of power, conditions knowledge and makes real comprehending possible. Setting up, then, as the phase requires, comprehension itself and the infinite self-development of comprehending subjectivity as the highest end, we consider this end first from the psychological, then from the ontological standpoint.

What the moment of power signifies, when the demand is for an ever richer development of a conceptual content infinite in itself, can be seen, from the standpoint of psychological subjectivity, from sense-perception and the sensible. There is in the sensible not merely an immediacy, but an antilogical immediacy; it is from this antilogical immediacy that the infinite wealth of determinations proceeds which alone can articulate thought and raise reflection to an actual comprehending. For what is it that prevents the many sensations from being reduced, in the name of the concept, to an insipid general determination, and this general determination in turn, since it is sensible, from being reduced to a logical reflex in the logical universal, like the \emph{Now} and \emph{Here} in the Hegelian phenomenology of spirit? Evidently the antilogical element! With sensation there is indeed given to us in actuality an infinite multiplicity of particular sensations; each such sensation has in itself the two-sidedness that, in the same instant in which it \opage{190} points to an antilogical immediacy that resists reflection, it precisely furnishes a reflex that determines and enriches reflection in a peculiar way. Think of sound and tones, think of light and colors, think of the point-for-point reciprocal determination of the aesthetic effects of colors and tones with the corresponding mathematical-physical functions; think of the swarm of rational determinations of knowledge contained in such mathematical-physical functions; think of the infinite satisfaction of the knower who, in developing all determinations of knowledge, develops himself, preserves unity with himself, and transforms the countless impressions of reality, determined by observations, into conceptual determinations in an ever-growing content of knowledge, and one will surely admit that the antilogical element of existence is decisive for a progressive development of concepts.

Should it be objected on behalf of idealism that pure comprehending must in this way be encumbered with so many elements of contingency, so many admixtures of sense-images and finite representations, that the idea as idea cannot possibly mirror itself in such coarse, muddied media, then one should consider well that the admixture of the contingent and the sensibly obscure by no means stems from there being an antilogical element here in comprehension, but rather from this antilogical element resting on a power that falls outside finite, psychological knowledge.

Seen from an ontological point of view, the moment of power is absorbed in an infinite manner into the moment of knowledge; the true, all-comprehending knowledge is a knowing power, an omnipotent knowledge. The original unity of knowledge and power in the form of knowledge here shows us comprehending subjectivity in a system of active general forms, overreaching universalia, which posit and presuppose res, since they everywhere posit and presuppose real media in which the light of the idea can refract its rays.

\opage{191} That res, the immediate existences in their singularities and infinite multiplicity, stand in power, will thus mean: 1) that things themselves become objects, not fleeting semblance, but actual objects for ideal knowledge; 2) that ideal knowledge essentially needs such actual objects, because the logical forms of comprehension must point for point be integrated by antilogical elements; 3) that comprehending subjectivity absolutely masters things, since the antilogical determination of power, which from the side of power is the reality of things, corresponds exactly to the logical determination of power, which, from the side of comprehension, is identical with ideality itself.

On these terms the phase is determined by the infinite energy of knowledge and of knowing subjectivity; the idea-movement goes from selfhood through otherness back to selfhood, from the subjective through the objective back to the subjective, from the universal through the singular back to the universal. But does this circular movement not then yield an empty circle, a fruitless repetition of the same and the same to infinity?

In abstract generality the idea-movement is indeed a repetition; but repetition is in the concept an infinite renewal, because comprehension itself is an infinite reciprocal determining of recollecting and producing, and ideal reciprocal determining is thus one with an infinite productivity. All that is new is old, ``there is nothing new under the sun!'' this is the idealess, trivial conception of the category of repetition, a conception that is always stricken with logical barrenness, since it is conceptless. In relation to such a conceptless repetition Lessing is right when he somewhere expresses himself roughly thus: ``if God held all truth in his right hand and the search for truth in his left, and bade me choose, I would choose the left.'' Were it possible to possess truth and the true in a dead being, a dogmatically closed form, then the idea would have to vanish, and the mocker would be \opage{192} right when he finds it tedious to be an omniscient one, since an omniscient one would never be able to come to know anything new. By virtue of ideal productivity, on the contrary, the old is constantly new; thus repetition is an essential moment in original comprehending, whose end point everywhere coincides with its new starting point\footnote{In his treatise ``Det astronomiske Aar'' J.\ L.\ Heiberg (Prosaiske Skr.\ vol.~9, p.~79 ff.) has, with particular regard to statements by S.\ Kierkegaard and Göthe, illuminated ``the dialectic of repetition'' in the following way. ``Now the complaint is heard that repetition is tedious because it constantly brings us the same, now the opposite complaint, that what it brings is fleeting and perishable, and that it must therefore continually put the new in place of the old. But this contradiction lies in repetition itself and constitutes its nature. What it brings must be perishable, for otherwise it could not be repeated; but it must also be imperishable, for otherwise the new could not be the same as the old. Whoever is inclined to see everything in a hypochondriac light can then, at nature's repetitions, both grieve that everything is changeable and be vexed that everything is unchangeable. But in this way one lets the changeable and the unchangeable fall apart in one's consideration, holds them separate as two independent properties unaffected by one another, and in that case one may well have reason to grieve over the one and be vexed at the other. But with a sound moral constitution one puts them in mutual connection, as opposite sides of the same concept, and now has equally valid reason to rejoice in and be enlivened by both \dots\ But what chiefly matters is that nature's objective repetition meets such a subjectivity in the apprehension as makes it its own and works it up into something new. For what nature bestows is like the talismans of the old fairy tales, which have wondrous powers but are morally indifferent, so that they become a blessing in the one's hand, in the other's a ruin.''}. For psychological subjectivity in its human limitation the conditions of knowledge are restricted; therefore \opage{193} ideal productivity can at this stage be so easily exhausted, and what should make the impression of novelty, that is, of originality, can look like a wearisome repetition of the old. For ontological subjectivity in its unity in principle with the divine idea the possibilities of productive knowledge are inexhaustible; here the past is essentially rejuvenated in the present; here the old is eternally a new --- under the sun.

\begin{center}{\bfseries β)\quad The originality of power: universalia post rem.}\end{center}
\addcontentsline{toc}{paragraph}{β) The originality of power: universalia post rem}

Only through singular determination can the universal be enriched to infinity; but singular determination is in its immediacy a determination of power: as many determinations of knowledge in the concept, so many determinations of power in the existence corresponding to the concept. If now the singular has, in real conceptual development, the same originality as the universal, then this contains a summons to consider the idea-movement from a standpoint opposed to idealism, to seek its beginning and end not in the universal but in the singular, not in knowledge but in power, and thereby to see how of itself it transforms itself into a movement of reality. The singular thus forms, in this connection, an opposition not to the composite, but to the universal.

The singular in its existential immediacy is only this singular, not the universal; and yet the singular, precisely this singular, is to be determined by the universal, apprehended as though it actually were the universal: here is the problem, the problem whose peculiar solution the nominalists of the Middle Ages designated by the formula universalia post rem. The nominalist standpoint is essentially psychological; here one begins with a finite separation of subject and object. Cognizing subjectivity stands as an existing individual in the midst of a world of existing objects; each object is a this-being, something individual and limited, like subjectivity itself. \opage{194} Cognition is conditioned by a plurality of singular objects being designated by a common name; this is a house, this and this likewise, etc. But although the many singular objects can be comprised under a common designation, yet the one object in its existential singularity has nothing in common with the other. It does indeed look as though there must be something common here, something universal that ran through the singulars; but that is an illusion: the common, the universal pertains only to the designation, not to the object, has only a nominal, not a real, only a subjective, not an objective validity. The existing objects --- res ipsæ --- are then, according to this view, not merely independent of and indifferent to, but also heterogeneous with, the universalia by which they are designated. That universalia are post rem means that our concepts of things are indeed abstracted from things, but that the thing itself is as indifferent to its concept as the body to its shadow. Far from belonging among realia, universalia themselves are only nomina, flatus vocis.

The contradictory element in the nominalist view nevertheless comes to light at once when the relation between the concept and its existence is considered a little more closely. The concept could not possibly be a common designation for certain determinate objects if there were nothing common in the objects themselves that could be held fast by the designation; the concept could not possibly be abstracted from the object if there were not something in the object itself that corresponded to the concept. With the recognition that the abstraction-concept must essentially correspond to the objects from which it is abstracted, that the universal which is brought out in the concept as what is common to the objects is also actually found in the objects, that singular existences, despite all their singularity, must necessarily fall under the universal, be determined and characterized by the universal, nominalism is transformed into empiricism.

\opage{195} Without entering upon any actual solution of nominalism's logical or metaphysical problem concerning the relation between universalia and res, the doctrine of experience proceeds straightforwardly from the presupposition that things must allow themselves to be determined by corresponding abstraction-concepts, provided, be it noted, that care is taken, by exact observations, that the concepts themselves are determined by things. The whole attention of cognizing subjectivity is thus directed to the peculiarity of things. The peculiar in things and the universal in concepts here furnish conditions for so rich a reciprocal determination of the subjective and the objective that a whole world of facts can be mirrored in a system of corresponding abstraction-concepts.

Empirical conceptual development does indeed begin theoretically, insofar as it begins by taking things as they are once given, and by directing its determinations of thought according to the determinations of existence that obtain; but the peculiar self-activity of subjectivity cannot fail to appear. It soon becomes apparent that the concepts which have been determined in a peculiar manner by things react determiningly upon things, and assert for the universal an ideal superiority over the singular. When empiricism has in this way developed into realism, the concepts have in reality become something quite other than subjective nominal values; they have become designations for substantial forms, for objective powers, for the unchangeable laws to which things are subject; the result is that the universal governs the singular. From the standpoint of realism, however, this cannot be understood as though the singular were to be transformed anew into a reflex of knowledge, a real moment in universal knowledge; on the contrary! realism lays the emphasis on res and thereby on the reality of the singular as singular. The movement of the concept is thus this: the singular comes, by means of the universal, to be determined in its singularity; by this determination of the singular in its singularity comprehending \opage{196} subjectivity gets the objective in an objective manner into its power: real cognition is a practical cognition.

``We learn,'' says the natural scientist, ``to master nature by our ever-growing insight into the connection of forces, by our discoveries of the laws of nature''; and it must be admitted that natural science has in the last century furnished a magnificent proof of the truth of the old proposition that knowledge is power. But the conceptual development taking place in natural science can in this connection also serve as proof of the impossibility of attaining complete insight into the relation between concept and existence so long as knowing subjectivity cannot give up its psychological standpoint.

If we note the manner in which empirical science on the whole works up its material and develops its real concepts, there presents itself here a natural distinction between the image, the description, the characterization and the concept proper. The image, the immediate designation of the singular in its singularity, represents the object, the real concept in its individual, antilogical form. Instead of producing the object itself, knowing subjectivity must, when we take the task from the theoretical side, thus restrict itself to imitating it by an image. As the images of several similar objects are taken up into the circle of representation and set in motion, they are worked upon by reflection and furnish elements for description. In the description of nature, which already requires a separation of the essential and the inessential, the transition to characterization is seen; characterization unites the presentation of the individual with the presentation of the universal in so peculiar a manner that comprehending proper can thereby be said to have been initiated. On what, then, does comprehending proper rest? Evidently on insight into the analogies and causal relations by which a unity of connection with its various ramifications can be determined.

\opage{197} That the concept as concept must in this way become poorer than the corresponding existence is self-evident. However far investigations are pursued, intellectual progress will always remain an approximation. And when the special sciences with their real investigations cannot get beyond mere approximation, then philosophy cannot do so either, as far as conceptual concretions are concerned. ``In its most finite reality the matter is an immediate existent, and this existent is, in relation to its concept, a specimen. In this existing stone we thus have a specimen of a stone, in this plant a specimen of a plant, in this animal a specimen of an animal, etc. In the matter as specimen we have the content of the concept, but not its form; in the matter as concept we have the form, but not the actual content. Between the matter as specimen and the matter as concept there now lies the example. In opposition to the specimen the example is the universal and to that extent akin to the concept; in opposition to the concept the example is the individual and to that extent akin to the specimen. If, then, there is to be a real unity of development of the matter and development of the concept, such a unity-in-difference must necessarily be conditioned by a development of examples.''\footnote{Forelæsn.\ over ``Phil.\ Propæd.'' 1861--62 p.~41.} But the development of examples is subject to an infinite approximation.

If we take the task from the practical side, a similar relativity will appear. What the construction-concept means, and how it is related to the antilogical judgment: \emph{the universal is the singular}, has already been treated in the foregoing. Here there is still to be added an indication of the intermediate determinations of relativity. If the abstraction-concept is imperfect, the construction-concept must of course be so too. The material to be used must indeed be found as something given; this given then has in its existential concretion a \opage{198} multiplicity of essential properties, of which only a few are brought out and held fast in corresponding conceptual determinations. The conditions and the more particular circumstances under which the production takes place are here subject to so many variations, alterations and changes that an infinite number of experiments would be needed before the properties of a single substance, even in a quite special circle of reciprocal actions, e.g.\ the chemical, could be fully brought to light. The production determined by the natural concept, e.g.\ the production of water, does indeed correspond to the concept itself, insofar as regard is here had merely to the law of existence that is fulfilled by the production under given conditions; but what a distance there is between the insight sufficient to bring oxygen and hydrogen to combine into water and the insight that would be required to see through the disposition either of oxygen or of hydrogen to chemical combinations with all possible substances under all possible conditions! The consideration of a single, limited example alone can almost make the understanding that strives after strict concepts giddy. ``With the number of the elements, with the number of the atoms that unite into a group,'' says Liebig, ``the directions of attraction are multiplied; for which reason the strength of attraction also decreases in the same degree as the multiplicity of the directions increases. A particle of common salt, one of the smallest particles of cinnabar, presents a group of no more than two atoms; a sugar atom, on the other hand, contains thirty-six, the smallest particle of olive oil several hundred single atoms \dots\ Without anything being added or taken away, we can imagine the thirty-six single atoms of the sugar atom arranged in a thousand different ways; with every change in the position of a single one of these atoms the compound atom ceases to be a sugar atom, etc.''

The concepts of natural science are concepts of facticity to such a degree that science at bottom prefers to avoid every attempt \opage{199} to explain to itself the inner how of the activities, content when it can merely demonstrate the factually valid laws; a proof that the general determinations themselves are nevertheless essentially apprehended as given facts.

If idealism is on the whole inclined to confuse its abstract knowledge of the divine with divine knowledge itself, then realism, with its widely ramified masses of information and its manifold knowledge of the finite, is inclined to deny the possibility of a true cognition of the idea. But both one-sidednesses will vanish as the relation between psychological and ontological subjectivity is more closely illuminated. It is the absorption of knowledge into power on which, in the present connection, everything depends. Just as the perfect knowledge-concept, in the first phase of comprehension, presupposed the perfect power-concept, so now the power-concept here, in the second phase of comprehension, presupposes its perfect knowledge-concept.

If, however, the power-concept were one with its knowledge-concept in such a way that in this unity there were no difference, then it would be a rather indifferent formality whether we preferred the designation that it is power that is absorbed into knowledge, or the converse, that it is knowledge that is absorbed into power. But here too it holds that the totalized unity consists in a totalized opposition. The infinite approximation by which the real concept, from the standpoint of human knowledge, can be perfected to ever greater agreement with its specimens by no means points to a completed comprehending that would presuppose what Hegelian logic so falsely presupposes, that the adequate concept immediately coincided with its existence. On the contrary! The gradual development of real concepts proves precisely that, e.g., the concept horse could be perfectly adequate in all determinations to the existence horse, and the opposition of concept and existence \opage{200} nonetheless persist, insofar as the concept as concept is logical, existence as existence antilogical.

In perfect comprehending, in the idea itself, the knowledge-concept, the concept as concept, thus forms a thoroughgoing opposition to the power-concept, the concept as actual existence; and yet the knowledge-concept is in the idea one with the power-concept. The unity-in-opposition must then develop itself in a movement. In the first phase of comprehension, where the antilogical became a moment for the logical, the movement was a pure idea-movement; here in the second phase of comprehension, where the logical becomes a moment for the antilogical, the idea-movement must transform itself into a real movement.

What confusion arises when the idea-movement and the real movement, following Hegel's immanent method, are without further ado lumped together under the name of ``movement of the matter'' has already been shown from more than one side: ``The mystification with the `movement of the matter', `the objective self-movement of the concept', stems in particular by no means, as so many have assumed, from the method's being dialectical, but rather from the methodical dialectic's being carried through only one-sidedly. The dialectic of concept and concept the method knows excellently, but the dialectic of the concept of being and that which is it knows only poorly: instead of developing the concept of the thing, the concept that is absorbed into the thing, the method everywhere develops the concept of the concept, the concept of the thing's concept, the concept in which the thing itself is not. A couple of examples will illuminate the matter better than lengthy explanations. For being to be, something that is must be; for something that is to be, there must be a this-being. Elements such as oxygen, sulfur, potassium, etc.\ are each something that is, in that they are, each by itself, a this-being. From the expression \emph{this-being} it is seen that that which is can itself only be pointed out, not stated. But what cannot be stated cannot, of course, be thought either. If now the unity of thought being and that which is itself
\opage{201} is nonetheless justified by the fact that the unthinkable becomes a negative element, a limit of thought, which everywhere belongs to determinate thinking, then it surely follows from this that the unity of thinking and being must be just as dialectical as the unity of being and that which is.''

``But if the dialectic of the concept of being and of that which is itself is different from the dialectic of the concepts of being among themselves, of being and non-being, of becoming and determinate being, then it surely also follows from this that the movement of thought must be different from the movement of being. The \emph{becoming} of which logic treats, when being unites with nothing so that determinate being becomes, is heterogeneous with the becoming of which chemistry treats when sulfur unites with oxygen so that sulfuric acid becomes, potassium with oxygen so that potash becomes, and sulfuric acid in turn with potash so that sulfate of potash becomes.''\footnote{Forelæsn.\ over ``Phil.\ Propæd.'' 1861--62 p.~55--56.}

By making clear the opposition between concept and existence as a peculiar opposition between the knowledge-concept and the power-concept, we learn to see what, from a realistic standpoint, necessarily belongs to a perfect comprehending. It is, in the first place, 1) the thing with its properties, next 2) the external and contingent in the nevertheless essential connection of things, and finally 3) the opposition between what is and what becomes, what rests and what is moved, on which attention must here be fixed.

1) As regards the thing with its properties, it must of course, insofar as it is seen from the point of view of the concept, be apprehended as a totality of determinations of existence. That this totality, when it expresses the power-concept, must not only present itself in a different light but also be produced in a quite different way than when it is to express the knowledge-concept, will be seen on closer reflection. \opage{202} If we set the concept of silver over against existing silver, then the identity of this opposition, if the comprehension is to be perfect, must necessarily be had in two forms: in the form of the concept as knowledge-concept, in the form of existence as power-concept. If we now look at the knowledge-concept, then mineralogy, by its presentation of the properties of silver, surely offers us contributions enough, at least for a thought-experiment that might indicate how the concept of silver would have to be developed if it were to satisfy the demands of knowledge in every respect. Whether we think of the whiteness and the luster, or of the ring, or of the form of crystallization, or of the weak attraction to oxygen, the strong attraction to chlorine, or of whichever of the properties of silver we may care to name then, from the standpoint of knowledge,
% Errata: (Trykfeil og Rettelser, p.~202 line~13 from top): ``nævne'' [name] should be corrected to ``nævne:'' --- here rendered as printed (without colon).
insight into each of these properties will be conditioned by the singular's being taken up as a moment in the determination of the universal. The luster of silver presupposes, in order to become a determination of the concept, a knowledge of what metallic luster is; but a knowledge of what metallic luster is acquires significance for comprehension only on the presupposition that there is already here a knowledge of what luster on the whole essentially is, hence a knowledge that could make the reciprocal action between light and bodies perfectly transparent, a knowledge that would only be complete when optics was complete. If we next take up affinity, then the attraction of silver to chlorine must necessarily reflect itself as a special determination of the concept in the determination of its attraction to sulfur, to bromine, to iodine, to fluorine, cyanogen etc. The determination of the individual attractions will thus in the concept coincide with a determination of the chemical attraction of silver in its concrete generality; but this concrete generality will in turn be a particular moment of universal chemical attraction in its thoroughgoing articulation, its many-sided concretion. A perfect insight into the attraction of silver to chlorine would presuppose a perfect insight \opage{203} into the whole system of chemical attractions, hence a perfect chemical knowledge.

Perfect knowledge thus proceeds --- for it is the law of knowledge --- as regards the concepts of existence, \emph{from the singular to the universal, through the singular to the concretely universal}, just as well as imperfect knowledge does. But this movement, which is so decisive for the knowledge-concept, does not fit the power-concept at all. In that knowledge relies on power, it presupposes existence; its logical determinations are together reflexes of the antilogical determinations of existence; the individualized universal concepts are, in this circle, universalia post rem. It is otherwise with the power-concept. Insofar as the determinations of power are originally identical with the determinations of knowledge, the power-concept has in the system of these determinations of knowledge a hidden knowledge, a logical content, and may well, with respect to its ideal origin, be said to begin with a system of universalia ante rem. But the character of the power-concept is, as we have already seen, recognizable precisely by the fact that it transforms the logical judgment: \emph{the singular is the universal}, into the antilogical one: \emph{the universal is the singular, indeed this singular}. Whereas the concept in the form of knowledge begins by presupposing the existing object as a subject of the judgment, and then develops by developing the system of all the predicates, the power-concept conversely begins by presupposing a system of predicates, and then comes into force by letting these predicates concentrate themselves in their subject as in an existing object. In the existing object the predicates are transformed into determinations of existence, into properties of the thing, into attributes, and into accidents of the substance. These accidents, these properties and attributes can now, in the sphere of the power-concepts, no longer be transformed into predicates, into logical universal determinations; their character is antilogical to such a degree that every attempt to apprehend them as universalia must lose itself in absurdity. As \opage{204} power-concept, silver exists in a system, not of logical predicates, but of actual properties, individual attributes, factual accidents, determinations of existence so antilogical that each one of the properties is just as immediately one with, as it is immediately different from, all the others. Here one cannot now think of the luster of silver in all generality, nor of metallic luster in particular generality, nor of silver luster in individual generality; no, here one cannot think at all of luster in any kind of generality whatsoever, but only of this luster of this existing silver. Existing silver with this determinate luster has, not a general attraction to chlorine in all generality, but under \emph{these} conditions \emph{this} attraction etc. In existing silver the determinate luster and the chemically determinate attraction have their factual identity; but this identity is so antilogical that no deity, still less any human being, shall undertake to deduce \emph{this} luster from \emph{this} attraction, or \emph{this} attraction from \emph{this} luster\footnote{That here, from the standpoint of science, there can of course be talk of --- at least to a certain degree --- explaining the one property by the other cannot come into consideration so long as what is aimed at here is exclusively a determination of the character of the power-concept. By a change of standpoint the properties are transformed back into predicates, and reflection appropriates a system of universalia post rem.}. To comprehend is then, in the present connection, not restricted to thinking or reflecting; to comprehend is: in hidden unity and factual understanding with an infinite thinking, an infinitely determinate reflecting, \emph{to objectify, to grasp, to master}.

2) From the relation between the knowledge-concept and the power-concept, the reciprocal determination of the external and internal, contingent and necessary relations of things is now explained in a very natural way. As existing objects, things come forward \opage{205} against one another in such a way that on a superficial view it looks as though the one thing did not concern the other at all; every thing is what it is; here the logical once more runs up against the antilogical. But it must constantly be emphasized that things as power-concepts are nevertheless at bottom concepts; that the content of the power-concepts is logical, although the form is antilogical; that the singular has surely not torn itself loose from the universal, but on the contrary presented itself as the universal's own individual determinateness, a determinateness without which the universal itself could in no way obtain reality in the objective: in spite of the sublated connection of things there is thus here nevertheless an actual connection. This connection of things now presents a double side, an antilogical and a logical one. Seen from the antilogical side the connection is \emph{external}; in the external connection with its changing circumstances contingency reigns: contingency is precisely the antilogical element of individual determinateness and of actual connection. From the logical side the connection is \emph{internal}; in the internal connection with its essential conditions necessity reigns, with necessity the laws, and with the laws the system of universalia that in the world of facts is cognized as objective powers. In the reality of the laws power has its ideality; therefore the laws are essential intermediate determinations between the knowledge-concept and the power-concept. Only in the form of knowledge is the law a law, for the law is the universal; only in the form of power is the law a law, for the validity of the law is its fulfilment, the fulfilment is conditioned by existence, existence by a comprehending power.

3) The opposition between the real movement and the pure idea-movement can then no longer be doubtful either. As the perfect knowledge-concept, the concept has ideal determinate being in a mighty reflecting; the mighty reflecting, comprehending all concepts, is a movement identical with rest, a movement in the pure idea. From the standpoint of human knowledge \opage{206} such a movement can be seen in a systematic development of the ontological categories, although such a development is of course infinitely far from being congruent with the conceptual movement in which the idea itself develops and determines its logical content. As a peculiar power-concept the concept is a thing in the series of things, and as an existing thing expressly determined for a finite opposition of rest and motion. To illuminate this finite opposition we must not restrict ourselves to considering the mechanical alone: every natural process is an example significant in its kind. The movement of reality presupposes space and time; what in the form of knowledge is united in pure reflection is laid out in the circle of power, by means of the forms of externality, through systems of actual transitions. For comprehension the task is to be able to cognize in these transitions the existential expressions of an at bottom essential and infinite reflecting.

``The conflict of existence with the ground is a contradiction which finds its resolution by the help of the forms of externality: space and time. If we ask what water as matter essentially is, then water is an essence reflected in the ground; for water is, as a product of oxygen and hydrogen, reflected in its constituents, and these constituents conversely reflected in their product. The product thus reflected is moreover --- in the ground --- burdened with the contradiction that it cannot exist insofar as its constituents have determinate being, and the constituents cannot exist insofar as the product has determinate being: product and constituents mutually presuppose and sublate one another. Furthermore, this product reflected in the ground is burdened with the contradiction that its form of aggregation is at once gaseous, liquid in drops, and solid, while each of these forms nevertheless mutually excludes the others \dots\ It is impossible to say what most characterizes the essence of water, whether its formation from, or its separation into, oxygen and hydrogen; for in the ground the separation must surely be reflected in the formation, and the formation conversely in the separation: separation and formation are, with their difference, in \opage{207} the ground one. It is otherwise with existence. In immediate existence one and the same product cannot possibly be at once both formed and dissolved; but what cannot happen at once can surely happen successively, therefore we have time; and what cannot happen at the same time in one place can surely happen simultaneously in different places, therefore we have space.''\footnote{Forelæsn.\ over ``Phil.\ Propæd.'' 1860--61 p.~173--74.} It is the power-concepts that fill up time in that they fill space.

\begin{center}{\bfseries γ)\quad The original unity of power and knowledge: universalia in re.}\end{center}
\addcontentsline{toc}{paragraph}{γ) The original unity of power and knowledge: universalia in re}

In its opposition to power, knowledge begins with the universal; the formula universalia ante rem is to that extent significant for the originality of knowledge. In opposition to knowledge, power begins by positing and presupposing the singular; the formula: universalia post rem, is in this respect significant for the originality of power. A look back at the two phases will, however, soon convince us that the reciprocal determination of the universal and the singular cannot simply be stated by scholastic formulas. If we consider the idea-movement that holds for knowledge, then it is evident that the universal in principle presupposes the singular, that in its unceasing self-determining it must constantly have the singular as middle term in the doublings in which it has its infinite wealth, that the universal, in other words, must just as well result from the singular as precede the singular, that thus universalia ante rem must originally presuppose universalia post rem. If we consider the real movement that holds for power, then the universal can surely only be abstracted from the singular on the presupposition that the universal itself is posited in and with this singular; in order to posit the singular as the existential expression of the universal, the power-concept must presuppose the knowledge-concept: \opage{208} universalia post rem is only a confirmation of universalia ante rem. But when the extremes: ante rem and post rem thus unceasingly presuppose and produce one another, then it is natural to seek the expression for the true in the uniting middle, hence to choose the designation universalia in re.

In a formal respect the designation universalia in re certainly also has its significance; but what is characteristic is that the logical significance, in this last designation just as well as in the two preceding ones, is precisely in its concept the opposite of what it is in the scholastic meaning. According to the scholastic meaning, the proposition: universalia sunt in re will evidently be understood in such a way that thought is to rest content with the being that universalia have in res, and hence to apprehend the unity of universalia and res not merely as a unity of being but as a unity that is. But here the misunderstanding shows itself again. The proposition: universalia sunt in re by no means belongs to the restful propositions, but on the contrary to the restless ones, haunted by all sorts of hidden contradictions; the designation in re surely posits at once the identity and the difference of universalia and res. What, pray, is res? The system of determinations of existence that are united in the object, the existing singular. And universalia? The system of predicates that in the judgment correspond to the determinations of existence holding for the subject. But now it has surely been shown in what precedes that the determinations of existence can no more be predicates than the predicates can be determinations of existence, since the predicates are logical, the determinations of existence antilogical; wherein, then, it lies: that universalia, insofar as they are in re, cannot be universalia, and, insofar as they are universalia, cannot be in re. The reciprocal determination of universalia and res thus demands a movement anew.

\emph{In the eternal reason knowledge and power are originally one.} It is remarkable how naturally and easily one gets scientific consciousness, especially if one paraphrases \opage{209} the determinations of knowledge as ``determinations of thought,'' to grant this proposition, even from the most different standpoints. ``The acts of nature,'' says H.\ C.\ Ørsted,\footnote{Naturlærens mechaniske Deel. Kbhvn.\ 1844. First Introduction.} ``conform to the acts of thought. The laws of nature can therefore also be called thoughts of nature. All the laws of nature together therefore make up a whole, just like all the thoughts in reason. Nature must thus be apprehended by us as a rational whole. The multiplicity in nature is therefore to be regarded as an unfolding of thought, its changes as a movement of thought. But just as in all thinking there is an activity, and the thought conceived without this is only a form of thinking, so in every existing thing there is an activity by which it asserts its existence and manifests it against other things. Without this activity it has no actuality. In thinking we call this active element will, in external objects force.\footnote{The categorial relation between force and power is in the present section only fleetingly touched upon, because the analysis of forces, insofar as it is given in Grundideernes Logik, belongs essentially under the circles of power. That power and force are related can be seen from linguistic usage. Heiberg has in his ``speculative logic'' (Anfr.\ Skrft.\ p.~72) stated the relation thus: ``Power is absolute necessity. Its moments are force and manifestation, of which the one is the limit of the other, and to that extent its quality. Force is the possibility of manifestation, and manifestation is the impossibility of force, i.e.\ the cessation of force, for in manifestation force is exhausted and ceases, as in its limit. But force, considered as possibility, is the \emph{condition} for manifestation, and manifestation is the \emph{existence} that comes forth where the condition obtains; whereas \emph{necessity} is the activity which, by sublating the condition, gives the conditioned existence, and (since the condition or the presupposition is sublated) leads what exists back to immediate being, not, as the highest logical principles do, by annihilating reflection, but on the contrary by completing it, so that it comes to its end and is sublated. The necessary is indeed necessary by means of an other, which is its condition, and this condition is itself existing or independent; but since the necessary does not itself come into existence before the existence of the condition is sublated, the necessary is just as well unconditioned, i.e.\ absolute, or: it is because it is.''} But just as willing and thinking \opage{210} are one, only considered from two sides, so also force and law of nature; they are manifestations of the divine will and the divine thinking.''

But when, in a logical sense, one is to give an account of the original unity of knowledge and power, and thereby of the possibility that ``the acts of nature'' can ``conform to the acts of thought,'' then a great problem shows itself, the problem with whose solution the logic of subjectivity must culminate. An attempt to answer the question what it really means that ``the acts of nature'' conform to ``the acts of thought'' will already make the problem recognizable. If we choose as an example of an act of nature the fall on the inclined plane, then it is in the first place evident that the acts of thought required in the observer in order to understand such an act of nature by no means coincide with an apprehension that immediately and directly conforms to the act of nature itself. Whole systems of logical determinations are united into antilogical unity in the existing inclined plane, in the falling body, and in consequence of this in the movement of falling itself; these systems united in the act of nature must be separated by distinct acts of thought and combined in a peculiar way if the act is to be seen to express determinate thoughts of nature.

``A body,'' it is said,\footnote{H.\ C.\ Ørsted. Anfr.\ Skr.\ p.~202.} ``needs as many times more time to fall down the length of the inclined plane than to fall through its height as the former is greater than the latter.'' Hereby, then, a law of nature is expressed as a thought of nature; but to what degree the acts of nature in the phenomenal movements by which the law is fulfilled deviate from the acts of thought whereby \opage{211} the fulfilment of the law is seen, shines forth when one casts a glance at the mathematical expressions\footnote{``If the height of the inclined plane is denoted by $a$, its length by $l$, then \dots\ the fundamental velocity for the fall on the inclined plane will be $g\frac{a}{l}$, and the motion is determined by the expressions
\[ h = \frac{gat}{l} \qquad r = \frac{ga}{2l}t^{2} \]
The duration of the fall on the inclined plane is found from the second expression by setting $r = l$, and if the corresponding time is denoted by $t_1$, one gets
\[ t_1 = l\sqrt{\frac{2}{ga}} \]
The time $t_2$ that a body would take to fall through the height of the inclined plane is'' --- here reference is made to something preceding --- ``determined by
\[ t_2 = \sqrt{\frac{2a}{g}} \]
and from this expression in combination with the preceding one we find
\[ \frac{t_1}{t_2} = \frac{l}{a} \]
or, in other words, the fall on the inclined plane requires longer time in the same ratio as the length of the inclined plane is greater.'' H.\ C.\ Ørsted, Naturlærens mechaniske Deel 3rd edition by C.\ Holten. 1859 p.~187.} that determine the fulfilment of the law from its various main sides. Certainly in all these main points ``the thought of nature'' meets the mathematical thought; but the act of nature is for that reason no more a mathematical operation than the mathematical operation is an act of nature; one does not fall in an equation, not in a quotient, not in a square root.

For the act of nature it holds that universalia are in re; of the act of thought, on the other hand, it holds that universalia, post rem as well as ante rem, are in mente. Every experiment presupposes a system of universalia ante rem; for the experiment acquires its significance precisely by the fact that in \opage{212} the given case a determinate question is put to nature for an answer, and the questioner must necessarily have a sum of already developed concepts in mente. Likewise every experiment presupposes that a system of universalia post rem can be developed; for the experiment is surely undertaken precisely with the aim of getting the concepts of existence translated into the form of the concept, and of keeping the universalia abstracted from the facts in mente. The acts of thought required for undertaking the experiment, and the sum of acts of nature in which the experiment is fully expressed, are, be it noted, as acts, not running parallel. If the act of thought could directly and in all immediacy follow the act of nature, then reflection would have to coincide with sense-perception. Nothing else, then, would be needed for understanding a natural event than that the observer should from all sides have seen what was here to be seen, heard what was here to be heard, in short: sensed what was here to be sensed; but just as little as sense-perception can directly follow thinking, so that the movement of sense should correspond point for point to the movement of thought, just as little can the act of nature directly follow the act of thought.

From the difference of character here demonstrated it follows, then, with irrefutable consistency, that act of thought and act of nature can just as impossibly be one in ``the eternal reason'' as in ``human cognition.'' It must certainly be granted that the eternal reason does not need, in order to cognize the fall on the inclined plane, such a cutting up and piecing together of preliminary knowledge, preparations, and deliberations as the human one does; one cannot possibly imagine the eternal reason busy adding, multiplying, extracting the square root, impressing upon itself formulas for the velocity, the space traversed, the time etc. But it nonetheless stands fast that the determinations of ideality which at the psychological standpoint are had in an imperfect way in mente humana must, \opage{213} at the ontological standpoint, as surely as there is an eternal knowledge, be had in a perfect way in mente divina.

That natural science as a special science cannot find itself called upon to undertake logical investigations concerning the peculiarity of the act of thought in opposition to the act of nature is understandable enough; but it is remarkable that the problem seems hitherto to have escaped even the attention of philosophy. ``The immanent method'' is in this respect so far from suspecting any difficulty that, on the contrary, as we have seen, in all dogmatic immediacy it proceeds from the presupposition that the movement of thought and the movement of the matter must everywhere coincide. ``The syllogism\footnote{J.\ L.\ Heiberg. Specul.\ Log.\ (Anfr.\ Skr.\ p.~93 ff.) (Prosaiske Skrifter, vol.~1, p.~303--304).} must consequently, no more than the judgment in general, be regarded merely as a human operation of thought; on the contrary, the syllogism must be regarded as the truth of all that exists, as the reason in all that is there; and everything is thus a syllogism. Or: everything is concept, and its existence is the difference between the moments of the concept, which difference consists in this, that the singularity of every thing stands under a particularity of properties or being outside itself, whereby it raises itself to its universal nature and takes this up into itself. Thus, e.g., the whole of human life is a syllogism. The individual is the singular. Its needs, its activity, its property, in a word everything by which it stands in connection with the external world, is the middle term between itself and the universal, or \emph{means} (i.e.\ middle terms, media) for ennobling and cultivating it, i.e.\ for uniting its individuality with what is universally valid. Regarded as human operations of thought, judgments and syllogisms are the expression of the subjective cognition of that fundamental reason in things, or the congruence of thought with actuality.'' One would thus, according to this statement, assume that the rational movement of thought \opage{214} was congruent with the movements in actual events; that there was thus in actuality an actual universal and an actual particular, just as well as an actual singular.

The ideal fusion of the subjective with the objective, which we already emphasized in setting up the problem of subjectivity, and have since sought to illuminate more closely from different sides and under different turns, must then in the present connection be subjected once more to a special test. It is the relation between the singular, the particular, and the universal on which attention is fixed. How, then, according to the speculatively immanent method, are the singular, the particular, and the universal determined by the syllogism?

``In the syllogism,'' says Heiberg,\footnote{Anfr.\ Skr.\ p.~93. (Pros.\ Skr.\ 1, p.~304).} ``the singular, the particular, and the universal are determined as \emph{individual}, \emph{species}, and \emph{genus} \dots\ If one denotes the singular by E, and the universal by A, then the judgment can be represented by the expression E --- A, which means that the singular and the universal are united, or that the singular is the universal. If now further the particular is denoted by S, then the complete judgment or the syllogism is represented by the expression E --- S --- A, which means that the particular is the middle term whereby the singular is connected with the universal, or that the singular is the universal by means of the particular \dots\ But\footnote{Anfr.\ Skr.\ p.~97--98. (Pros.\ Skr.\ 1, p.~313--314).} since the singular in that first proposition is determined as the universal, consequently as that which is common to both the other terms, it itself becomes the connecting middle, which gives the second syllogism, S --- E --- A, whose meaning is that the particular is the universal by means of the singular, since this is itself both particular and \opage{215} universal, according to the first proposition. --- Since further, according to the first proposition, the singular is determined as the universal, and according to the second the particular is likewise determined as the universal, the universal is the link that connects the singular with the particular, which gives the third syllogism, E --- A --- S, whose meaning is that the singular is the particular by means of the universal. And since now finally, according to this last syllogism, the singular is the particular, and according to the preceding one the particular is the universal, the particular again becomes the link, and produces the first syllogism, E --- S --- A, which had at first appeared without a preceding syllogism. Hereby the system of syllogisms is completed, since each of the judgments is now mediated.''\footnote{As an example of what is rational in the fact that everything ``is a system of three syllogisms,'' the following remark is appended (Anfr.\ Skr.\ p.~99): ``The individual, which by its sensuous nature connects itself with reason, stands under the form of the first syllogism, E --- S --- A. But the individual is itself the union of sensuous nature and reason; this gives the second syllogism, S --- E --- A. And finally it is reason in which both the singularity of the individual and its sensuous nature are grounded; this gives the third syllogism, E --- A --- S. Many disagreements and confusions in science have their ground in the fact that one keeps to the single syllogism instead of to the whole system. Hereby different, indeed opposed, assertions come about, of which each is correct insofar as it is one side of the truth, but incorrect insofar as its truth is one-sided.''}

If it were now a matter of carrying out the formal side of this systematics further, then it would here be an easy matter to show how the syllogism E --- A --- S essentially corresponds to universalia ante rem, the syllogism S --- E --- A to universalia post rem, and the syllogism E --- S --- A to universalia in re; but the formal side is already so predominant that \opage{216} the systematics, even to a doubtful degree, approaches schematism. If ``the system of the three syllogisms'' is to acquire any true significance in the determination of the relation between the act of thought and the act of nature, in other words, in the determination of the relation between the idea and actuality, then each of the system's three syllogisms must in turn furnish the type for a sphere of systems in which the peculiar movement can be seen to be determined by the nature of the term whose development the movement serves. If, then, the development and articulation of the universal is set as the task, according to the schema: E --- A --- S, we get an ideal system; if conversely the development and articulation of the singular is set as the task, according to the schema: S --- E --- A, we get a real system; if finally the reciprocal determination of the universal and the singular, or the development of the particular, is set as the task according to the schema: E --- S --- A, then it will be the teleological unity of both one-sided systems that articulates itself in a complete system.

1) \emph{The ideal system}. The singular is the universal, E --- A, this judgment is logical, hence expressly an act of thought. The essential difference between subject and predicate rests, according to what we have said before, not on any difference in the completeness of the determinations, for the same totality that is found in the subject can also be found in the predicate; but the opposition, the decisive and thoroughgoing opposition between the subject and the predicate, rests on this, that the determinations of the predicate are logical, those of the subject antilogical. In its logical form the totality of the predicate will, as the universal, express unity, while the totality of the subject in its existential, antilogical form will be able to repeat itself to infinity, and thus express multiplicity. On the side of multiplicity as well as on the side of unity we here have the absolute itself; namely, on the side of unity, hence in the form of the universal, the absolute-absolute; on the side of multiplicity, \opage{217} hence in the individual form of the singular, the relative-absolute: under both forms the concept, insofar as it is a knowledge-concept, is had in its absolute content as idea.

The singular in its existential immediacy the divine knowledge --- unless it should be able to see the square as a circle --- can no more see through immediately than human knowledge can: what exists is and remains the object of knowledge, and as object of knowledge a barrier for knowledge. On the other hand the divine knowledge has the infinite advantage over the merely human that the knower sees the existing singular in its adequate idea. The individual idea, the adequate idea of the singular, thus repeats itself in multiplicity with the singular; thus the judgment E --- A acquires the significance, decisive for the doctrine of ideas, that the same totality, the same absolute, the same idea, can at once be seen in the light of unity and of multiplicity, of absoluteness and of relativity. In one and the same idea knowing subjectivity can thus have countless determinate ideas, and in the countless ideas one single determinate idea.

It will, as a consequence of this, then also be no hindrance to perfect comprehending that here, instead of uniform repetition of the same singular, there enters an existential difference of totality between the singular existences among themselves, so that the singulars divide into two or more series. On the contrary! The difference of totality is surely precisely a peculiar expression of the fact that power serves knowledge, in that the knowledge-concept is articulated by the varying power-concept. With the difference of the totalities of existence the particular comes forth. But just as the particular in its conceptual unity always expresses doubling, so that where there is one particular there are also two, so the origin of the particular shows itself at once double-sided: it is the particularity in the singular that makes the universal into a particular, which can here be denoted by the judgment: E --- S. \opage{218} But since the power-concept, as concept, can only transform the form of the knowledge-concept, not its content, particularity cannot possibly come forth in the singular except on the presupposition that it is already in the universal; which can here be denoted by the judgment: S --- A. The unity of the universal idea with the individual ideas is thereby articulated into a system of particular ideas; in each particular idea the whole system is transparent.

In the sphere of ideas particularity is not merely a middle term between the universal and the singular, but at the same time the absolute unity of both. Here we have the syllogism: E --- S --- A, and Heiberg is to that extent right when he lays weight on the fact that the singular, the particular, and the universal are determined in the syllogism as \emph{individual}, \emph{species}, and \emph{genus}. When the particular is species, then the universal is genus; but in this connection what matters above all is to determine the movement as act of thought, and to present \emph{genus} and \emph{species} not as merely formal categories but as ideas. It is the universal itself at whose articulation the ideal system aims from all sides. Whether it is otherwise a question here of the individual or of the species or of the genus, it is nonetheless nothing but determinations of knowledge, or alterations, formations, and enrichments of the universal, that are concerned: hereby, and only hereby, is the system of movements peculiarly characterized as a system of acts of thought.

From these acts of thought it is expressly seen that the totality of determinations which in an existential way corresponds to the individuals in a certain series furnishes precisely the antilogical moments whereby the category \emph{species} is integrated into an expression of unity for a relative absolute, and becomes a particular idea. Further it will be evident from these acts of thought that the particular determinations of totality whereby the species belonging under the same genus differ from one another among themselves do not fall outside of, but coincide with, the system of determinations that characterizes the genus as genus. As long as the \emph{genus}
\opage{219} is taken merely as a category, it is naturally poorer in determinations than the species that fall under it; but as the absolute expression of unity for the different peculiarities distributed among the species, the genus is idea. ``Here a generic concept is not formed because different species occur here; rather, different species occur here because the idea of the genus comes to existence only by differentiating itself into species-ideas, just as the species-idea comes to existence only by differentiating itself into individual ideas.''\footnote{Forelæsn.\ ov.\ Phil.\ Propæd.\ 1860--61 p.~324--25.}

Finally, from the acts of thought decisive for the ideal system it will be possible to understand that the universal which in a determinate connection is articulated and determined as \emph{genus} in opposition to \emph{species} is eo ipso determined as \emph{species} in opposition to the \emph{family}, which in the logical sense is only a genus of higher order, etc., so that in this way the system can be articulated into systems of systems ad infinitum. The movement from the individual ideas through the species to the genus, through the genus-species to the genus of genera in the absolute, content-rich unity of the universal idea, is to be taken as subjectivity's deepening of itself in the content of its knowledge. ``The law of concretion, identical with the law of selfhood and decisive for the development of ideas, is already determined once the idea of the genus has been reached. Instead of determining the passage through the individuals to the species, through the species to the genus, through the genus to the family, etc., as a transition from the more concrete to the more abstract, we must then, according to the law that has been demonstrated, conversely determine this passage as a deepening into ever richer, ever more inward concretions.''\footnote{Forelæsn.\ ov.\ Phil.\ Propæd.\ 1860--61 p.~325.} --- Grounding reflection already gave us a self-deepening; but its one-sidedness showed itself in this, that it led to ferment and brooding. This obscurity is now completely overcome: \opage{220} in the ideal system the absolute unity of the ideas is just as clear and determinate as the infinite multiplicity of the relative ideas; for the unity is everywhere seen and posited together with the multiplicity by the same act of thought of the same thinker.

2) \emph{The real system}. The universal is the individual, A --- E; this judgment is antilogical, and hence designates a natural act. In its pure universality the universal is always a thought; but in this system the thought is identical with power: the universal is a mighty thought, a power of thought, and hence a law. But the law is actual only in its fulfillment; the fulfillment is in the individual, determined by the individual and one with the individual. Everything that actually exists exists by virtue of a fulfillment of law, and just for that reason exists as an individual: what rests, rests in accordance with a fulfillment of law; what moves, moves by a fulfillment of law. The same holds of all becoming, arising, passing away, etc. The law, the universal power, does not proceed from the world of eternal thoughts in such a way that it intervenes immediately in the finite and individual; no, the law, the mighty universal, is on the contrary determined by the reciprocal determination of the individual properties of individual things. It is everywhere the peculiar, punctual reciprocal action of individuals with individuals, down to the very smallest particulars, on which the satisfaction of the demands of the universal must here depend: the fulfillment of the law is precisely determined by the connection of things which it itself determines. From this standpoint the antilogical judgment A --- E is to be regarded as an expression for fulfillment of law in general; it is to that extent a logical symbol of all actual and possible natural acts.

From the unity of the universal with the individual in the fulfillment of law it follows that the fulfillment of a single law can in concreto be united with a fulfillment of several mutually heterogeneous laws. It is not merely a one-sided abstract determination of universality that from this standpoint is seen to become one with an abstract \opage{221} determination of existence; here whole systems of universalia make up the logical content which the individual existence, res, the immediate concept of power, expresses in antilogical form. The concretion of the determinations of existence in their ideality gives us once more the universal in the essence of things under the form of categories such as \emph{ground} and \emph{possibility}; yet with the heightened peculiarity which is an expression of the fact that the forms of reflection are now integrated by the antilogical moments of actuality. Substantial reflection gave us the category of powers with its existential significance; but the category of power in logical reflection is one thing, and a system of determinations of power whose union factually expresses the factually given existence is another. The necessary unity of actuality and real possibility has already been demonstrated from the standpoint of the concept as a necessary unity of the objective with the subjective in ontological objectification; but in the present connection the unity in its generality is not enough; the whole emphasis falls on the determinateness decisive for the actual, namely that it is each time just \emph{this} unity under \emph{these}, just \emph{these} circumstances, on \emph{these}, just \emph{these} conditions. If the law as law was the universal, then the determining ground, the real possibility, is now the particular: in order to be fulfilled in the individual, the law must essentially be determined by the particular. The categorical movement from concept to existence thus proceeds from the \emph{universal} resting in the concrete ideality of power (the moving ground, the actual possibility) through the \emph{particular coalescence} of the determinations of power (this ground, sufficient on these conditions; this possibility, completing itself into a disposition in these reciprocal actions), and can to that extent be designated by the judgment: A --- S.

The judgment A --- S, however, first obtains its proper existential validity when it is seen in connection with the ground's issuing in the determinate existence, the possibility's \opage{222} conversion into the individual actuality, the execution of the disposition, and hence with the particular's transition into the individual; this transition we designate by the antilogical judgment: S --- E. Natural acts are categorical acts of power; they presuppose the idea and satisfy the idea, but do not coincide with the acts of thought. Thus, in relation to natural production, the idea of the genus is the possibility of the species, the idea of the species the possibility of the individual; but the production itself shows that it is not the idea as idea, but the individuals, that actually accomplish the natural act. Since the concept of power everywhere has the individual, the individual being, for its reality, it has laid claim to all determinations of knowledge to such a degree that they must all find their logical downfall in the antilogical determinateness of existence. That this downfall is nonetheless one with a continual reawakening, that ideality in categorically determined forms pervades and sustains reality, so that at all points we have universalia in re, is a plain and necessary consequence of the fact that what is objective in objectivity would have to vanish as a phantom if it were not ontologically objectified. Ontological objectifying, as a natural act --- for the natural act is not an abstractly objective event that carries itself out in an abstractly objective way --- is just as little an immediate intuiting as a pure thinking or general reflecting; it has all logical presuppositions, all universal and particular determinations of the concept, present in the moment of activity and application. But these logical determinations are not set free into logical ideality; on the contrary! they are bound, spent, and exhausted in the real production of the objectified object: S --- E --- A.

That the natural act corresponds to the act of thought will then mean, in other words, that it has the act of thought at once both as immanent and as hovering above it. Insofar as the concept of knowledge is determined, and by the determination transformed into a concept of power, the act of thought is immanent; that is to say: it is not an \opage{223} act of thought, but a real production which at all points expresses the thought by sublating the thought, and thus corresponds to the demand of thought. Insofar as the act of power categorically actualizes a system of determinations of knowledge, the act of thought must necessarily hover above it. It would indeed be impossible for the concept of power to be a concept if it were not a real expression of a comprehending power; but a comprehending power must absolutely be a knowing power; but a knowing power is a mighty knowledge, for real comprehending presupposes the logical, and the logical rests on knowledge. If the ideal system with its acts of thought is an empty semblance when the real is lacking, then the real system with its acts of power is a meaningless fact when the liberation of ideality necessary for the logical development of concepts is thought away.

In the real system, too, the difference between the idea as category and the category as idea makes itself felt. Seen from the empirical point of view, the idea can show itself only as category. The thing itself, the concrete, the actual itself, is the individual. ``Nature,'' it is said, ``everywhere gives us only individuals, not species, still less genera.''\footnote{Cf.\ Forelæsn.\ over Phil.\ Propæd.\ 1860--61 p.~250.\ ff.} Genera and species are then, on this way of looking at things, merely concepts of abstraction, universalia post rem, which have nothing to do with the real productions themselves. But how thoughtless such a view is can be seen already from the fact that it makes the natural act itself thoughtless. It is through the identity of the logical and the antilogical in the form of knowledge that the category becomes idea; it is through the corresponding identity in the form of power that there is a reality, and that this reality expresses the idea. If there were nothing that bore within itself the disposition to reality, then nothing could be realized either, and then there would be no reality; but empty thoughts, abstract concepts, subjective fancies, and \opage{224} arbitrary representations do not bear within themselves the disposition to reality, and therefore, strictly speaking, cannot be realized either. Only when the concept in its logical essence so masters the antilogical moments that, by means of this ideal power, it has the disposition to reality within itself, can it --- in a natural act --- be realized; but such a concept is more than a category, it is an idea. It is just as little the formal concepts as the concept-less elements that, either each by itself or when they are added together, have the disposition to reality; no, it is in the original, absolute unity of concepts, elements, and forces, hence in the idea, that the disposition must essentially be sought: only the ideas have the disposition to reality; only the divine ideas can be realized in the proper sense\footnote{That the foolish and confused can nevertheless be realized to a certain degree rests on the fact that the moments of the idea can be parceled out in the finite, and the disposition to reality can to a certain degree be confused.}.

3) \emph{The teleological unity of the ideal and the real system}. To what degree the Hegelian logic, as soon as one disregards the mystical twilight of ideality in pregnant expressions, sinks down into a mere doctrine of forms is especially conspicuous when one considers its treatment of ``teleological activity.'' ``In mechanical as well as in dynamic activity'' --- it says\footnote{Heiberg.\ Specul.\ Logik.\ (Anfr.\ Skr.\ p.~107.\ ff.; Pros.\ Skrifter, 1st vol.\ p.~333 ff.)} --- ``the concept, as pure objectivity, is external to itself, there as the immediately outer, here as the relative and mutual externality of the positive and the negative. But in teleological activity, which is the activity of the concept (just as the mechanical is that of immediacy, and the dynamic that of reflection), it returns to itself, and since the concept is both subject and object, it presents the same opposition anew, and sets itself, as subjective, against the objective external world, \opage{225} but in such a way that it regards this as the opposition which is to be sublated in itself. At the subjective standpoint the concept, in opposition to the objective world, is determined as the intention or plan which has indeed been conceived but not yet carried out, i.e., as \emph{subjective end}. In this syllogism, in which the subjective end is determined as the individual, and the fulfillment of the end as the universal with which the individual is to be concluded together, the \emph{means} is the particular or the connecting link that unites the intention or plan with its execution. The result of the syllogism is thus the executed intention or plan, or the fulfilled end, i.e., \emph{the end as objective}; or the result is the judgment in which the subjective is concluded together with the objective: E --- S --- A. The second syllogism, S --- E --- A, yields as its result the judgment that the means itself is the objective, i.e., that the intention is already fulfilled in the means. And in the third syllogism, E --- A --- S, this is expressed still more determinately by the resulting judgment that the subjective end has its objectivity in the means.''

If one considers the juxtaposition of the categories mechanical, dynamic, and teleological activity exclusively from the side of form, and keeps to the agreed designation according to which ``the subjective end'' is to coincide with ``the individual,'' and ``the objective'' with ``the universal,'' then it cannot be denied that ``teleological activity'' has here been illuminated in such a way that the concept in teleology is clearly seen to stand higher than the concept in the mechanical and dynamic form. But when it is then to be investigated how it actually comes about that the concept, although it seems to remain at each station where it is, in the mechanical as ``mechanical object,'' in the dynamic as ``dynamic object,'' --- nevertheless manages to raise itself upward from stage to stage and, despite the many ``primal divisions,'' not to say parcelings-out, to preserve its unity with itself, then we are once again standing before formalism. \opage{226} It is indeed a strong personification, an anthropomorphism that encroaches upon so ethereal a being as the pure logical concept, when it is related --- for the exposition here comes very close to a narrative --- of the concept that it ``sets itself, as subjective, against the external world, but in such a way that it regards this as the opposition which is to be sublated in itself.'' The ``regarding'' concept, which has here taken its stand ``at the subjective standpoint,'' comes, mirabile dictu! into the bargain to regard itself ``as the intention or plan which has indeed been conceived but not yet carried out''; what wonderful suppleness such a logical concept must possess!\footnote{So that the reader shall not think that Hegel himself speaks more clearly, the following is set down here for comparison. ``Der Zweck ist nämlich der an der Objektivität zu sich selbst gekommene Begriff; die Bestimmtheit, die er sich an ihr gegeben, ist die der \emph{objektiven Gleichgültigkeit und Aeußerlichkeit} des Bestimmtseyns; seine sich von sich abstoßende Negativität ist daher eine solche, deren Momente, indem sie nur die Bestimmungen des Begriffs selbst sind, auch die Form von objektiver Gleichgültigkeit gegen einander haben \dots\ Der Zweck ist in ihm selbst der Trieb seiner Realisirung; die Bestimmtheit der Begriffs-Momente ist die Aeußerlichkeit, die \emph{Einfachheit} derselben in der Einheit des Begriffs ist aber dem, was sie ist, unangemessen und der Begriff stößt sich daher von sich selbst ab. Dieß Abstoßen ist der \emph{Entschluß} überhaupt, der Beziehung der negativen Einheit auf sich, wodurch sie \emph{ausschließende} Einzelnheit ist; aber durch dieß \emph{Ausschließen} \emph{entschließt} sie sich, oder schließt sich \emph{auf}, weil es \emph{Selbstbestimmen}, \emph{Setzen seiner selbst ist}.'' Wissenschaft d.\ Logik 2ter Theil p.~219--20.}

The fault is not that the self-moving concept is spoken of here in a pregnant sense; for if the expression concept were expressly chosen and held fast in such a way that it could actually designate the unity of essence that is active in all things and comprehends all things within itself, why then should one not just as well be able to speak of the all-active concept as of the all-active reason? No, the fault, the fundamental fault, highly questionable for a logic, is precisely \opage{227} this, that the concept itself is not comprehended in such a way that the development of the concept justifies the independence ascribed to it in the pregnant turns of phrase. To judge by the turns of phrase, one would think that the concept of logic was actually a divine being, indeed God himself --- for the concept of God is of course already present in pure being ---; but the development of the concept proves clearly enough that the concept from first to last has its becoming, existence, and movement in a human thinking, in a thinking that is not merely human but human to such a degree that by abstracting from itself it even imagines it can make itself divine. The fundamental fault is that the act of thought without further ado becomes a natural act, the logical development of forms a real development of the thing, the ideal system a real system, and conversely.

If the emphasis is laid on the opposition between the ideal and the real system, the former will one-sidedly present freedom, the latter necessity. Logical freedom then rests on this, that every category, every determination of the concept, and every particular concept is developed and clarified in the light of the whole infinite system of concepts. Thus, e.g., the concept being can just as little be a concept in divine as in human thinking unless it is lifted out by itself, by a special act of thought, in opposition to non-being; thus the concept ground cannot be a concept unless the ground is expressly comprehended in its unity-in-difference with the concept existence; in short: from this side the ideal system can be regarded as a divine logic, a system of thoughts, concepts, and ideas that move in eternal thinking and rest in eternal knowledge.

In opposition to logical freedom, real necessity rests on the immediate unity of the concept with its existence, hence on its transformation from a concept of knowledge into a concept of power: the real system is a system of categories in antilogical form; here there is in the objective sense nothing but \opage{228} individual existences and their reciprocal actions. These reciprocal actions, with the corresponding formations, transformations, and changes, are partly successive, partly simultaneous; here there are movements in time and space; here, finally, there is the opposition of motion and rest.

Since the aim of the ideal system is a development of concepts as concepts, and the aim of the real system an actualization of concepts in existence, the opposition between the ideal and the real system is a teleological opposition; the former is a realm of knowledge, the latter a realm of power. The teleological opposition rests in the teleological unity; this unity is the all-embracing idea of subjectivity. Existence could not possibly be the realized concept if there were no concept that could be realized. But in what exists the concept is sublated; the concept of power, the factual concept, is something existing, and this existing thing a concept which is no concept. In order that the concept can be realized, it must be; in order that it can be, it must be comprehended: the ideal system with its concepts is means, i.e., essential medium for the real system. Conversely, no concept can be developed, articulated, or posited as a concept of existence unless in its development it has originally taken up that sum of antilogical reflexes by which alone its existential peculiarity can be determined: the real system is in this respect means, that is, actual medium for the ideal system.

The unity of real necessity and logical freedom is, owing to the infinite reciprocal determination of the essential and the actual medium, so thoroughgoing that necessity at every point expresses freedom, and freedom explains necessity. For what, in this connection, is freedom? The subjective essentiality of ontological objectification. Therefore all existences and existential reciprocal actions, although in actual existing they are anticategorical, must nevertheless be apprehended in categorical form, determined as facts that satisfy categorical demands. \opage{229} Freedom in the ideal system is without arbitrariness; here no other thoughts are thought, no other concepts developed, no other forms clarified, no other ideas held, than such as --- in the infinite mediation of the mediate with the immediate --- are demanded partly mediately, partly immediately, in order that existence may proceed from existence by proceeding from the ground, actuality from the actual by proceeding from possibility, and that the natural act may be determined as a categorical act of power. Necessity in the real system is more than contingent. Contingent or merely hypothetical necessity (if a is, then b is) already shows the reciprocal determination of category and fact: if the presupposition is factual, then the consequent is categorical. The actual connection of existing things rests on an essential connection of categories of existence; the system of things is essentially a system of grounds, for the thing is the existential expression of the sufficient ground; the system of actualities is in a corresponding way a system of possibilities, the system of individuals a system of species and genera, etc. But the systems of grounds, possibilities, genera, and species, with their real as well as ideal connection, are clarified in the ideal system, hence in freedom.

In the ideal system all natural acts are reflected as mighty acts of thought whose movement is logical; in the real system all acts of thought are reflected as categorical acts of power whose movement is antilogical. The unity-in-opposition, the articulated reciprocal determination of act of thought and natural act, a reciprocal determination in which the antilogical is means when the logical is end, and conversely, is \emph{teleological}. The clear, teleologically transparent unity of the realm of knowledge and the realm of power is the idea itself, a light of truth in subjectivity.

\opage{230}\begin{center}
{\bfseries c)\quad The Logical Ideal: Comprehending and Knowledge.}
\end{center}
\addcontentsline{toc}{subsubsection}{c) The Logical Ideal: Comprehending and Knowledge}

That in logic, as in every other science, it essentially depends on trustworthy starting points and reliable marks, if investigations are to be undertaken at all, goes without saying. It is psychological subjectivity, with its presuppositions and conditions, from which we must constantly proceed in the logic of fundamental ideas, and to which we must again return. In the logical analyses everything depends on how the starting points given as facts are used: to remain in finitude does not satisfy; to abstract from everything finite and let the thoughts hover in the empty infinite does not satisfy either. In the treatment, within logic, of the relation between perfect knowledge and the psychological factors of knowledge there is something here that recalls the treatment of the so-called ``indeterminate forms'' in mathematics; none of the finite presuppositions can simply be retained as a fact, but every fact can and must be used as a moment: the task is like calculating with infinite magnitudes.

Immediately and straightforwardly no one can say how much $0 \cdot \infty$ is worth, or what value one must assign to $\infty - \infty$; likewise no one can immediately and straightforwardly decide how a real concept issues in its existence, and what movements the acts of thought must call forth in the pure idea. Not by a leap, but by a transition corresponding to the peculiarity of the functions, must comprehending subjectivity in each case move from the finite to the infinite, if it is to succeed in finding the way back from the infinite to the finite. To illuminate the advances which the human spirit, despite all its missteps, errors, and one-sidednesses, has made in this respect too, the phenomenology of spirit together with the history of philosophy will constantly continue to furnish new and essential contributions. Trivial considerations such as these: that philosophy is still \opage{231} in its childhood, or that from old age it is entering its second childhood, that it has not yet led to anything certain, since one philosopher thinks one thing and another another, when the one thinks that one should hold to the principle of contradiction and subjective logic, the other that only dialectic and objective logic are able to express the absolute; when the one thinks that it is empirical inquiry, the other that it is speculation, which alone leads to the cognition of truth, etc. --- science can hardly find it worth the trouble to refute\footnote{``Dear Jensen! You are surprised that I, notwithstanding the continuing disputes of the philosophers, ascribe reality to philosophy. Do you think, I would ask you in reply, that any philosopher could come into dispute with you? You are surely too modest to expect it. Do you not think, then, that the disputes of the philosophers must be grounded on a mutual understanding, an agreement; that their disputes are facts which presuppose that they have a common ground to stand on?'' (Poul Møller.)}.

Modern philosophy, as regards the problem of subjectivity, is oriented by an entirely different standard than the ancient could ever have been. With the many-sided illumination of the problem of subjectivity there naturally follows a corresponding illumination of the problems of objectivity. When a Fichte in all categories sees only the subjective, while a Hegel in all of them essentially sees the objective, it is not thoughtlessness but on the contrary persistent thinking, carried through from a peculiar starting point, that produces such one-sidednesses. It is at bottom one and the same analysis of the infinite that has here been undertaken from two opposite sides. Or how would it be possible to let the universal concepts of knowledge stand forth in the light of subjectivity if a subjective thinking, reflecting, and acting did not in an infinite manner permeate them all? How would it be possible to unfold the categories in so objective a system that it looks as \opage{232} vividly as if the concept of the thing were the thing itself, and the movement of thought one with the actual movement of the thing, if these a priori forms of objectification of our knowledge were not, in an infinite sense, as it were shot through with flashes of actuality from the objective world?

When now, further, in sharp opposition to the objective as well as to the subjective carrying-through of the analysis of the infinite, finitude and the finite are accentuated anew, --- whether this otherwise happens in a finer way with the help of the ``reals'' and mathematical discreteness or in a coarser manner with the help of atoms and other materialistic props --- then such a reaction on behalf of the finite facts is just suited to bring forward, in a refreshing and beneficial way, the points of actuality and articulating intermediate determinations which those analyses of the infinite, concluded too early and guided by a one-sided principle, have not yet been able to assimilate. On this occasion too we must recall Poul Møller's words, ``that true philosophical inquiry is a work cohering through several centuries, an unbroken chain whose later links presuppose the earlier as their necessary condition.'' --- ``The goal, absolute knowledge, or spirit knowing itself as spirit'' --- thus Hegel concludes his ``Phänomenologie des Geistes'' --- ``has for its path the recollection of the spirits as they are in themselves and accomplish the organization of their realm. Their preservation, in respect of their free existence appearing in the form of contingency, is history, but in respect of their comprehended organization it is the science of phenomenal knowledge; both together, comprehended history, form the recollection and the place of judgment of absolute spirit, the actuality, truth, and certainty of its throne, without which it would be the lifeless, the solitary; only from the chalice of this realm of spirits does infinity foam forth for it.''

In the closest connection with the task of clarifying, out of the psychological, \opage{233}the ontological in knowing subjectivity stands another, equally essential task, namely that of determining the relation between the human and the divine in ontological knowledge itself. This task has its difficulties in this, that every uncertainty concerning the reciprocal determination of the finite and the infinite within the human consciousness's own bounds at once brings with it a corresponding uncertainty in the determination of the difference between human and divine knowledge.

The logical ideal, the ideal of knowledge, in its pure universality must then necessarily be misunderstood and botched, from the one-sided standpoint of finitude as well as from that of infinity. By remaining seated in the midst of finitude, subjectivity may miss the ideal in two ways: partly by a thoughtless and self-contradictory conception of the perfections of divine knowledge, partly by a narrow-minded holding fast to the psychological conditions of human knowledge. In the first case the ideal becomes meaningless, in the latter it is rejected. If knowledge is not possible except on human terms, if every knowing subjectivity must absolutely have a brain, like man, five senses, like man, etc., then it follows from this that knowledge by its nature is and must be relative, that there can be only a difference of degree between lower and higher stages of knowledge, but no absolute or perfect knowledge. If, on the other hand, the knowing subjectivity, with the help of abstractions and speculative reductions, is done with the finite too easily, and in pure infinity comes to confuse its empty beholding with ideal knowledge, then in mystical enthusiasm it will all too easily mix up the divine with the human and take the ideal in vain. ``Speculation is everything, is to behold in God, to contemplate that which is. Science itself has value only insofar as it is speculative, that is, contemplation of God as he is. But the time will come \opage{234} when the sciences will more and more cease, and immediate cognition will enter. Only in the highest science does the mortal eye close, since it is no longer man who sees, but the eternal seeing itself is seeing in man.'' (Schelling).

What these misunderstandings signify in a logical respect can best be illuminated in what follows when the ideal of knowledge is characterized with particular regard to \emph{intuitive} and \emph{discursive knowledge}.

\begin{center}{\bfseries α)\quad Intuitive Knowledge.}\end{center}
\addcontentsline{toc}{paragraph}{α) Intuitive Knowledge}

In the present connection --- to explain the ideal of knowledge --- it will not be without logical interest to cast a glance at the strictly orthodox dogmaticians' conception of God's omniscience, and in this respect to acquaint oneself with a testimonium spiritus of an older date. The success with which the more recent speculative dogmatics has known how to gather and use philosophical grains of gold --- and when it is only a matter of coating the old dogmas with a thin gilding, one can make do with a little gold --- has of course already been acknowledged in our foregoing; here, in confirmation of how the Holy Spirit, at the time when the ecclesiastical system of doctrine still kept itself pure of all philosophical contamination, must have attended to the ``dogmatic enlightenments,'' we shall cite a few further features.

For the determination of the ideal of knowledge the dogmatics of the 17th century offers interesting main points. Omniscience'' --- it says\footnote{Hutterus Redivivus. Leipzig 1836. p.~138--139.} --- ``is, with respect to its object: 1) \emph{a necessary} knowledge (scientia necessaria, qua Deus semetipsum atque in se ipso omnium rerum necessitatem perspicit), 2) \emph{a free knowledge} (scient.\ libera [repræsentatio visionis s.\ intuitionis], qua Deus omnes res præter ipsum vere existentes vere novit), 3) \emph{a conditional} \opage{235} \emph{knowledge} (scient.\ media [sc.\ de futuro conditionato s.\ futuribili, hypothetica, simplicis intelligentiæ], qua Deus perspicit omnia, quæ, positis quibusdam conditionibus, evenire potuissent). With respect to the perfection of the mode of cognition, God's knowledge is: 1) \emph{intuitive} (intuitiva [pura, immediata] i.\ e.\ sine sensu, imaginibus, abstractione, discursu et ratiocinio), 2) \emph{simultaneous} (simultanea), 3) \emph{distinct} (distinctissima), and consequently 4) \emph{perfectly true} (verissima).''

If we now subject these perfections of knowledge, assembled under one survey, to a closer examination, it becomes quite evident to us that ``the dogmatic talent'' has come to set up these perfections not in accordance with any real insight into the functions of knowledge, but by an immediate leap from finite to infinite. That this leap, logically understood, has taken place in all possible thoughtlessness, for that we have the testimonium spiritus --- in thus ascribing to the testimonium spiritus a ``theoretical significance'' --- that the one of the perfections set up directly annihilates the other. ``With respect to its object'' the divine knowledge is to know the same things in three different ways. ``God sees through himself and in himself the necessity of all things.'' Here the dogmaticians seem to have had in mind what Proclus teaches in his ``Platonic Theology'': ``God cognizes the divided undividedly, the temporal timelessly, the non-necessary necessarily, the changeable unchangeably, and in general all things more excellently than according to their order.'' But what now does the \emph{free knowledge} mean? Answer: a knowledge of free, i.e., independent things existing outside God. But how can things outside God be known by God to be with freedom when God in himself nevertheless knows that they are with necessity? If freedom in the divine knowledge is without further ado the same as necessity, why then does one use two such different expressions as freedom and necessity? If freedom is one thing and
\opage{236} necessity another, how then does it come about that the same things can be both free and necessary? If one says: ``Things are necessary in their ground and grounding, but contingent in their external and finite existence, and free in this contingency,'' then we must ask again how contingency and freedom are then related to the ground? For if in what is contingent to human eyes there really is something of which God knows that it is at bottom detached from necessity and posited solely by his free will (libertas), then the things existing outside God are by that very fact loosened from the ground of necessity; the \emph{necessary} knowledge, the knowledge by which God in himself sees through the necessity of all things (omnium rerum necessitatem), must therefore be an untrue knowledge.

In determining ``free knowledge,'' dogmatics has expressly proceeded from the presupposition that there is a world of actual things existing over against God (res præter ipsum vere existentes), and that God cognizes these things in their finitude, in that he cognizes them as they actually are (vere novit). This regard of the divine knowledge for finitude and its conditions comes out still more clearly from the standpoint of hypothetical knowledge; here there then arises a finite separation not merely between possibility and actuality, but within possibility itself between those possibilities which, under given conditions, become actual, and those possibilities which do not, because the conditions do not occur. In all this there is a certain naturalness. The dogmatician takes actuality straightforwardly, as it appears to human eyes, brings out its various sides, and then it of course goes without saying that God, who is omniscient, must know all possible and all actual things from all possible and actual sides. To this there then also corresponds the subdivision falling under ``free knowledge'': remembrance (reminiscentia), sight (visio), and foreknowledge (præscientia).

\opage{237} But everything that is thus set up under the two rubrics, ``free'' and ``conditioned knowledge,'' must be torn down again when the ``mode of cognition'' in and for itself is to be more closely determined. Here ``intuitive knowledge'' is now the most important category of all: how then is ``intuitive knowledge'' characterized? It is, we are told, altogether immediate. The immediacy is then here to be understood more precisely thus, that ``intuitive knowledge'' is, in the first place, without sense-perception and without images (sine sensu, sine imaginibus); but already these negative determinations must awaken misgivings. If the divine knowledge is without sense and sensing (sine sensu), how then can it be determined by a seeing, a sight (visio)? If the divine beholding is to be immediately one with what is beheld, so that there is no difference here at all, then one must necessarily assume one of two things: either what is beheld must vanish in the pure beholding, which then becomes a beholding that beholds only itself, a beholding of a beholding, an infinitely empty beholding; or else the beholding must be absorbed in what is beheld in such a way that the object of knowledge steps into the place of knowledge, and knowledge vanishes. Let us try it once more with the intuitive knowledge of a stone: either the stone must vanish entirely for the intuition, or the intuition for the stone; thus: if the stone is, then the intuitive knowledge is not; if the intuitive knowledge is, then the stone is not. What holds of the stone naturally holds of every finite and actually existing thing; thus: if the existing thing is, then the intuitive knowledge is not; if the intuitive knowledge is, then the existing thing is not. But when the divine knowledge is intuitive, and an intuitive knowledge of existing things is a contradiction, how then is one to understand a proposition as dogmatic as this one: scientia libera est repræsentatio visionis, qua Deus omnes res præter ipsum vere existentes vere novit?

If any meaning is to be connected with what dog\opage{238}matics teaches about God's intuitive knowledge, then one must necessarily presuppose that there is a difference between beholding and what is beheld, between knowledge and the object of knowledge, and that this difference has immediate validity; one thing, then, is the beholding, another that which is beheld by it; one thing is knowledge, another is the object of knowledge. Under this presupposition it is then possible, at least to a certain degree, to connect a meaning with the assumption that God has an intuitive knowledge of all existing things. The model is again human. We humans can, after all, in a way also have an intuitive knowledge of certain easily surveyed objects (of this stone, this plant, this house), when they are placed in a favorable light. The difference between human and divine knowledge is naturally this, that the former in each individual case confines itself to apprehending a single object from a single side, while the latter apprehends at once (scientia simultanea) all possible objects equally clearly and equally distinctly (scientia distinctissima) from all possible sides. But do we not now get into difficulty with the dogmatic declaration that God's intuitive knowledge is without sense and sense-perception (sine sensu), without image and representation (sine imaginibus)?

If intuitive knowledge absolutely presupposes a unity-in-difference of knowledge and the object of knowledge, then an intuitive knowledge without sense and sense-perception is a meaningless contradiction. What indeed is sense-perception other than an immediate apprehension, hence an immediate knowledge of something that immediately determines this knowledge and thereby constitutes its object? That divine sense-perception must be essentially different from human sense-perception goes without saying; but if dogmatics does not dare to ascribe sense-perception to God because the divine sense is so infinitely different from the human, how then does dogmatics dare to ascribe to God \opage{239} any knowledge at all, since the divine knowledge after all is and must be infinitely different from the human?

Dogmatics teaches that God in intuitive knowledge cognizes things without images (sine imaginibus). But those who without further ado put forward such an assertion seem not so much to have been prompted by some testimonium spiritus sancti as rather, by an obscure feeling of theology's scientific need, to have been tempted to bear a theoretical testimony that at bottom they have thought neither about the conditions of human nor about those of divine knowledge. How in all the world could a beholding of an existing object be thought, if it may not be thought in such a way that an image of the object forms itself in the beholder by means of this beholding? If one gives up this doubling, and lets God's intuitive knowledge intuit the object as object without an intervening image, then one falls back with inescapable necessity into the absurdity that, as already remarked, it must be either knowledge that in intuition devours the object, or the object that absorbs knowledge.

The absurdity in the declaration put forward can moreover, from the standpoint of dogmatics itself, be seen in several ways. What becomes---not to speak of the divine remembrance (reminiscentia)---of God's foreknowledge (præscientia), if in the divine knowledge there is no difference, and that an essential difference, between the image of the thing and the thing itself? If the thing coincides absolutely with its image, then one may turn and twist as dogmatically as one likes: one has plainly denied the divine foreknowledge; for foreknowledge rests precisely on the sight of the thing, that is, the image of the thing, going before the thing itself.

If it is really a truth and not an empty solemn figure of speech that the Holy Spirit---``in that we thus ascribe to the testimonium spiritus sancti not merely practical but also theoretical significance''---has helped the older as well as the \opage{240} newer dogmaticians with their scientific conception of ``God's intuitive knowledge,'' ``the all-penetrating seeing,'' then to us frail humans, who, without being able to appeal to supernatural inspirations, must see the matter with merely human eyes, it undeniably comes to look as if the Holy Spirit were somewhat careless with its theoretical testimonies; indeed it sometimes looks as if, in the testimonies it bears in recent times, it had, to put it plainly, forgotten a good many of the testimonies it held fast with the utmost strictness in an earlier period. This holds in particular of the determination of the relation between things and the images of things. If God's ``intuitive knowledge'' was in older days sine imaginibus, without imagination and consequently without imaginings, then it has now in these latter days, by way of compensation for those deficiencies, conceived by ``speculation'' and got itself a ``pleroma, a paternal pleroma'' of filial-speculative imaginings. ``The paternal pleroma, which is revealed in the Son as a necessarily originating realm of ideas, is transfigured by the free artistic activity of the Spirit into a Kingdom of Glory (δόξα), where the eternal possibilities play before God's face as magical actualities, as a heavenly host of visions, of plastic prototypes for a revelation ad extra, to which they, as it were, long to be released.''

Similar difficulties force themselves upon us when, instead of looking at what is intuitable in sight and image, we keep specifically to the side of the understanding, to what concerns pure thinking, comprehending, and insight. Here dogmatics now wants abstraction and discursive knowledge simply excluded from the intuitive; the divine intuition is sine abstractione, sine discursu, sine ratiocinio. What is intended by this is clear; the divine knowledge cannot proceed in the same finite and mediate way as the human. ``Our knowledge is piecework''; the piecing-together is conditioned by abstraction, by what is discursive in our investigations and deliberations, \opage{241} what is mediate in our inferences. But it is one thing that the divine knowledge does not use the mediate in the same finite way as the human, and something quite different that in its immediacy it should exclude the mediate. If God's intuitive knowledge really is without all abstraction, what then follows from this? That it must be without all thinking, hence thoughtless! All thinking rests to such a degree on abstracting and on integrating what has been abstracted that thought must necessarily come to a stop when these purely intellectual activities cease.

Without abstracting and thinking, the immediacy of intuition is an animal immediacy. To intuit an object in such a way that the intuition is lost in animal immediacy is what is called gawping. The animal gawps at the strange object; the cow gawps at the new gate, etc. Animal gawping is characterized precisely by this, that its immediacy lacks the reflection necessary for all thinking and understanding. The animal sees the thing without being able to raise itself to a seeing of this its seeing; its seeing is immediate in the sense that it cannot free itself from immediacy. On what condition, then, can such a liberation take place? On that of abstraction, only on that of abstraction! for abstraction conditions negation.

How many abstractions must not be presupposed in order that one who contemplates an object may be able to say to himself: ``here I see a plant''; ``here I see a stone.'' What breathing is for our organic life, abstraction and its sublation are for all thinking and comprehending; all thinking and comprehending---in this Hegel is unconditionally right---goes through mediations; but mediation is an infinite mediating, in the course of which the mediate comes forth in order to be thoroughly formed, sublated, and to vanish into a higher immediacy. Wherever there is talk of an immediate knowledge, one must take care to distinguish between the immediacy with which one begins and the immediacy with which, \opage{242} through many mediations and intermediate determinations, one completes. The more mediations a knowledge has run through, the higher in turn is the immediacy that will result from it.

That this law, since it is a law of knowledge, must hold for divine as well as for human knowledge, can be made clear without difficulty. If we assume with our dogmaticians that ``God's intuitive knowledge'' is a clear and distinct knowledge (scientia distinctissima) of all individual existing things in all their particulars; if we further assume with the dogmaticians that God's intuitive knowledge of each individual existing thing is immediate (pura, immediata), what then follows from this? Seen from the standpoint of dogmatics itself, it will necessarily follow that in God's intuitive knowledge there must be just as many intuitions, just as many immediate glances, in other words, just as many divine eyes, as there are existing things and individual determinations of the existing things in heaven and on earth. Whether dogmatics now prefers to express itself more biblically and speak of God's eye, hence of the eye God ``\emph{has},'' or more speculatively of the eye God ``\emph{has not},'' the result is nonetheless this, that in the one divine eye there must be an infinite multiplicity of eyes, in the one universal intuition an infinite multiplicity of separate intuitions, in the one all-penetrating divine glance an infinite multiplicity of divine glances.

If one now wants quite thoughtlessly to hold fast the multiplicity of separate intuitions in immediate opposition to the immediate unity of intuition, or---since the contradiction is equally glaring whichever ``dogmatic elucidation'' one chooses---to let the immediate unity of intuition be absorbed in the immediate multiplicity of separate intuitions, then one blasphemously transforms the divine eye into a faceted eye, a monstrous cockchafer's eye, and confuses the immediacy of God's intuitive knowledge with the immediacy \opage{243} of an animal gawping. Should immediacies of this nature still to this day be repeated, inculcated, and admired as scientific inspirations, then a testimonium spiritus may very easily transform itself into a testimonium stupiditatis. Such testimonia stupiditatis---in that we thus ascribe to testimonia stupiditatis ``not merely practical but also theoretical significance''---plainly stem not from religious profundity but, especially as regards the new rehashes of old scholastic subtleties, from logical frivolity.\footnote{``Whoever wants to find favor with the literary rabble must not rise to the region of pure speculation, nor crawl upon the prosaic earth either, but be like a bird with half-clipped wings, which flies up a little now and then, but nonetheless keeps plopping down again.'' This ``rule of prudence,'' taken from Poul Møller's Strøtanker, ought perhaps especially to be recommended to the ``dogmatic talent'' that honestly strives, not exactly to fathom a truth, but---by making use of sundry testimonia---to maintain an opinion.}

\begin{center}{\bfseries β)\quad Discursive Knowledge.}\end{center}
\addcontentsline{toc}{paragraph}{β) Discursive Knowledge}

Strong by virtue of its strong faith and above all borne up by its testimonium, the theological-dogmatic consciousness easily comes, as if by a breath of ``inspiration,'' to a kind of climactic hovering in the high, star-rich vault of heaven, where the ideal of knowledge appears roughly like the white shimmer of an infinitely blue tone; by way of compensation we then have, at the sensuous-realistic level of consciousness, a standpoint firmly on the ground.

The sensuous-realistic way of viewing things can, as regards the ideal of knowledge, be referred to the following three points: 1) the peculiar character of our senses and their relation to our thought; 2) the peculiar character of our finite and mediate knowledge corresponding to the peculiar relation between our sense-perception and thought; 3) the impossibility of conceiving a knowledge that should be of a higher kind than our own human knowledge.

\opage{244} 1) Now as regards, in the first place, the peculiar character of our senses and their relation to thought, sensualism can appeal to the fact that the immediate must from the beginning be apprehended in an immediate way. ``What a sensation is we learn by sensing; what a sight is, by seeing; what a hearing is, by hearing. But from this it then follows that our knowledge of what sensation is must wholly and entirely conform to the properties of human sensations. What sensations an animal, e.g.\ a snail, may be able to have, of that we are able to form a representation only insofar as we presuppose that there must be an essential likeness between the sensations of the animal and those of man. The same holds of sight and hearing; we are compelled to judge all other seeing by our own seeing, all other hearing by our own hearing. How unconditionally we are bound to what is immediate in each peculiar sense is further seen from the fact that each sense apprehends its own; what is seen cannot be heard, what is heard cannot be felt, etc. However much one trains, sharpens, and arms the individual sense, it is nevertheless not possible to bring about a smooth transition from one sense to another: a sharpened hearing does not approach becoming sight, nor a sharpened sight becoming hearing. Certainly one can in a way make the audible visible, tones for example by sound figures, but this happens only in that the same thing presents two heterogeneous sides, of which the one corresponds to the other, the audible to the visible, or conversely. That the senses, which are so heterogeneous with one another, are also heterogeneous with thought is equally settled. Psychologists may repeat long enough the assurance that the understanding is a sense; they will nonetheless not be able to show us what sort of sensible thing it is that the understanding senses. To make the understanding into a sense for the non-sensible, hence into a sense for what cannot be sensed, is about the \opage{245} same as renouncing it: every sense must have its sensible; a sense for something non-sensible is a contradiction. One thinks about the sensible, indeed one can and must think by means of the sensible, but the sensible itself cannot be thought. As thought is, so must the understanding be, and conversely. An understanding that should be able either to do without the sensible, or to make itself one with the sensible, or to penetrate the sensible, is so alien and incomprehensible to our nature that we must flatly deny its reality. As a consequence of this we must then also deny that there should be conditions of knowledge of another kind than the human, or any knowing subjectivity of a higher order than the psychological.''

2) On the peculiar character of the senses' relation to one another and to thought rests---and this is the next main point---the peculiar character of our knowledge. ``If one wants without further ado to call immediate knowledge intuitive, then one may well say that our knowledge at all points begins with the intuitive; but intuitive knowledge then in this way comes merely to form the presuppositions for the discursive; thus there must always be something immediate that can furnish the foundation for the mediate. The mediate presupposes middle terms; every middle term is a means, but the means, considered by itself, is immediate. Now as the mediate is related to the immediate, so, as regards our knowledge, the discursive is related to the intuitive. If the object of knowledge is a totality, then our knowledge of the whole is mediate insofar as it comes about through an immediate but successive apprehension of the individual parts and their mutual connection: the partial and mediate here coincides with the discursive. Conversely, there can also be a certain immediate apprehension of the whole, namely a total impression; the mediate and discursive then comes to rest on the partial investigation of the particulars.''

\opage{246} ``The relation between the intuitive and the discursive can accordingly be determined by the opposition between analysis and synthesis; if the synthesis is immediately given for intuition and our knowledge is to that extent intuitive, then the analysis is discursive; if the individual element brought out by analysis is immediately present for intuition and thus to be apprehended intuitively, then the synthesis, the partial combination of the parts, the successive unification of the elements, gives us the discursive. The same can be seen from the relation between the premises and the conclusion in every syllogism. Insofar as a premise is evident by itself, it is apprehended intuitively; but precisely the circumstance that here there is a premise, that one must therefore go either from premise to premise or at any rate from premise to conclusion, shows the discursive. Now since the discursive is everywhere recognizable by the plurality of the different points of departure, by the partial, the successive, and the mediate, and the intuitive with all its immediacy is likewise partial, successive, and mediate---namely insofar as the partial and successive immediacy must always be partially conditioned---it is seen from this that the intuitive, far from having any independent significance over against the discursive, is on the contrary to be apprehended as the discursive knowledge's own element, an element that constantly points beyond itself and first gets its significance through the heterogeneous contributions illuminating one another. We therefore do not have, when we look at the totality, an intuitive knowledge by itself over against or alongside a discursive one; rather all our knowledge, with its heterogeneous elements of immediacy, is essentially discursive.''\footnote{In Kant the separation of sensibility and understanding, of intuition and concept, corresponds in a peculiar way to the separation between intuitive and discursive cognition. ``Reflectiren wir auf unsere Erkenntnisse in Ansehung der beiden wesentlich verschiedenen Grundvermögen der Sinnlichkeit und des Verstandes, woraus sie entspringen, so treffen wir hier auf den Unterschied zwischen Anschauungen und Begriffen \dots\ Dieses ist der \emph{logische} Unterschied zwischen Verstand und Sinnlichkeit, nach welchem diese nichts, als Anschauungen, jener hingegen nichts, als Begriffe liefert \dots\ Auf den hier angegebenen Unterschied zwischen \emph{intuitiven} nnd \emph{discursiven} Erkenntnissen, oder zwischen Anschauungen und Begriffen gründet sich die Verschiedenheit der \emph{ästhetischen} und der \emph{logischen Vollkommenheit} des Erkenntnisses \dots\ Die logische Vollkommenheit des Erkenntnisses beruht auf seiner Uebereinstimmung mit dem Objecte; also auf \emph{allgemeingültigen} Gesetzen, und läßt sich mithin auch nach Normen a priori beurtheilen. Die ästhetische Vollkommennheit besteht in der Uebereinstimmung des Erkenntnisses mit dem Subjecte, und gründet sich auf die besondere Sinnlichkeit des Menschen.'' Kants Logik von G.~B.~Jäsche. 1800. Werke von G.~Hartenstein. Leipzig 1838 vol.~1, p.~360--61. --- But if one wants in this way to hold fast an intuitive knowledge over against a discursive one, then one must at least admit that intuitive knowledge is without all insight and concept. ``Es finden daher bei der ästhetischen Vollkommenheit keine objectiv- und allgemeingültigen Gesetze Statt, in Beziehung auf welche sie sich a priori auf eine für alle denkende Wesen überhaupt allgemeingeltende Weise beurtheilen ließe.'' Even when, as in the Kritik der Urtheilskraft, an account is given of what beauty is, the cognition becomes discursive.}

\opage{247} 3) ``If our knowledge in its total development is conditioned by nothing but partial developments and is thus essentially discursive, then the same must hold of all possible knowledge. We grant that on other globes in the vault of heaven there may be rational beings whose knowledge is far more perfect than the human; but the perfection then rests on a difference of degree, not on a difference of sphere. If our knowledge is essentially discursive, then the discursive belongs originally and necessarily to the essence of knowledge, and then the higher, highest, and very highest knowledge of the higher, highest, and very highest rational beings must necessarily be discursive. What this means can easily be elucidated \opage{248} by a couple of examples. Among the many imperfections of our knowledge is, among other things, this, that on the whole we learn so slowly, and easily forget one thing over another. How long does it not take us to go through a more comprehensive course, be it mathematical or logical; what exertions does it not cost us when we are to work our way through great systems, and unite insight into the particular with a survey of the whole! Now since experience nonetheless teaches that within the circle of human individuals there can be a great and remarkable difference with respect to intellectual endowment, there is nothing here to prevent the assumption that on other globes there may be rational beings with far greater and more excellent mental powers than those of men. These beings might then be endowed with a sight surpassing the human in strength, sharpness, and fineness; but they could not possibly be endowed with a sight that made it possible for them at one time to see in all directions and to see all things in one and the same instant equally distinctly. Let them, as regards mental powers, be so richly equipped that in a day, indeed in an hour, they can master systems whose appropriation would cost the greatest human genius half a lifetime; they would nevertheless not be able to think two thoughts, let alone two theorems or two proofs, equally distinctly at one time. Let them be ever such great physicists, let their sensing, representing, and thinking accompany one another to such a degree that they are close to counting the vibrations of the air as easily as they hear the tones, and the vibrations of the ether as quickly as they see the colors: a simultaneous hearing, a simultaneous seeing, a simultaneous counting will nonetheless be impossible for them. To see the whole surface of a board with one glance is an impossibility; yet the succession, even under the conditions of our sight, takes place so quickly and easily that, when the surface is not too large, we have a sensation as if we saw the whole at once. If the surface is, e.g., red, then the homogeneity will even contribute to heightening the impression \opage{249} of simultaneity. But if we now imagine that each of the red rays were also to be apprehended by itself, then the dizziness begins; if we further demand that the glance follow, and the thought count, the vibrations in each ray (thus for each ray 477 billion vibrations per second), then everything goes black before our eyes; if we then finally add the demand that the one who sees the colored surface and counts the vibrations of the ether shall at the same time hear the Overture to The Marriage of Figaro and note each tone and count the vibrations of the air, then we have the meaningless.''

``What can be derived from this is something very important; for from this it can be derived that not merely human knowledge, but every knowledge, however perfect otherwise, is and must be limited and, in consequence of its limitation, imperfect. A perfect knowledge, a knowledge without one-sidedness, without bounds and limitations, is a triangle without boundaries, a circle with infinite radius.''

The denial of the ideal of knowledge thus motivated is well suited to make us attentive to the peculiar reciprocal relation in which the theory of cognition and logic must necessarily stand to one another. If we abstract from the theory of cognition, it will be impossible to develop the logic of subjectivity, since we must then lack the points of departure that are determined solely by the psychological conditions of subjective knowledge; if we keep exclusively to the theory of cognition, hence to the psychological presuppositions of subjective knowledge, we come consistently to a denial of the absolute validity of knowledge and thereby of truth. What a difference it makes whether the problems of cognition are treated as in an introduction to the logic of ideas, or are set up on the foundation of the ontological self-development of subjectivity, is easily seen by a look back.

The peculiar character, pointed out under the first point of view, of our mutually heterogeneous senses and their relation to thought would, if we kept exclusively to the theory \opage{250} of cognition, certainly be incompatible with the concept of a perfect, all-embracing and penetrating knowledge. But if we recall what has been developed at the stages already traversed of \emph{immediate}, \emph{reflecting}, and \emph{comprehending} subjectivity, the matter takes on a different aspect. While \emph{immediate subjectivity} had an eye only for the psychological peculiarity of the senses, the \emph{reflecting} saw clearly that this psychological peculiarity must at all points have its corresponding ontological determinateness. It is now this determinateness that for \emph{comprehending subjectivity} is expressly articulated through the reciprocal determination of the logical and the antilogical. It is not the psychological facts of sensibility, but the ontology of sense-perception and of the senses, on which the emphasis must now fall.

What has been brought out here under the second point of view to show that all our knowledge is and must necessarily be discursive likewise continues to have incontestable validity so long as we keep to the finite separation of the subjective and the objective. But already the way in which \emph{reflection of the understanding} gradually passed over into \emph{essence-reflection} made it evident, after all, how the finite in the relation dissolved point for point into peculiarly determined infinity; the infinitely determined reciprocal relation between subjectivity and the objective was thereupon made clear, at the standpoint of comprehending subjectivity, by a demonstration of the difference between psychological and ontological objectification, a difference which, when the question concerns the relation between intuitive and discursive knowledge, is of decisive significance. The theory of cognition knows only psychological objectification; objective logic ``lets the subject determine itself to object,'' but otherwise seems to have no inkling that there is, in the ontological, that is, in the subjective-objective logical sense, a problem of objectification. So \opage{251} long as the problem of objectification has not yet come into being, there will certainly be required partly a mystical contemplation, partly a dogmatic testimonium, in order to raise consciousness from discursive cognition up to the divine region of intuitive knowledge, of blessed beholding.

Finally, as regards the third point, it will continue to be an impossibility for the theory of cognition as well as for the phenomenology of spirit to determine the relation between imperfect and perfect knowledge. In order to treat the question of the possibility of an absolutely perfect knowledge, subjectivity, in the course of a progressive logical self-development, must not merely have learned what transformations finite knowledge with its heterogeneous cognitions can and must undergo; it must, by looking in all the tasks of knowledge at the principle of knowledge, also have gained such an insight into the relation between the original and the derived that in its derived knowledge it can everywhere see a reflected splendor of the original.

Whether consciousness, as at the standpoint just discussed, keeps exclusively to the finite and sees in the ideal an unclear fancy, a contradiction; or whether, as at the theological standpoint, it proceeds from an immediate and consequently dualistic opposition of the human and the divine, and so comes partly to an abstract and vacuous, partly to a thoughtless and self-contradictory conception of omniscience, in either case the ideal of knowledge is violated. But the ideal is then also violated by the philosophy of the absolute giving itself the air of having attained an absolute knowledge of the absolute.\footnote{``In dem Wissen hat der Geist die Bewegung seines Gestaltens beschlossen, insofern dasselbe mit dem unüberwundenen Unterschiede des Bewußtseyns behaftet ist. Er hat das reine Element seines Daseyns, den Begriff, gewonnen \dots\ Indem also der Geist den Begriff gewonnen, entfaltet er das Daseyn und Bewegung in diesem Aether seines Lebens und ist \emph{Wissenschaft} \dots\ Wenn in der Phänomenologie des Geistes jedes Moment der Unterschied des Wissens und der Wahrheit und die Bewegung ist, in welcher er sich aufhebt, so enthält dagegen die Wissenschaft diesen Unterschied und dessen Aufheben nicht, sondern indem das Moment die Form des Begriffs hat, vereinigt es die gegenständliche Form der Wahrheit und des wissenden Selbsts in unmittelbarer Einheit. Das Moment tritt nicht als diese Bewegung auf, aus dem Bewußtseyn oder der Vorstellung in das Selbstbewußtseyn und umgekehrt herüber und hinüber zu gehen, sondern seine reine, von seiner Erscheinung im Bewußtseyn befreite Gestalt, der reine Begriff und dessen Fortbewegung hängt allein an seiner reinen \emph{Bestimmtheit}.'' Phänomenologie des Geistes. Hegel's Werke. vol.~2. Berlin 1832 p.~609.}

\opage{252} That the ideal of knowledge is violated by the last of these ways of viewing as well as by the first is evident. However much we exert ourselves toward an intuitive knowledge of the absolute, the discursive, the conditioned and relative nonetheless assert themselves everywhere, not merely in the sciences but in science. If the knowledge that philosophy has attained at a given stage of development is to be passed off as an absolute knowledge in the sense that the ideal itself has been attained, then two consequences harmful to science spring from this. In the first place, the disproportion between the actual achievement and the ideal demands of knowledge---even if one imagined a philosophy as complete as that which the greatest philosophical genius might one day be able to develop toward the end of days---will nevertheless be so conspicuous that it must tempt acuteness to accuse philosophy of an unseemly lack of self-criticism. If there is any science that is strictly called upon to take note of the limits of human knowledge, it is surely philosophy.

Secondly---and this drawback is by no means lesser than the preceding one---the unclarity that always arises when philosophy confuses its knowledge of what
\opage{253} which belong to a perfect knowledge, with the perfect knowledge itself, have a dulling and stupefying effect on philosophical inquiry. Instead of veiling what is deficient and imperfect in all our knowledge, the inquirer is on the contrary called upon to bring the problems forward. The progress of our knowledge is not such that the remaining problems should become fewer, easier and less significant the nearer we approach the ideal; on the contrary! the more clearly the ideal is seen, the higher the problems become, and the deeper runs the stream of investigation.

\begin{center}{\bfseries γ)\quad The Ideal of Knowledge.}\end{center}
\addcontentsline{toc}{paragraph}{γ) The Ideal of Knowledge}

When intuitive knowledge is in an immediate manner conceived and determined as the ideal of knowledge, then perfect knowledge must exclude all comprehending; for no comprehending is thinkable sine abstractione, sine discursu, sine ratiocinio: the All-knowing One of the dogmatists knows everything, but comprehends nothing. If, on the other hand, knowledge is to be developed into a comprehending, then it becomes discursive; but when discursive knowledge posits the intuitive as something subordinate, draws it down into finitude and uses it as an element, then comprehension itself becomes so finite, mediate and restless that the ideal vanishes. What, then, is here the universal, knowledge or comprehending? It is evidently knowledge; for there can be a knowledge without comprehending, but not, conversely, a comprehending without knowledge. But when knowledge is the universal and comprehension the particular, then one must investigate more closely how this particular then relates to its universal. As what is particular with respect to knowledge, comprehension forms an intermediate form between the two extremes of knowledge: the beginning and the consummation. With respect to the beginning, comprehension stands higher than knowledge; with respect to the consummation, knowledge in turn stands higher than comprehension. That knowledge stands higher than comprehension thus by no means signifies that comprehension is excluded, but rather that the mediate is elevated and transfigured in a higher immediacy. ``Should anyone, however, \opage{254} say: our cognition is thus only immediate because it is so infinitely mediate, such an utterance would surely easily be taken as proof of hair-splitting and dialectical wordplay. Yet the validity of the proposition advanced will become evident precisely immediately, once the understanding of it has first been mediated. In a chain of connected links it certainly holds that, while the first and second links are in immediate connection, the first and third are only in a mediate, the first and fourth in a still more mediate connection. But with such a parceling out of immediate and mediate the concept `inner infinity' is incompatible. In intellectual connection one does indeed also speak of the connection of extreme terms by middle terms; but the deeper insight goes, the more the finite middle terms, as determinations of concepts, are dissolved in the intimate unity of the extreme terms. In order not to be dazzled by the mediacy in the middle terms that come in, one must look at the action of light and make clear to oneself the difference between the transparent and the opaque.''\footnote{Philosophie og Mathematik.\ Kjøbenhavn 1857 p.~41.}

In order to make clear the difference between the transparent and the opaque, we must above all note that transparency is determined not merely by the object that is to be seen, but also by the sight. If there are different kinds of sight, then there must also be able to be different kinds of transparency. Thus something may very well be transparent to thought and to pure insight without its therefore being transparent to the eye. When the botanist contemplates a tree, he is surely just as little able either to see through the trunk or to see its root through the earth as any other observer; nonetheless the whole organism of the tree will be transparent to the botanist in an entirely different way than to the uninformed; the same holds when the mineralogist contemplates a mineral, etc. But such relative \opage{255} differences are, however many gradations they may run through, insufficient to give us a concept of the knowledge that bears everything that is. The fundamental question is: what transformations must the knowledge known to us undergo in order that the ideal can be thought realized, and on what conditions are we able to see the possibility of such transformations? The answer to this question will depend essentially on how we more precisely conceive and determine what we have above called the ontology of sense-perception and of the senses.

What ontological significance, then, can be ascribed to sense-perception? Sense-perception is the necessary and therefore universally valid intermediate determination between the pure ideality of knowledge and the immediate reality of the object. That immediate existence is grounded in reflection, and reflection in reflecting subjectivity, has in the foregoing been illuminated from so many sides that in this connection we can hold fast to the result obtained as an indisputable fact. Subjective reflecting and the objective immediacy of existence relate to one another as the logical to the antilogical, and to that extent furnish heterogeneous extremes. The middle term that is to unite these extremes by expressing the natural, i.e.\ immediate, unity-in-difference of the heterogeneous factors must therefore double itself in such a way that it immediately and at one and the same time expresses the antilogical in logical and the logical in antilogical form; here we have sense-perception and the sensible. Whichever of our senses we choose as an example, we are convinced that the matter stands thus. In existential immediacy the object of sight stands outside sight; only the image of the object is in sight. Here there is in sight, as in every sense-perception, a peculiar relation between essence and actuality. Essentially the one who sees sees only the image, not the object; but in actuality he sees the object, not the image: the object is visible, and yet this visible thing is not the object. In order \opage{256} to see how this contradiction is resolved, one must see how it arises. It is an unreflected reflection, an immediate reflecting, on which all sense-perception rests. Every sense-perception is determined immediately insofar as it is determined by the impression; the reality of the sense-impression lies in the antilogical: in this respect it holds that what is seen is the object itself, not its image. But since sense-perception takes place in subjectivity and rests on an immediate, and therefore unconscious, reflection, the sense can be determined by the impression only to do the determining itself; the ideality of sense-perception lies in the logical: in this sense it holds that what is seen is not the object itself but its image. On the immediate conversion of the antilogical into the logical, and conversely, rest both the contradiction and its resolution. From what is peculiar in this conversion it is at the same time explained that the senses can be so heterogeneous among themselves: the heterogeneous stems, namely, from the immediate determinateness that is inseparable from the antilogical element.

We have here a point of departure for judging the transformations of sense-perception. That this world, in its multiplicity of sensible things, can have actual validity only insofar as it has validity in and for the Idea, is evident. But if there is now a system of sensible determinations of existence which in the Idea of Knowledge has validity for the knower, then it follows from this with irrefutable consistency that in the Idea of Knowledge there must also be a system of peculiar sense-perceptions and determinations of sense-perception that has validity for the knower; for sense-perception and the sensible are categories that stand and fall together.

From the logical standpoint of idealism one has not hesitated to ascribe thinking, comprehending and knowledge to the Idea: ``Die absolute Idee allein ist Seyn, unvergängliches Leben, sich wissende Wahrheit und ist alle Wahrheit\footnote{Hegel Die subjektive Logik.\ Werke vol.~5.\ Berlin 1834 p.~238.}''; further, ``Die \opage{257} Methode ist daraus als der sich \emph{selbst wissende}, sich als das Absolute, sowohl Subjektive als Objektive, zum \emph{Gegenstande habende Begriff}, somit als das reine Entsprechen des Begriffs und seiner Realität, als eine Existenz, die er selbst ist, hervorgegangen.'' But if one is in earnest about placing an absolute knowledge into the Idea as Idea and letting this knowledge have its object in existential actuality, then it will appear on closer consideration that knowledge in the form of thinking and comprehending cannot possibly apprehend actual existence unless sense-perception forms a mediating transition from the antilogical reality of the object to the logical essentiality of pure knowledge. If one dare not grant sense-perception ontological validity because, according to our own immediate experience, it is bound to psychological conditions, then it is mere phrases and inflated turns of speech when, on the ground of the ontological validity of thought, one would acknowledge ``the knowing truth,'' ``the concept knowing itself in everything''; for if one holds exclusively to immediate experience, then thought surely presupposes its psychological conditions just as much as sense-perception does.

What sort of transformation of the psychological conditions, then, does the ontology of sense-perception essentially presuppose? It is, according to what we have already shown in an earlier connection, the finite separation of knowledge and power, upon whose sublation the emphasis must here fall. On the psychological standpoint the different sense-perceptions are essentially conditioned by different ``affections.'' These affections all point to a dualism of knowledge and power. The power lies in the object that affects us; the knowledge is determined by the affected sense. In this respect it must be granted that sense-perception and sensibility belong to what is lower in our nature, thought with its ``functions'' and concepts to what is higher.

But on the ontological standpoint this dualism is \opage{258} removed. For ontological subjectivity knowledge is power and power is knowledge; but sense-perception is by no means thereby sublated; on the contrary! inasmuch as the unity of power and knowledge at the same time expresses an opposition, sense-perception designates what is immediate in this unity-in-opposition, seen from the standpoint of knowledge. What intermediate determinations ontological sense-perception here furnishes can easily be shown. By itself the determination of power expresses the antilogical reality of existence; thus the single atom, considered by itself, is an existence-point; this existence-point is furthermore a material, a sensible point, hence a sense-perception-point. If the sensible existence-point were not at the same time a sense-perception-point, then there would be here no unity of power and knowledge; but in consequence of the unity of power and knowledge it holds: as many sensible existence-points, so many ontological sense-perception-points. The existence-points are furthermore in various relations to one another (gravity, cohesion, etc.); but it has shown itself in the foregoing from many sides that all objective relations rest on subjective reflections. The immediate unity of relations and reflections is then in this connection determined from the side of power by a system of sense-perception-points, from the side of knowledge by a system of sense-perceptions; to each sensible relation (a pressure, a movement of falling) there corresponds an ontological sense-perception. That it is not sensible terms of relation which, in their mutual relations, may be assumed to sense one another or to have sense for one another, is understandable from what has been treated under ``reflecting subjectivity'' concerning the reciprocal determination of relation and reflection. It is ``the hovering-over reflections'' which here, in accordance with the antilogical factor taken up on the standpoint of the concept, have received so peculiar an increment of existential immediacy that precisely from the side of this immediacy they are to be conceived as ontological sensations.

We can illuminate the matter from yet another side. Among the philosophers who have come closest to \opage{259} discovering the ontology of sense-perception and of the senses we must mention Leibnitz in particular. As is well known, Leibnitz lets matter consist, seen objectively, of ``naked monads,'' seen subjectively, of obscure, confused representations. It is evident that Leibnitz in the obscure perceptions (perceptions petites) has come to touch upon the ontological unity of sense-perception and the sensible; this touching upon it is here all the more to be emphasized as it is punctual and, by its punctuality, fitted to furnish presuppositions for more special analyses; but two misunderstandings grounded in the doctrine of monads come along with it and render the hint unusable. The first misunderstanding is that matter, as a product of the obscure, confused representations, should be a kind of result of the impotence of ideality and knowledge, that matter, as a consequence of this, should have no validity for the divine knowledge; this is the old misunderstanding that has hitherto repeated itself at all stages of spiritualism and idealism. The second misunderstanding is that the obscure representations corresponding to the sensations should be, as it were, manifestations of force of the existence-points themselves. Such a view must lead to the greatest absurdities. Were the existence-points to be self-representing representations, then it would surely follow from this that every object consisted of representations and that everything represented; but this proposition is just as absurd as the proposition ``that every object consists of thoughts and that everything thinks.'' The difference between manifestations of force and ontological sense-perceptions is in immediacy what the difference between relations and reflections is in reflection.

The existence-point, the sense-perception-point determined by power, is, as the object of knowledge, a limit for knowledge; all knowledge, even that which in the Idea is divine, has a limit insofar as it has an object. There must in all knowledge be something that is known; if what is known is only knowledge itself, then knowledge is empty, and empty knowledge is in the act of transforming \opage{260} itself into non-knowledge. In order that knowledge can have reality, what is known must have reality; but when what is known is to have reality, there must here be at one and the same time a \emph{content} of knowledge and an \emph{object} of knowledge. A closer investigation of this will lead us to an insight into the ontological reciprocal determination of sense-perception and reflection, and thereby to a clearer understanding of the relation between comprehending and knowledge.

On the necessary unity of sense-perception with reflection rests ontological objectifying. This objectifying has, as was indicated in an earlier connection, two main sides: the idealistic, corresponding to the \emph{ideal system}, and the realistic, corresponding to the \emph{real system}. From the idealistic side objectification is logical and free of all limits; every limit is in this sphere transformed into a peculiar moment of determination in the infinitely transparent content of knowledge; here in this circle there are no objects: the adequate ideas everywhere and from all sides take the place of the objects. In this sphere there are then also no sense-perceptions; for the particular, diverse and peculiar sense-perceptions vanish everywhere and from all sides as illuminations, alterations and shadings by which the ideality of the content is articulated in inner infinity. It is otherwise when we look at the real system. From the realistic side objectification is antilogical and unfree, bound and determined by countless limits. In this sphere knowledge has not an immanent content, a system of clear ideas, but an infinite aggregate of real oppositions, a world of objects, hence a world of limits; instead of the pure reflections we here get everywhere and from all sides nothing but manifestations of sense, sensations of reason, ontological sensations.

The ontological sensations relate to the pure reflections as the moments of the object to the essential moments of the content. Through each moment of content knowledge receives an ideal determinateness, through each moment of the object a real limit. Just as the content consists in a system of \opage{261} moments of content, so the object consists in a system of moments of the object. As the moment in question is, so is the corresponding totality: just because the moment of content is transparent in its ideal determinateness, the content itself is transparent; conversely! just because the moment of the object is and, in accordance with its character, must be opaque, the object itself, as object, is opaque. To extol an ideal of knowledge that in its kind should be so exalted that the knower immediately saw through every object of knowledge is to extol one's own thoughtlessness.

``But,'' it will perhaps be objected, ``when the actual objects in their actuality cannot do other than furnish limits even for the divine knowledge, then there surely remains here a multiplicity of determinations of the object, hence a multiplicity of determinations of being, which must fall outside knowledge; then the ideal is surely violated after all, ontological subjectivity denied, and the principle of knowledge abandoned.'' This objection rests on a misunderstanding that can be removed without difficulty. Certainly every object, indeed every moment of the object, must set knowledge a limit; thus every material atom must set a limit even to omniscience itself; but one must not forget that the restriction which knowledge thus receives through each moment of the object is a negation that belongs within real knowledge itself, and precisely thereby satisfies one of the essential demands of knowledge. The limit is a determination of power; but the power which thus sets knowledge limits is, in unity-in-opposition with knowledge, a categorial expression of the very reality of knowledge. Every determination of the object by which knowledge is restricted is identical with a corresponding determination of content by which knowledge is extended: that the restriction of knowledge at all points conditions its extension is an eloquent testimony to the infinitely overreaching ideality of real knowledge.

Insofar as the moment of the object in its immediacy is \opage{262} a limit for knowledge, matter with its atoms and masses will surely be impenetrable to the one who in the Idea knows everything; but insofar as every moment of the object is a moment of power, and as moment of power a real negation, a determination of essence in and through the divine knowledge itself, dark, impenetrable and opaque matter can nevertheless in the Idea be seen through by powerful knowledge, cognized and comprehended by knowing power.

The contradiction that the opaque is nevertheless transparent finds its resolution at all points in the teleological reciprocal determination of the ideal and the real system. From this reciprocal determination it follows, namely, and that, be it noted, not merely in abstract universality but also in everything special, single and punctual, that what in the antilogical form of existence is moment of the object and limit is, in the ideal form of knowledge, moment of content and clarifying determinateness.

The opposition between the transparent content of knowledge and the object impenetrable to knowledge receives yet a new and richer illumination when we look at the shadings, alterations and transitions that the peculiar reciprocal determinations of reflections and sensations furnish for the equalizing of every possible tension between the demands of knowledge and of power. If the categories of knowledge are bound and to that extent hidden in the determinations of power, then the categories of power, on the other hand, are liberated and clarified in ideality as transparent determinations of knowledge; but between the extremes of perfect ideal clarity and complete material obscurity there now falls the system of intermediate determinations that is articulated precisely by the transitions of reflections and sensations into one another.

The more sensible the relation is, the more immediately does the sensation of reason, ontological sense-perception, coincide with the manifestation of force; the more elementary the manifestation of force is (e.g.\ the cohesion of solid bodies), the more obscure is \opage{263} the corresponding sensation: the sensations can, as regards such a system of material relations and movements of form, lose themselves in differential effects of different order. Seen from the standpoint of the real system, the light in the divine knowledge can thus run through countless gradations, all the way from the complete darkness determined by the impenetrability of the object, through the twilight of the sensations and the daylight gleam of the sense-reflections, to the highest clarity of ideality.

Hereby the unity of comprehending and knowledge is now attained. The real system in this sense establishes the essential significance of \emph{discursive knowledge} for the ideal of knowledge itself. From the standpoint of the ideal system, on the other hand, the stated oppositions of light and shadow have only the effect that they endlessly bring new determinations of knowledge to the realm of ideas, namely such as develop everywhere and from all sides through the existential formations. There is thus in the circle of power, in the world of things, not an object, indeed not a moment of an object, that does not have in the ideal system its corresponding determination of knowledge in perfect clarity. Since now the ideal system, with its whole content of categories, concepts, moments of ideas and ideas, is perfectly transparent to knowing subjectivity, it is seen from this that the ideal of knowledge also in itself assumes the character of \emph{intuitive knowledge}.

In the teleological reciprocal determination of intuitive and discursive knowledge the ideal of knowledge has infinite reality, and with a logical demonstration of this reality the logic of subjectivity must conclude.

\begin{center}---\end{center}

\opage{264} \begin{center}
{\large\bfseries Chapter Two.}\\[3pt]
{\Large\bfseries Functions of Subjective Knowledge: The Logic of Objectivity.}
\end{center}
\addcontentsline{toc}{section}{Chapter Two. Functions of Subjective Knowledge: The Logic of Objectivity}
\markboth{The Idea of Knowledge}{The Logic of Objectivity}

\noindent From the very beginning the question of objectivity has been drawn into the treatment of the problem of subjectivity; thus there runs through the whole of the previous chapter a development of the objective corresponding step by step to the development of the subjective. The opposition between this and the preceding chapter is, however, determined by the fact that, just as in the logic of subjectivity the objective has furnished only element and point of transit for the analysis of the subjective, so now in the logic of objectivity the subjective is to furnish element and point of transit for the analysis of the objective. In this there is, so far, nothing strange. On the other hand, one might perhaps find it strange that the logic of objectivity is referred to functions of a subjective knowledge; one would surely expect that a logic of the objective would have to rest precisely on functions of an objective knowledge. The matter, however, stands thus. As long as the objective is merely treated in general, its peculiar character must necessarily be volatilized; for the universal is not, as has all too one-sidedly been emphasized in opposition to the psychologically subjective, the objective, but on the contrary what is in its essence originally subjective: to treat the objective in general is to treat it in a subjective manner, is, in other words, to treat it not as something objective but as something subjective. If, e.g., Idea and actuality are set against each other, then the Idea, although it is itself the unity of the subjective and the objective, nevertheless immediately passes over to the subjective side, while actuality remains standing as the particular expression of the objective. The reason for this is easy to see. The distinguishing mark of the objective is everywhere otherness, singularity in its immediacy, the externally manifold, the diverse, etc. Now there can surely be said in general \opage{265} a good deal, even a good deal that is apt, about the manifold, the diverse and the special; but the manifold in its multiplicity, the special and the diverse with their singularities and peculiarities, one cannot possibly dispose of by a treatment in general. The manifold demands to be determined in a manifold manner, the special in a special manner: here, in the logical sense, lies the difficulty with the objective.

The distance between a presentation of the objective in general and an exhaustive presentation of what it essentially is in its existential peculiarities is so great that an intermediate sphere must necessarily be sought here between the extremes: in this intermediate sphere lies precisely the logic of objectivity. From this standpoint we must make an essential distinction between the logic of \emph{objectivity} and that of \emph{objects}. By the logic of objects we designate the outermost limit to which a special development of the special can be carried without thought's thereby lacking the clear inwardness, the insight into the essential unity of the manifold, the transparency of the moments, which the logical can never renounce. In the logic of objects the gaze is directed punctually to the different forms and shapes that the objective, what is brought forth by objectification, can assume in the different spheres of the objective world; the functions of knowledge that correspond to this are then also to be determined as functions of objective knowledge. In the logic of objectivity, on the other hand, attention dwells on the peculiar points of view and the particular, not arbitrary but necessary, acts of subjectivity under which the objectification of the manifold and diverse takes place; the functions of knowledge that correspond to this point in all particulars from the objective to the subjective, and are for that reason to be determined as functions of subjective knowledge.

What difference it makes whether it is the objects in their singularities, or the acts of objectification \opage{266} valid for a particular circle of objects, to which attention is directed, can easily be illustrated by a few examples. The ellipse is a mathematical, feldspar a mineralogical object; such objects are now treated in the sciences concerned with such exclusive regard to the peculiarities that obtain that there is not even reflection on what it properly means that they are objects. For a special science can very well treat its objects logically without its therefore needing to occupy itself with treating their logic. The logic of objects expressly presupposes that what makes the objects objects becomes in particular an object for reflection. Since now the objects, in their capacity as objects, can be of different kinds and belong to highly different circles of existence, it becomes a task for the logic of objects to show the peculiarities by which the diverse objects belonging under different circles enrich the concept \emph{object} with new and essential determinations; in this respect the mathematical object is then of an entirely different kind from the mineralogical.

That there must now be a connection between the logic of objects and that of objectivity is self-evident, and can likewise be seen from the examples cited. In order to grasp the objective difference between the mathematical and the mineralogical object, one must necessarily give oneself an account of the subjective difference that lies in the manner in which objects of such different kinds are objectified; but the investigation of the special relations in which sense-perception, imagination and thought must work together in the objectification of the diverse objects belongs precisely under the logic of objectivity.

But just as we see the logic of objects and that of objectivity in their inner union, we see at the same time the difference and opposition that assigns each of the tasks named its place in the logical system. The logic of objects, namely, is not content with a general demonstration of the properties that distinguish the objects insofar as they belong to different spheres; \opage{267} it goes into the particular and investigates the differences by which objects of the same sphere are opposed to one another, as well as the relations, series and transitions that they essentially form with one another in accordance with their mutual kinship. Thus the logic of objects, if the ellipse is to be chosen as an example, will by no means confine itself to characterizing the ellipse as a mathematical object in logically objective opposition to feldspar as a mineralogical object; no, the logic of objects will expressly show the relation between the closely related and yet mutually different objects, such as, as regards the mathematical: the ellipse, the hyperbola and the parabola; as regards the mineralogical: the diamond, graphite, etc. The interest rests, to keep to the mathematical, naturally not on the specifically mathematical problems, for logic is not mathematics and can never become it; but the specifically mathematical determinations here furnish the elements through whose working-up logic becomes able to articulate the development of the concept of the object in the sphere concerned. The logic of objectivity cannot enter upon such analyses of singulars\footnote{The manner in which the proposition: ``When the understanding operates on the basis of the imagination, mathematics comes about,'' is illustrated in my lectures on ``Phil.\ Propæd.'' from the academic year 1861--62 p.~29--34, can in this connection be said to belong under the logic of objectivity; whereas the investigation of whether, if the formation of a planetary orbit were to depend on accidental circumstances, the greatest probability would be for a parabola or a hyperbola or an ellipse (Anfr.\ Skr.~95--98), must be referred to the logic of objects. How precarious it is when a philosophy that has not been properly attentive to what a logic of objects demands would treat the objective can be seen from Hegel's objection to the sudden reversal in centrifugal force. ``What Hegel adduces as the mathematical view, that the centrifugal force should decrease from aphelion to perihelion, is precisely the opposite of what mathematics proves, namely that it increases from aphelion, the point at which it has its minimum, to perihelion, where it precisely reaches its maximum; of a sudden transition --- `ein plötzliches Umschlagen' --- there can here be no question at all.'' (Anfr.\ Skr.\ p.~64--65). Had Hegel, before he undertook to judge of the changes in centrifugal force, first investigated what sort of mathematical object the ellipse is, he would soon have discovered his misunderstanding.}; since it, itself \opage{268} still belonging under subjective knowledge, is to form the transition to an objective knowledge of the objective, it keeps to an intermediate stage, so that it uses the objective for a more penetrating specification of the subjective, while nevertheless at the same time, conversely, it applies the subjective to a more far-reaching articulation of the objective.

What this means can be illustrated here by a retrospect; for the logic of objectivity is also to be determined as a continued development of what we in the previous chapter have called ontological objectification.

The ontology of objectification is built on the following foundation. The finite separation between the subjective and the objective that characterizes the psychological standpoint vanishes in the logic of the fundamental ideas; but the opposition does not vanish. The Idea as Idea is indeed ``the unity of subject and object,'' but since the unity is a unity-in-opposition, the opposition is preserved. For the psychological conception the sensible falls outside sense-perception, the thinkable outside thought, the objects outside objectification. Therefore from the psychological standpoint it looks for all the world as if there could very well be a world of sensible things without anyone sensing them, a multiplicity of existing objects without anyone objectifying them. It is otherwise in ontology: here a sensible that is not sensed, a thinkable that is not thought, an objective that is not objectified, are only hovering possibilities, subjective intimations, empty representations without reality. So far is it from being the case that an objective world should \opage{269} be able to have actuality without an objectifying subjectivity that it would not even be possible unless every object by itself were expressly determined by a corresponding objectification, and such an objectification corresponding to each object would not be possible unless there were in the act of objectification itself just as many particular determinations of subjectivity as there are objective sides in what is objectified.

It is only superficiality and half-measures when in logic one rests content with an explanation such as this, that the Idea, ``the knowing truth,'' objectifies everything to itself in general by a general thinking; even if, in order to make some concessions to natural sense, one adds intuition and lets ``the Absolute'' itself intuit everything in a general intuiting, it is emptiness nonetheless. In the Idea itself, the concrete Idea, objectification has, as was said, the same peculiarities in the subjective form that the object has in the objective. The sensible object is not objectified by a general sense-perception, for a general sense-perception is just as unactual as a general sensible; no, the sensible, this determinate sensible, which unites in itself the many different and diverse sensible properties, is expressly objectified by a sense-perception that unites in itself equally different and diverse sense-perceptions: not a single alteration can come forth in what is objectified without a corresponding alteration having to take place in the objectification.

That the objects in actuality come forward and force themselves upon us without our noticing the least thing of the ontological objectification is in this connection not strange. How much self-observation and reflection is not required in order that the sensing individual, even merely from a psychological standpoint, may become aware that in all feeling, seeing and hearing there is an objectifying? Yet another example. How easily and naturally can the visible objects divert our attention from the light in which we see them; it \opage{270} seldom occurs to us that in visible things we have to do not immediately with the things, but only with the visual images. The apprehension of what is nearest of all, of what is most immediate, here demands the most strenuous reflection. ``If the ether in the whole of cosmic space were at rest,'' says the physicist, ``complete darkness would reign everywhere''; so important a role, then, does the vibrating ether play in cosmic space; and yet no mortal, not even the most sharp-sighted physicist, can immediately see either the ether or its vibrations.

To become aware of the objects is an easy matter; but to become aware of ontological objectification there is needed a reflection guided from many sides and stimulated by many logical difficulties. The Hegelian logic is a great example of the fact that there can be an acute observation of objectified movements of concepts without attention being yet in the least directed to the logical in objectification itself. And yet it makes an essential change, a thoroughgoing difference, whether or not one looks, in what is objectified, at objectification itself. If we overlook objectification, then what is objectified surely shows us, from its side, how existence proceeds from its ground, actuality from its possibility, the effect from its cause, etc.; but how does it go with insight? Seen from the point of view of what is objectified, the ground is in existence, the possibility in actuality, the cause in the effect, the disposition in the execution; and yet there is here an opposition, an essential opposition between ground and existence, between possibility and actuality, between disposition and execution. What, then, does this opposition signify? It is an opposition of the ideal and the real, which constantly comes forth through, and essentially rests on, a doubling in what is objectified itself. Existence, this existent, has the ground in itself, that is, existence is at one and the same time both ground and existence: \emph{existence}, insofar as the existent has sublated its ground; \emph{ground}, insofar as
\opage{271} the existent in turn goes to ground and thereby grounds another existent. In a corresponding way the actuality in this actual thing is at once both possibility and actuality, and so forth. How, then, is this doubling really to be thought of as taking place? Through an incessant conversion, a transformation! Actuality, itself a transformation of possibility, an expression of the consummation in which the real possibility vanishes, is surely actual only insofar as it is precisely in the midst of a transformation whereby it becomes possibility again. And on what, then, does the transformation rest? On the immanent negation! Here the problem is at once recognizable; for the immanent negation, when we look at what is objective in objectivity, is a speculative fantasy: the more one urges, in an existential sense, the immanence of negation and negativity, the more high-flown the negation becomes, the more it betrays its subjective origin. If one sees this, while one nevertheless still goes on staring at what is objectified and overlooking the objectification, where does one then end up? Why, then one arrives at such fine explanations as this one, that the categories ground, possibility, disposition, etc.\ are not merely subjective, but subjective in the sense that they concern only our psychological consciousness of things, without in the least concerning the things themselves. Instead of proceeding from the ground and being grounded, the one existent then comes, without further ado, to proceed from the other existent; instead of arising from its possibility, the one actual thing must now arise from the other actual thing, and so forth. Now since the thought that the one existent should be able to proceed from the other without an intervening ground is just as unthinkable as that the one actual thing should be able to arise from the other actual thing without an intervening possibility, one is compelled, in order to avoid ``the immanent negation,'' to reject the proceeding of the existent from the existent, the arising of the actual from the \opage{272} actual; and where do we then end up? At local displacements, merely external changes! Yet here too one cannot rest content; for it has been amply shown in what has gone before that even a spatial movement, a change in what is external, is impossible in the objective sense when one abstracts from the objectification.

Only in the ontology of objectification is the living unity of knowledge and power cognized, in which alone the reciprocal determination of the subjective and the objective in its inner infinity can be grounded. Objectification is in this respect determined equally by necessity and by freedom. In the real system, in the sphere of power, necessity reigns: objectification is determined by the objects, partly coming to be, partly existing; in the ideal system, in the sphere of knowledge, freedom reigns: objectification is determined by the categories, moments of the idea, and ideas that develop and move within knowing subjectivity, which objectifies everything in knowledge. To objectify the concept as existence presupposes necessity: hereby knowledge perishes in power; to objectify existence as concept presupposes freedom: hereby power is transformed into ideality and transfigured in knowledge; to objectify the concept as existence in order, under this presupposition, to objectify existence as concept, and conversely, presupposes the teleological: hereby freedom passes in an infinite way over into necessity, and the latter into the former.

These results, won by the analyses in the preceding chapter, thus furnish the foundation; the further development, the development whereby the logic of objectivity is specifically determined, now depends on the way in which the specifications are carried through. At this, as at every turning point, it is a matter of choosing the right point of view and finding out the procedure corresponding to that point of view. The task is to master multiplicity, to let the distinct acts of objectification come forth in such a way that the essence of otherness, the \opage{273} world of objects, can be made clear to the objectifying selfhood. The appearance in which the objective is determined point by point subjectively, and the subjective objectively, then gives us objectivity. Objectivity is thus two-sided; the two-sidedness presents a double point of departure, namely a point of departure in the objective, which conditions a further development of the subjective, and a point of departure in the subjective, which in turn conditions a punctual analysis of the determinations by the objective.

If we now designate the distinct act of objectification, be it an act of thought or an act of nature (the act of nature, be it noted, considered from the standpoint of subjectivity), by the category \emph{function}, then in the points of departure indicated we have the following directions, characteristic for the treatment of functions. The objective point of departure lies in the real system and leads to a logical development of the essence of functions and of their diverse peculiarities, that is, to a \emph{logic of functions}; the subjective point of departure belongs to the ideal system and leads to a special exposition of the intellectual acts necessary for the objectifying activity of subjectivity, in other words, of the \emph{logical functions} in the proper sense. Finally, since the unity-in-opposition of the real and the ideal system likewise requires its punctual development, the logic of objectivity must close with an exposition of the totalized reciprocal action of the objective and subjective moments in the \emph{systematics of functions}.

As regards the \emph{logic of functions}, it is of decisive importance for the analysis that the objective determinations, antilogical in themselves, are all of essentially subjective, hence of logical, significance. This is by no means the case when attention fixes on the objects. Every specification of the peculiarities of an existing object will always lead to a conflict between an essential and an inessential, a necessary and a contingent, a universal and an individual, whose \opage{274} prolixity can easily lead logic astray, since, insofar as it does not then by its forbidding character tempt to untimely abstractions, it only all too easily leads it away from total cognition. It is otherwise when we consider functions. The specific peculiarities that essentially characterize a function are, however much the punctual determinacies may concern the antilogical, nevertheless of logical significance even in the finest details; for the universal predominates. As regards \emph{the logical functions}, it is of importance that an articulation of the objective shine through at every point; only on this condition will it be possible to determine the distinct intellectual acts as peculiar functions, and only by what is peculiar in these functions possible to demonstrate the overreaching power of objectifying subjectivity over the objectively manifold. Finally, as regards the \emph{systematics of functions}, the logical endeavor must aim at working together what is diverse, demonstrating the ideal reciprocal determinations of the correlate functions, and thus leading the developed multiplicity of objectivity back to the principle itself, the principle of knowledge, which is to be determined from all sides as the original unity of the subjective and the objective.

\begin{center}
{\bfseries A.}\\[2pt]
{\large\bfseries The Logic of Functions.}
\end{center}
\addcontentsline{toc}{subsection}{A. The Logic of Functions}

In the determination of the objective a tinge of the subjective shows itself everywhere; however much subjectivity is affected and determined by objects, it still lies in its essence that it can be determined only to determine itself. It is not merely a sign of the limitation of human subjectivity that it finds itself so disposed to treat objects as subjects and to attribute a certain selfhood to what is selfless; it is at the same time a sign of its divine origin, its inner \opage{275} infinity, its overreaching power: it appropriates everything, sets its stamp on everything, and would cognize its own in everything. When the child personifies its footstool, caresses it, and converses with it as with a relative; when a gifted people of antiquity populates rivers and forests with nymphs and dryads, and decorates the starry heavens with images of ``gods and heroes,'' one may of course say that so naive an activity of fantasy does not concern logic. But notwithstanding, the matter still has a side that must, from a logical point of view, attract attention. It is not fantasy alone, but also the understanding, that expressly demands that objects be regarded as acting, suffering, doing subjects. ``In order to be able to speak and judge about everything, we must be able to make everything the subject in a logical judgment. But everything that is posited as subject becomes hypostatized. In sentences such as these: `The stone has killed the man'; `Brutus has slain Caesar'; `the universal requires an example'; `the man demands what is owed him,' the subjects assume the very same hypostatic character, and it is, according to the nature of language, altogether impossible to avoid personifications. So long as the talk turns on finite objects, the misunderstanding is easy to avert; but when the inquiry oversteps the boundary of immediate experience, the separation is so difficult that one finds the most utterly non-self-subsistent hypostatized, and the most utterly self-subsistent transformed into a reflex. And yet in every case it must be possible to determine with certainty whether we are dealing with a merely sensuous substance (substrate), a provisional personification, or an eternal hypostasis, if insight is otherwise to be clear and distinct.''\footnote{Den propæd.\ Logik.\ p.~182.}

If the involuntary personification is to vanish as a moment transparent to consciousness in objectification proper, then there must be had in reflection an expression that \opage{276} at once indicates equally both the subjective and the objective side of the matter. If we say, e.g., that things act, the personification is conspicuous. One may indeed speak of ``an act of nature'' just as well as of ``an act of thought''; but the figurative element in such designations is unmistakable; thought, strictly speaking, can no more act than nature can: action is in the proper sense reserved to free will. Less prominent is the personification in expressions such as: things suffer; the building has suffered much from damp, from lightning, from earthquakes and the like; and yet the category \emph{suffering} is evidently a category that belongs to the subjective. On the other hand, it can be said of things in the proper sense, hence without semblance of personification, that they are moved, changed, act, are acted upon, and so forth; for the categories movement, change and activity belong just as much to the objective as to the subjective. If now an expression of unity could be demonstrated for such categories, and if this expression of unity could further assume such a form that it thereby became fit for a widely ramified articulation, and if this articulation could finally be laid out in such a way that it went in two directions, the subjective and the objective, now the one, now the other becoming decisive, then such a category would surely have to play a considerable role in any logic laid out on the development of the idea. But just such a category we have in the expression, applicable in so many ways and significant for so many turns, transitions and determinations: \emph{function}.

How many-sided and yet peculiar the category function is becomes evident at once when we bring out some characteristic senses that usage has already sanctioned.\footnote{``The word \emph{function} occurs in a very comprehensive meaning. Wherever there can be talk of acting, suffering, doing, function designates a special law-determined activity. It is said of an official that he is in function; one speaks of the function of the eye, of the function of a nerve, etc.\ \ldots\ Pure magnitudes, to be sure, cannot exercise real activity, but by being placed in various connections they can nevertheless essentially alter their mutual relations. A magnitude $x$ can, so to speak, be appointed to various offices; it can, as addend $(a + x)$, as factor $(ax)$, as exponent of a power $(a^{x})$, as logarithm, etc., acquire a very different influence, and to that extent appear in different functions. --- `If a magnitude $y$ depends according to a certain law on other magnitudes $x, z, t, u \ldots$, so that every value of $y$ is determined by the values arbitrarily assigned to $x, z, t, u \ldots$, then $y$ is said to be a function of these magnitudes, all regarded as variable, $y$ dependent, the others independent. This is commonly designated: $y = f(x, z, t, u \ldots)$.' C.\ Ramus. --- According to this mathematical determination of the concept, changes would thus proceed from the independents and exercise their lawful influence on the dependent.'' Phil.\ og Mathem.\ p.~5--6.}

\opage{277} Thus, as regards activity, function is used not merely of activities in concreto, as when it is said that an official is in function, or when the function of a nerve is discussed; but also of the most abstract activities, as when Kant teaches that concepts rest on functions, and by ``function'' understands the unity of the act of ordering different representations under a common representation. Indeed, so comprehensive is the meaning of function that the distinction which Kant on this occasion would establish between ``affection'' and ``function'' proves to be vanishing. Whereas concepts rest on functions, intuitions, in Kant's opinion, were to rest exclusively on affections; but on closer consideration it becomes clear that affections themselves must be determined as functions. For by functions we may, according to usage, think not only of law-determined activities but also of law-determined changes. From this point of view the finite difference between the active and the passive, between acting and suffering, vanishes, so that every suffering is just as much taken up into a corresponding function as \opage{278} every acting. Indeed, we are even led involuntarily a step further; for since an affection, a suffering, can be a law-determined change only in that it expressly corresponds to a determinate acting, every suffering must surely be precisely a function of a corresponding acting. The expression of unity lying in function, the categorial determination that comprises the abstract as well as the concrete, the active with the passive, the universal with the special, the subjective with the objective, is dependence, \emph{law-determined dependence}. In how intimate a connection the concept \emph{function} stands with the concept \emph{dependence}, law-determined dependence, can then be seen expressly from the mathematical meaning of function. ``One cannot say,'' it is stated in mathematics, ``that $2$ is a function of $3$ because $2 = 3^{2} - 7$, for $2$ is also $= 3 - 1$, or $= 3^{3} - 25$, etc., so that $2$ can depend on $3$ in manifold ways. But one can indeed say that $y$ is a function of $x$ by virtue of the equation $y = x^{2} - 7$, and that therein $y = 2$ corresponds to $x = 3$.'' Thus: the condition for $y$'s showing itself dependent on $x$ is that different values be assigned to $x$; the condition for the dependence cited designating a function is that the manner of dependence be more precisely fixed. But the condition for the more precise determination of the manner of dependence is in turn, when we look more closely, this: that the universal function must with inner necessity divide itself into distinct functions: $f(x) = y$ first acquires meaning by being posited as $f(x) = ax$, $f(x) = x^{a}$, $f(x) = a^{x}$, etc.\footnote{On ``the concept of function'' cf.\ Mathem.\ og Dialekt.\ p.~4--27.}

If we now draw a parallel between the mathematical meaning of function and what we may in this connection call its real meaning, the expression of unity, \emph{law-determined dependence}, points expressly to the reciprocal determination of presupposition and consequence. Insofar as an action is to be de\opage{279}termined as the consequence of another action, it is a function of the same; the same holds of action and suffering. Thus not merely can one activity be a function of another activity, but an acting of a suffering, a change of a change, etc. That function in its essence unites the subjective with the objective shines forth from its significance for comprehending and knowledge; for what, pray, is it to comprehend something, if not to see through the function or functions by which the something in question is determined? To comprehend equilibrium is to know the system of functions on which equilibrium rests; to comprehend chemical formations is, were it possible, to reflect through the infinitely intricate system of the functions of affinity; to comprehend the existence of the garnet is to account for the functions of its constituents, which in turn presupposes an insight into the way in which isomorphous substances function in one another's place. Insight into one function is conditioned by insight into another.

But although function according to its concept comprises both the subjective and the objective moments of the matter, the opposition of the subjective and the objective will nevertheless continue to hold for function itself, and that not merely from the standpoint of finitude but also from that of infinity. As regards finite knowledge, we have in ordinary observational psychology and in empirical physiology examples that show how the function of a nerve can be apprehended partly so one-sidedly subjectively that reality, the objective, is lacking, partly so one-sidedly objectively that ideality, the essentially subjective, is lacking. The explicative psychology of observation, remark and assurance makes the matter easy for itself when it passes over the functions of the brain and nerves with the remark that the brain and the nerves are organs, means and conditions for all the many spiritual activities that come into activity in these organs. Despite all as\opage{280}surances about the necessity of the nerves as conditions for the life of consciousness, and whatever else may be remarked about the actions and reactions ``which are hereby to be remarked,'' the subjective explanation of the phenomena of consciousness nevertheless continues to lack any scientific point of support in the objective, so long as not the slightest attempt is made to have the reciprocal determination of the ideal and the material, the psychical and the bodily, investigated in the very essence of nerve function. The subjective one-sidedness and scientific superficiality is recognizable in this, that there is talk far and wide of the many psychical and bodily functions without anything further determinate being thought by a function in general, let alone that any effort should be traced here to deduce the distinct functions in their multiplicity from one or more fundamental functions.

If descriptive psychology holds one-sidedly to the subjective, insofar as, in faith and trust in the originality of the functions of the soul, it holds to the conscious ``living'' in the ``circle of representations,'' with the assurance that nerves and ganglia are conditions, but only conditions, for the psychical functions, then one finds in the already extensive literature of empirical physiology examples enough that show what it means to apprehend a nerve function from a one-sidedly objective standpoint. Without troubling himself further about what is psychical, and hence essentially subjective, in sensibility, a physiologist can then begin by establishing it as a fact that ``sensibility is an inherent property of a ganglionic tissue,'' and thereupon, by a sharp separation between property and function, introduce investigations of what in the matter is externally objective. The property of a nerve can thus be determined as a peculiarity that belongs to it insofar as it is a nerve, but that is not determined by its special connection with some other organ. The property of a nerve will thus ``depend on its structure, just as \opage{281} the elasticity of steel depends on the structure of steel''; whereas the function of a nerve will rest on the use that is made of the property, and depend on the connections in which the nerve stands with other parts of the body. ``The property,'' the physiologist then says, ``is unchangeable so long as the structure persists: sameness of structure everywhere entails sameness of properties; whereas the functions are just as changeable as the character of the organic connections.''\footnote{Cf.\ Forelæsn.\ over ``Phil.\ Propæd.'' academic year 1860--61, p.~367 ff.}

Without dwelling in this connection on the question whether physiology is served by restricting the meaning of the word function solely to activity and action, or not, we can nevertheless readily see from the features indicated that a development laid out in this way will lead to one-sidedly objective results. That we find ourselves on the preserve of the objective is seen at once from the multiplicity of facts that here presents itself for consideration: ``if this nerve is cut, feeling will cease; if, on the other hand, this one is cut, paralysis will set in; are there, then, nerves that according to their `property' are only sensory nerves, or are all nerves, according to their `property,' equally fit to be both sensory nerves and motor nerves, so that the difference that obtains rests solely on the function?'' These are, in their sphere, important investigations, and important results to which such investigations can lead; but even if one imagined all these investigations concluded, all questions belonging here answered in a manner satisfactory to empirical physiology, we would still be infinitely far from the untying of the knot so long as the concept of a nerve function is held fast in so one-sidedly objective a meaning that the subjective moment is more and more lost from sight. To what degree this is the case is seen from the highly disquieting fact that physiology, especially in the very latest time, has set \opage{282} about apprehending the organism as ``a living machine,'' without considering that there [is] a difference of sphere between a machine and a mechanical system of nature, and more than one sphere's distance between a merely mechanical and an organic function.

Yet it is not merely in finitude, with its partitionings and one-sidednesses, that an opposition between the subjective and objective extremes of function shows itself; from the standpoint of infinity too the opposition is seen, but here, naturally, in another and higher light. If the finite presuppositions and conditions are transformed into moments in a thoroughgoing reflection, then both that which in the relation of dependence is determined as the dependent and that which is presupposed as the independent will transform themselves into function. We then have not merely what mathematics, in a more restricted meaning, presents as a function of a function, $y = f(\varphi(x))$, but functions of functions to infinity. Seen from this side, function belongs to the ideal system, and its form will in that respect not be separable from its content. If, on the other hand, the presupposition with its conditions is held fast in finite determinacy, yet in such a way that the whole of actuality is at the same time regarded as subject to the power of the idea, we obtain, as regards the requisites of function, a significant distinction between the something that is in function, the something that is a function of another, and the other of which something is a function; this formal distinction is in turn determined by real connections of diverse functions, existential knots, complexes of functions in a material unity settled by power. From this side, then, the real system belongs to the objective sphere; in this sphere the form of function will in various ways enter into opposition to its content, and this opposition will precisely condition the development of the exact determinations and punctual particularities that are significant for the objective world.

``If one develops, as Hegel in his direction has done, the \opage{283} categories of quantity and measure in logical-dialectical form in such a way that the moments of operation are only fleetingly touched upon, the development of concepts must become one-sided and, at several points, even as regards the universal, untenable. Concepts that have their scientific roots in the soil of exact analyses are stunted when the exact moments must be missing from their dialectic. That a dialectical development of the concepts of quantity does not make up for the lack of insight into the fundamental relations of the operations is seen from this, that while there is indeed a transition from exact analysis to dialectical analysis, there is not, conversely, a transition from general-logical and dialectical analysis to exact analysis: `Non datur metaphysica in mathesin via.'\,''\footnote{Mathem.\ og Dialekt.\ p.~25.} That mathematics, although an abstract a priori science just as much as logic, has quite another grip on finitude, on actuality, hence on objectivity, than ``the speculative logic'' with all its categories of actuality and objective concepts, is explicable when one considers that logic, as this science has hitherto been treated, must always let go of the individual in order to grasp the universal, whereas mathematics knows how, with all punctiliousness, to determine the universal in such a way that it can grasp the individual in it. For a philosophical exposition of the concept the word is used, and the word speaks one-sidedly of the universal; for an exact exposition of functions the language of signs is used, and the language of signs, the language of functions, is precisely suited to take up the individual into the universal, the antilogical into the logical.

Finally, the concept function, which is especially important for a logic of functions, will itself be able to furnish the principle of division according to which the different stages of the development follow one another essentially and naturally. The law-determined dependence by which function is recognized is, as we have seen, specifically designated by the relation between the presupposition and its condi\opage{284}tioned consequences. The presupposition by itself, in its immediacy, is not the function, but an element, a requisite essential for the concept; that the consequence, taken by itself, cannot then be the function either is self-evident when we recall that the consequence is consequence only in reflection of its presupposition, just as the presupposition for its part is presupposition only in reflection of its consequence. But then is it perhaps the manner of dependence, the conditioned relation between presupposition and consequence, in which function must properly be said to consist? This assumption is not satisfactory either; for the manner of dependence cannot possibly be determined when one abstracts from the dependent and the independent. A one-sided emphasis on necessity, the necessity with which the consequence sets in when the presupposition is given, would in any case give us only the law of the function, not the function itself. Function, when it is determined many-sidedly, is at once both the law and the fulfillment of the law; the concept of function is in its unfolding a true concept of totality: presupposition and consequence, necessity and contingency, universal and special, being and becoming, acting and suffering, in short, the most diverse conceptual determinations are taken up as moments in the many-sided development of this concept, a development which for that very reason must coincide with a development of the objective unity of the subjective and the objective, that is, with a development of objectivity.

From the subjective side it has already been shown in our preceding chapter that the presupposition, reflected in negative immediacy, is \emph{ground}, and that the ground, not in the psychological but in the ontological sense, is essentially subjective; in opposition to what is subjective in the ground we have, in the conditions indispensable for grounding, the immediately objective; if grounding is now posited specifically in one of its moments, namely the condition, then a peculiar expression offers itself for the objective unity of the subjective and the objective: the \opage{285} functions belonging hereunder are all to be determined as \emph{functions of the ground}.

It has further been demonstrated at the standpoint of causality what it means that the ground developed into actuality is \emph{cause}. If the presupposition corresponding to the function is now apprehended as ``thing-cause,'' then we certainly have, in the cause itself as thus existing, the one-sidedly objective; but the one-sidedly objective that lies in existence and actuality then also expressly presupposes that the cause is actual. But the actuality of the cause presupposes a possibility from which it has proceeded; in the real possibility of the cause we again have the moment of subjectivity. Since, however, the moment of subjectivity is, as regards the causal relation, otherwise determined by its objective than was the case with grounding in general, the expression for the objective unity of the subjective and the objective assumes so special a stamp that the functions belonging hereunder must be characterized as \emph{functions of the cause}.

Finally, from the point of view of totalizing reflection it has been made clear what it means that a reflected presupposition can in its originality contain a whole totality of ideal-real causal moments, that a presupposition thus totalized in itself is, in other words, a principle, and the principle, as the identity of ground and cause, a cause of causes, hence: a \emph{ground-cause}. As ground-cause the principle has an overreaching power over its conditions and therefore gives its real possibilities an inner peculiar determinacy in a peculiar \emph{disposition}. In the disposition as disposition there is had a clear reflection of the ontological energy of subjectivity; but the disposition has reality only in its execution; with the execution, the peculiar development of totality, the actualization of the disposition's demands, comes the objective, and with this objective a widely ramified system of presuppositions and consequences, relations of dependence and peculiar modes of \opage{286} acting, which give the element of objectivity a decisive significance for the determination of the subjective. Thus apprehended, then, the principle as ground-cause will yield the richest contributions to an essential determining of the objective unity of the subjective and the objective: the functions belonging hereunder, for whose concrete formative activity the functions of lower order furnish the elements, we then designate in what follows as \emph{functions of the ground-cause}.

\begin{center}
{\bfseries a)\quad Functions of the Ground.}
\end{center}
\addcontentsline{toc}{subsubsection}{a) Functions of the Ground}

Founding reflection has taught us in the foregoing what it means that all determinations of ground are determinations of reflection; every such determination of reflection is as a semblance of other in other, a semblance determined by reflected semblance, a pure reflex that cannot be held fast immediately. In reflection itself the ground is one with essence: essence is the formal ground of things. It is the pure, immanent conversion of negative and positive moments, of identity in difference and of difference in identity, that on the whole characterizes the \emph{functions of the formal ground}. In the real ground, on the other hand, immediacy is seen to permeate reflection; the immediate in the ground itself is indeed negatively reflected; but to the negative determinations in the ground correspond the positive ones in existence; from the standpoint of existence the ground thus has essential reality, and can to that extent assimilate to itself all the real moments of actuality. It is by the assimilation of such real moments that we shall in what follows recognize the \emph{functions of the real ground}. Now if, as must plainly follow from this, all actuality-reflections have an essential side from which they show themselves as functions of the ground, and if it is, as remarked above, characteristic of the functions of actuality that they are found united in existential knots in which the most diverse determinations meet, then such an existential knot must, as regards grounding, necessarily \opage{287} point to a peculiar union of diverse grounds. If we now call such a union of diverse grounds a mixed ground, then it will be the diverse determinations of ground that, in the development before us, are to furnish the distinguishing mark of the \emph{functions of the mixed ground}.

\begin{center}{\bfseries α)\quad Functions of the Formal Ground.}\end{center}
\addcontentsline{toc}{paragraph}{α) Functions of the Formal Ground}

So long as the task was to assert the validity of ontological subjectivity against an abstraction in which it was volatilized, we had to become aware from many sides that the part of Hegel's logic which he himself calls the subjective lacks subjectivity; now, by contrast, when the task is to uphold the right of objectivity against a development of thought that dissolves its real systems into abstractly subjective forms, we must with the utmost emphasis recall that the part of logic which Hegel calls the objective lacks objectivity.\footnote{It is remarkable to what degree the subjective in Hegel runs together with the objective; one would almost be tempted to assume that logic, in his opinion, must be the science that from beginning to end varied the ambiguity that the subjective is objective and the objective subjective. We find the ambiguity already in the Introduction under ``Allgemeiner Begriff der Logik,'' where (Anfr.\ Udgv.\ p.~35), concerning the relation between the phenomenology of spirit and logic, it is said, among other things: ``Der Begriff der reinen Wissenschaft und seine Deduktion wird in gegenwärtiger Abhandlung also insofern vorausgesetzt, als die Phänomenologie des Geistes nichts anderes als die Deduktion desselben ist. Das absolute Wissen ist die \emph{Wahrheit} aller Weisen des Bewußtseyns, weil, wie jener Gang desselben es hervorbrachte, nur in dem absoluten Wissen die Trennung des \emph{Gegenstandes} von der \emph{Gewißheit seiner selbst} vollkommen sich aufgelöst hat\ldots\ Die reine Wissenschaft setzt somit die Befreiung von dem Gegensatze des Bewußtseyns voraus. Sie enthält den \emph{Gedanken}, \emph{insofern er eben so sehr die Sache an sich selbst} ist, oder \emph{die Sache an sich selbst}, insofern sie \emph{ebenso sehr der reine Gedanke} ist. Als \emph{Wissenschaft} ist die Wahrheit das reine sich entwickelnde Selbstbewußtseyn, und hat die Gestalt des Selbsts, daß \emph{das an und für sich seyende gewußter Begriff, der Begriff als solcher aber das an und für sich seyende} ist. Dieses objektive Denken ist denn der \emph{Inhalt} der reinen Wissenschaft\ldots\ Die Logik ist sonach als das System der reinen Vernunft, als das Reich des reinen Gedankens zu fassen. \emph{Dieses} Reich ist die \emph{Wahrheit, wie sie ohne Hülle an und für sich selbst ist}. Man kann sich deswegen ausdrücken, daß dieser Inhalt die \emph{Darstellung Gottes ist, wie er in seinem ewigen Wesen vor der Erschaffung der Natur und eines endlichen Geistes ist}.'' But even if one is willing to overlook what is ``Ueberschwängliche'' in the last declaration, and to refer the whole to the sphere that we have in the foregoing designated by the \emph{ideal system}, the whole explanation nevertheless founders on the opposition between the ideal and the real system, an opposition that logic is not to veil but, on the contrary, to unveil.}
\opage{288} In no part of logic does the opposition of the objective and the subjective come forward more sharply and in a more conspicuous way than precisely in the section with which the Hegelian logic begins, namely in the very first section, which deals with \emph{quality}. Quality is immediate determinateness; but just because the determinateness is immediate, the separation of its objective and subjective, its antilogical and logical side, is also immediate. The abruptness, the immediacy in the separation, must be negatively recognizable in the functions. It is, according to the characterization given above, the functions of the formal ground that here come into consideration.

In what relation, then, must the doctrine of \emph{quality} place itself to the doctrine of functions, when the matter is seen from the objective standpoint? Objectively determined, quality in its pure immediacy is so antilogical that it is even incommensurable with language; its immediacy in the sensible is \emph{this} sensible, and this sensible is the unsayable. Quality \opage{289} and the predicate can in this respect be said to form extremes that have the property as their middle. Insofar as the property falls under reflection, it passes over into the universal and identifies itself with the predicate; insofar as the property, as an immediate determination of existence, indicates what is peculiar to the individual thing, it withdraws from reflection and coincides, in antilogical immediacy, with quality. Since each quality is a determinateness immediate for itself, the one quality is without essential connection with the other; the one may well be simultaneous with the other and follow upon the other; but since, as long as each individual quality is apprehended by itself, it is not possible to demonstrate any determinate relation of dependence whatsoever in which the one should stand to the other, it will naturally not be possible either to present the one as a function of the other. By transforming experience into a multiplicity of isolated sense-impressions, and letting each sense-impression count as a determinateness immediate in itself, a quality by itself, the Humean skepticism can be defended. This skepticism may properly be said to rest on the loose ground of the abstractly objective qualities.

Instead of attending to what is immediately objective in quality, Hegel, without himself noticing it, has gone to the opposite extreme, and, by apprehending quality as concept, has treated it only in logical generality, has posited it only as a pure determination of thought. In order to justify the proposition that being and nothing are the same, Hegel must remind us most emphatically that here one may by no means think of any determinate being or any determinate nothing. Quality, immediate determinateness, thus designates not anything immediately determinate, but the thought that there is, in general, and thus in all indeterminacy, something determinate. Once quality has been changed into a determination of thought, being is quite rightly one with nothing; that is to say: negation is immanent, for thought's expression for the qualitative limit is \opage{290} precisely negation itself. Insofar as thinking, pure, abstract thinking, is to decide the matter, Hegel can calmly challenge sound common sense to name a single example of anything actual in which the unity of being and nothing should not be found.\footnote{``Da nunmehr diese Einheit von Seyn und Nichts als erste Wahrheit ein für allemal zu Grunde liegt, und das Element von allem Folgenden ausmacht, so sind außer, dem Werden selbst, alle ferneren logischen Bestimmungen: Daseyn, Qvalität, überhaupt alle Begriffe der Philosophie, Beispiele dieser Einheit. --- Aber der sich so nennende gemeine oder gesunde Menschenverstand mag auf den Versuch hingewiesen werden, insofern er die Ungetrenntheit des Seyns und Nichts verwirft, sich ein Beispiel ausfindig zu machen, worin Eins vom Andern (Etwas von Grenze, Schranke, oder das Unendliche, Gott, wie so eben erwähnt, von Thätigkeit) getrennt zu finden sey. Nur die leeren Gedankendinge, Seyn und Nichts, selbst sind diese Getrennte, und sie sind es, die der Wahrheit, der Ungetrenntheit beider, die allenthalben vor uns ist, von jenem Verstande vorgezogen werden.'' Anfr.\ Skr.\ p.~82.}
% Errata (p. 290): the errata list corrects ``sind außer,'' to ``sind, außer''; kept as printed.
Now since, by virtue of the immanence of negation, Nothing emerges from Being, Other from Something, Many from One, and so forth, every subsequent member undeniably becomes dependent on its predecessor, and must, essentially seen, relate itself to it as consequence to ground; and so there can here to that extent be talk of determining such relations of dependence by corresponding functions.

In what relation, then, must the doctrine of quality place itself to the law of functions, when the matter is seen from the ontologically subjective standpoint? If we cast a glance at the categories of quality, as they arrange themselves easily and naturally in our language in accordance with the Hegelian type, we obtain, following Heiberg's Outline, the following schema: 1) \emph{Absolute being}: a) Being and Nothing, b) Becoming, c) Determinate being; 2) \emph{Determinate being}: a) Something, b) Other, c) The One; 3) \emph{Unity of being}: a) One, b) Many, c) Quantity. Now since the meaning of \opage{291} each such categorial expression obviously depends on the peculiar operation of thought from which it has emerged, and such an operation of thought can always be regarded as a function of the understanding, one may indeed say with a certain justification that every category is a function of the place it occupies in the system. But the difficulty is that this general truth cannot --- which is precisely what testifies to its subjective one-sidedness --- get beyond the purely general; the difficulty is that the form of the functions of thought coincides with the content itself in such a way that the reciprocal determination of form and content characteristic of the function cannot be specified punctually; the difficulty is that the objective element is lacking. According to the system, the category Quantity is a function of the reciprocal determinations One and Many, and a peculiar function at that, insofar as this reciprocal determination is peculiar; but that the reciprocal determination One and Many is peculiar can be seen, among other things, from the fact that it is different from the reciprocal determination Something and Other, and this in turn different from the reciprocal determination Being and Nothing. Nevertheless it will be impossible to determine what is peculiar to these functions in such a way that the determination of the functions' peculiarity could be separated from the functions themselves. What lies in the act of thinking whereby One becomes Many and the Many become One can be seen only from this act of thinking; the same holds of the logical operation whereby Something passes over into Other, and Other into Something.

The intellectual constituents that must necessarily step apart from one another when a function is to be determined for itself, namely presupposition, condition, consequence, find themselves, under the treatment of ``the speculative logic,'' in such fluid transitions that each of the function's separate constituents must itself become a function. Thus the One is a function of Something and Other only by virtue of the fact that Something and Other, the moments \opage{292} of the One, alternately come forth as abstractions of the common concrete, hence as functions of the One: Something is just as much a function of the One as the One is a function of Something. Further, the One is a function of Something only insofar as One is a function of the One. ``When Something is designated as the One, it has taken its Other up into itself, and has regained its unity. This unity is itself the unity of absolute being and determinate being, or the not abstractly-immediate, but concrete and resulting \emph{unity of being}.''\footnote{Heiberg.\ Anfr.\ Skr.\ p.~17. (Pros.\ Skr.\ vol.~1, p.~138.)} Thus: the operation of thought whereby the category is determined is a function; the category determining the operation is a function, namely of the act of thinking whereby it was itself determined; the new category resulting from the operation of thought is then finally also a function: thus arise functions of functions to infinity. The functions penetrate and shine through one another in so many illuminations, alterations, transformations and fleeting transitions that the only thing by which the objective determinations of difference are replaced is the words of language: the quality Something has immediate determinate being in the word Something, just as the quality Other has it in the word Other. The immediate determinatenesses are to such a degree separated from all actual existence and grown together with linguistic expression that the philosophy of pure qualities can be regarded as a pure philosophy of words.

To the abstractly objective qualities the concept of function cannot be applied, because the subjective is lacking; as regards the concepts of quality, the abstractly subjective determinatenesses, the peculiarity of the functions cannot be more closely determined, because the objective is lacking. In both respects, then, objectivity must be missed. He who fixes his gaze exclusively on what is objective in quality misses objectivity just as much as he who puts the concept of quality in the place of quality; with his loose multiplicity of immediate \opage{293} object-impressions Hume was infinitely far from cognizing an objective world. That he missed objectivity made itself known at once, as we saw, in the fact that the category of function became unusable.

An abstractly objective doctrine of quality there is not; but is there not, then, an objective doctrine of quantity? To answer this question we shall again cast a glance at the Hegelian schema. As the great trilogy we can then, following Heiberg's direction, have: 1) \emph{Absolute quantity}, 2) \emph{Quantitative determinate being}, 3) \emph{Unity of quantitative being}. If this schema is now filled in, the course of the categorial development is as follows:

1. \emph{Absolute quantity}. a) ``One and Many are the moments of quantity. But One is continuity of itself; whereas the Many, which are mutually separated, are discreteness.'' b) ``The transition from continuity to discreteness, and conversely, is the \emph{becoming} of quantity. This becoming is \emph{quantum} or \emph{magnitude}.'' c) ``The unity of discreteness and continuity thus arising is \ldots\ the resulting one, whereby they become the moments of quantity. Thus quantity is \emph{the determinate quantum} or \emph{the determinate magnitude}.''

2. \emph{Quantitative determinate being}. a) ``The determinate magnitude is quantitative \emph{unit}.'' b) ``But as unit it is seen merely from the side of continuity. Insofar as discreteness is also its moment, it is at the same time \emph{amount}.'' c) ``But the resulting unity of that immediate unit and the amount is \emph{number}.''

3. \emph{Unity of quantitative being}. a) Number, as unit in amount, is a \emph{sum}. To produce such a sum from immediate numbers (nothing but 1s) is to \emph{numerate} or \emph{count}. But to produce it from numbers that are already themselves sums is to \emph{add}, which consists in counting sums. Since the sum is number itself, unit and amount are its moments.'' b) ``When the identity of both moments develops into \opage{294} difference, the \emph{product} arises. To produce this is to \emph{multiply}, which consists in adding a certain amount of sums which are themselves equal, and consequently constitute a \emph{unit}.'' c) ``The regained identity of unit and amount is the \emph{square}. To produce this is to \emph{square}, which consists in taking a number as many times as it itself contains amount, or in multiplying the number by itself.''

It will be seen from this that the operations of thought that come into application in the doctrine of quantity are different from those that hold for the determination of the pure qualities: to \emph{count}, to \emph{multiply}, to \emph{square} are acts of the understanding in quite another sense than the reflections whereby Something passes over into Other and Other into Something. The mere fact that on the foundation of pure quantity there has arisen a scientific edifice of imposing grandeur, namely mathematics, whereas no doctrinal edifice has ever been erected on the foundation of pure quality, must arouse attention: quantity, pure quantity, plays quite a different role in our cognition of the objective than pure quality. The task, then, is to show where the indicated difference of sphere comes from and in what it essentially consists.

Hegel presents the difference between quality and quantity thus: ``Die Qvalität ist die erste unmittelbare Bestimmtheit, die dem Seyn gleichgültig geworden, eine Grenze, die ebenso sehr keine ist \ldots\ Weil das Fürsichseyende nun so gesetzt ist, sein Anderes nicht auszuschließen, sondern sich in dasselbe vielmehr affirmativ fortzusetzen, so ist das Andersseyn, insofern das \emph{Daseyn} an dieser Kontinuität wieder hervortritt, und die Bestimmtheit desselben \emph{zugleich} nicht mehr als in einfacher Beziehung auf sich, nicht mehr unmittelbare Bestimmtheit des daseyenden Etwas, sondern ist gesetzt, sich als repellirend von sich, die Beziehung auf sich als Bestimmtheit vielmehr an einem andern Daseyn (einem fürsichseyenden) zu \opage{295} haben, und indem sie \emph{zugleich} als gleichgültige, in sich reflektirte, beziehungslose Grenzen sind, so ist die Bestimmtheit überhaupt \emph{außer sich}, ein sich schlechthin \emph{Aeußerliches} und Etwas ebenso Aeußerliches; solche Grenze, die Gleichgültigkeit derselben an ihr selbst und des Etwas gegen sie, macht die qvantitative Bestimmtheit aus.''\footnote{Anfr.\ Skr.\ p.~209.} Now had Hegel actually brought everything that germinates in this excellent transitional determination to complete development, he would absolutely have had to become aware, as regards objectivity, that quantity designates a sphere of its own.

The very way in which externality, finitude, diversity, multiplicity in quantity alternate with a formal internality, infinity, uniformity and streaming unity must strike the logician. That Hegel from his standpoint has characterized the conflict between ``the forms of number'' and ``the determinations of thought'' is not to be denied; on this occasion he even comes, in certain places, to express himself with a certain logical passion, as when it is said: ``Wenn nun aber die Denkbestimmungen durch Eins, Zwei, Drei, Vier für die Bewegung des Begriffs, als durch welche er allein Begriff ist, bezeichnet werden, so ist dieß das Härteste, was dem Denken zugemuthet wird. Es bewegt sich im Elemente seines Gegentheils, der Beziehungslosigkeit; sein Geschäft ist die Arbeit der Verrücktheit.''\footnote{Anfr.\ Skr.\ p.~249.}

But instead of seeing that what here repels thought is precisely the antilogical element, which is inseparable from objectivity and the objective, Hegel holds exclusively to the logically universal, believes he has grasped true objectivity in it, and so on the whole gets no further than a polemic, by no means unjustified for that matter, against using geometrical figures, circle, triangle \opage{296} etc., as mere \emph{symbols} (the circle of eternity, the triangle of the Trinity). We shall gladly grant him that philosophy cannot immediately borrow formulas from mathematics in order to express concepts in them, since it is precisely philosophy's business to decide how far the mathematical formulas express thoughts and concepts; but then we must also make the demand of philosophy that it take along all the elements of the concept, and not fob us off with something merely subjective when it has promised us the objective.

As long as quality is held fast by itself, the extremes must alternate: on the one side a discreteness without continuity, on the other side a continuity without discreteness; these extremes dissolve in quantity. The dissolution takes place in that the objective qualities, mutually excluding one another, yield reflexes, fleeting determinations of difference, for the development of the categories of quantity. The qualitative difference between Something and Other, which, so to speak, poured itself out in the stream of alteration, vanished as a subsidiary determination in the result that Something and Other are one and the same. Instead of a stream of Something, Other, Other, \ldots\ we thus get a stream of One, One, One, \ldots\ to infinity. This stream has now dissolved discreteness into continuity, for the one One is like the other; it is unity that continues itself in every One: the continuing unity is continuity itself. But the stream has at the same time broken continuity and posited discreteness; for each single One is the whole unity by itself; it is unity that, in the repetition of the One, in \emph{plurality}, has come outside itself and become discrete: discreteness is unity posited as plurality.

What it means that One and Many are the moments of quantity, that quantum, the determinate magnitude, is a determinate unity of continuity and discreteness, etc., is excellently developed in Hegel's Logic, and can be elucidated in many ways from a Hegelian standpoint. But what is not developed in \opage{297} Hegel's Logic, and cannot be illuminated from a Hegelian standpoint either, is precisely the peculiar relation in which quantity as quantity evidently stands to objectivity. Just as continuity on the whole designates the subjective, so discreteness, wherever its immediate validity is acknowledged, points to what is objective in itself; as the formal unity of discreteness and continuity, quantity is a categorial expression for formal objectivity. But the significance of this can show itself only when quantity is determined in unity with function.

If it is a lack of logic when one immediately sets up a concept of quantity separated from, or at any rate in merely finite opposition to, the concept of quality, then it is also a lack of logic when one develops the categories of quantity in so abstract a generality that the concept of function is overlooked. So far is the doctrine of functions from standing in a merely accidental relation to the doctrine of quantity, that it is precisely quantity, as mathematics for its part amply shows, in which function must seek its exact determinations: the logic of quantity is inseparable from a logic of functions. We convince ourselves of this by a glance at the \emph{concept of function}.\footnote{Cf.\ Mathematik og Dialektik. En philos.\ Afhandl.\ by R.\ Nielsen, Kbhvn.\ 1859.}

For the peculiar determinateness of the mathematical function there is required: 1) a constant and a variable magnitude, 2) two variable magnitudes, an independent and a dependent one\footnote{That a function such as $y = f(x, z, t, u \ldots)$ contains several independents is left out of consideration in this context, where it is merely a matter of a conceptual development.}, and 3) a relation between the function's dependent and its independent that can be determined as a peculiar fundamental relation. On whichever of these three requirements we reflect, the reciprocal determination of quantity and function will be conspicuous. \opage{298} How could the difference between discreteness and continuity be expressed more clearly than by an exact separation of the limit that is and the limit that becomes: by virtue of the limit that is, quantity is quantum; by virtue of the limit that becomes, quantum proceeds from a quantifying. If the quantifying is sublated in its determinate quantum, then the limit is a limit that is, and the magnitude is determined as \emph{constant magnitude}: continuity has been absorbed into discreteness. If, conversely, the determinate quantum passes over into quantifying, then the limit is a limit that becomes; with the limit that becomes, the magnitude is \emph{variable}: discreteness is absorbed into continuity. ``In the expression $f(x) = ax$, $a$ is constant, $x$ variable; in the expression $f(a) = ax$, on the other hand, $x$ is constant and $a$ variable. Now if every magnitude, according to its concept, can be now constant, now variable, according to circumstances, then it can already be seen from this how easily the axiom of mathematics, `every magnitude is equal to itself,' can be attacked by dialectic. That every finite magnitude must, insofar as it is constant, be equal to itself, $a = a$, is incontestable; but how do matters stand with the variable magnitude? \ldots\ That the determination \emph{variable} belongs to magnitude just as essentially as the determination \emph{constant} is seen from the fact that magnitude is defined precisely as that which can be increased and diminished \ldots\ The axiom set up, `every magnitude is equal to itself,' then dissolves into two propositions, of which the first is a tautology, the second a contradiction. For it is not magnitude as magnitude, but magnitude as constant magnitude, that is equal to itself; the emphasis falls not on being magnitude, but on being constant magnitude. And what, then, is a constant magnitude? A magnitude that is equal to itself. So the axiom says: a magnitude that is equal to itself is equal to itself; a pure tautology! If, next, magnitude is apprehended as that `which can be thought increased or diminished,' then the emphasis falls not on being magnitude, but on being variable \opage{299} magnitude. The variable must show itself in varying. The variable magnitude, which according to the axiom was of course also to be equal to itself, can thus resemble itself only in constantly becoming unequal to itself''; but to have one's self-equality as an inequality with oneself is a contradiction. ``These difficulties are removed when we look at the function. The concept of function is, dialectically seen, the true concept of magnitude, since the two-sided determination of magnitude, to be constant, i.e.\ equal to itself, and yet variable, i.e.\ unequal to itself, is here posited in unity. Even if one cannot immediately say $x = x$, insofar as $x$ is variable, one can nevertheless say with truth $f(x) = ax = f(x)$, since the same function must, as function, be equal to itself under all changes. The tautology being removed, the contradiction is at the same time sublated. The unequal no longer stands as abstract negation over against the equal, or the variable over against the constant; the function lets precisely the equal reflect itself in the unequal, the variable in the constant, the moment of difference in identity.''\footnote{Mathem.\ og Dialekt.\ pp.~5--7.}

The difference between the constant and the variable magnitude is thus a difference within unity, which is developed in and with the function. In the equation $f(x) = y = ax$, the difference between $a$, the constant, and $x$, the variable magnitude, is expressly designated by $f(x)$, which indeed shows precisely that the changes of the magnitude corresponding to the function depend on the changes in $x$; whereas the unity-in-difference of $a$ and $x$ is designated by $y$, since $y$ shows that from the way in which $x$ enters into combination with $a$ a new magnitude results. To be sure, the mathematical function must be regarded as a function of the formal ground, just as much as the abstractly logical one; but the difference is nevertheless striking. The abstractly logical function lacks the element of punctuality to such a degree that it eludes every exact articulation; its articulations consist in clear, sharp determinations of reflection; it is to infinity a reflex of reflexes: so wholly does it belong to ideality and subjectivity. \opage{300} The mathematical function is, as regards both thought and language, the pasigraphic sign-language, expressly laid out for punctuality and exact analyses. Without giving up the universal, it is nevertheless able to specify and determine the individual; its peculiarity expresses itself precisely in the articulating interplay between the universal and the immediately individual. By this articulating interplay the function becomes able to express a system of thoughts in a system of laws of existence, and thus to make the objective distinct by means of the subjective. ``How function and law coincide in the fundamental concept is seen from innumerable examples. That $F(x, y, z) = 0$ is the equation of a surface means, in other words, that this function expresses the law for the changes of coordinates, corresponding to a given system, whereby the points in the surface in question are determined. The number of points is infinite; an infinite number of points cannot be determined singly; in the choice of the points one wishes to determine singly there is thus something arbitrary and accidental; but in accordance with the function, the determination of the innumerable points is reflected precisely in the way in which the individual point is determined: thus the finite everywhere enters under the infinite, the changeable under the permanent, phenomena under the law, the accidental under the necessary, the thing under the concept, the particular under the universal.''\footnote{Mathm.\ og Dialekt.\ p.~20.}

``An idealism that really concerns itself only with what is actual in the finite and material sense, insofar as it everywhere aims at dissolving this actual as a configuration of the universal concept, as a moment vanishing in the pure idea, has at bottom no essential need to engage with exact analyses. In Hegel's Logic, accordingly, the treatment of the mathematical, relegated to a couple of remarks, stands in rather loose connection with the dialectical development of the categories \opage{301} of quantity; that the Hegelian method must consistently lead away from the mathematical, as it leads away from the physical analyses of actuality, is shown by his philosophy of nature, and in particular by his explanation of the law of falling bodies.''\footnote{Anfr.\ Skr.\ p.~147.}

The most telling proof of the logical connection between the mathematical function and formal objectivity is the fact that mathematics, although an a priori science just as much as logic, must nevertheless furnish a foundation indispensable to physics. ``The history of science,'' says Ramus, ``presents the fact that the greatest discoveries in physics and astronomy, which have intervened so powerfully in both the spiritual and the social life of mankind, have been partly conditioned by, partly most closely connected with, the discoveries of the most abstract truths of mathematical analysis. If the study of higher mathematics is long and difficult, its results in turn afford an all the more beautiful reward. Fruitless, on the other hand, are the endeavors of those philosophers who would shorten the way by mystical and metaphysical speculations.''\footnote{Anfr.\ Skr.\ p.~148.}

\begin{center}{\bfseries β)\quad Functions of the Real Ground.}\end{center}
\addcontentsline{toc}{paragraph}{β) Functions of the Real Ground}

It is an unmistakable merit of the Hegelian logic that from its standpoint it shows us both the transition of quality to quantity and the return of quantity to quality in such a way that an advance is thereby made to the determination of measure (das Maaß). Clearly and sharply Heiberg, albeit in a way somewhat deviating from Hegel, has designated the transition thus: ``The square, as the identity of unit and amount in the higher circle of quantitative unity, \ldots\ or as unity itself in quantitative being, is thus more than quantity; for it is at the same time the mode of \opage{302} quantitative being, or \emph{modality}.''\footnote{Anfr.\ Skr.\ p.~22. (Pros.\ Skr.\ vol.~1, p.~150.)} To this is added in a remark: ``It might seem strange that the category of quantity concludes with the square, and that the other powers do not come into consideration. But power is a concept that is not purely quantitative; on the contrary, it is the beginning of modality. The square is indeed a power, but in this capacity it is seen from a new standpoint \ldots\ For power expresses the intensive magnitude or the quality of the quantitative; it is by it that quantity begins to return to quality, out of which it has developed.''

But from this it then also follows that, if a misunderstanding lies hidden in the transition from quality to quantity, the same misunderstanding must repeat itself when the question is of a transition from quantity to ``measure'' or ``modality.'' ``The return of quantity to quality, or the resulting unity of quality and quantity, is expressed in general by \emph{modality} or the \emph{mode} of determinate being. For the mode in which something is, is the result of its quality and quantity, or has these as its presuppositions. E.g.\ the mode in which a state exists is conditioned partly by qualitative considerations, such as its climate, its products, etc., partly by quantitative ones, such as its extent, population, etc. Likewise the mode in which a human being lives (his mode of life) is qualitatively determined by his position, his sphere of activity, his aptitudes, his inclinations, etc., but quantitatively by the extent of those qualitative determinations, and further by his fortune, his income, etc. When one says that in all things it comes down to the mode, there lies in this the correct meaning that the thing must not be regarded one-sidedly from the quantitative standpoint, since the qualitative in the thing modifies the quantitative.'' This determination of measure as modality is in and for itself quite \opage{303} excellent; but the question is: where does the qualitative come from, how do matters stand with the transition?

From quality in its objective immediacy, as we have seen, quantity cannot possibly be derived; the emergence of quantity, not from quality itself but from the category of quality, rests precisely on the fact that this category, when thought through, is really quantity: quantity can from this side be determined as quality identical with its category. But from this it follows that quantity has in itself from the very beginning all the qualitative that it needs for the specification of its quantitative determinations, indeed all the qualitative that it is at all able to take up in its capacity as pure quantity. For quantity has this qualitative element of its own as a subsidiary determination indispensable to discreteness. Discreteness is posited by continuity's being broken; but continuity is broken precisely by arbitrary limits being posited. Only by arbitrary limits can quantity become quantum; but the arbitrary limit is, as a limit that is, the categorial determinateness of quality. In the series designating plurality, One, One, One \ldots\ the difference between One and One must absolutely be immediate and qualitative; if the difference were not immediate and qualitative, the equation ``One $=$ One'' would annihilate discreteness and thereby make every further determination of quantity impossible.

But from this it follows further that the whole transition from quantity to modality, insofar as it is to be made, following Hegel's direction, in a logically immanent way by a development of the categories of quantity themselves, is a fictitious movement. Or should the category \emph{product}, which after all points precisely to a particular mode of addition, not lie just as near to modality as the category \emph{square}, which after all designates only a particular mode of multiplication? If it depends solely on there being a mode here, then the first mode can surely represent modality just as well as the latter. In the examples adduced to illustrate the categorial significance of modality \opage{304} and mode there is not a single feature that points to the square's expressing precisely the unity of the quantitative and the qualitative. ``Das Qvantum,'' says Hegel, ``in seinem Andersseyn sich identisch mit sich setzend, sein Hinausgehen über sich selbst bestimmend, ist zum Fürsichseyn gekommen. So qvalitative Totalität, indem sie sich als entwickelt setzt, hat sie zu ihren Momenten die Begriffsbestimmungen der Zahl, die Einheit und die Anzahl; die letztere ist noch im umgekehrten Verhältnisse eine nicht durch die erstere selbst als solche, sondern anderswoher, durch ein Drittes bestimmte Menge; nun ist sie nur durch jene bestimmt gesetzt. Dieß ist der Fall im Potenzenverhältnisse, wo die Einheit, welche Anzahl an ihr selbst ist, zugleich die Anzahl gegen sich als Einheit ist.''\footnote{Logik 1ster Theil.\ Anfr.\ Udg.\ pp.~391--92.}

One may readily grant that the concept of quantity gains in formal respect by such a development; but it must be most decidedly denied that a real transition from the sphere of quantity to that of modality could be made by a merely formal development of the pure concept of magnitude. By overlooking what is essential in the matter, an otherwise acute and persevering thinker easily comes to make something all too significant of the inessential. Thus Hegel, in his conception of the relation of powers, has not only been dazzled by an ambiguity in the logical; this ambiguity has even borne unhealthy fruit in his philosophy of nature. How has the concept of the square of time not been mystified in the remarkable speculative interpretation of the law of falling bodies,\footnote{``Die der Einheit, als der Form der Zeit, entgegengesetzte Form des Außereinander des Raums, und zwar ohne daß irgend eine andere Bestimmtheit sich einmischt, ist \emph{das Qvadrat}: die Größe \emph{außer sich kommend} in eine zweite Dimension sich setzend, sich somit vermehrend, aber \emph{nach keiner andern als ihrer eignen} Bestimmtheit --- diesem Erweitern sich selbst zur Grenze machend, und in ihrem Anderswerden so sich nur auf sich beziehend.''} in which Galilei's discovery was to be glorified! \opage{305} By considering the connection between the determinations of quantity and the mathematical functions, one easily convinces oneself that the function expressed by the fundamental equation $f(x + y) = f(x) f(y)$ is just as far from satisfying modality's demand for an element qualitative in the proper sense as the function expressed by the equation $f(x + y) = f(x) + f(y)$. How, then, do we reach modality? Evidently only on the condition that we, in all immediacy, take up the immediate determinateness of quality in its unity with the corresponding quantity. Only in quality does quantity have existence. The abstraction by which we come from quality to quantity cannot sublate itself by continued analysis of pure magnitudes; it can only be integrated by quantity's being seen again in its existence, i.e.\ in its determinate quality. ``The moments of the square, unit and amount'' --- it is said\footnote{Heiberg.\ Anfr.\ Skr.\ p.~24. (Pros.\ Skr.\ vol.~1, p.~153.)} --- ``are, when they are regarded as the new immediacy in the higher circle, also moments in modality, so that the unit designates the intensive, the amount the extensive magnitude.'' From the formal side there is nothing to object to in this. The opposition between intensive and extensive magnitude is certainly an opposition of modality; but intensive magnitude cannot possibly arise by a conceptual development of pure quantity. What ``the speculative logic'' would assert, namely that ``the root in the power is the extensive, the exponent the intensive magnitude,'' has no warrant, either in mathematics or in physics or in ordinary usage or in the thoroughly thought-through concept. Far more significant, then, in this context is the Kantian distinction: ``all phenomena,'' says
\opage{306} Kant, ``are, since they can be apprehended only under the forms of space and time, extensive magnitudes; all phenomena are, because every sensation can be intensified and weakened to a determinate degree, intensive magnitudes.'' This distinction is admittedly one-sided, insofar as it is still seen only from a psychological standpoint; but this much has at least been gained, that the intensive magnitude has been determined by a manifestation of force, since it has been determined by the degree of strength that a sensation apprehending the increase can have. Difference in intensity is difference in strength; but strength presupposes a force, and force an inner. Only through the opposition between inner and outer, between force and manifestation, does the opposition between intensive and extensive magnitude acquire its full significance.

According to the view that the exponent is the intensive magnitude when the root is the extensive one, the becoming of modality is to be designated by an infinite series of powers: $n, n^{2}, n^{3}, n^{4}, n^{5} \ldots$ The terms of such a series are then to be regarded as intensities of the 1st, 2nd, 3rd degree etc. The infinite series of powers is, further, to sublate itself when it comes to a power in which the exponent is the magnitude itself: $n^{n}$. ``In the series of powers,'' it is said, ``there is no determinate relation between the root and the exponent, since the latter is subject to an infinite alteration. But in intensity the exponent forms a ratio to the root, and this ratio is that of unity. Namely, in $n^{n}$, $n : n = 1$. In the series of powers the extensive (the root) is the unchangeable, and the intensive (the exponent) the changeable. But in the ratio the intensive (the exponent) is the unchangeable, and the extensive (the sides of the ratio) the changeable. Or: in the power the exponent is the other of the root, hence the changeable in the power; but in the ratio the sides are the other of the exponent, hence the changeable in the ratio. When intensity, or the unity in the power-ratio, is regarded as the immediate beginning of the ratio, it is the \emph{equal-sided} \opage{307} ratio. The exponent of this ratio is $= 1$, or unity. This exponent remains unaltered, while the sides of the ratio are indifferent quanta which, solely under the presupposition of mutual equality, pass over into infinite series without the unity of the ratio being disturbed. $1$ designates the ratio of $1 : 1$, of $2 : 2$, of $3 : 3$ etc.\ to infinity. The exponent is the qualitative determinateness in this alternating externality. But just as unity conditions number, so the equal-sided ratio conditions the infinite series of unequal-sided ones. Here the exponent passes over into the infinite series and is designated successively by $2, 3, 4, 5 \ldots$ But to each of these exponents there corresponds in turn an infinite series of ratios in which the sides, solely under the presupposition of proportional inequality, are indifferent quanta, while the exponent, as the qualitative determination, remains unaltered. E.g.\ $2$ designates the ratio of $2 : 4$, of $3 : 6$, of $4 : 8$ etc.\ to infinity.''\footnote{Heiberg.\ Anfr.\ Skr.\ p.~24--26. (Pros.\ Skr.\ vol.~1, p.~156--57).}

Since the conceptual development here given is meant to justify the assertion that intensity is ``the unity in the power-ratio,'' it cannot possibly avoid taking account of the functions by which the magnitudes are determined. There is, then, no pure magnitude or unit of magnitude which, when regarded by itself alone, can be called either intensive or extensive. The magnitude $27$ is in and for itself neither intensive nor extensive; whereas the magnitude $3^{3}$, according to the explanation given, must be said to stand at the transition. For if it occurs in a series of powers such as $3, 3^{2}, 3^{3}, 3^{4} \ldots$, then it is intensive with respect to $3^{2}$, but extensive with respect to $3^{4}$; only insofar as $3^{3}$ is expressly posited as the last term of the series is it decisively intensive. But although the whole train of thought, so laid out, evidently rests on the principle that the power as a function is higher than the product and the product \opage{308} higher than the sum, it nevertheless comes into contradiction with this principle of its own. For by regarding ``the unity in the power-ratio'' as ``the immediate beginning of the ratio'' and thereby determining ``the equal-sided ratio,'' it makes a turn so speculative that the matter, seen from a mathematical standpoint, must become utterly meaningless. How in the world could one justify an assertion so strange as this, that the exponent of the ratio, $1$, in the equal-sided ratio should have anything to do with $n^{n}$? But the contradiction is best seen in the following way. Both ``the equal-sided'' and ``the unequal-sided ratio'' are determined by the function \emph{quotient}; but the quotient and the product are inverse functions --- $\varphi(\psi(x) = x$, so that $\psi(x) = \frac{x}{a}$ when $\varphi(x) = ax$ ---, hence correlative concepts, concepts which in the logical sense stand at the same stage of development. Now if, as ``speculative logic'' itself so emphatically stresses, the power is a higher concept than the product, then it must of course also be a higher concept than the quotient; in going from ``power'' to ``ratio,'' logic, in the belief that it really went a step forward, has then evidently gone a step backward. But from this step backward something more is to be seen; for from it one sees expressly that speculation, in spite of all its exertions, has made no progress, that it still stands at the very first beginning of the analysis of quantity and has not moved from the spot. If the equal-sided ratio is expressed by the equation $\frac{x}{x}$, hence $1 = f(x)$, and the unequal-sided ratio by the equation $a = \frac{x}{y}$, hence $a = F(x,y)$, then the reciprocal determination of determinateness (quality) and indeterminateness (quantity) that one wishes to attain by this can be attained just as well by employing sums and differences: $a = \pm x y$.
% Errata: the list of misprints and corrections (p. 308, l. 3 from bottom) gives a correction for this formula; the printed ``a = ± x y'' is kept as in the original (the recorded correction ``+x5'' → ``x + 7'' is itself uncertain).

Yet there is one circumstance here that might seem to speak for the assumption that the intensive could be derived from \opage{309} quantity itself, and that is the fact that the category \emph{degree}, a category that stands in the closest connection with the category of intensity, is undeniably used as an expression of geometrical magnitude. How, then, does ``speculative logic'' deduce ``degree'' from ``ratio''? --- ``The unity of the unequal-sided and the equal-sided ratio is \emph{degree}. By the first of these moments degree is quantitative and indeterminate; by the second it is qualitative and determinate. Regarded as magnitude, it is therefore indeterminate in itself; or the determinateness by which it is magnitude falls outside it, namely in another magnitude. E.g.\ a degree of a circular arc is, as magnitude, not determinate in itself or absolutely, for on the larger circle its absolute magnitude is greater than on the smaller; but it is determined by the angle of which it is the measure. Degree thus has its determinateness not in itself but outside itself, or it is indeterminate in its absolute but determinate in its relative magnitude; this follows directly from its being the unity of the equal-sided and the unequal-sided ratio. Both are moments in degree; the latter is its relation to the whole, the former its identity with the external determination. E.g.\ a degree of the circle stands to the whole circle in the unequal-sided ratio $360$; but its identity with the corresponding angle constitutes its equal-sided ratio, which is designated by $1$. This identity, by which degree has its determinateness outside itself, constitutes its qualitative side.''\footnote{Heiberg.\ Anfr.\ Skr.\ p.~26. (Pros.\ Skr., vol.~1, p.~157.)}

Now if the circle could really be adduced as an example of a pure magnitude, freed from all quality, then the above interpretation might perhaps seem to furnish a kind of proof that a magnitude, merely by being determined on all sides as magnitude, gave itself intensity and thus passed over into modality. But the circle as a geometrical magnitude already has a quality; it is expressly qualified as extensive magnitude. \opage{310} So far is the quality ``\emph{extensive}'' from being derived here from pure magnitude that, conversely, it is magnitude that in this case has extension as its qualitative foundation. Add to this that the degree of a circular arc has nothing whatever to do with the intensive. No mathematician will concede that the degrees on an arc of greater radius have greater intensity than the degrees on an arc of smaller radius; if one could speak here of the intensity of the degree, one would also have to be able to speak of the intensity of the angle, and to determine the intensity of the angle as a function of the radius.

It is not with a degree of the circle as it is, e.g., with a degree of heat. The degree of heat is an intensive magnitude; for with each degree the heat rises in strength. The 10th degree of heat unites in itself the strength of 9 degrees with the strength of the first degree. As a unit of strength growing in arithmetical progression, degree is here an intensive magnitude, and as intensive magnitude a qualified unity of continuity and discretion. Degree is each time the whole qualified unit: the second degree of heat does not subsist beside the first but has the first within it. Insofar as each succeeding higher degree has all the preceding lower ones within it, continuity is realized and unity preserved. This is seen, then, also from the fact that the transition from one degree to another passes through gradations that even out every finite inequality; by the continuity preserved in the stream of gradations, degree expresses the quantitative in quality. But as intensive magnitude, degree also expresses realized discretion: each degree is a unit of strength qualified in itself; as the unit continues, it is transformed in such a way that at every new stage it becomes a distinctive determinate being for itself. Thus the 50th degree of heat is a quite different heat from, e.g., the 10th, the 10th a different one from the 9th. The different units of strength are different realities; each of these realities expresses in its \opage{311} own way quality in the quantitative. The functions by which degree as intensive magnitude is determined are all functions of the real ground.

The categories \emph{degree} and \emph{ground} stand, as far as objectivity is concerned, in a far more intimate connection than one would suppose from the Hegelian logic. The succeeding degree relates, insofar as it then has any intensity, to the preceding ones as existence to ground; when the 10th degree of heat is reached, it presents itself as an existence for which the 9 lower degrees furnish the ground; the 10th degree is itself 10 degrees; its existence has the 9 degrees as its ground. The objective is expressly determined by the fact that the reciprocal relation between ground and existence can be articulated by exact specification of gradations and degrees. As long as one judges changes of temperature by sensation alone, or at any rate by an immediate estimate, the gradations will be lost in the indeterminate, and the apprehension will be one-sidedly subjective; if on the other hand one has the thermometer with its determinate divisions to go by, the determination becomes exact and the apprehension objective.

The transition from degree to measure is stated by ``speculative logic'' in the following way: ``Both moments of degree, the unequal-sided and the equal-sided ratio, or the quantitative and the qualitative moment, can thus be regarded as the absolute and the relative. But the absolute moment is indeterminate, and what is determinate is only the relative. The sublation of this difference is \emph{measure}, which therefore consists in the moments' obtaining the same common determinateness. Thus measure is both the return of the ratio to absolute modality and the return of quantity to quality. --- In measure the relation is sublated insofar as its magnitude is not merely relative, like that of degree, but at the same time absolute. An ell, a pound etc.\ are absolute magnitudes, although they at the same time designate a relation to that of which they are measures \ldots\ Degree can be extended to infinity, but measure can be \opage{312} full, and is then concluded\footnote{Heiberg.\ Anfr.\ Skr..\ p.~27. (Pros.\ Skr.\ vol.~1, p.~159.)}.'' --- The emphasis that in this deduction of measure from degree is laid on the reciprocal determination of the equal-sided and the unequal-sided ratio leads one to suppose that there must here be in the concept some immanent negation or other, which would show the necessity of the one of these ratios integrating the other. But where is this negation? In the transition from something to other, from one to many, from continuity to discretion, an immanent negation can be demonstrated; but a negation that leads from the equal-sided ratio to the unequal-sided is not to be discovered. It is said, as quoted above in the doctrine of the ratio, that ``just as unity conditions number, so the equal-sided ratio conditions the infinite series of unequal-sided ones.'' But on what ground, then, can a series such as $2 : 4$, $2 : 6$, $2 : 8 \ldots$ be said to be conditioned by the series $1 : 1$, $2 : 2$, $3 : 3$? In such a juxtaposition there is not even a glimmer of necessity that one must first have a constant exponent before one can think a variable one; indeed it is even quite incomprehensible how one gets from the former to the latter. For in the unequal-sided ratio the exponent is of course not merely constant, but so wholly restricted to expressing unity $= 1$ that one cannot at all see how such an unchangeable exponent should ever pass over into becoming changeable. But if the changeable exponent is merely set up in an external way beside the unchangeable one, what becomes then of the deduction, the immanent deduction? If ``the equal-sided'' and ``the unequal-sided ratio'' themselves stand in an entirely external, indifferent and contingent relation to each other, then they cannot, in the stricter sense, be said to constitute the moments of degree either; and when they do not in the most essential sense constitute the moments of degree, then the emergence of measure from degree, supported on those moments, is an empty formalism.

\opage{313}What is decisive for degree and gradation is that quantity shows itself as a determination by real quality. The real determination of quality does not surface little by little out of the ``equal-sided'' and ``unequal-sided ratio'' of quantity; it announces itself at once when the abstraction of quantity is to be integrated. If the abstraction of quantity is not integrated by immediate elements of intuition, then neither is it integrated by continued analyses of magnitude; for one can go through pure mathematics in its whole extent without thereby in the least approaching real qualities. The real determination of quality lies as near to summation as to potentiation; and when real quality is involved in quantification, degree and measure also emerge without any artifices with ``equal-sided and unequal-sided ratios.'' Only with degree and measure do the relations of dependence determined by the functions become real relations of ground: objectivity becomes real.

How matters stand in this connection with real objectivity can be seen from determinate functions. The effect of gravity, e.g., stands, according to Newton's law, in inverse ratio to the squares of the distances; if this relation of dependence is stated by the equation $T = F(x) = \frac{a}{x^{2}}$, then we have here a function of the real ground, and in this function conditions for an exact determination of degree and measure. By $a$ the constant unit of attraction is designated; this unit of attraction is an intensive magnitude, a unit of strength, which presupposes a force determined by the mass, a force which in turn, under given conditions of distance, manifests itself with differing strength: to the unit of strength, intensity as intensity, there thus corresponds a determinate unit of distance. Intensity is at one and the same time both constant and variable; as constant it is designated by $a$, as variable by $T$; as constant it is essentially in the ground, as variable in the manifestation.

That the relation between $a$ and $T$ is a real relation of ground \opage{314} is now easy to see. If $x = 1$ is set, then $a = T$; but notwithstanding this, $a$ is nevertheless essentially different from $T$; $a$, the force by which a determinate body, e.g.\ the earth, is able to attract another body, e.g.\ the moon, at a distance $x = 1$, is here to be apprehended as the inner, as possibility, in opposition to $T$, the attraction taking place, which thus designates the manifestation of the force, the actual, the outer. Now insofar as the force $a$ is the earth's force of attraction, a force that has an independent existence, since its existence is the terrestrial globe itself, $a$ is from this standpoint certainly to be apprehended as cause, not as ground; but insofar as reflection here is directed, not to the existential unity of the earth and the force, but to the essential relation between the force in the earth and its effect on the moon, the force has only a determinate being negatively reflected in its manifestation; the intensity $a$ is then the real ground, and the manifestation $T$ the grounded phenomenon.

The function of the real ground, $T = F(x) = \frac{a}{x^{2}}$, then shows us how the relation of dependence expressed in the fundamental law is articulated by the condition falling between the ground and its consequences, the influence of distance entering as $\frac{1}{x^{2}}$. Only through this articulation does the expression of dependence become function in the stricter sense. Every existence can certainly be said, quite generally, to be dependent on its ground; when the ground changes, the corresponding existence must also change: existence is in this respect always a function of its real ground. But the exact determination of the function, the determination by which its significance for objectivity first fully comes to light, can be attained only by its assuming a mathematical form. What this means is seen clearly enough from the example adduced. The same ground has here a series of highly different consequences. Seen superficially, the change in the consequence certainly seems \opage{315} to depend, not on the ground, since the ground is of course essentially the same, but on the condition: the condition $\frac{1}{x^{2}}$ shows us $x$ as the independent variable. A closer consideration, however, convinces us that the union into which the independent enters with the constant, $\frac{a}{x^{2}}$, expresses precisely the union into which the condition, as fundamental condition, enters with the ground. In the equation $\frac{a}{x^{2}} = F(x) = T$ the first term, $\frac{a}{x^{2}}$, indicates the difference, the last term, $T$, the unity, and the middle term, $F(x)$, the unity-in-difference of the ground and its condition.

Now since not only the difference and the unity, but also the unity of difference and unity, are punctually expressed in the form of the equation, all changes in the phenomenon $T$ can be determined by changes in $x$, and since the changes in $x$ can be exactly specified, the relation of dependence designated by the function can at the same time be articulated to infinity. Under this articulation the reciprocal determination of the intensive and the extensive, of degree and measure, becomes clear. As extensive magnitude, $x$ is precisely suited for exact measurement. The standard of measure, the qualified unit, is a unit of length, since $x$ itself is a length; the qualitative is in this case immediately posited together with the quantitative. That the categories of measure and standard of measure by no means always presuppose the categories of degree and intensity our example can illustrate, since it expressly shows that the intensive magnitude is measured by the extensive, but not conversely. If the category of measure were to correspond precisely to the development of concepts according to the Hegelian method, then it would surely have to be discernible from logic that measure and standard of measure could be taken only from such magnitudes as had already taken up the intensive and degree into their concept; but this logic can in no way justify. On the contrary, there is of course a measure for the extensive that is quite \opage{316} independent of the intensive; whereas there is not, conversely, any measure for the intensive independent of the extensive.

What it signifies that functions of the real ground must condition our insight into real objectivity is seen clearly enough from the relation between analytical mechanics and mechanical physics: the exact determinations in pure mechanics are related to the determinations of existence in physics as possibility to actuality. In the functions of the real ground the subjective is articulated by the objective, for the essence of the functions is logical; but the objective is then also at the same time made clear by the subjective, for the real changes corresponding to the function are, in their immediate reality, antilogical. The functions of the real ground are, from the logical side, functions of knowledge, from the antilogical side functions of power, expressions of unity articulated in the ontological sense, of the objective and subjective fundamental determinations of objectivity.

\begin{center}{\bfseries γ)\quad Functions of the Mixed Ground.}\end{center}
\addcontentsline{toc}{paragraph}{γ) Functions of the Mixed Ground}

If the functions were merely logical, so that the determinatenesses by which they are articulated rested either on arbitrary assumptions or at any rate on reflexes of existence which were held apart from existence itself by abstraction, then the inner connection of the functions would also be determinable in an abstractly rational way. The one function would then, in every distinctive totality of functions, have to be derivable with logical necessity from the other, or, if they were coordinate with one another, from one and the same principle; furthermore, the subordinate totalities would have to be both referable to and derivable from superordinate ones, so that there would be formed a system transparent in innumerable systems, an ideal system of exact functions. But in the actual world, in the sphere of acts of power, the matter stands, as has been shown in the foregoing, quite otherwise. Here in the real system the categorical determinations of knowledge are converted into \opage{317} categorical-anticategorical acts of power. The consequence of this is that instead of totalities of functions in logical transparency we get complexes of heterogeneous functions, factual totalities, unities that are determined as existing objects. Attention has already been drawn in an earlier connection to how many diverse properties a single substance such as, e.g., silver can unite: to each of these properties there naturally corresponds, under given conditions, a particular function. That the properties, although in the substance silver they have a substantial unity, cannot with logical necessity be derived either from one another or from some general category of silverness has likewise been demonstrated; but when the properties united in the \emph{substance} as material substance (attributes and accidents) either cannot at all, or can only in a very restricted sense, be derived from one another, then the same must hold of the corresponding functions. When, then, we single out one such function by itself and consider it in its distinctive mode of acting, while at the same time we are convinced that this distinctive mode of acting is conditioned by the function's connection with the other functional dispositions belonging to the substance, we have here a \emph{function of the mixed ground}.

The Hegelian logic's treatment of measure, of the relations of measure and their nodal line, certainly belongs among the most excellent things it has achieved at all. Real objectivity here makes itself known in many main features. But in two points the one-sidedness of the method is nonetheless seen, as regards this section too: 1) the absolutely immanent deduction will not, as already shown, suffice; 2) the presupposed unity of quality with quantity is too abstract and uniform. It is the latter point that we shall illuminate in this connection. ``Alles, was da ist,'' says Hegel,\footnote{Logik.\ 1ster Theil.\ Anfr.\ Udg.\ p.~403--404.} ``\emph{hat ein Maaß}. Alles Daseyn \opage{318} hat eine Größe, und diese Größe gehört zur Natur von Etwas selbst; sie macht seine bestimmte Natur und sein Insichseyn aus. Etwas ist gegen diese Größe nicht gleichgültig, so daß, wenn sie geändert würde, es bliebe was es ist, sondern die Aenderung derselben änderte seine Qualität. Das Qvantum hat als Maaß aufgehört Grenze zu seyn, die keine ist; es ist nunmehr die Bestimmung der Sache, so daß diese über das Qvantum vermehrt oder vermindert zu Grunde ginge.''

As regards the statement that ``quantum has as measure ceased to be a limit that is no limit,'' it must be remembered that, for a quantum to become measure, all that is required is that it be expressly posited as a constant magnitude. The limits of the constant magnitude are always qualitative; this holds even to the extent that not only finite but even infinitely small magnitudes can become measures. If $y = f(x)$, then $dx$ is just as qualified a standard of measure for $dy$ as a mile for the length of a road, or an ell for a piece of linen. That ``something is not indifferent to its magnitude'' is to be conceded; but when the quality inseparable from the magnitude is apprehended as a qualitative determinateness for itself, separated from other qualities, how then is one, by mere quantifying, to get beyond the infinity of gradations to the fixed limit of measure? Let the quality be, e.g., red. This red is, as quality, not merely in unity with quantity; but the unity is even so two-sided that the quantity is at one and the same time both extensive and intensive. That red as extensive magnitude can quantify to infinity is evident: the red quality of course sets no limit to the increase or diminution of its extension. As intensive magnitude, the red color can furthermore be intensified and weakened through innumerable gradations without any fixed or final limit immediately being demonstrable. There are certainly limits here to intensity, just as there \opage{319} are limits to the magnitude that lies between $0$ and $1$; but since the gradations in intensity are just as fluid as the alterations in pure quantifying, the limit of measure is not yet posited by this quantifying.

If the determinate limit of quantity is to be expressly qualified as a limit of measure, then the determination of the limit must itself be a function of several co-operating qualities. Heat, immediately apprehended, is a quality that has cold as its limit, just as light is a quality that has darkness as its limit; with respect to quantity these qualities are, moreover, intensive magnitudes. But heat as a quality for itself can of course quantify to infinity and run through innumerable gradations without any measure or fixed limit being seen in this. By what, then, does the fixed limit come about? Evidently by the quality heat acting in unity with other qualities, with the qualities in the bodies that are heated and altered by the heat. ``Heat,'' says physics, ``expands bodies and governs the forms of state''; in both respects the effect of heat is a function of the mixed ground, and in both respects it acts to determine measure. In pure gradation the qualitative becomes a moment for the quantitative, so that the limit can constantly be exceeded, and the change, like the change of heat, can go on to infinity. In measure (μέτρον), on the other hand, the qualitative has quantity as its moment; thereby a limit is posited that, despite gradation, cannot be exceeded. Such qualitative limits are seen, under the influence of heat, in cohesion. For under the influence of heat solid bodies expand according to a quite different standard of measure than liquid ones, and these in turn according to a different standard than gaseous ones. If the activity of heat has its measure in the expansion of solid bodies, then the differing capacity for expansion of different solid bodies has in turn its measure in the activity of heat. Now since one and the same raising of \opage{320} the degree of heat produces a different expansion in the different solid bodies, these can be said, as regards expansion, to have their measure in heat, in one another mutually, and each for itself in itself.\footnote{``Um ein Beispiel anzuführen, so ist die \emph{Temperatur} eine Qvalität, an der diese beiden Seiten, äußerliches und specificirtes Qvantum zu seyn, sich unterscheiden. Als Qvantum ist sie äußerliche Temperatur und zwar auch eines Körpers als allgemeinen Mediums, von der angenommen wird, daß ihre Veränderung an der Skale der arithmetischen Progression fortgehe und daß sie gleichförmig zu- oder abnehme; wogegen sie von den verschiedenen in ihr befindlichen besondern Körpern verschieden aufgenommen wird, indem dieselben durch ihr immanentes Maaß die äußerlich empfangene Temperatur bestimmen, die Temperatur-Veränderung derselben nicht der des Mediums oder ihrer untereinander in direkten Verhältnisse entspricht. Verschiedene Körper in einer und derselben Temperatur verglichen, geben Verhältnißzahlen ihrer specifischen Wärmen, ihrer Wärme-Kapacitäten. Aber diese Kapacitäten der Körper ändern sich in verschiedenen Temperaturen, womit das Eintreten einer Veränderung der specifischen Gestalt sich verbindet. In der Vermehrung oder Verminderung der Temperatur zeigt sich somit eine besondere Specifikation \ldots\ Näher betrachtet würde daher das Verhältniß eigentlich nicht als Verhältniß von einem bloß qvantitativen und einem qvalificirenden, sondern von zwei specifischen Qvantis zu nehmen seyn. Wie sich das specificirende Verhältniß gleich weiter bestimmen wird, daß die Momente des Maaßes nicht nur in einer qvantitativen und einer das Qvantum qvalificirenden Seite einer und derselben Qvalität bestehen, sondern im Verhältnisse zweier Qvalitäten, welche an ihnen selbst Maaße sind.'' Hegel Logik.\ 1ster Theil p.~410--411.}

But what can be the ground why one solid body, other conditions being equal, does not expand as much as another? Indeed, what is the ground why the expansion of gaseous bodies proceeds according to a different law than the expansion either of liquid or of solid ones? Existence's need for measure and limit! It is a mixed ground that motivates the satisfaction of this need. In \opage{321} the first place, the peculiarity of cohesion, as regards the different bodies, is a function of the quality of the cohering parts, but the quality of the parts is in turn a function of their position, figure, mutual distance etc. Secondly, as regards the forms of state, each state is grounded in a system of distinctive conditions; by its system of inner and outer conditions the one form of state stands in qualitative opposition to the other. Now since the conditions vary, and the grounding is a function of the varying conditions, it is evident from this that diverse presuppositions, conditions and circumstances can enter as moments into every particular grounding and thus give the sufficient ground the character of \emph{a mixed ground}.

``When a body,'' says physics, ``passes over from the solid to the liquid state, one observes two remarkable phenomena. First, it continues to be solid up to a certain determinate degree of heat, which is unchangeable for the same body, and at which melting can first begin; and second, the degree of heat does not change, however much heat penetrates into the body. The unchangeability of the melting point and the binding of heat are two essential conditions for melting.'' Hereby the qualitative opposition between the solid and the liquid form of state is expressly designated. As long as cohesion is strong enough to withstand the influence of heat, the body remains in its solid state; but the force of cohesion has its qualitative limit: if the limit is exceeded, the parts are separated, and a new form of state is formed. During the separation of the parts, heat penetrates into the spaces formed in such a way that it thereby as it were hides itself in the body and merges with the new form. Instead of acting on feeling or on the thermometer, hence \emph{outward}, heat acts only to maintain the separation of the parts that has been brought about, hence \emph{inward}. The latent or bound heat is thus the measure of \opage{322} heat, the capital of heat, that is needed to hold the parts in the order corresponding to the new state, that is, to maintain the quality of the form of state.

That conditions of different kinds intermix---the intermixture being itself grounded---in the grounding of the difference of the forms of state is manifest. Sulfur, e.g., melts at $111^{\circ}$ into a thin and yellow liquid; heated still more strongly, it becomes, at $160^{\circ}$, viscous, dark, reddish-brown, and at about $200^{\circ}$ quite stiff and tough. If the molten sulfur is cooled slowly, it crystallizes, but in a crystal form of its own. At about $400^{\circ}$ it boils and is thereby changed into a yellow vapor that occupies a volume 5--600 times greater than the solid sulfur; if the sulfur vapors are mixed with cold air, they condense in the form of a fine yellow powder and form the so-called flowers of sulfur; if the sulfur is precipitated from a salt solution, it forms a still much finer powder (milk of sulfur), etc.

The cohesion valid for each form of state thus has not merely a strength determined once and for all, which heat must overcome when the form changes; but the strength of cohesion can itself change during the raising of heat. Such a change, such a sudden reversal, as when the cohesion of sulfur, which was too weak to withstand a heat of $111^{\circ}$, again became so strong that it could oppose a heat of $160^{\circ}$, and even offer vigorous resistance to a heat of $200^{\circ}$, would surely be altogether inexplicable if the change of cohesion were not grounded in distinctive co-operating influences stemming from other changes, effects which only arise at a given degree of heat and mix themselves into the matter.

This node-forming mixture, this---one might be tempted to call it---articulating ``mixture'' of grounds, comes clearly to the fore in chemical proportions; measure is here sharply designated.
\opage{323} For the combination of substances is not restricted to a single ratio: combinations take place according to multiples of equivalents; but these multiples are on the whole to be referred to simple ratios of measure. If we disregard organic bodies, then it holds as a rule that 1 equivalent of the one substance enters into combination with 1, 2, 3 equivalents of the other, or 2 equivalents of the one with 3, 5, 7 equivalents of the other. Here the limit of measure is qualified. ``How products that are formed from the same constituents acquire, precisely through measure and limit, a character corresponding to the ratios of equivalents, we can best illustrate---since oxygen combines with all hitherto known elements except fluorine---by the so-called degrees of oxidation. A substance combined with oxygen is an oxide; thus manganese, $\mathrm{Mn}$, a grayish-white, slightly lustrous and very brittle metal, combines with oxygen in 5 different ratios: $\dot{\mathrm{M}}\mathrm{n}$, $\dddot{\mathrm{M}}\mathrm{n}_{2}$, $\ddot{\mathrm{M}}\mathrm{n}$, $\dddot{\mathrm{M}}\mathrm{n}$ and $\ddddot{\mathrm{M}}\mathrm{n}_{2}$, and thus forms 5 different oxides of manganese. But why, now, can substances that stand in affinity not enter into combinations in all possible ratios, as Berthollet in his time assumed? What does it mean that combinations take place at intervals and as if by leaps? --- The craving of existence for measure and limit! The difference between the 5 degrees of oxidation cited is more than a difference of degree; it is a difference of character determined by measure.''%
\footnote{Lectures on ``Phil.\ Præpæd.''\ 1860--61. p.~127--128.}

In what sense the general doctrine of functions can be said to form the connecting link between the logic of subjectivity and that of objects will be evident at this stage of the development. So long as we, with Hegel, kept exclusively to the ontological determinations of quality, and in general of the categories at large, it was impossible to get the concept of \apage{324}function articulated; the transitions were too fluid, the forms too mobile, the stream of continuity washed away every immediate mark of difference: the reason for this was that the professed objectivity was really subjective. In the present context the relation is reversed. Here there is no lack of determinations of discretion, of determinately present marks of difference and fixed limits. On the contrary! the fixed limits, the qualitative marks of measure, press themselves upon us here in such a multiplicity and with so decisive an immediacy that it is now precisely for this reason impossible to keep the concept of function in continued development. For just as little as the concept of function permits continuity to absorb discretion, the singular, the punctual, to be volatilized in the universal, just as little can the concept of function permit discretion to supplant continuity, the singular, the special and externally different, to shoot up with such luxuriance that it overgrows and hides the universal. Where the latter is the case, the difficulties with the concept of function will soon make known that the logic of objectivity is about to founder on the rocks and reefs of finite, material objects.

We see this from our present example. Our five different oxides of manganese, of which the first, $\dot{\mathrm{M}}\mathrm{n}$, is a protoxide, the second, $\dddot{\mathrm{M}}\mathrm{n}_{2}$, a sesquioxide, the third, $\ddot{\mathrm{M}}\mathrm{n}$, a peroxide, the fourth, $\dddot{\mathrm{M}}\mathrm{n}$, an acid, the fifth, $\ddddot{\mathrm{M}}\mathrm{n}_{2}$, a peracid, are five peculiar objects, five different chemical characters. Insofar as the difference of character of the five objects is determined by the different quantities of oxygen, the chemical peculiarity of the manganese oxide can, to be sure, be said to be a function of the quantity of oxygen. If we denote the manganese oxide quite generally by $M$ and the quantity of oxygen by $m$, we thus obtain what seems a very exact expression in the equation: $M = f(m)$. But closer consideration will soon convince us that the exact \opage{325} notation is only a semblance. According to the notation one would indeed assume that $m$ was the independent, $M$ the dependent variable. Were this so, then the constant, finite or infinitely small, increment of $m$ would have to furnish an exact measure for the changes in $M$; but this, when we keep to the matter itself, is not at all the case. In the given example the qualitative changes can in no way be derived from the quantitative: a proof that the law of dependence, and with it the universal, is still veiled. The determination of quantity is unfree, immediately bound to the quality: the five quantities of oxygen have been found in a purely empirical way.

If the function were to be said to be determined, then it would have to be possible, when one knew one of the oxides, e.g.\ the protoxide, to derive the other formations from it, thus possible, without any experiments, to say at once how many and which oxides of manganese there can be at all. Why, now, is such a prediction not possible? That the character of the manganese oxide depends on the quantity of oxygen is certain; but on what, then, does the quantity of oxygen depend? An attempt to answer this question will at once direct us to the \emph{mixed ground}. To get an impression of how different the presuppositions and conditions are under which the particular degrees of oxidation are formed, discovered and prepared, or, in other words, to form a representation of how heterogeneous the functions must be which fundamentally cooperate in determining the process of oxidation from which the characteristic oxide must each time emerge as product, one need only cast a glance at the information that the general chemistry of substances communicates in this respect. The formation of manganese protoxide is then described, among other things, thus: ``Formed by violent ignition of manganese sesquioxide hydrate, or of whole crystals of manganese peroxide, in hydrogen gas, it is obtained in light-green pseudomorphs. When heated in air it glows into brown manganese protoxide-sesquioxide. Manganese protoxide is the base in the \apage{326}manganese protoxide salts; from these it is obtained, precipitated by potash, as a white hydrate, which in air oxidizes further to blackish-brown sesquioxide hydrate.'' The quantity of oxygen of which the manganese protoxide, according to the equation $M = f(m)$, was supposed to be a function, is thus a function of a process of oxidation; the process of oxidation, or still more exactly of deoxidation, is in our special case a function of the violent ignition of the manganese sesquioxide hydrate; but the peculiarity of this ignition is itself a function of the peculiarity of the sesquioxide hydrate, and this peculiarity again a function of the functions in which the formation of the hydrate is grounded, etc.

The similarity between the abstractly metaphysical and the concretely physical is, with respect to the concept of function, most striking. From logic's ``objective,'' that is, ontological-subjective and abstract-metaphysical standpoint, we saw functions of functions of the functions, whose logical transparency made every determination of punctuality impossible. From the concrete-physical starting point---in our arbitrarily chosen example, the specifically chemical one---we again meet functions of functions of functions etc., without discovering any prospect whatsoever of a closer determination of the exact nature of the individual functions. The similarity pointed out nevertheless has its real ground in a thoroughgoing dissimilarity; if the functions in the abstract-metaphysical form could not be punctually determined because discretion vanished in pure transparency, then the functions in the concrete-physical form cannot now be articulated for the opposite reason, namely that their essential and necessary connection is hidden in the mixed ground, in the night of grounds, in the logical opacity of material activities.

Since the functions of the mixed ground must all, by the nature of the case, lend themselves to development in exact form, inasmuch as they are all specifications of a general law of nature, while they nevertheless hide the inner necessity behind \opage{327} the existential immediacy of the facts, it must be recognized as a scientific task of a higher order to discover methods by which it becomes possible to master the mixture and to let the fundamental activity that governs the heterogeneous elements clearly come forth. A great step toward the solution of this great task has, as far as chemistry is concerned, certainly been taken in the very latest times in the ``thermochemical theory'' set forth by Professor Julius Thomsen.%
\footnote{Tidsskrift for Physik og Chemie. 1st year, 3rd, 4th, 6th number. Copenhagen 1862.}
From this theory we single out, in the present context, the following fundamental features characteristic for logic. ``To speak in general of stronger and weaker affinity is too indefinite; here a measure must be shown for the absolute magnitude of affinity; this measure is the development of heat. In the same chemical process an equally great quantity of heat is always developed when the external conditions are the same. The development of heat that accompanies the formation or the decomposition is always equally great, whether the compound is formed or decomposed all at once or by degrees. In the decomposition of a chemical compound just as much heat is produced as is developed in the formation of the compound.''

One sees from these features how the theory knows how to discover, in the mixed, the special and the intricate, what runs throughout, the simple and universal. Only when it has secured presuppositions in the universal does it turn toward the particular and grasp the special. Thus, with express reference to the relation between the metals and sulfuric acid in its different states of concentration, the following question is raised: ``\emph{Which metals decompose dilute sulfuric acid with evolution of hydrogen?}'' The universal consideration is this. ``When one wishes to perform a work, the quantity of work must be greater than, or at least as great as, the work. When one wishes to decompose a chemical \opage{328} compound by the quantity of work (quantity of heat) which appears in the formation of another chemical compound, this quantity of work must be greater than that which is required for the decomposition of the first compound. To answer the question posed, one must therefore know 1) the quantity of heat which is bound in the decomposition of water, and 2) the quantity of heat which is set free in the formation of the sulfate from its constituents: metal, oxygen and sulfur. The quantity of heat which appears in the formation of 1 equivalent of water, or which is bound in the decomposition of the same quantity of water, amounts to $4340^{\circ}$.'' For a metal to be able to decompose dilute sulfuric acid, the following condition, expressed in the language of signs, must therefore be fulfilled:
\[ (\mathrm{R},\ \mathrm{O},\ \ddot{\mathrm{S}}\ \mathrm{Aq}) > (\mathrm{H.\ O}) \]
Since the decomposition of water requires a quantity of heat $= 4340^{\circ}$, we obtain:
\[ (\mathrm{R},\ \mathrm{O},\ \ddot{\mathrm{S}}\ \mathrm{Aq}) > 4340^{\circ} \]
as the condition for the possibility of the decomposition of water. As a remarkable example of the theory's victory over an experience erroneous in a particular point, the following case deserves special mention. According to the condition set forth in the theory, lead
\[ (\mathrm{Pb},\ \mathrm{O},\ \ddot{\mathrm{S}}\ \mathrm{Aq}) = 4690^{\circ} > 4340^{\circ} \]
should indeed also be able to decompose dilute sulfuric acid with evolution of hydrogen. Experience, however, had hitherto by no means confirmed this fact: on the contrary! the decomposition of water by lead was something unknown to experience. Here, then, there was, as it seemed, a conflict between the a posteriori and the a priori. But the theory itself already contained a hint of how easily experience could have erred on this point. The affinity of oxygen for hydrogen and for the constituents of the salt gives, in the case of lead, a difference of only $350^{\circ}$, whereas with the affinities in the case of iron and zinc one gets a difference of $1410^{\circ}$ and $2230^{\circ}$ respectively''.
It is thus \opage{329}natural that the decomposition of water, when lead is used, must proceed incomparably more slowly than with iron and zinc; but, as a consequence of this, it was then also natural that experience, so long as it had not yet been made attentive by the theory, could have stopped short at an incomplete investigation. Carefully conducted experiments then also showed this to be the case. ``The first disagreement between experience and theory thus ended in the latter's favor, and this,'' says the author, ``was certainly \emph{the first example of a chemical process that has been found according to the method of the exact sciences.}''

The fundamental thought of the theory is that the chemical force in each substance must, at the same temperature, always have the same intensity. If the intensity is now denoted by $i$, and the heat-tone, a common term for development of heat and binding of heat, by $W$, then the heat-tone appearing in the formation of a compound, a function of the intensity in the constituents that form the compound, can be expressed by:
\[ \mathrm{W} = \mathrm{f}\,(i,\ i',\ i''\ \ldots). \]%
\footnote{J.\ Thomsen. Bidrag til et thermochemisk System. Videnskabernes Selskabs Skrifter. 5th series. vol.~3. Copenhagen 1853. p.~119--20.}
The function of the mixed ground has thus regained the universal and thereby acquired the exact form.

\begin{center}{\bfseries b)\quad Functions of the Cause.}\end{center}
\addcontentsline{toc}{subsubsection}{b) Functions of the Cause}

It would be a great misunderstanding if the difference between what we here call functions of the cause and functions of the ground were taken to mean that the former fell outside the latter. There is, to be sure, not always a cause where there is a ground; but since the cause, as we have seen in the preceding chapter, is the existing ground, the ground active in actuality, there must naturally always be a relation of ground where there is a relation of \apage{330}cause. The modes of dependence that we have cited above as examples of functions of the real and of the mixed ground can therefore, from an altered point of view, also be taken as examples of functions of the cause. With causality as starting point, apprehension again goes in two directions: the \emph{mechanical} and the \emph{dynamic}. What the mechanical means has already been illustrated from one side in what precedes; but what, now, does the dynamic really mean? In the ``speculative logic'' the following distinction is made between ``mechanical and dynamic activity.''

``The mechanical object is already in itself a plurality or an aggregate of objects, of which the one is external and indifferent to the other, i.e.\ immediately being. The mechanical activity is itself equally immediate, and determined only as the one object's pressure or \emph{impact}, which is \emph{communicated} to the other. Thus it is \emph{indifference}, and indifference is itself the mechanical \emph{force}. The sublation of indifference by means of impact (the sublation of the indifferent force in its manifestation) is \emph{motion}. The result of motion is the object's return to indifference, but to one that is no longer immediate but mediated, i.e.\ \emph{rest} or \emph{centrality}. With this return of the object to indifference, the activity of immediacy has thus completed its circuit, and that of reflection begins. Indifference is thus just as much \emph{difference}. For centrality or mechanical rest is a rest in motion, and consequently just as much unrest. This difference is, because it denotes the sphere of reflection, the relation between the positive and the negative. The activity in this relation is \emph{dynamic}, and difference is itself the dynamic force.''

The development of dynamic force thus coincides here with a development of what Hegel understands by the ``chemism'' of the concept. ``Difference,'' it is said, ``is striving for its own sublation, namely the striving of the positive for \opage{331}union with the negative, and conversely. This striving is \emph{affinity}. The result of affinity is the object's return to difference, but to one in which affinity is sublated. Thus difference is neutrality by means of affinity.''

The universal-logical element in these categorial determinations is now, on the whole, supposed to be recognizable in this, that the opposition between mechanical and dynamic activity holds not only in the bodily but also in the spiritual. ``In the spiritual world the mechanical force (indifference and inertia) shows itself in all mechanical actions; in all knowledge that is merely \emph{communicated} without being apprehended; in all memory-work, in all instilled faith that is not cognition, in all symbolism and ceremony that becomes habit; in all isolated maxims, principles, aphorisms that are torn out of their dialectical connection and thus not apprehended in their truth'' \ldots\ ``Dynamic drives are \ldots\ those which are sublated in their satisfaction, or are sated. The insatiable, on the other hand, are mechanical. For they present a motion into the infinite which even in its rest is only a renewed circuit about an unattainable center; whereas affinity, on the other hand, has its goal, in which it is neutralized.''

The chosen point of view is, however, both too special and too general; too special insofar as by the mechanical it aims all too one-sidedly at motion about a center, and, as regards the dynamic, at chemical activity; too general, on the other hand, insofar as it gives, neither in mechanical nor in dynamic respect, the least hint of how and on what conditions the activities belonging here can be specified. The task is thus a double one: partly to show the foundation on which the opposition between the mechanical and the dynamic rests in its widest extent; partly to show the principle of specification. Now as regards, first, the foundation, it evidently receives its \opage{332} widest extension when it is determined by bringing the opposition between discretion and continuity under the relation of cause. Insofar as causality is posited in the moment of discretion, the mechanical prevails; insofar as it is posited in the moment of continuity, the dynamic prevails. Now since the relation of cause itself, in its difference from the general relation of ground, is characterized by discretion, the mechanical will be recognizable by this, that all determinations of ground are punctualized by finite determinations of cause, while the dynamic will be recognizable by this, that the finite determinations of cause are transformed into immanent determinations of ground. So determined, the opposition is visible in the history of philosophy, even from the beginning; for the doctrine of transformation in the Ionian school is evidently dynamic, in opposition to the doctrine of alteration in Anaxagoras and Democritus, which is mechanical. That the doctrine of transformation is older than the doctrine of alteration, the dynamic view older than the mechanical, cannot surprise us when we consider that speculation on the whole is older than exact natural science, and that speculation to this very day feels far more attracted by the dynamic than by the mechanical view of nature.

What aversion speculative philosophy of nature on the whole must have for the exact analyses of mechanics can be seen from manifold examples. ``The Newtonian theory, according to which light is supposed to propagate itself in lines, or the wave theory, according to which it is supposed to propagate itself in waves, like the Eulerian ether or like the vibration of sound, are,'' says Hegel, ``material representations, `die für die Erkenntniß des Lichts nichts nutzen'.''%
\footnote{Phil.\ Propæd. Copenhagen 1857. p.~52.}
By this abstracting from all mechanical and therewith from all actual, in the proper sense objective, starting points, one-sided speculation shows plainly enough that, with respect to the logic of objectivity, it has no \opage{333} inkling of what a principle of specification means. So long as one does not feel prompted to investigate the objective in the activity of light, so long as one in general lets the supposedly higher rational hover in complete indeterminacy above the exact specifications, one may still in some measure rest content with an explanation such as this: ``the first qualified matter is matter as pure identity with itself, as the unity of reflection-into-self; insofar it is only the first, itself still abstract manifestation. Determinately present in nature, it is the relation to itself as independent over against the other determinations of the totality. This existing, universal \emph{self} of matter is \emph{light}; as individuality, \emph{the star}; and this as a moment in a totality, \emph{the sun}''; --- for it is certain that such an explanation is essentially subjective, although it is meant entirely objectively. Lack of objectivity is recognizable here, as everywhere, by lack of reliable specification, of a differentiation corresponding to actuality. Yet without any specification whatever objectivity cannot be thought; if, now, philosophy of nature cannot appropriate the principle for the determination of the discretions valid in actuality, then it must help itself in another way: then it must appeal to ``intellectual intuition,'' and try to poeticize its way into ideal specifications. ``The universe,'' says Schelling, ``is to be thought under the image of a line, in whose midpoint falls $A = A$, at whose ends there falls, on the one side, $\overset{+}{A} = B$, i.e.\ a preponderance of the subjective, on the other side $A = \overset{+}{B}$, i.e.\ a preponderance of the objective, yet in such a way that relative identity takes place even in these extremes. The one side is the real or nature, the other the ideal. The real side develops itself according to three potencies, potency denoting a determinate qualitative difference of subjectivity and objectivity: the first potency, $A$, is \opage{334} matter and the force of gravity---the greatest preponderance of objectivity; the second potency, $A^{2}$, is light, an inner---as gravity is an outer---intuiting of nature. Light is the higher manifestation of the subjective. It is the absolute identity itself. The third potency is the common product of light and the force of gravity, the organism, $A^{3}$. The organism is just as original as matter.''%
\footnote{Cf.\ Philos.\ Propæd. Copenhagen 1857. p.~50--51.}

It is, as regards the specifying principle, again the category of function about which everything comes to turn. By speculating one-sidedly in the dynamic, idealism must constantly volatilize all moments of actuality and put in the place of the efficient causes a series of subjective explanatory grounds, which it then passes off as objective grounds of reason. In opposition to this, empiricism goes the reverse way; without brooding over the deeper grounds, it everywhere presupposes the efficient causes. According to Bacon's words: vere scire est per causas scire, the empirical method has a sharp eye for actuality, thus for causal discretion, thus for the mechanical function, thus for objectivity. The mechanical functions are related on the whole to the dynamic as the objective side of objectivity to the subjective, as reality to ideality; it is therefore also on the basis of the mechanical that the dynamic is articulated, and we can in this respect grant Kant that there is as much science in natural science as there is mathematics in it.

\begin{center}{\bfseries α)\quad Mechanical Functions.}\end{center}
\addcontentsline{toc}{paragraph}{α) Mechanical Functions}

The objective in objectivity is the mechanical; but the objective is in itself antilogical: the mechanical must therefore reflect itself in the antilogical. Antilogical determinateness, insofar as it is apprehended logically, is exact; in every numerical value, every external standard of measure, every immediately posited limit \opage{335} an antilogical reflex is seen. In being articulated by antilogical reflexes, reflection is transformed into intellectual operation. Intellectual operation in turn expresses the unity of the mathematical and the abstractly mechanical. There is an essential and thoroughgoing difference between a development of concepts that proceeds according to universal-logical reflections and a development of concepts that is specially conditioned by mathematical operations. In mathematical operations the concept is at every point, as it were, in the process of having its immanent determinations of essence converted into external singular determinations; this antilogical element in punctuality itself Hegel certainly noticed, but what this dotted, punctual externality really has to do with the concept as concept, of that he has had no inkling.
% Errata: the print reads "det i Mechaniske"; the book's errata list corrects this to "i det Mechaniske," which is the reading translated here.
From every side he sees in the mechanical, the mathematical, the exact, with its peculiar operations and language of signs, something merely subordinate, a lower, a foreign thing, something ruinous to genuine comprehending: ``Es bewegt sich im Elemente seines Gegentheils, der Beziehungslosigkeit; sein Geschäft ist die Arbeit der Verrücktheit.'' Hegel's misjudgment of the significance of mathematical operations for the logical-antilogical development of concepts is a striking proof of his total misjudgment of the objective in objectivity.

``When Hegel shows the fruitlessness of wanting, like Leibnitz, to make the syllogistic figures the object of a combinatorial calculus, and censures a Ploucquet, who by mathematical determinations of the logical believes he has brought matters so far: posse etiam rudes mechanice totam logicam doceri, ut pueri mathematicam docentur, then he is in the right; to transform logic into a mechanics of the understanding is spiritless \ldots\ But with all this Hegel, in his judgment of the mathematical mode of expression, has overlooked a side which, as regards thought and the concept, is the real main side, namely the punctual \ldots\ Spoken language, which cannot by the word point out the singular in its immediate determinate being'' --- the antilogical! --- \opage{336} ``also cannot by the word make punctual determinateness distinct; spoken language therefore has, when thought demands the punctual, a deficiency that must be made good by the mathematical language of signs \ldots\ What velocity is one can indeed in a way express in words, but the concept of velocity demands, when it is to be made distinct with its exact determinations, exact analysis. It makes a change of sphere in cognition when one, after having explained in words what velocity is, applies the mathematical notation, calls space $s$, time $t$, velocity $v$, and sets $v = \frac{s}{t}$ as the definition of velocity; indeed, in order that the definition shall fit not merely uniform but also non-uniform motion, goes further, puts instead of the space $s$ the element of space $ds$, instead of the time $t$ the element of time $dt$, and thus defines velocity by $v = \frac{ds}{dt}$. He who has the concept of velocity only insofar as it can be presented and made distinct in words has it only in hovering generality, and thus in a superficial way. The category \emph{velocity} belongs to a circle of concepts which have the peculiarity that they cannot be made distinct by general reflection, but demand for their essential and punctual development the intellectual act, the mathematical operation: all concepts of function are concepts of operation.''%
\footnote{Mathem.\ og Dialektik p.~23--25.}
Such categories as \emph{inertia}, \emph{force}, \emph{moving force}, \emph{living force} etc.\ are to such a degree peculiar operational categories that he who knows their meaning only from linguistic usage will not be able to form any distinct concept of the role they play in mechanics. According to the ``speculative logic's'' presentation of mechanical motion, one would have to take inertia, synonymous with indif\opage{337}ference, for a mechanical force. ``It is,'' says Heiberg, ``to be regarded as an effect of the invincible power of the concept that the physicists, who are precisely wont to regard force as independent over against its manifestation, have nevertheless been compelled to adopt the true determination of abstract force, namely as the force of indifference, vis inertiæ, for whose sublation in manifestation the communication of an impact is required.'' This proof of ``the invincible power of the concept'' is, however, so extremely weak that it is on the point of turning into a proof of the mechanical impotence of the concept, namely of the speculative concept. For it is true, to be sure, that Newton and the older writers used the expression vis inertiæ; but this expression is so far from
% Errata: the print repeats "at" across the line break on p.~337 ("saa langt fra at / at være adæqvat"), a dittography not reproduced in the English.
being adequate to the mechanical concept that physics in recent times---which is precisely to be regarded as an effect of ``the invincible power of the concept''---has found itself called upon to exchange it for the expression lex inertiæ. The law of inertia is then stated in the following propositions: ``A body that is at rest begins no motion without an external cause; a moving body continues in its motion with unchanged velocity and unchanged direction, insofar as no external causes act upon it.'' Motion is, on this view, a state, just as much as rest.

If one were now able to say with ``the speculative logic'' that inertia is sublated in motion, then inertia would have to be absolutely one with rest; for it is indeed rest that is sublated in motion. But motion is a state that is in turn sublated in rest; rest is then also a manifestation of inertia, and inertia is thus one with motion. To resolve the contradiction that the manifestation, which was supposed to be the sublation of the force, coincides immediately with the force itself---that, in other words, the sublation of the force is not its sublation---``the speculative logic'' appeals to ``centrality.'' But so long as it is here a question only of
% Errata: the print reads "Inertie, ved Stød og Bevægelse,"; the book's errata list corrects this to "Inertie og Bevægelse ved Stød," which is the reading translated here.
inertia and motion by impact, there can here as yet be no question of centrality. That ``the single \opage{338}planet has the same centrality as the sun,'' or that ``the single globe is itself centrality by means of motion (E --- S --- A),'' that ``its motion'' insofar ``has returned to the indifference or rest which is no longer immediate,'' one may readily grant, without therefore being able to grant that herein is found any trace of something that could resemble a proof that inertia should be force.

What, then, is inertia? A logical expression for the antilogical, a negative determination for the objective in objectivity, an extreme abstraction in which otherness, determinately present selflessness, is held fast. The positive corresponding to the negative determination is matter. Matter as matter is a res iners, that is, indifferent to motion and rest. Inertia, or the indifference of matter, is by no means the same as sluggishness; the German designation ``Trägheit'' is in this respect misleading. If inertia were sluggishness, then matter would, by means of inertia, have to be less disposed to motion than to rest; but inertia is supposed precisely to denote that it is equally disposed to both. Far from being any force, inertia is on the contrary a categorial expression for lack of force, insofar as it is an expression for matter's lack of self-activity. And yet this lack of self-activity is by no means the same as forcelessness or inactivity: by inertia is denoted precisely that property of matter by which it presents an effect that has its cause outside itself, transforms every such effect into a state, and preserves it as a state after the cause has ceased. Logically seen, the effect is inseparable from its cause; but inertia posits separation, finite, real separation between effect and cause. The doctrine of inertia is the doctrine of external causes, and precisely thereby a doctrine of the objective world, for the objective world is a world of external causes.

It is so far from being the case that ``the communication of an impact'' \opage{339} should denote ``the sublation of inertia in manifestation,'' that the motion corresponding to the impact is precisely an expression for the existence of inertia, insofar, namely, as it is an expression for the fact that the effect can continue after the cause has ceased. When, then, the effect continues unchanged by means of inertia, the motion is uniform. Uniform motion, or, what is synonymous with this, constant velocity, is therefore in mechanics an exact expression not for force but for inertia. If the space $s$ is a function of the time $t$, then the first differential coefficient $\frac{ds}{dt} = a$ or, since from this it follows that the second differential coefficient must become $0$, $\frac{d^{2}s}{dt^{2}} = 0$, is a notation for inertia.%
\footnote{``Insofar as $s$ is a space already traversed in the time $t$, one cannot yet know whether this space has been traversed with variable or invariable velocity. To determine the velocity, one must give the time an increment, that is to say: provide time for continued motion. If the increment of the time is now called $h$, we have the series:
  \[ f(t+h) = s + \frac{ds}{dt}\,\frac{h}{1} + \frac{d^{2}s}{dt^{2}}\,\frac{h^{2}}{1\cdot 2}
     + \frac{d^{3}s}{dt^{3}}\,\frac{h^{3}}{1\cdot 2\cdot 3} + \cdots \]
  If the increment $h$ is a finite magnitude, and the velocity constant, the whole series will reduce to $f(t+h) = s + \frac{ds}{dt}\,\frac{h}{1}$; if $h$ is infinitely small, equal to $dt$, the same will hold even when the velocity is variable. On the other hand, when $h$ is finite and the velocity variable, the terms of the series can be increased to infinity. Seen from the standpoint of the concept, the series indicates a totality of objective determinations of reflection. The individual terms can, $s$ signifying space, $v$ velocity and $\varphi$ force, be rendered by different mutually corresponding expressions:
  \[ \begin{aligned}
     f(t+h) &= s + \frac{ds}{dt}\,\frac{h}{1} + \frac{d^{2}s}{dt^{2}}\,\frac{h^{2}}{1\cdot 2}
        + \frac{d^{3}s}{dt^{3}}\,\frac{h^{3}}{1\cdot 2\cdot 3} + \cdots\\
     &= \ldots + v\,\frac{h}{1} + \frac{dv}{dt}\,\frac{h^{2}}{1\cdot 2}
        + \frac{d^{2}v}{dt^{2}}\,\frac{h^{3}}{1\cdot 2\cdot 3} + \cdots\\
     &= \ldots + \ldots + \varphi\,\frac{h^{2}}{1\cdot 2}
        + \frac{d\varphi}{dt}\,\frac{h^{3}}{1\cdot 2\cdot 3} + \cdots
     \end{aligned} \]
  The categories \emph{motion}, \emph{velocity}, \emph{force} and \emph{change in force} are thus determined by objective reflection. Immediately, $s$ denotes a length of space; but in the equation $s = f(t)$ this length of space is seen in reflection of a length of time; it owes, indeed, its determinate being to time, insofar, namely, as it is traversed in a given time. Immediately, $v$ denotes a velocity, but in the expression $\frac{ds}{dt}$ the concept of velocity is determined by objective reflection of space and time in the given ratio. Immediately, force is denoted by $\varphi$; next, by the reflection of velocity and time in the determinate ratio $\frac{dv}{dt}$; finally, since $v = \frac{ds}{dt}$ resolves into the reflection of space and time, by a doubled reflection in the ratio $\frac{d^{2}s}{dt^{2}}$. In this way the changes of force, with respect to the ratio of space and time, can be determined as reflections of reflections to infinity.'' Phil.\ Propæd. p.~165--66. Cf.\ Mathem.\ og Dialktk. p.~78--88.}
\opage{340} Admittedly, one could here, from the standpoint of ``speculative logic,'' raise the objection that there must surely be force in inertia, since mechanics itself ascribes to it the maintenance of a motion once given. For if the maintenance of motion is a manifestation of inertia, then, as it seems, inertia must also be the force whose sublation is seen in the manifestation. But such an objection holds only so long as, from a one-sided interest in the logical, one overlooks the antilogical. Were inertia to be conceived as mechanical force, it would have to be conceived as cause; but inertia is not cause, it is ground. By putting the law of inertia in place of the force of inertia, physics has, with fine tact and in an acute manner, indicated a categorial difference between ground and cause. According to the law of inertia, it is not an internal cause, but on the contrary the lack of an internal cause, that has as its consequence that the body's transition from rest to motion must be determined by an external cause: such a lack \opage{341} of cause, which nevertheless seems to act as cause, such a cause negatively reflected and sublated in reflection, shows precisely what it means that something is not cause, but ground. As ground, inertia, when more closely determined, is not a real but a formal ground; its phenomenon is the mechanical tautology: rest is rest, motion is motion; its essence is reflected in the antilogical identity: this is this. As the categorial expression of antilogical identity: this rest is what it is, namely this rest; this motion is what it is, namely this motion, inertia can be designated just as exactly by an immediate rest,%
\footnote{A caused rest, a rest designated by the fact that the resultant of the cooperating forces is $0$, cannot come into consideration here, where abstraction is made precisely from all external causes.}
as by a uniform motion. The equation $\frac{ds}{dt} = C$ is, insofar as $C$ is a constant, an exact designation of inertia; if $C$ has a finite value, then inertia is designated by uniform motion, motion with unchanged velocity; if $C = 0$, then inertia is designated by rest, the rest that continues according to the judgment: rest is rest, and that would therefore continue for a whole eternity, if an external cause did not come along and set what is at rest in motion.

A striking proof that inertia, although it is not an acting cause, is nevertheless a determining ground, is precisely the peculiar way in which, in mechanics, it unites with mechanical force. If force, the acting cause, sublated inertia, the second differential coefficient $\frac{d^{2}s}{dt^{2}}$ could not possibly represent force. But since the moved mass, on account of inertia, retains the velocity imparted to it, the velocity that force brings about in each equal element of time becomes a constant increment, an increment by which \opage{342} the velocity already attained is steadily increased; thus a uniformly increasing motion arises. If $\frac{ds}{dt} = v$, and $v = f(t) = gt$, then $gt$ means that the force $g$ brings about the velocity $v$ in the time $t$. By acting in the time-element $dt$, $g$ produces the velocity-element $dv$, so that $dv = g\,dt$, and $\frac{dv}{dt} = g$. But since $v = \frac{ds}{dt}$, we have
\[ \frac{dv}{dt} = \frac{d.\,\dfrac{ds}{dt}}{dt} = \frac{d^{2}s}{dt^{2}}. \]

Seen from the point of view of mechanics, force is the cause of motion, and motion therefore the effect of force. In a certain respect one must now indeed grant ``speculative logic'' that ``physicists regard force as independent of its manifestation,'' insofar, namely, as they keep to the effect without asking about ``the unknown cause.'' But since mechanics nevertheless has the cause reflected in the effect, and on the whole can take an interest in the cause only insofar as it makes itself known in the effect, one can with equal right say that mechanics does not regard force as independent of its manifestation, but sees force in the manifestation, and designates precisely its identity with the manifestation by positing the second differential coefficient of the space traversed, and hence the differential coefficient of the velocity, which measures the manifestation, as the exact expression for force. The assurance in abstraction with which mechanics, from its point of view, knows how to remove all irrelevant considerations when the question concerns mechanical force, or, still more precisely, what is mechanical in force, shines forth from its ``principle for the measurement of forces.'' For this principle presupposes ``that any forces whatever, even if of different physical origin, never modify one another's magnitude by their simultaneous action on the same point. When, e.g.,'' (following a definition already given, namely that two forces are equal when, acting on the \opage{343} same point in opposite directions, they are in equilibrium) ``$P = Q$ and $Q = R$, one must also always have $P = R$, i.e., when the forces $P$ and $Q$ acting on the same point in opposite directions are in equilibrium, and likewise $Q$ and $R$ are in equilibrium, then equilibrium must also necessarily obtain between $P$ and $R$ acting against each other. Likewise, when $P = P'$ and $Q = Q'$, one must also in all cases have $P + Q = P' + Q'$, i.e., $P$ and $Q$ acting simultaneously on a point in the one direction must give equilibrium against $P'$ and $Q'$ acting simultaneously on the same point in the opposite direction, if separately $P$ and $P'$ give equilibrium, as well as $Q$ and $Q'$. Herein consists the law of the \emph{mutual independence of mechanical forces from one another}, which shows itself in nature by the fact that all of nature's results concerning the transformation and equilibrium of forces are on the one hand necessary (mathematically deduced) consequences of the principle stated for the measurement of forces, while on the other hand they agree exactly with experiments.''%
\footnote{Ramus. Analytisk Mechanik. Kbhvn. 1852 p.~4.}
Mechanics has thereby clearly shown that it not only lets forces reflect themselves in their manifestations, but even regards this reflection as the true expression of their mechanical essentiality. If here, in the measurement of forces, no account is taken of their ``different physical origin,'' then here, in the determination of force itself, no account is taken either of the individual physicist's subjective representation of what may possibly lie behind the manifestation: scientifically seen, mechanical force is as the manifestation, the manifestation as the force.

As the differential coefficient of velocity, when velocity is a function of time, as the second differential coefficient of the space traversed, force reflected in its manifestation is expressly determined as accelerating force (force accélératrice). Accelerating force can in turn be either constant or \opage{344} variable. If the force is constant, this must be seen from the manifestation: the manifestation of constant force is constant increment of velocity. If the velocity that the force has produced by acting in the time $t$ is determined by the equation $v = gt$, so that $t$ is the number of units of time, e.g.\ seconds, elapsed since the beginning of the action, then $g$, the velocity the force imparts in one second, is precisely the measure of intensity, and hence the manifestation of the force is identical with the force itself. If, on the other hand, the accelerating force is variable, then the increment of velocity, the manifestation of the force, must naturally also be variable. But although the force $\varphi$ in the equation $v = \varphi t$ is itself variable, insofar as it is a function of time, $\varphi = f(t) = \psi t$, $\psi$, as cause of the change in $\varphi$, may nevertheless very well be constant. The cause of the change of force is itself force; if this force $\psi$ is constant, then its manifestation, namely the change of force $d\varphi$, must likewise be constant. The force $\psi$, reflected in its manifestation and changing the force $\varphi$, is then determined by the equations:
\[ \psi = \frac{d\varphi}{dt} = \frac{d^{2}v}{dt^{2}} = \frac{d^{3}s}{dt^{3}}. \]
In the matter itself there is nothing here to prevent the force $\psi$ that changes the force $\varphi$ from in turn being changed by a force $\omega$, while the change of velocity must nevertheless continue to be the measure of the accelerating force, however many changes it may undergo; just as $\frac{dv}{dt}$ is a manifestation of $\varphi$, so $\frac{d^{2}v}{dt^{2}}$ is a manifestation of $\frac{d\varphi}{dt}$, which in turn is a manifestation of $\psi$; we have here, therefore, an example showing how mechanics can let a force reflect itself in a manifestation of a manifestation of a manifestation to infinity.

Reflected in velocity, force is an operational concept that articulates the subjective by means of the objective, the logical by means of the antilogical. The essence of the change of velocity and of force is in reflection, hence from the subjective side of the matter, \opage{345} one and the same essence; as regards the phenomenon, the immediately objective, on the other hand, the extensive becomes an extensive measure of the intensive, and the line, the path of motion, an image of force. ``A system of forces acting on a point $M$ in any directions whatever in space can be represented by the system of lines $\mathrm{MA}_{1}$, $\mathrm{MA}_{2}$, $\mathrm{MA}_{3}\ldots$, whose directions coincide with those of the forces, while at the same time their lengths are in proportion as the magnitudes of the same forces. One can take each of the lines, by reference to the chosen unit of length, to be determined by the same number as the corresponding force acting along the direction of the line, and one then says that \emph{the force is equal to the line}.''%
\footnote{C.\ Ramus, Analyt.\ Mechan. p.~4.}
Already from this it is seen how the objective side of objectivity reflects itself in the subjective.

For immediate intuition, pure space must necessarily present itself as the pure conflict between the subjective and the objective. As the proper external form of externality, as the schema of otherness, the medium of antilogical relations, the extreme abstraction of everything objective, space, pure space, must itself be objective. In this its objectivity it accordingly shows itself as an eternal being without becoming, an infinity resting in itself, an existing emptiness that is indifferent to everything that fills it. The space-filling objects must of necessity presuppose space, but so objective is pure space that it seems not to presuppose the space-filling objects at all. Abstraction, which can remove all the objects and retain space, cannot, on the other hand, remove space and retain the objects. But the abstraction that was to prove the absolute objectivity of space, of pure space, turns round at its decisive turning point and furnishes precisely a proof of the absolute subjectivity of space. Abstraction is, as we have seen from the very beginning, \opage{346} a general act of subjectivity, an activity of thought, a function of subjective knowledge. When abstraction has taken away all intuitable, space-filling objects, intuition is empty, and empty intuition, deprived of everything intuitable, now intuits in empty space its own emptiness. Since the emptiness is brought about by abstraction, and abstraction is a subjective act that has significance only for the subjective, it is naturally pure imagination if anyone supposes that, besides the subjective emptiness, there should also be an objective emptiness here. On the contrary! an objective emptiness is the same as an objective nothingness, a determinately existing nothing, a subjective phantom: pure space is absolutely subjective.

Only when the concept of space has been made clear from the subjective side by reflexes of what is objective in the phenomena can the objective validity of space be cognized. The punctual starting point of the development is already given by the fact that space is seen as a function of time. For it is not possible to reflect this function through without becoming aware of the thoroughgoing difference between a unit of space and a unit of time. If we go back once more to even or uniform motion, then, in the equation $s = vt$, the time $t$ is the number of units of time, the velocity $v$ the number of units of space traversed in one unit of time, the space $s$ the number of units of space that corresponds exactly to the number of units of time in $t$. From all sides, therefore, the function is determined by the relation between the units of space and the units of time. Now if these units of space and time were not qualitatively different, the function itself would be impossible. But if we lay stress on the qualitative difference, it is seen at once that the unit of time, by its nature, is and must be subjective, namely ontologically subjective.

That the unit of time is essentially subjective means, in other words, that it is negative, that it is in a ceaseless coming and vanishing, that its being is becoming, its essence \opage{347} reflection. But if the unit of time is negative, then the unit of space, qualitatively different from the unit of time, must necessarily be positive; if the unit of time is logical, the unit of space must be antilogical; if the unit of time is subjective, the unit of space must be objective. This opposition is made clear when we consider what it means that, after the lapse of a number of units of time corresponding to $t$, a number of units of space corresponding to $s$ has been traversed. For what has become of the elapsed units of time? That they have elapsed means that they have vanished, are irrevocably lost; no retrograde motion can here restore the vanished immediacy; the moments of the past belong to memory, they are sublated in subjectivity.

It is otherwise with the units of space; these elements have been covered, but not therefore sublated; they have been traversed, but have not elapsed: they did not come into being with motion, and therefore could not cease with motion. Here, in the unit of space, is an antilogical being that resists the reflection prevailing in becoming. The space $s$ covered in the time $t$ with the velocity $v$ is a determinately existing length of path, a well-preserved line whose finite parts are equally immediately in being. A proof of this, accordingly, is also the fact that the same, the very same space can be traversed countless times by countless forward and backward motions. So long as $v = v$ and $s = s$, one will naturally also have $t = t$; but while $t = t$ never is or can be the same $t$, $s = s$, on the other hand, can be the same $s$. That the traversed units of space, in qualitative opposition to the units of time, preserve an immediate determinate being is a decisive expression of the objectivity of space. The objectivity of space, in opposition to the subjective essence of time, is accordingly also made intuitable in the three dimensions of space. Why, one asks, has time only one dimension, while space has three? The answer is: because time is essentially subjective, space objective. The ideal \opage{348} length of time, valid only for reflection, has no simultaneous points at all; it would indeed also be an absurd contradiction if there really could be two times at one time. Now since the line of time cannot have immediate determinate being, it cannot be thought to move and describe any surface either; time, in other words, cannot have any breadth: time cannot be squared. It is otherwise with space. The line of space, which has all its points in simultaneous determinate being, can---in objectification---move and describe a surface; the surface, which likewise has its parts in simultaneous determinate being, can in turn move and form a mathematical body: the mathematical body is an antilogical form, an exact expression of the objective essence of space. By means of the three dimensions infinitely many directions of motion are now possible; but for that reason all possible directions can also be determined by the three dimensions. Mechanics has known how to put this to use, in that it has employed the coordinates of the so-called trihedral angle of analytic geometry. If we take these coordinates to be $x$, $y$, $z$, then an equation such as:
\[ \left(\frac{d^{2}x}{dt^{2}}\right)^{2} + \left(\frac{d^{2}y}{dt^{2}}\right)^{2}
   + \left(\frac{d^{2}z}{dt^{2}}\right)^{2} \div p^{2} = 0 \]
can show how the different directions in which accelerating forces bring about different motions can be exactly specified. A single glance into such a system of motions and manifestations of force must be sufficient to convince us of the objectivity of space.

And yet the system of objective determinations here indicated is still only a shadow-outline of the objective, not the actual objective itself. That the original relation between space and time is grounded in the existence of matter can be seen from the identity, namely the mechanical identity, of inertia and force. How different inertia is from force has been shown above; but the difference is a difference within unity, and here the emphasis now falls on the unity. Immediately, space and \opage{349} time could never enter into relation with each other; their relation in the pure concept of velocity, $\frac{s}{t} = v$, is an abstraction that presupposes matter. It is therefore not space and time, but matter, that in reality is the original.

Now it is not the case that only inertia is in matter, while force, the force that acts on matter from without, should be outside all matter: the force that acts \emph{on} a matter is itself in a matter. The contradiction that matter at one and the same time has the acting force in itself and yet outside itself is resolved in the determination of realized otherness. Realized otherness is the existence of matter; therefore the one atom is outside the other, the one mass outside the other, the one acting cause outside the other. The acting causes are in this connection all of them material causes; and yet it looks as though matter were constantly acted upon without itself acting. Now this follows directly from the fact that the force that is in a given mass never acts on the mass itself, but always on another mass, or, what is the same thing, the force that acts on a given mass never proceeds from the mass itself, but from another mass; realized otherness, actualized selflessness, has its mechanical expression in the system of external causes, a system that is articulated exactly by reciprocal determination of inertia and force.

Force, which by virtue of inertia is an accelerating force, is then by virtue of matter also \emph{a moving force} (force motrice). If velocity is posited in abstracto as the manifestation of force, $g = v$, then there is a double presupposition here that shines through. In the first place it is presupposed that the force $g$ acts in one unit of time, for all acting takes place in time, and makes time actual. The equation $g = v$ is to that extent only a special case of the general: $v = f(t) = gt$. In the second place it is presupposed that the force \opage{350} acting in the unit of time has a determinate matter, a determinate atom, on which it can act; hereby the equation $v = gt$ is then transformed into $mv = mgt$. Here there now enters a peculiar reciprocal relation between what we shall call \emph{mechanical products and mechanical concepts}. Insofar as the mechanical concept is presented as a product, its qualitative-quantitative moments step apart from one another in an exact manner, such that each moment becomes a factor: the logical reflects, in discretion, the antilogical. Insofar as the mechanical product is expressly conceived and determined as a concept, the particular factors, the antilogical elements, are transformed into reflexes and subsidiary determinations in the logically articulated unity; the mechanical concept is, by virtue of the product, an operational concept. In the product as mechanical real concept, the function presents a transparent unity of mathematical, mechanical, and logical operation. What this means, a consideration of particular products can easily show us.

Let us first examine somewhat more closely $gt$, the product of force and time. Seen from a general-logical point of view, it would almost seem absurd that two such heterogeneous factors as force and time should be capable of being multiplied by each other. One cannot possibly multiply the quality tone by the quality color; how does it come about, then, that in mechanics one can multiply the quality force by the quality time? It is a fundamental unity consisting in qualitative units that matters in this connection. In order to manifest itself, force must originally presuppose time; the determination of time therefore does not come from outside, but belongs as essentially to the manifestation itself as the subjective belongs to the objective. Conversely, the manifestation for its part does not enter into time either, as though an empty time could be held fast by itself over against that by which it happened to be filled: time has existence only in and through what fills it, just as the subjective has existence only in and through the objective. The one \opage{351} of the two qualitative units can therefore, according to the fundamental unity, be posited as a quale, while the other is determined as quantum. If, e.g., the unit of force, $g = g$, is posited as quale, the unit of time can be posited as quantum. That the unit of time is posited as quantum makes itself known by the fact that $t$, the independent variable of the function, is assigned values $1, 2, 3, 4\ldots$ The products $g$, $2g$, $3g$, $4g\ldots$ thus indicate the exact difference between manifestations of a constant force insofar as it acts in $1, 2, 3, 4\ldots$ units of time (seconds). If, on the other hand, one posited the unit of time $t = t$ as quale, the unit of force would have to be transformed into quantum. Instead of $v = f(t) = gt$ we would thus get $v = f(g) = tg$; $g$, the independent of the function, could then receive the values $1, 2, 3, 4\ldots$ From this it is seen that a force $= 4$ must in one and the same unit of time accomplish 2 times as much as a force $= 2$, 4 times as much as a force $= 1\ldots$ From both sides, then, it becomes evident that a product such as $gt$ is something more than a simple number, that it is a mechanical real concept. In the equation $v = f(t) = gt$, $v$ presents the concept in its immediacy, $gt$ the moments of the concept, and in such a way that the quantifying falls on time; to $g\,dt$, the time-element qualified by force, there accordingly corresponds $dv$ as velocity-element.

That the product $gt$, however, suffers from an abstraction incompatible with actuality can be seen from the fact that the manifestation $v$ corresponding to the essence of the product is still held in pure quantifying, and therefore lacks qualitative determinateness. If velocity is to be qualified, it must expressly be referred to a matter that has the velocity. When velocity is then determined as a velocity of mass, qualified velocity shows itself as a \emph{quantity of motion}.

From the side of the concept, it goes with the product of mass and velocity as with the product of force and time. When velocity, as the manifestation of force, is qualified in the product by mass, it turns out that the product in which time was \opage{352} qualified by force, namely $gt$, still lacks a factor essential to qualification, namely matter. If abstraction is made from all matter, the product $gt$ is meaningless; force, which acts in time, cannot act on time: it is not time, but matter, that furnishes the points of application (points d'applications) without which force cannot possibly manifest itself as force. If the product $mv$ is to be conceived as a concept, and the concept is to have reality, then the abstract product $gt$ must be transformed into the concrete $mgt$, a product which signifies that the force $g$ acts, or rather has acted, on the mass $m$ in the time $t$. Here, then, we again have the equation given above, $mv = mgt$, whereby the force $mg$ is determined as moving force.

Moving force, that is, the quantity of motion that the accelerating force imparts to the moved mass in a determinate unit of time, has as its measure the product of the velocity produced and of the mass it carries; but the mass that is moved forward in time with a certain velocity is of course at the same time moved forward in space: in every spatial motion there must naturally correspond an element of space to every element of time. If $\frac{ds}{dt} = v$, then $dt = \frac{ds}{v}$; the quotient $\frac{ds}{v}$ is then here an exact expression of the well-known proposition: time is equally long, but not equally useful; for $\frac{ds}{v}$ is constant, since $dt = dt$ is constant; now since $v$ varies, $ds$ must of course also vary, but $ds$ is precisely an indication of what can be accomplished with the velocity $v$ in a moment $dt$: in such a moment, namely, the force $g$ can move the mass forward an element $ds$ in space. Instead of $mg\,dt$ we then put $mg\frac{ds}{v} = mdv$; yet this is only a transitional determination, which is expressly seen from the fact that the moments $v$ and $dv$ belonging to \opage{353} the concept of velocity fall on opposite sides of the equation. If the separated moments are united, we obtain $mgds = mvdv$ and by integration $mgs = \frac{mv^{2}}{2}$, and hence at one stroke two peculiar products, and hence two new concepts of force, namely \emph{working force} $mgs$ and the so-called \emph{living force} (force vive); these forces stand in the exact relation that the living force $mv^{2}$ is equal to twice the quantity of work $2mgs$.

Inertia, force, space, time, motion, velocity: all these mechanical determinations are moments of one and the same fundamental concept, namely matter. In matter power has posited its definitive expression for what is objective in objectivity: in matter the antilogical has immediate reality; but this reality conditions precisely an infinite reciprocal play with the logical. In opposition to matter, inertia is logical; for it is indeed ground, ground for matter's always remaining in the state in which it is put by force. In opposition to force, on the other hand, inertia is antilogical; for it holds to tautology until force, by converting rest into motion and by changing velocity, introduces difference. In opposition to varying velocity, force is logical; for force as force is the inner, and velocity in mechanical motion is its manifestation, hence its immediate outer. In opposition to working force, moving force is logical, but antilogical in opposition to accelerating force; for what is concrete in finite relations, individualized by special conditions, is always antilogical in opposition to the abstract and to that extent universal. In countless ways, then, the reciprocal determination of the logical and the antilogical can be conceived; but the mechanical function, the function of the external cause, has the peculiar significance that it holds the moments of the concept that belong together in sharp separation, and \opage{354} precisely through this discretion lends itself to the punctual development of objectivity.

\begin{center}{\bfseries β)\quad Dynamic Functions.}\end{center}
\addcontentsline{toc}{paragraph}{β) Dynamic Functions}

If the mechanical function is determined as the function of the external cause, then the dynamic function must be determined as the function of the internal cause, a categorial expression of causal continuity. The opposition between mechanical and dynamic activity is generally conceived in such a way that the one falls outside the other: either an activity is mechanical or it is dynamic.

Many misunderstandings harmful both to life and to science have their origin here. Separated from the dynamic, the mechanical takes on the appearance of being something merely finite, pieced together, materialistic, devoid of idea. The thinking that holds fast to the ideal must therefore look down on the mechanical and seek truth in the dynamic; but separated from the mechanical, the dynamic becomes something hovering, indeterminate, and fantastic. Hegelian logic has quite rightly seen that the mechanical and the dynamic must be two sides of the same matter, two concepts of the same concept; but on this point too it becomes apparent how little the immanent development of the concept satisfies the demands of actuality. As the concept, in its progressive self-movement, raises itself from mechanical to dynamic objectivity as from the lower to the higher stage, the accessory representation gains the upper hand that the dynamic should constitute a higher sphere of concepts, lifted above the mechanical; the dynamic then continues, as before, to be abandoned to speculative fantasy-intuition. That the misunderstandings which thus germinate in logic must have serious consequences for a development of the concepts in more concrete form is understandable. The Hegelian philosophy's total misapprehension of the independence, the peculiar aims, and the procedure of the special sciences can on the whole be said to \opage{355} be grounded in its total misapprehension of the relation between the mechanical and the dynamic.

Natural science, which must everywhere look to the acting causes and their conditions, chooses mechanical points of departure and avails itself of discretion; thus astronomy speaks of centripetal force, tangential force, centrifugal force, and normal force as of different, determinately existing forces; thus optics speaks of particles of ether, waves, wave-breadths, oscillations, etc.; even physiology seeks, and must seek, a mechanical foundation for itself,%
\footnote{Cf.\ Forelæsngr.\ over ``Phil.\ Propæd.'' academic year 1860--61. pp.~177--200.}
in that it regards the organism and its individual parts as ``a machine or a similar physical-chemical apparatus,'' lets the heart count as ``a pump-works furnished with suitable valves,'' the arteries as ``elastic pipelines,'' etc. What is finite, merely hypothetical, and devoid of idea in this is what speculation now wants to attack; it wants, as it says itself, to vindicate the rational; but it does so in a wholly irrational manner. For no procedure can surely be more irrational than the one Hegel chose in his philosophy of nature, when he attacks Newton's mathematical treatment of the Keplerian laws.%
\footnote{Cf.\ Forelæsngr.\ over ``Phil.\ Propæd.'' academic year 1861--62. pp.~57--66.}
In order to sublate the one-sidedness in the mechanical partitioning of cooperating moments, one must above all seek to enter into the presuppositions of exact analysis and thereby show how the finite dissolves in the infinite, the conditions in the unconditioned, the thing-causes in the ground-causes. But instead of doing this, Hegel, from the speculative height of his logic, has simply rejected the exact-mechanical way of viewing things, and has accompanied the rejection with a critique in which he betrays a striking unfamiliarity with what he wants to criticize. Instead of securing the real concepts in their exact \opage{356} origin, he botches them with irrelevant reflections.

``The methodological fundamental error, that the original mode of formation of the various real concepts is botched from the beginning, is most intimately connected with another equally essential error, namely: letting speculative grounds step in place of acting causes, so that all the intermediate determinations of finitude vanish as by a conjuring trick before the logical Idea, which everywhere acts alone, is alone mighty, is capable of everything. Instead of an attempt to dissolve dialectically the material presuppositions, conditions, and intermediate causes of mechanical physics and, by the dialectical dissolution, transform them into causal moments for the Idea-cause acting in the totality, we meet in Hegel a direct, hence undialectical, rejection of the presuppositions of mechanics and an immediate, hence likewise undialectical, appeal to the all-moving Idea.''%
\footnote{Forelæsngr.\ ov.\ ``Phil.\ Propæd.'' 1861--62. pp.~73--74.}
``Die Centrifugalkraft als das Entfliehenwollen in der Richtung der Tangente, soll läppischer Weise den Himmelskörpern durch ein Werfen auf die Seite, eine Schwungkraft, einen Stoß zukommen, den sie von Haus aus erhalten hätten. Solche Zufälligkeit der äußerlich beigebrachten Bewegung, wie wenn ein Stein an einem Faden, den man schräg wirft, dem Faden entfliehen will, gehört der trägen Materie an. Man muß also nicht von Kräften sprechen \ldots''%
\footnote{Hegel Vorlesung.\ üb.\ d.\ Naturphil. Werke 7ter Bd. 1.\ Abtheil. Berl. 1842. p.~97.}
It must constantly be held fast that the misdirection has its source in logic; only thus is it explicable that any attempt whatever at a philosophy of nature in the Hegelian direction must fail, even if the palpable errors that can be demonstrated in Hegel's own work from the standpoint of the special science were corrected ever so carefully: the fundamental error rests \opage{357} on a superficial rejection of the mechanical and a correspondingly superficial predilection for the dynamic.

Hegel has, and rightly, reproached the Schellingian philosophy of nature for its empty schematism, its arbitrary symbolism, and its lack of true development of the concept. When the Hegelians polemicize against the Schellingians, one would almost believe from their utterances that they held to the fact, to the actual, to what accords with the understanding; thus Bayrhoffer remarks, in an introduction to what he calls ``the theory of natural development,''%
\footnote{Beiträge zur Naturphilosophie. Leipzig 1839.}
that empiricism and speculation, each from its own side, must both agree in opposing the boundless arbitrariness that began to gain entry into science with the Schellingian philosophy of nature. ``In letzterer Beziehung,'' it says with regard to speculation, ``war es \emph{Hegels} Riesengeist, welcher den \emph{Begriff} an die Stelle der \emph{Phantasie} setzte, und ohne die tieferen Anschauungen der Polarität zu verkennen, sie der Willkühr und Spielerei zu entreißen, die Natur selbst als einen vernünftigen Organismus zu erfassen suchte. Er hat auch hier \emph{den Grund} des absoluten Begriffes gelegt \ldots''%
\footnote{Anfr.-Skr. Indledning. p.~14.}
How the author himself wants this statement to be understood can best of all be seen from the way in which he applies ``the dynamic determinateness.'' ``Zunächst,'' he says of ``the planetary organism,'' ``muß die \emph{dynamische} Bestimmtheit demgemäß hier \emph{die Wiederholung} der allgemeinen Formen des kosmischen Organismus, aber unter dem Exponenten \emph{der Besonderung} und somit durchweg schon selbst der \emph{physikalischen Bestimmtheit} sein, mithin concrete Polarisirung in sich, so daß die Kräfte jetzt nicht mehr unmittelbar kosmische, sondern von \emph{diesen getragen tellurische} und tellurisch specificirt sind. Daher ist: 1) \emph{die Schwere} nun \emph{specificirt} \opage{358} in Einheit mit dem \emph{physikalischen} begriffsmäßigen Bestimmen, Moment an ihm, in ihrer \emph{freien Gestalt} erscheinend als \emph{Magnetismus}, \emph{engegengesetztes} Anziehen an dem \emph{Einen} Körper, identisch mit dem Körper als \emph{Cohäsion}, und deren Erzittern in sich, \emph{Klang}; 2) eben so das Licht \emph{specificirt}, \emph{frei} erscheinend als \emph{Electrismus} (und Galvanismus), als \emph{polarisirtes} Licht in der \emph{Differenz} körperlicher Formen (vergl. Hegel, Encyclopädie § 324. Anm. S.~328), dann aufgehoben in der \emph{reellen} Form erscheinend als \emph{Farbe}, in ihr verzehrend als \emph{Wärme} und \emph{Feuer}, immanente Negativität; 3) Diese dynamischen Momente resp. Stufen sind aber \emph{realisirt an den reellen} Formen und gehen in ihnen, zunächst im \emph{Galvanismus}, zugleich \emph{fort} zu den durch und durch qvalificirten Stoffen und \emph{deren} realen Spannung und Vermittlung \ldots'' As regards the chemical foundation, the following may serve to characterize the specification: ``Indem \ldots\ \emph{das solare} Princip sich die Materie des Planeten \emph{einfach} unterwirft, oder dieselbe sich solar \emph{bestimmt} als Reflexion aus der Sonne in sich, erzeugt sich das chemische Element des \emph{Sauerstoffs}, der chemische Lichtstoff, das oxydirende und Anderes verbrennende, \emph{allgemein in sich reflectirte Element}, der in den \emph{Gegensatz} eingegangene Lichtäther, welcher daher nicht mehr das Leuchten an ihm selbst, sondern das \emph{Befeuern} des Andern ist.''%
\footnote{Anfr.\ Skrift. pp.~92--94.}
To complete the picture, we cite a couple more features of the ``dynamic foundation'' that bears the living organism. ``Die allgemeinen \emph{dynamischen} Momente nun zunächst sind wohl wieder die Reproduction von Schwere, Licht und deren concreter Einheit, wie dieser im tellurischen Organismus specificirten Momente. Aber sie stehen nun unter der Form der \emph{Seele}, der negativen durchdringenden Einheit, des Fürsichseins. So ist nun: 1) die Unmittelbarkeit der Schwere als selbst zur \emph{Seele} \apage{359} reflectirt die fürsichseiende, idealisirte, das Centrum in sich \emph{selbst} tragende, nicht in Anderem habende Einheit, das \emph{reflectirte Insichsein} der Totalität, und frei hervortretend, in dem animalischen Organismus, \emph{die Sensibilität}, die \emph{Empfindung}; 2) das Scheinen des Lichts nur das Herausgehen der Sensibilität, des Insichseins gegen \emph{Anderes} und die hiermit eintretende Erregung und Bewegung in sich und gegen Anderes, als \emph{frei} erscheinend, in dem animalischen Organismus, die \emph{Irritabilität}; 3) die Rückkehr aus dem Andern in sich und so \emph{die Reproduction} \ldots''%
\footnote{Anfr.\ Skr. p.~129.}
If one compares this specification with Schelling's cited above, one will without doubt find that the difference between them is by no means as great as one should have expected, given the opposition appealed to between ``concept'' and ``fantasy.'' We find it speculative that $A$ in the first power is gravity, $A$ in the second power light, $A$ in the third power the organism; but we find it no less speculative that \emph{gravity}, which ``under the exponent of the particular'' becomes \emph{magnetism} and \emph{sound}, can ``under the form of the soul'' become \emph{soul}, \emph{sensibility} and \emph{feeling}; that light, which under the exponent of the particular specifies itself into \emph{electrism}, \emph{galvanism}, fire and \emph{heat}, under the form of the soul specifies itself into \emph{irritability}. The Schellingian ``fantasy'' and the Hegelian ``concept'' have, in the respect of the philosophy of nature, a highly striking mutual resemblance; the difference is only this, that Hegel has a greater wealth of logical forms at his disposal. As regards the things of nature, what is objective in objectivity, Hegel's ``concept'' is just as subjective as Schelling's ``intellectual intuition.'' If intellectual intuition is required in order to intuit water as ``a wet flame'' and the plant as ``the first organic water on earth,'' then intellectual \apage{360} intuition is surely also required in order to intuit how ``the solar principle simply (einfach) subjects to itself the matter of the planet,'' or how the planet, in determining itself solarly as reflection of the sun into itself, begets the chemical element of oxygen, the chemical light-substance, the light-ether that has entered into opposition, etc.

It is a misunderstanding if anyone supposes that such specifications might perhaps be justified by means of special explanations; the basic features of a speculative intuition would never conform to special explanations, but would rather make the special explanations nicely conform to ``the concept.'' It is a misunderstanding if anyone supposes that dynamic specifications of this kind could be corrected with regard to the researches of natural science; on the contrary! between the exact researches of natural science and the mystical-abstract speculations of the philosophy of nature every connecting link is lacking. It is, finally, a misunderstanding if anyone supposes that the determinations of concepts in the philosophy of nature should have a reliable corrective in the immanent method of Hegelian logic. Logic seeks only its own and demands no more than what it has itself given. If it is given that the universal is objective, and that the objective universal must go on specifying itself until it is in one way or another able to become subjective, then logic must find it perfectly rational that gravity, the objective universal, not only specifies itself into magnetism and sound, in order to become the particular, but, as a proof of how the objective itself becomes subjective, even specifies itself over into the singular and becomes soul. The soul is thus concretely subjective gravity, and gravity the abstractly objective soul.---Hegelian philosophy of nature stands and falls with Hegelian logic. If logic is scientifically tenable in its basic features, its principle, and its method, then the corresponding philosophy of nature is so too; if, on the other hand, the philosophy of nature stands in open conflict with actual research into nature, a conflict that does not raise
\opage{361} from individual mistakes and misunderstandings, but from principle and method, then the same must hold of logic.

The error in the speculative philosophy cannot possibly be that it brings forward the dynamic, for ``the dynamic activity'' certainly has reality just as much as the mechanical; but the fundamental error is that, under the pretense of taking the matter immanently and objectively, it apprehends the dynamic in subjective generality. The general thus becomes in this way a logical veil that covers over what is subjective and arbitrary in the specification; therefore, at this standpoint, however much talk there may be here of ``dynamic determinatenesses,'' ``dynamic foundations'' and the like, there can nevertheless be no talk, in the proper sense, of dynamic functions.

In philosophy as in the empirical sciences, the scientific use of the category \emph{dynamic} is essentially conditioned by the corresponding activity's allowing itself to be articulated. Now if the dynamic is a categorial designation for causal continuity, and if continuity cannot possibly be articulated without discretion, then it follows directly that the dynamic can be articulated only by the mechanical, that in other words the dynamic function can be apprehended only as a mirroring of the mechanical, determined by reflection. But the exact form of the mechanical function is mathematical; the transition from the mechanical to the dynamic must therefore be demonstrable in the mathematical: it is infinite discretion on which the transition rests.

As long as finite discretion is opposed to continuity in all immediacy, so long will the mechanical also oppose itself in a finite way to the dynamic; but when finite discretion dissolves itself into infinite discretion, it is transformed into a middle determination essential for the extremes; for infinite discretion is at one and the same time both continuity and discretion. That the doctrine of infinite discretion must essentially coincide with the doctrine of functions, and thus \opage{362} also with ``the analysis of the infinite,'' is easy to show. The equation $s = f(t) = vt$ presupposes, namely, not only that different values can be assigned to $t$, but that $t$ can vary by infinitely small increments $dt$. Insofar as different values are assigned to $t$, the discretion is finite; insofar as $t$ varies by constant increments of $dt + dt + \ldots$, the discretion is infinite. ``To the constant velocity $v = a = \frac{ds}{dt}$ there corresponds a simple continuity, a continuity which, since it is designated precisely by the differential coefficient of the first order being constant, we shall here call \emph{continuity of the first order}. If simple continuity can be comprehended only by a corresponding simple discretion that goes on to infinity, then the continuity doubled in itself cannot be comprehended either without a corresponding doubling of discretion. If the velocity, altering with each moment, gives us the series $v_{1}, v_{2}, \ldots v_{n-1}, v_{n}$ in $n$ successive moments, so that $v_{p} - v_{p-1} = dv$, whereby the spatial increment $ds$ in each $dt$ is successively transformed into $ds_{1}, ds_{2} \ldots ds_{n}$, then, since $v_{n} \mathrel{\substack{\textstyle> \\ \textstyle<}} v_{n-1}$ and $ds_{n} \mathrel{\substack{\textstyle> \\ \textstyle<}} ds_{n-1}$, there is no continuity to be discovered in these two series, insofar as they are considered \emph{immediately}. Continuity, which demands that the same must repeat itself without interruption in each moment, first shows itself when we discover, in the velocity as well as in space, the possibility of an element unchangeable under the changes. If it is thus given that the force $\varphi$ is constant, then the change of velocity $dv$ brought about thereby in each moment, $dt = dt$, also becomes constant; but the constant spatial element corresponding to this is not $ds$; it is, according to $dv\,dt = d^{2}s$, precisely the element $d^{2}s$. The continuity articulated by the functional law $s = at^{2}$, which presupposes that the differential coefficient of the second order is constant, we accordingly call, with regard to this, a \emph{continuity of the second order}. If doubled continuity demands a doubled \opage{363} discretion, then in a corresponding way the tripled continuity, namely that which corresponds to $s = at^{3}$, demands a tripled discretion, a spatial element of the third order. For here it is not the force itself, $\varphi$, but the cause of the change of force, $\psi$, that is constant. The constant $d\varphi$ now brings about in each moment, $dt$, a constant change in the increment of velocity, which in turn brings about a constant change, $d^{3}s = d^{3}s$, in the increment of the spatial increment; $d^{3}s$ is therefore the element of discretion in space corresponding to this law of continuity. It is the third differential coefficient that is constant here. \emph{The continuity is of the third order}.''

``From the articulations thus demonstrated it is seen that continuity is by no means interrupted by more and ever more changes being drawn in under the law of continuity. Should continuity ever break under the weight of the multitude of changes, then for the rising order of the elements an insurmountable barrier would at last have to show itself, without this, however, arising from any constant differential coefficient; but since this is a contradiction, since the order of the elements can indeed rise and rise to infinity, the continuity of the $n^{\mathrm{te}}$ order corresponding to a constant $n^{\mathrm{te}}$ differential coefficient, $\frac{d^{n}s}{dt^{n}} = C$, must hold, also for $n = \infty$. If $s$ is a fractional, irrational or transcendental function of $t$, e.g. $s = \cos.t$, then there is indeed here no
% Errata: (Trykfeil, p. 363 line 9 from below): printed "cos.t," --- the errata list corrects this to "cost,". Transcribed faithfully as printed.
constant differential coefficient; but this by no means signifies that the continuity is broken, but on the contrary that it is so infinitely articulated that one cannot demonstrate in isolation the atom of discretion corresponding to such a continuity.''%
\footnote{Forelæsn.\ over ``Phil.\ Propæd.'' academic year 1860--61 p.~77--79.}
--- If the articulating reciprocal determination of discretion and continuity is made clear in this way, then the relation between causal discretion and causal continuity, thus between the mechanical and the dynamic, is \opage{364} at the same time made evident. For rational mechanics does not stop at the fact that finite masses attract one another in a finite way, but interprets the attraction in such a way that every smallest part in every mass must, in that the attracting is attracted, be at one and the same time both point of departure and point of application for an acting force. Insofar as the question here is not of this or that actual thing, which is always subject to these or those conditions, but of the essence of phenomena, the principle of the actual, the finite is consistently dissolved into the infinite as into its ground, and this infinite is in turn apprehended as the real possibility that is always, under given conditions, actualized in a corresponding finite expression.

``If, to refer to an example, the material volume $\varrho v = M$ vanishes into an extensive element $dM$, then the force (weight) $P$ at the same time vanishes into an intensive element $dP$; a finite force (weight) is thus possible only through a correspondingly finite volume of the mass. If $p$ denotes the weight of a unit of volume, then the weight of a volume element, expressed in polar coordinates, becomes $dP = p\,r^{2}\,dr\,\sin\alpha\,d\alpha\,d\theta$. If this equation is now integrated, one obtains finite weight through finite volume. If this is, e.g., a homogeneous sphere, so that $p$ is thus constant, one obtains, for $r = a$, by the integration:
\[ \begin{aligned}
P &= p \int_{0}^{a}\!\!\int_{0}^{\pi}\!\!\int_{0}^{2\pi} r^{2}\,dr\,\sin\alpha\,d\alpha\,d\theta\\
&= 2p\pi \int_{0}^{a}\!\!\int_{0}^{\pi} r^{2}\,dr\sin\alpha\,d\alpha
   = 4p\pi \int_{0}^{a} r^{3}\,dr = \tfrac{4}{3}\,p\pi a^{3}.
\end{aligned} \]
With every variation of $a$, even the very smallest, a corresponding variation must take place in $P$, since $p$ is constant, and conversely; that is, every change in the material, the extensive, even the very slightest, is at the same time a change in the force, the intensive. From the analysis thereby indicated there follows, then, as an incontrovertible result, that matter presupposes force to infinity, just as force in turn presupposes matter to infinity, that matter \opage{365} can thus as little be without force as force without matter, that the determinate being of matter and of force is essentially one and the same being. Since force is possible only by means of matter, it is impossible to derive matter from force; since conversely matter is possible only by means of force, it is impossible to derive force from matter; the \emph{dynamic} view, which makes force the principle, is just as untenable as the \emph{materialistic}, which makes matter the principle: the principle is neither force as force nor matter as matter, but the dialectical unity of matter and force.''%
\footnote{Mathm.\ og Dialekt. p.~127--129.}
How the dynamic is articulated by the mechanical becomes evident, then, especially from the way in which the intensive in our cited example is articulated by the extensive. As the extensive magnitude $M$ vanishes into $dM$, the intensive magnitude $P$ vanishes into $dP$, and both are identical in their rudiment: the rudiment of $M$ is inseparable from the rudiment of $P$. As $dM$ is integrated to $M$, and $dP$ to $P$, they are both identical in their actuality: the actualization of $M$ is at the same time an actualization of $P$. Here, in the actual sphere, as regards the relation between its mass and weight, there is not an intensive magnitude by itself, connected with an extensive one, but an intensive magnitude whose reality consists in the extensive, an extensive magnitude whose reality is posited in the intensive. That the intensive has reality by means of the extensive, and conversely, means in other words that the dynamic has reality by means of the mechanical and conversely; but the intensive is indeed the inner, the extensive the outer; in reality the dynamic is therefore related to the mechanical as the inner to the outer.

The question of the conversion of the mechanical into the dynamic stands in the closest connection with the question of the relation of the discretion actualized in matter to material continuity, thus in the closest connection with the question \opage{366} of the reality of atoms. By placing the essence of matter in pure extension, Cartesius made space-filling matter a continuum to such a degree that he denied the existence of an empty space. Against this Gassendi maintained that there must necessarily be empty interstices, partly of smaller extent (vacuum disseminatum), partly of larger (vacuum coacervatum): ``Were there no empty interstices between the parts of matter, then in the first place all \emph{motion} would be impossible. Were space completely filled, then indeed not a single body could move from its place. Where then should such a body go? Since it cannot occupy the same place as another body, it would have to drive the other away from its place; but where is this other body to go when all space is occupied? The expelled body would then in turn have to expel another; but in this way the same difficulty returns again and again. Furthermore, all \emph{rarefaction} and \emph{condensation} would be impossible. Air, e.g., is indeed very strongly condensed. Now if the whole of space were filled with air, so that every part of space were occupied by its corresponding part of air, then an increasing condensation would have to entail that more and more parts of air came to occupy one and the same place; before the condensation one and the same part of air would then have to be in two or more places at the same time, etc.'' These and similar difficulties can all be removed when one lets infinite discretion furnish the middle determination between finite discretion and continuity. We think of a body enclosed in a bounded space, whose finite spherical parts have a diameter equal to $\mu$ and are at a distance equal to $\sigma$. The greater or smaller the distance between the spherical parts, the denser or rarer, naturally, is the body. We assume $\sigma = q\mu$, thus: $\frac{\sigma}{\mu} = q$; thus: the greater $q$, the less the density. Without the density itself being changed, finite discretion can however approach \opage{367} infinite discretion; the mass with the arbitrary density can, in other words, approach a continuum. If $\frac{\sigma}{\mu} = q$, then also $\frac{d\sigma}{d\mu} = q$; but by putting $d\sigma = q\,d\mu$ instead of $\sigma = q\mu$, continuity has evidently stepped into the place of discretion. In that the finite part $\mu$ has been divided into infinitely many infinitely small parts, $d\mu$, $d\sigma$ has replaced $\sigma$, and $d\sigma = q\,d\mu$ is then an expression for the distance between the infinitely small parts having become infinitely small. Even if there can be no question in actuality of an absolute continuum, there can nevertheless be question in actuality of a matter approaching infinitely near to becoming a continuum.

The doctrine of the approximation of material discretion to continuity is of the very highest importance when it is a matter of determining the relation between \emph{alteration} and \emph{transformation} in the material. Mere alteration rests on finite discretion, transformation on infinite discretion; alteration is in its essence mechanical, transformation dynamic. With respect to alteration, the atom has immediate independence as a point of matter; with respect to transformation, the atomistic point of matter has reflected independence as a point of function. ``That the atom is a point of function means, then, in other words: 1) that its peculiar determinate being is determined by its peculiar acting, 2) that by its activity it both produces and itself undergoes a series of law-determined changes, 3) that under all changes it keeps itself unchanged, not in material immediacy, but in essential form, --- a reflected constant! By apprehending the atom as a constant of reflection, philosophy can appropriate chemistry's exact interpretation with all the punctuality belonging to it \ldots If the atoms, as material parts, are to fill space, their mutual reciprocal action must naturally be conditioned by a special distribution in space; chemistry is thus in this respect justified \opage{368} in laying particular weight on the arrangement of the parts in all formations and transformations. But to want to explain everything by mere rearrangements is superficial. It is true that attempts have been made to make intuitively clear the different ways in which single atoms can group themselves into compound ones; but the attempts belonging here are too poor, arbitrary and insufficient even merely to explain a single one of the transformations that fall, e.g., under isomerism. In metameric phenomena, such as $\mathrm{C}_{4}\,\mathrm{H}_{6}\,\mathrm{O}_{3} + \mathrm{C}_{2}\,\mathrm{H}_{6}\,\mathrm{O}$ as opposed to $\mathrm{C}_{2}\,\mathrm{H}_{2}\,\mathrm{O}_{3} + \mathrm{C}_{4}\,\mathrm{H}_{10}\,\mathrm{O}$, each of the elementary atoms is indeed, despite the rearrangement, like itself, namely as a point of function: $\mathrm{C} = \mathrm{C}$, $\mathrm{H} = \mathrm{H}$, $\mathrm{O} = \mathrm{O}$; but under the conditioned reciprocal action each single $\mathrm{C}$ in $\mathrm{C}_{4}$ nevertheless acquires, by means of $\mathrm{C}_{4}\mathrm{H}_{6}$, a different character, a different influence and significance than that which it can obtain in $\mathrm{C}_{2}$ by means of $\mathrm{C}_{2}\mathrm{H}_{6}$, just as each single $\mathrm{H}$ in $\mathrm{H}_{6}$ also plays a different role in relation to $\mathrm{C}_{2}$ than in relation to $\mathrm{C}_{4}$.''%
\footnote{Forelæsn.\ over ``Phil.\ Propæd.'' 1860--61 p.~155--156.}
In objectivity itself, therefore, the dynamic corresponds to the mechanical as the intensive to the extensive, as the inner to the outer, the ideal to the real, the logical to the antilogical, the subjective to the objective; yet it must constantly be held fast that each of these main sides can in its peculiar way reduce its opposite to a moment and thus present the whole of objectivity from a different side. In the compounds $\mathrm{C}_{4}\,\mathrm{H}_{6}\,\mathrm{O}_{3} + \mathrm{C}_{2}\,\mathrm{H}_{6}\,\mathrm{O}$ and $\mathrm{C}_{2}\,\mathrm{H}_{2}\,\mathrm{O}_{3} + \mathrm{C}_{4}\,\mathrm{H}_{10}\,\mathrm{O}$ it is the different positions of the atoms that express the mechanical, whereas their different significance as moments in the compounds expresses the dynamic. As moments of the cause, the atoms are reduced to real determinations in ground and possibility; thus the dynamic comes to designate partly the ground, partly the possibility, and in both cases ideality, the ontologically subjective. Considered by \opage{369} itself, $\mathrm{C}_{6}\,\mathrm{H}_{12}\,\mathrm{O}_{4}$ is only a chemical possibility. The apprehension of such a possibility is abstractly dynamic; the dynamic apprehension shows that the mechanically existing atoms, in their immediate independence, are without chemical influence; the chemical possibility is, in its immediate determinate being, only a mixed mass, an aggregate. In order to acquire chemical significance, the atoms must give up their independence and become moments. Only in the opposition between the compound $\mathrm{C}_{4}\,\mathrm{H}_{6}\,\mathrm{O}_{3} + \mathrm{C}_{2}\,\mathrm{H}_{6}\,\mathrm{O}$ and the compound $\mathrm{C}_{2}\,\mathrm{H}_{2}\,\mathrm{O}_{3} + \mathrm{C}_{4}\,\mathrm{H}_{10}\,\mathrm{O}$ is the dynamic demand fulfilled. The possibility, the chemically subjective, is here actualized in two compounds, two chemical objects. The chemical objects that satisfy the dynamic satisfy at the same time the mechanical; if the position of the atoms is conditioned by the dynamic peculiarity of the moments, then the peculiar significance of the moments is in turn conditioned by the position of the atoms: the reciprocal determination reflects itself in the principial unity of the act of thought and the act of nature; the expression of the unity is the chemical formation. From the subjective side, the formation is an ontological objectification, an act of thought determined by moments of power; but the objective side of the objectification is the formation itself, an act of nature determined by moments of knowledge.

\begin{center}{\bfseries γ)\quad Aesthetophysical Functions.}\end{center}
\addcontentsline{toc}{paragraph}{γ) Aesthetophysical Functions}

In that the abstract-speculative philosophy of nature volatilizes the mechanical in a fantastic way and holds to the dynamic, it commits what is, in logical respect, a threefold blunder. In the first place it comes to tear what is ideal in objectivity loose from its material foundation; in the second place it deprives itself of every means of investigating the exact reciprocal determination of the ideal and the material; in the third place it entangles itself in the illusion that the ideal itself should be what is objective in the matter, although the ideal is precisely the subjective. These misunderstandings, so baneful for the logic of objectivity, have their essential origin in an unclear \opage{370} apprehension of the relation between the mechanical and the dynamic main side of the functions. As long as one wants merely to hold to the abstractly general, it is easy to see that dynamic activity must presuppose mechanical activity, and conversely; but when the question is of what is peculiar in dynamic activity, when sound, light, heat, magnetism, electricity, etc., are each to be derived, with what is dynamically special to it, from corresponding mechanical presuppositions, then an infinite chasm opens between what is dynamic in the matter and the mechanical, then a dualism shows itself here. The fundamental error in the speculative philosophy of nature, a fundamental error which evidently stems from its ambiguous logic, is now this: that that dualism is not allowed to come forward sharply and clearly. The first condition for a dualism's being sublated must, however, undeniably be that it is allowed to come forward: if it is not allowed to come forward, then it smolders, then it is veiled, and then we have mystification. Before we proceed, then, to an attempt to develop and resolve the dualism cited here, we shall prepare ourselves by a closer consideration of the way in which the doctrine of sound and light, of tones and colors, has been treated by speculation.

In resonance and tone the speculative philosophy of nature recognizes, and rightly so, an utterance of the ideality of matter. ``Beim Ton der Körper,'' says Hegel,%
\footnote{Naturphil. p.~207.}
``fühlen wir, wir betreten eine höhere Sphäre; der Ton berührt unsere innerste Empfindung. Er spricht die innere Seele an, weil er selbst das Innerliche, Subjektive ist. Der Klang für sich ist das Selbst der Individualität, aber nicht das abstract Ideelle, wie das Licht, sondern gleichsam das mechanische Licht, nur als Zeit der Bewegung an der Cohärenz hervortretend.'' But how does he now determine the relation between the material and the ideal, the mechanical and the dynamic? \opage{371} Ambiguously and mystically! ``Zur Individualität,'' it says, ``gehört Materie und Form; der Klang ist diese totale Form, die sich in der Zeit kund giebt, --- die ganze Individualität, welche weiter nichts ist, als daß diese Seele nun mit dem Materiellen in Eins gesetzt ist, und es beherrscht als ein ruhiges Bestehen.'' Such an explanation is unscientific. The categories form, resonance, individuality, soul are used here as univocal words; the figurative, the half-figurative and the non-figurative are mixed in ingenious immediacy: it is a handling of categories that is dubious for the strict logician.

That an ideality such as ``this soul'' is posited in unity with ``the material'' explains nothing, as long as no attempt is made here to show \emph{how} and \emph{why} just ``this soul'' is to be united with this material. By way of elucidation is added here: ``Was sich hier zeigt, dem liegt nicht Materie zu Grunde; denn es hat nicht seine Objektivität in einem Materiellen. Nur der Verstand nimmt zum Behuf der Erklärung ein objektives Seyn an, indem er von einer Schallmaterie, wie von Wärmematerie spricht.'' But such an elucidation elucidates nothing. One may readily grant that there is no sound-matter existing beside or within the sounding matter; but what is gained thereby? Every sounding matter is indeed itself a sound-matter; there does, therefore, lie a matter at the ground of sound: sound must necessarily have ``its objectivity in something material.'' This Hegel has indeed himself in a way presupposed, when he says: ``Der Klang gehört dem Reiche des Mechanismus an, da er es mit der schweren Materie zu thun hat.'' But instead of taking seriously the task of developing the mechanics of resonance by showing how all determinations of difference in resonance and sound are articulated by the dissolution of mechanical functions into dynamic ones, he disposes of the whole thing with general remarks. ``Die Form, als sich dem Schweren entreißend, aber ihm noch angehörend, ist somit noch bedingt: die freie physicalische Aeuße\opage{372}rung des Ideellen, die aber an das Mechanische geknüpft ist, --- die Freiheit in der schweren Materie zugleich \emph{von} dieser Materie.'' Everywhere that things begin to become serious with objectivity and the objective, the speculative philosophy runs up against a dualism that it can indeed mystify, but by no means master. It is, then, a physical and a metaphysical explanation that here stand over against one another. Speculation cannot deny or directly reject what is physical in the matter, the mechanical --- for that the fact is too insistent; but neither can it engage in thorough analyses of the physical and mechanical --- for that its consideration is on the whole too high-flown, and its trust in the logical-metaphysical too blind. So it happens that the two factors, so mutually at odds, weaken each other, that the physical development is hampered by regard for the metaphysical, and the latter in turn by regard for the former. ``Solche Erscheinung, daß ein Insichseyn als physikalisch zur Existenz kommt, kann \emph{uns} nicht in Verwunderung setzen; denn in der Naturphilosophie liegt eben dieß zu Grunde, daß die Gedankenbestimmungen sich als das Wirkende zeigen.''%
\footnote{Hegel. Anfr.\ Skr. p.~208.}
According to this intimation one would therefore expect to see the pervasive unity of the physical and the metaphysical developed clearly and on all sides; but instead we meet the dualism of the a priori and the empirical in the most naive form. The a priori is seen by itself as the necessary; the empirical, on the other hand, in its contingency, becomes the object of intimation, description, etc. ``Das Nähere der Natur des Klanges ist nur kurz anzugeben, indem diese Gedankenbestimmung empirisch durchzugehen ist.''%
\footnote{Ibid.}
In the remarks that then follow on resonance in metal, glass, etc., every regard for the air is passed over, although the air, \opage{373} as the real medium of resonance, surely must also have its share in resonance. It says: ``Luft und Wasser klingen \ldots nicht für sich selbst, wenn sie auch der Mittheilung des Klanges fähig sind''; but this explanation is unclear. If air and water cannot resound in and for themselves, then metal and glass in airless space cannot resound in and for themselves either. One misses an elucidation of how it is possible that air and water should be able to communicate resonance without themselves resounding.

Finally, then, the admission is brought forward here too that the relation between the physical and the metaphysical in resonance is almost impossible to comprehend. ``Die Geburt des Klanges ist schwer zu fassen. Das specifische Insichseyn, von der Schwere geschieden, ist, als hervortretend, der Klang; er ist die Klage des Ideellen in dieser Gewalt des Andern, ebenso aber auch sein Triumph über dieselbe, indem es sich in ihr erhält. Der Klang hat zweierlei Weisen seiner Hervorbringung: α) durch Reibung, β) durch eigentliches Schwingen, Elasticität des Insichseyns.'' Now an utmost effort is made to articulate the reciprocal determination of the physical and the metaphysical, an attempt which, in order to round off the picture of the whole mode of explanation, we shall still add. ``Bei der Reibung ist auch dieses vorhanden, daß, während ihrer Dauer, eine Mannigfaltigkeit in Eins gesetzt wird, indem die verschiedenen außer einander seyenden Theile momentan in Berührung gebracht werden. Die Stelle eines jeden, somit seine Materialität, wird aufgehoben; sie stellt sich aber ebenso wiederher. Diese Elasticität ist es eben, die sich durch den Klang kund giebt. Aber wird der Körper gerieben, so wird dieses Schlagen selbst gehört; und diesem Tone entspricht eher das, was wir Schall nennen. Ist das Erzittern des Körpers durch einen äußeren Körper gesetzt, so kommt das Erzittern beider Körper zu uns; Beides greift ineinander, und läßt keinen Ton rein. Die Bebung ist dann nicht sowohl selbstständig, sondern gegenseitig gezwungen; das nennen wir dann Geräusch.''%
\footnote{Anfr.\ Skr. p.~209.}
\opage{374} With the attempt to loosen the ideal from its material foundation another attempt stands in the closest connection, the speculative attempt to let the material itself transform itself into the ideal; an example of such an attempt in the higher mystical art of transformation we have in the doctrine of light. ``Die Materie, als die unmittelbare, in sich zurückgekehrte, freie selbstständige Bewegung'' --- here mécanique céleste is meant --- ``ist einfache, sich selbst gleiche Gediegenheit.'' This is now explained somewhat more precisely thus: ``Indem die Bewegung in sich zurückgegangen ist, so hat die himmlische Sphäre ihr selbstständiges ideales Leben in sich vollendet und beschlossen; das vollkommene Insichseyn ist eben ihre Gediegenheit \ldots Die Materie, wie sie erkannt worden als diese Unruhe des Wirbels der sich auf sich beziehenden Bewegung und als die Rückkehr zum Anundfürsichseyenden, und dieß Insichseyn, welches da ist gegen das Daseyn, ist das Licht. Es ist die in sich verschlossene Totalität der Materie, nur als reine Kraft, das sich in sich haltende intensive Leben, die in sich gegangene himmlische Sphäre, deren Wirbel eben diese unmittelbare Entgegensetzung der Richtungen der sich auf sich beziehenden Bewegung ist, worin, in dem Heraus- und Hineinströmen, aller Unterschied sich verlöscht; \ldots Das Licht ist diese reine daseyende Kraft der Raumerfüllung, sein Seyn die absolute Geschwindigkeit, die gegenwärtige reine Materialität, das in sich seyende wirkliche Daseyn, oder die Wirklichkeit als eine durchsichtige Möglichkeit \ldots''%
\footnote{Anfr.\ Skr. p.~130--131.}
If we consider the whole utterance from the standpoint of ideality and judge by the total impression, then it cannot be denied that there are beautiful and profound features in the explanation cited. To the mechanical totality, the solar system, there corresponds a dynamic totality whose moments are united in an ideal way as the pure force in the all-manifesting light. That the same existing totality is thus posited at one and the same time both as \opage{375} mechanical and dynamic, both in the form of reality and in that of ideality, is truly profound. But the whole explanation, although it is meant objectively, is altogether subjective. How subjective the explanation is can be seen from the fact that it is not even possible to subject it to any scientific corrective whatsoever. When someone asserts that light is ``die in sich verschlossene Totalität, die in sich gegangene himmlische Sphäre,'' then it will surely be just as impossible to adduce scientific proofs \emph{against} as \emph{for} such an assertion.

If an account is to be given here, as regards phenomena such as those of light and sound, of the relation between the ideal and its material foundation, then it is altogether necessary that weight be laid on the corresponding media, thus on the difference between the air and the ether, and, as regards each of these media in turn, on the punctual presuppositions and conditions to which exact analysis with its functions must consistently lead. In that natural science proceeds from the presupposition that the activity of light, like every other activity in nature, must necessarily have its material foundation, it does not stop at such an assumption in all generality, but investigates the matter in the particular and then assesses the general presupposition solely and exclusively according to its applicability to, and agreement with, the special cases. From this standpoint ``the natural science of light'' can open peculiar insights for us into the articulation of matter; from this standpoint the relation between the material and the ideal becomes clear; from this standpoint the dispute between the emission theory and the undulation theory acquires scientific significance, also for logic.

The theories named are indeed about equally old. ``According to the one, the \emph{emission} or \emph{corpuscular theory} put forward by Newton, light was a substance. It consisted of small bodies which the sources of light cast out with an extraordinarily great velocity, and which, on striking the retina, \opage{376} produced the impression of light; they were not affected by gravity and would therefore fly through celestial space in straight lines, \emph{rays of light}. Now when space was traversed in all directions by these `corpuscles' or grains of light, it might seem that collisions would have to take place every moment, and a change in the direction of motion occur very frequently. One would then every moment get such stray little bodies into the eye and see flashes from all sides; but this difficulty can well be avoided. For if a luminous point sent, e.g., 100 grains of light into the eye every second, the eye would receive a steady impression of light, but since light traverses 40{,}000 miles in a second, there would be 400 miles between two successive grains of light. Another objection was that the impact of these grains would give a quite perceptible mechanical effect, even if they were very small. If a grain of light had the same mass as $\frac{1}{2000000000}$ grain, the impact would make just as great an impression as a ball of 25 grains' weight fired from a rifle; but even in the focus of large converging lenses and concave mirrors one has never been able to discover any trace of mechanical effect. But one is thereby only led to take their masses to be many times smaller than the above-mentioned magnitude.''%
\footnote{C.\ Holten, Lysets Naturlære Kjøbhvn. 1861 p.~163--64.}
The Newtonian conception of what is material in the nature of light forms, as one sees, a striking contrast to the Hegelian; while Hegel aims at transforming luminous matter into pure ideality, a logical essentiality, an immaterial matter, Newton is close to regarding the immaterial as something merely material, insofar as he apprehends light without further ado as a substance. On behalf of ideality one would therefore be tempted to regard Hegel's view as far more scientific than Newton's. Nevertheless the matter stands the other way round. Even if we apply the measure of the idea, we must \opage{377} concede the decisive advantage to the Newtonian mode of consideration. The point is that Newton's hypothesis indicates determinate points of departure for fruitful investigations; it thus not only gives directions for reliable correctives, but even contains, dialectically understood, the germ of its own self-sublation. ``From this hypothesis,'' it says, ``the general phenomena can be explained quite easily. That the illumination is inversely as the square of the distance follows quite simply from the number of grains of light that will strike the same surface at different distances from the point of light. The law of reflection could be derived quite simply from the impact of the grains of light against bright elastic surfaces; but here it must be remarked that, with respect to bodies as small as these grains of light, no surface is even and smooth. Newton therefore assumed that the parts of the bodies that reflect light exert a repulsion on the grains of light, which is perceptible at a finite though exceedingly small distance.''%
\footnote{C.\ Holten. Anfr.\ Skr. p.~164.}
Only insofar as the hypothesis demands ``that light shall travel $\frac{4}{3}$ times faster in water than in air,'' while in actuality it can be shown that precisely the opposite happens, does it come into conflict with a fact.

Before, however, we develop what is thereby indicated, we shall first cast a glance at the undulation or wave theory put forward by Huyghens. ``Huyghens assumes celestial space to be filled with a substance that combines an extraordinarily slight density with an extraordinarily high degree of elasticity. This substance, which he calls the \emph{ether}, penetrates transparent bodies, although in them it may have a different density and elasticity than outside. This ether is now set by luminous bodies into a wave motion, and it is these ether waves which, when they strike the retina, produce impressions of light.''%
\footnote{C.\ Holten. Anfr.\ Skr. 166--167.}
\opage{378} Now if the dispute for and against the two theories were to be conducted in general terms, so that assumption was set against assumption, reasons against reasons, the dispute would never come to an end. A reasoning from reasons is inexhaustible, and an endless dispute is meaningless. It is from this standpoint that speculative philosophy has viewed the matter, and for that reason has not only felt itself raised above any further consideration of these hypotheses, but has even, by reasoning from reasons, seen itself able to refute them both. ``Jene Physik von Lichtpartikeln,'' says Hegel,%
\footnote{Hegel. Naturphil. p.~140.}
``ist um nichts besser, als das Unternehmen desjenigen, der ein Haus ohne Fenstern gebaut hatte, und das Licht nun in Säcken hineintragen wollte.'' By this forceful argument, then, Newton's hypothesis is completely refuted. As regards the propagation of light, the following reasoning must weigh heavily upon optics. ``Von der Vorstellung der \emph{Optik} aber, daß von \emph{jedem Punkte} einer sichtbaren Oberfläche (den jede Person an einem andern Orte sieht) \emph{nach allen Richtungen} Strahlen ausgeschickt, also von jedem Punkte eine \emph{materielle Halbkugel} von unendlicher Dimension gebildet würde, wäre die unmittelbare Folge, daß sich alle diese uendlich vielen Halbkugeln (wie Igel) \emph{durchdrängen}. Statt daß jedoch hierdurch zwischen dem Auge und dem Gegenstande eine verdichtete, verwirrte Masse entstehen und die zu erklärende Sichtbarkeit vermöge dieser Erklärung eher die Unsichtbarkeit hervorbringen sollte, reducirt sich damit diese ganze Vorstellung selbst eben so zur Nichtigkeit, als die Vorstellung eines concreten Körpers, der aus vielen Materien so bestehen soll, daß in den Poren der einen die andern sich befinden, in deren jeder selbst umgekehrt alle anderen stecken und circuliren \ldots''%
\footnote{Anfr.\ Skr. p.~138--39.}
That so speculative a critique strikes the undulation theory just as much as the \apage{379} emission theory is obvious, and it is moreover expressly emphasized, since pure manifesting is here spoken of as ``ein Verhältniß, aus dem, als der in sich verhältnißlosen Reflexion-in-sich, alle die weiteren Formen von \emph{Vermittlungen}, die ein Erklären und Begreiflichmachen genannt zu werden pflegen, Kügelchen, Wellen, Schwingungen u.\ s.\ f., so sehr als Strahlen, d.\ i. feine Stangen und Bündel, zu entfernen sind.''

These objections against the material basis of the action of light are, of course, not one hair better than those brought forward against the conditions of the action of gravity, in particular against ``ein plötzliches Umschlagen'' of centripetal and centrifugal force.%
\footnote{To what degree those loose, untenable, overbearing reasonings must in their time, when speculative philosophy was a power, have offended the experts and contributed to bringing philosophy into discredit, is understandable when one compares what is inflated and confused in such explanations with the clarity, certainty and precision that distinguish the exposition of the special sciences. --- ``By a ray of light one understands'' --- according to the wave theory --- ``any line along which the light propagates itself \ldots The waves spread out as concentric spherical surfaces around the luminous point; at a considerable distance from this, one can regard two neighboring rays as parallel, and the wave surface then becomes a plane which stands perpendicular to the direction of the ray of light. All ether particles in the same wave surface are always in the same state of vibration.'' ``According to this theory the impression of light will be proportional to the living force of the vibrating ether particles, taking account for the present only of the greatest velocity which they attain during their to-and-fro vibrations. The motion which proceeds from the radiating point will, after a certain lapse of time, have reached a spherical surface $a$, in which all the ether particles, whose number we may denote by $n$, had the same velocity of vibration $h$. Later the same motion reaches another spherical surface $a_{1}$, in which the ether particles, whose number be $n_{1}$, have the velocity of vibration $h_{1}$. If now the mass of an ether particle be denoted by $m$, the sum of the living forces in the two spherical surfaces will be $n\,mh^{2}$ and $n_{1}\,mh_{1}^{2}$; but since it is the same original impulse that has produced both effects, and the ether must be regarded as perfectly elastic, one has
  \[ n\,mh^{2} = n_{1}\,mh_{1}^{2} \]
  If the radii of the spherical surfaces are denoted by $r$ and $r_{1}$, one has
  \[ \frac{n_{1}}{n} = \frac{r_{1}^{2}}{r^{2}} \]
  which in connection with the above gives
  \[ \frac{mh^{2}}{mh_{1}^{2}} = \frac{r_{1}^{2}}{r^{2}} \]
  The living force of the ether particle, or in other words its effect, the illumination, is therefore inversely proportional to the square of the distance.'' C.\ Holten. Anfr.\ Skr. p.~167--68.}
\opage{380} Since speculation, however much it accentuates in its ``objective logic'' matter as the bearer of all formal activity, does in actuality volatilize the material basis of the mechanical-dynamic functions, it has of course barred for itself the access to clear insight into the interior of actual matter, and thereby to a closer understanding of what is objective in objectivity. For such an insight, such a closer understanding, is reached only with the help of mathematical hypotheses. However dissimilar the two hypotheses under discussion otherwise are, they nonetheless both agree in this, that they take seriously what we may call a mechanics of light. The significance of mechanical functions with exact presuppositions, conditions and consequences here appears in its true light. The material determinations that furnish the presuppositions by no means depend immediately upon an immediate sense-perception; it is rather consequences, and so determinations of an intellectual nature, that thought must here transform into actual presuppositions, while the sensibly given presuppositions are transformed into actual consequences. Since rationes cognoscendi and essendi thus incessantly alternate with one another and reflect \opage{381} themselves in one another, it cannot fail that the material determinations must correspond point for point to the intellectual ones, and these to the former.%
\footnote{The ground (ratio cognoscendi) for Newton's letting the parts of the bodies which reflected the light exert a repulsion upon the particles of light was evidently the equality of the angle of incidence with the angle of reflection; but this sensibly given presupposition, of which the assumption of repulsion was an intellectual consequence, at once transforms itself into an actual consequence, a consequence for which the assumption of repulsion, that first intellectual consequence, furnishes the actual presupposition, the proper ground of existence (ratio essendi). What \emph{repulsion} is for the bodies that reflect light, \emph{attraction} is for refracting bodies. From attraction as ground of existence (ratio essendi) there then results the intellectual consequence:
  \[ \frac{\sin i}{\sin b} = \frac{h_{1}}{h} \]
  where $i$ signifies the angle of incidence, $b$ the angle of refraction, $h$ the velocity before and $h_{1}$ the velocity \emph{after} refraction. If now the consequence thus obtained is to be comprised in a ground of cognition (ratio cognoscendi), it must be proved that the ratio itself is independent of the angle of incidence. The proof can be carried out in consequence of attraction as ground of existence (ratio essendi). For if the attracting force is called $Y$, one will obtain, according to propositions of rational mechanics (cf.\ C.\ Holten. Anfr.\ Skr. p.~164--65):
  \[ h_{1}^{2} - h^{2} = \int_{0}^{y_{0}} Y\,dy, \]
  which, since $y_{0}$ and $h$ are constant, must make $h_{1}$ constant. But when one must have:
  \[ h_{1}\sin b = h\sin i \]
  whatever the angle of incidence may be, then the consequence breaks forth, the consequence at once intellectual and yet actual, which overthrows the theory.}
\opage{382} In punctuality objectivity is seen again. Were punctuality to fall away, general considerations would at once step into the place of facts, and explanations back and forth, for and against, would bear witness that what is subjective in objectification had separated itself from what is objective. Now, on the contrary, since the consequence in each single case is just as punctual as the presupposition, the conflict between theory and fact in a single determinate case may be of so decisive a nature that it overthrows the theory. Thus the one circumstance that light, in inescapable consequence of the corpuscular theory, should travel $\frac{4}{3}$ times faster in water than in air is sufficient to overturn this whole theory.

Now as regards the undulation theory in particular, it at once points us to a remarkable analogy between three material media: \emph{water}, \emph{air} and \emph{ether}. What is common to these media is wave motion; but what is universal in wave motion is again to be determined in a particular way in each of the particular media, so that water waves are waves of another kind than air waves, and these again of another kind than ether waves.%
\footnote{According to the undulation theory, ``the velocity of the ether waves would then have to be 41{,}500 miles per second, 919{,}000 times as great as that of sound in air, and if we were to transfer this difference solely to density, the density of the ether would have to be 845{,}000 million times less than that of air, and one cubic mile of ether would contain no more mass than $1\tfrac{1}{2}$ pounds.'' C.\ Holten. Lysets Naturl. p.~167.}
With the exact analysis of sound waves and light waves the mechanical first receives its full due; but now there also appears the dualism between the mechanical cause and the dynamic effect in all its strength.

Through the mechanical, what is active steps out, so to speak, mimically into local determinations (wave motions, vibrations); these local determinations, as mechanical \apage{383} causes, have not the remotest likeness to the corresponding dynamic effects. Or what likeness could one well demonstrate between sound and air waves, between light and ether waves? Even if one were to try to explain the subjective difference between sound and light from the objective difference between air and ether, one could nonetheless not possibly suppose that the subjective difference between light and heat should allow itself to be explained from an objective difference between two kinds of ether; and yet the subjective difference between light and heat is just as great as the difference between light and tone.

The dualism of the material and the ideal thus appears also in natural science as a dualism of the objective and the subjective: to what is objective in air waves and ether waves corresponds what is subjective in our hearing and in our sight. So sharp in this matter is natural science's separation of the objective and the subjective functions that ``common sense'' cannot quite agree with itself as to whether one dare admit such a separation or not. It is true, to be sure, that the brightness of light can be apprehended only by one who sees, just as sound and tone can be apprehended only by one who hears; when one closes one's eyes it is dark, whether the light shines or not; when one is deaf there is always silence around one, however noisy those who hear may find the surroundings. But such an objectivity there nonetheless seems to be in sound as well as in light, that the natural understanding can no more reconcile itself to the sun's losing its shine, even if all eyes closed, or to tones themselves vanishing when no one hears them, than it can reconcile itself to there being neither solid nor fluid bodies if there were no one who could observe them and give them the predicates solid and fluid. It is again the ontological problem of objectification that we meet here.

From the objective side the functions are physical, from the subjective side aesthetic; now since the aesthetic, in the narrower as well as in the wider sense, is decisive for any possible \opage{384} apprehension of sound and light, the activities in question are in this connection to be determined as \emph{aesthetic-physical functions}.

Now if the mechanical-material forms the objective extreme in the functions, the aesthetic-ideal the subjective, then the dynamic, insofar as it has its physical side in original unity with the mechanical and, on the other hand, its aesthetic side in unity with the ideal, will furnish the middle determination.

As regards, then, in the first place the unity-in-difference of the mechanical and the dynamic, the logic of objectivity has a corrective to apply which neither the natural science of sound nor that of light will reject. The mechanical function is, as we have seen, a function of the cause. If one now reflects here exclusively on the mechanical cause, one is driven to so antilogical-objective an extreme that it is not possible to descry any inner essential connection between what is material in the cause and what is ideal in the effect. In acoustics it is said: ``A tone is produced by a series of vibrations which were equally long and follow one another so quickly that the ear cannot apprehend them separately. The faster the vibrations follow one after another, the \emph{higher} the tone is said to be; the slower, the \emph{deeper}. The \emph{pitch} of the tone is therefore determined when one knows how many vibrations it performs each second. The \emph{strength} of the tone, on the other hand, depends on the velocity the air particles attain during their vibrations; for the greater this is, the more powerful the impression the motion will make on the organ of hearing.''%
\footnote{H.\ C.\ Ørsted. Naturlærens mechaniske Deel. Third edition by C.\ Holten. Kbhvn 1859. p.~294.}
To this corresponds, then, what is imparted in optics: ``Just as rays of light can differ in the velocity the ether particles attain during their vibrations, whereby the strength of the illumination is determined, so they can also differ in the ether particles' \opage{385} period of vibration, and thereby the color is determined, the red rays being those with the greatest, the violet those with the least period of vibration. The more and less refrangible kinds of light therefore stand in the same relation to one another as higher and deeper tones.''%
\footnote{C.\ Holten. Lysets Naturl. p.~170.}
The highest that one reaches is an ingenious analogy between air waves and ether waves, together with, on both sides, an articulated parallelism between the causes and their effects. But the mechanical function is, as we have seen, also a function of the ground. Complete grounding is a motivation: here the corrective shows itself. By stopping at what is mechanical in the material cause alone, we did indeed obtain a punctualized parallelism of cause and effect, but no grounding. Or what ground-relation could one well demonstrate between the color red and the vibration number 477 billion per second, between the color violet and the vibration number 733 billion per second? If, on the other hand, we take the matter not merely according to the cause but from the ground up, it will appear that the mechanical is itself motivated by the dynamic. If there is a peculiar ground-relation between sound and the ear, between light and the eye, then there must also be a peculiar ground-relation between sound and the medium of air, between light and the medium of ether. While the causal relation shows us the punctual, the exact and mechanical, the ground-relation shows us what is universal in the matter, what is continuous and dynamic in the transitions.

The very circumstance that air and ether are here to be apprehended as media shows how the determination of the universal involuntarily comes to the fore. If what is peculiar about air waves is cause of sound, then what is peculiar in the activity of sound is in turn the ground of the air waves' \opage{386} peculiarity. Just as the ideal, although it can neither be arbitrarily connected with the material nor be represented as a transformed matter, nonetheless can and must both be itself determined by and in turn show itself determining for the material, so too the ideality of sound can and must both be itself determined by and show itself determining for the corresponding air waves. And the same holds, of course, of light and the ether waves. Between red and the vibration number 477 billion, between violet and the vibration number 733 billion, there is here a mechanical causal relation only insofar as there is here a dynamic ground-relation. The vibration number 477 billion is in the mechanical sense cause of the red, insofar as the vibrating particles are ether particles; but in each ether particle the dynamically ideal corresponds to its material as the inner to its outer: the vibrations themselves, the local changes of the ether particles, express the outward figurations of the activity, but not its essence. The essence of the activity is the ideal; the vibrations are only the material conditions without which the ether particles could not develop their ideal nature; the ideal, what is dynamic in the effects of light, is itself the ground of the ether waves' becoming the cause of the effects of light.

If we next consider the functions from the subjective, and so aesthetic, principal side, then the difference between the merely psychological and the ontological objectification comes into its full right. In the organs of sight and hearing the mechanical becomes a moment for the dynamic; but as long as the psychological dualism persists, what is aesthetic in the functions will constantly remain as though in a suspension between contingency and necessity. The aesthetically subjective comes in as something contingent, for the air and ether waves are independent of our human hearing and seeing; and yet the aesthetically subjective is something absolutely necessary for sound and light; for natural science itself impresses upon us the truth that sound, the \opage{387} audible, can no more be without a hearing than light with all its visible can be without a seeing.

The dualism of the subjective and the objective in the determination of such material idealities as sound and light by no means derives from these idealities themselves, but, like dualism in all other respects, solely and exclusively from the psychological mode of consideration. It is only an illusion when one assumes that the objectivity in sound and tones, in light and colors, should, ontologically seen, be more subjective than the objectivity in gravity, cohesion, hardness and the like. All objective relations rest upon subjective reflections; all objects and objective determinations presuppose an ontological objectification. The difference is only this, that the moment of subjectivity which otherwise conceals itself in the abstract-physical existences steps clearly forth and asserts its right in the aesthetic-physical functions.

\begin{center}{\bfseries c)\quad Functions of the Ground-Cause.}\end{center}
\addcontentsline{toc}{subsubsection}{c) Functions of the Ground-Cause}

The essential connection between the logic of objectivity and the logic of functions is unmistakable from what has already been developed. In the general content of thought that constitutes subjectivity's own logical essentiality, the function already points to the lacking objectivity, insofar as it points to a lack of punctual determinations. In the most abstract form the punctual determinations are exact determinations of magnitude; the antilogical reflexes, whereby the universal-logical concepts transform themselves into operational concepts and punctuality is introduced, here indicate the basic features of a manifold determining of the manifold, the first lineaments of objective existence: the functions of the formal ground thus led, in the exact form, to an adumbration of relations of objects and thereby to objectivity as \emph{formal objectivity}.

When the abstractions with which the formal ground is \opage{388} encumbered are integrated, the pure determinations of form transform themselves into determinations of existence, so that the content of the ground acquires existence and the ground becomes real. The reality of the ground is peculiarly determined by the unity of quality with existent quantity, and so by modality. What is objective in modality shines forth from the specific transitions of measures and relations of measure, from the widely ramified systems of real oppositions. The intensive, which comes to existence in its extensive; force, which has a conditioned existence in its manifestation; the inner, which on so many sides breaks out into its outer --- these fundamental factors accordingly brought about such a multiplicity of distinctive relations of dependence and peculiar modes of acting that the objectivity which was articulated by functions of the real ground acquired relative independence and was to be determined as \emph{essential, real objectivity}.

That the moments of existence which make up the content of essential objectivity belong to actuality, and consequently are determinations of power which in a categorical-anticategorical manner give the existent objects objective peculiarity, followed next from the mixed ground and from the reciprocal actions in which the causes of things have their subsistence. To be sure, the mixing of the grounds cannot take place in an external and contingent manner, as, e.g., the mixing of gases; the conditions which are externally determined by the things are at the same time internally determined by the grounds: the ground receives conditions only by positing conditions; the mixed ground is itself a unity, the unity in which the mixing is grounded. Nevertheless the mixing of the grounds did become a categorical expression for the finitization of essence in otherness, in the anticategorical, what is objective in objectivity. The reciprocal determination of logical and antilogical articulated itself in the relations of measure. The antilogical was that the limit of the measure was posited by a qualifying act of power; the logical was that the qualitative determinacies were motivated by a mixing of \opage{389} grounds grounded in the mixed ground. The real moments of the ground stepped forth, through the outburst of power, as causal moments and existing causes: the functions of the ground transformed themselves into functions of the cause. In time and space, in masses and forces, in changes of state of acting and suffering, motion and rest, causal objectivity then received so exact an analysis, so punctual a development, that it expressly assumed the stamp of \emph{actual objectivity}.

Actual objectivity is dualistic; the dualistic came to light in this, that the functions separated into two dissimilar, yet parallel-running systems, namely the mechanical and the dynamic. In the mechanical functions actual objectivity reached the uttermost limit of the objective; but the antilogical elements, which even in the form of the concept seemed to oppose all essential comprehending, insofar as they asserted the factual cause even to such a degree that the ground and the ground-relation disappeared, were transformed again in the dynamic functions into logical real moments: the analysis of the infinite, with its differentials and differential effects, then showed us the way to the leveling of the chasm that otherwise separates the mechanical from the dynamic.

The relation between the mechanical and the dynamic further presented itself as a relation between what is real and what is ideal in matter, and so as a relation between the objective and the subjective side of objectivity. In the aesthetic-physical functions the opposition of the objective and the subjective factor of objectification, and therewith the dualism of objectivity, was clearly expressed; but the aesthetic-physical functions at the same time articulated the reciprocal action of the opposed terms in such a way that the dualism perished in its consequences.

With the dissolution of the dualism the functions of the ground and of the cause transform themselves into \emph{functions of the ground-cause}; it is this transformation that is now to be more closely determined.
\opage{390} Insofar as the relative caused causes are moments for the principle itself, the content-rich self-cause reflecting itself in all things, all particular functions, all diverse acts of objectification, must be apprehended as activities of one and the same overarching subjectivity. The objective, which thus at once both resists and yet is subordinated to subjectivity, then shows itself, in accordance with the sublated dualism, both from a material and from an ideal side. The task, therefore, is to make clear the relation between material objectivity and subjective knowledge. This relation would be without any concept, and so incomprehensible, if power did not form the middle between the two extremes: matter in the object and knowledge in subjectivity. If one abstracts from power, one has only the choice between a one-sided idealism and an equally one-sided materialism. Idealism sacrifices matter to its ideal of knowledge, materialism sacrifices the ideal of knowledge to matter; only power, as the middle of the extremes, can master what is one-sided in the extremes.

In the first place, then, power is able to vindicate the reality of matter; for what is matter, after all, but the immediate outburst of power, the antilogical element of objectivity, the mighty otherness which not even the Idea itself in its divine reflection can transform into semblance? But on the other hand power would not, as was shown in the preceding chapter, be power if it were not able to master material reality and posit the substantial term as moment, as real negation for a perfect knowledge. Only in that matter with all its reality is transformed into moment and real negation for a divine knowledge is materialism overcome and idealistic one-sidedness sublated.

A thoroughgoing proof of the reality of matter we have in the mechanical functions. If the idealists could really bring themselves to study the analyses and open for themselves an insight into the infinite multiplicity of punctual determinations \opage{391} on which objectivity mechanically rests, they would soon give up every attempt to explain matter away and to volatilize the material. A thoroughgoing proof of the ideality of matter we have in the dynamic functions. That what is material in matter is not itself ideal, that an immaterial matter is a contradiction, must be admitted; but if one has fully seen what it means that atoms as function-points have no immediate stuff-like independence, that as constants of reflection they keep themselves unchanged only in their own changeability, then one will necessarily see in the material the basis for what is ideal in objectivity. If the materialists cannot, at the present stage of science, give an account of how the mechanical leads over with inner consequence into the dynamic, then their assertions and assurances about the absolute reality of matter are without any scientific interest; but if an account is to be given of how the mechanical and dynamic functions reflect themselves in one another, it will be impossible to acknowledge the material side of objectivity without at the same time acknowledging the ideal.

That objectivity is material and mechanical only insofar as it is ideal and dynamic means, in other words, that the objective is only insofar as it is objectified. In ontologically objectifying subjectivity power is identical with knowledge, and the material is to that extent identical with the ideal. But identity in turn has difference, and the unity-in-difference is, in the relation of the objective to the subjective, a unity-in-opposition. Upon the unity-in-opposition of knowledge and power there rests, then, the unity-in-opposition of material and ideal objectivity, and upon this in turn the ontological unity-in-opposition of sense-perception and thought.

That sense-perception, seen in the light of objectivity, has ontological validity just as much as thought can be clearly seen at this turning point. Functions of pure thought without supplementary determinations from sense-perception and the \apage{392} sensible are, as we have seen, empty and meaningless; in the purest reflections, indeed even in speculation, there is a sheen and intimation of transformed moments of sense-perception. But functions of the sensible without sense-perception, and of sense-perceptions without thought, are then also meaningless; for functions rest on relations, relations on reflections, reflections on thought.

The question of the reality of sense-perception stands in the closest connection with the question of the reality of matter and of material objects. If matter, seen in the Idea, has no reality, if its reality is only an expression of the impotence of knowledge, then sense-perception, ontologically seen, naturally has no reality either; then its psychological validity is an expression of the impotence of limited subjectivity. But now the converse holds too. If matter has, for the Idea itself, the reality of otherness, if this reality is a decisive expression of the unity-in-opposition of power and knowledge, then the material object must necessarily, by means of power, determine knowledge, and the immediate determination of knowledge correspond to ontological sense-perception. Just as psychological sense-perception is a categorical mark of the impotence of knowledge, so the ontological is an essential expression of the unity of knowledge with ideal power. By means of power knowledge can determine matter, for power is ideal; by means of power matter can determine knowledge; for power is real.

That idea-sensations, ontological sense-perceptions, must have reality is seen from the mighty opposition between real and ideal power. Upon the idea-sensations there rests, then, also immediacy, the real-ideal immediacy which conditions reflection and makes comprehending and knowledge possible. Or how, after all, should knowledge be conceivable without reflecting and comprehending? Let the knowledge that is perfect in the Idea be ever so intuitive, ever so speculative: our investigations of the ideal of knowledge have nonetheless convinced us that an intuitive knowledge without discursive \apage{393} presuppositions, a speculative knowledge without any reflecting, a divine knowledge without thinking, abstracting and comprehending, is an idle fancy. But no reflecting is possible without immediacy, no discursive presuppositions without finitude, no actual comprehending without antilogical reflexes: whence, then, come these conditions decisive for an unconditioned knowledge? Answer: from what is material and sensible in objectivity? But how do these sensible moments of immediacy come into the Idea? Answer: through the ideal-material affections of subjectivity, through idea-sensations, through ontological sense-perceptions! Or still more clearly: the sensible moments of immediacy come from actuality, from the realm of objectivity, into the Idea as knowing Idea, because they are originally in the Idea, namely in the Idea of Power. It is through the idea-sensations, the ontological sense-perceptions, that subjectivity, by means of its objective other, stands in immediate relation to itself. Take away these sense-perceptions, these idea-sensations, these impressions of subjectivity from material determinations of power: what then becomes of the infinite reflecting? It stops! And when reflection with its real negations has stopped, what then becomes of the infinite comprehending? It vanishes! But when the infinite comprehending with its ideal content and real objects has in this way vanished, what then is the Idea, ``the truth that knows itself'', with all its thinking, comprehending and knowing, but a semblance, a mystical essentiality, a logical dream?

By ontological sense-perception matter is determined in principle as the sensible. From the standpoint of finitude it looks as if the material were equally independent whether it was sensed or not; but in the principle sensibility is itself the substantiality of matter: to exist and to be sensed are for matter one and the same. Now if matter has its reality in sensibility, and sensibility in ontological \opage{394} sense-perception, then the question of the Idea's sense organs, of a sensorium of subjectivity, falls away of itself. Far from needing any special sensorium in order to sense the sensible, the sensible is on the contrary itself the Idea's sensorium. At the psychological stage the sense organ is one thing and the object of sense-perception another. Only the object is sensed; the organ itself is not sensed. At the ontological stage the sense organ is itself object and the object organ. Now if there are countless objects in the real system, and if every sensible object is itself an ontological sense organ, then it goes with subjectivity and the sense organs as with the ground-cause and the causes of things: the many different causes of things are moments and points of passage for the ground-cause; the many and different sense organs are moments for the subjectivity that senses in the Idea.

Now insofar as subjectivity as the principle of sense-perception itself, in particular limitation, forms the sensible through which it senses, the corresponding functions are \emph{organic}; insofar as the principle of sense-perception, as reflecting subjectivity, sets its intellectual acting against the sensible presuppositions by which this acting is conditioned, the corresponding functions are ideal-mechanical, functions of the \emph{mechanics of spirit}; but the mechanical conditions of spirit are taken up, as serving determinations, into subjectivity's appropriation of its objective content, transparent in the Idea, while the sensible continues to hold, and the functions are determined according to their \emph{energy}.

\begin{center}{\bfseries α)\quad Organic Functions.}\end{center}
\addcontentsline{toc}{paragraph}{α) Organic Functions}

``To assume a principle of life valid for organic individuals without presupposing a fundamental principle valid for the whole of existence is thoughtless. But how, after all, are we to convince ourselves that the external actual world, with its things, causes of things, conditions and laws, has so ideal an origin, \opage{395} that this whole is held together and borne by so spiritual a power, that this manifold is pervaded by a single, omnipresent principle?''%
\footnote{Lectures on ``Phil.\ Propæd.'' academic year 1860--61 p.~200.}
--- When in speculative cosmologies there is talk far and wide of an all-constituting, all-organizing activity of reason, it might indeed, on a preliminary consideration, seem as though the significance of such an all-organizing activity and ``activity-process'' were so evident in itself that no further logical preparation would be required for insight into the organic functions of reason. But if we look more closely, the concept of an all-organizing activity of reason is, after all, not exactly so easy to make clear.

The difficulty with an all-organizing activity of reason is, at the present stage of development of objectivity, by no means this, that the Idea itself appears as cause, that the eternal reason is able to penetrate, move and form matter; for the material energy of the principle has already been illuminated from so many sides that no doubt is left here as regards this point. In its generality the matter is clear; but with the particular and the singular the difficulties come. For on the one hand it is uncertain whether the all-organizing reason acts as an unconscious principle of nature or as self-determining spirit, as knowing subjectivity; on the other hand it is undetermined what is properly to be thought by an all-organizing activity, as long as the categories \emph{lay out}, \emph{order}, \emph{constitute}, \emph{articulate}, \emph{organize} are used indiscriminately, so that they almost count as synonymous words.

Now as regards, in the first place, the question whether the all-organizing reason is an unconscious natural being or a knowing spirit, Grundideernes Logik has, in the present connection, presuppositions enough for its answer. For it will be evident from all that precedes that the emphasis here cannot fall on the category \emph{reason},
\opage{396} since the eternal reason indeed acts at once both in the sphere of knowledge and in that of power; from which it then follows that its activity must be able to take a reflected as well as an immediate, a subjective and conscious as well as an objective and unconscious form. No, the decisive stress falls on the category of \emph{organization}; and from this it then follows that, if the category \emph{organizing} belongs within the sphere of power, hence within the real system, the reason that organizes everything must necessarily act as an unconscious principle of nature. The same holds of the activity that orders everything, constitutes everything, articulates everything.

But from the fact that the organizing activity belongs to the same comprehensive sphere as the constituting and articulating activity, it by no means yet follows that the functions named should have the same value; the real system itself indeed divides into lower and higher spheres, corresponding to totalities of lower and higher order. Thus a totality of lower order may very well be laid out, ordered, constituted and articulated without therefore being \emph{organized} in the proper sense of the word. The totalizing reflection has already shown us in the preceding chapter how essential the difference is between the mechanical, the mechanical-dynamic and the organic totality; each of the totalities named is in its own way constituted and articulated, but only the organic one is organized. From the standpoint of objectivity, hence with express regard to the corresponding systems of functions, we shall then in this connection consider the difference indicated somewhat more closely.

The abstract-mechanical, the in concreto mechanical-dynamic and the organic totality have, with all their diversity, an essential mutual likeness, indeed a common ideal-material basis. By what is ideal in the basis the totalities belong to the ideal system, by what is material in the ground they belong to the real system. How the disposition is formed, what a constituting essentially means, what moments the articulating \apage{397} activity contains, are here the decisive questions.

The formation of the disposition necessary for each totality takes place through a peculiar co-operation of ideal and material energy. If the energy is placed one-sidedly in the idea and the ideal, then matter becomes the element that is in itself passive. The formation of the disposition, the original constituting of the whole, then comes to rest exclusively on the peculiar self-activity of the idea-filled subjectivity. To an immediate intuition the constituting self-activity of the idea, when it is to be more closely determined, can present itself in different ways. On the one hand it may seem as though subjectivity itself, with abstract freedom, hence with pure arbitrariness, acted upon matter as upon a disordered mass, a pliant, compliant stuff, in which the forming reason can stamp its conscious thoughts and realize a plan of its own wisdom. On the other hand the matter can, when in order to veil the arbitrariness one introduces an ideal middle term, be represented in such a way that it is not, to be sure, immediately subjectivity itself, but a peculiar idea proceeding from the realm of ideas, on whose principial activity the movement of matter and the formation of the disposition essentially rest. But whatever modifications the idealistic way of viewing things undergoes, it will not be possible to conceal that the ideal energy is lost when the material energy is denied: the functions are ideal only by being material.

If the energy is placed one-sidedly in matter and the material, then the idea with its knowledge and self-determining must either vanish as a phantasm, a deuteroscopic vision, or at any rate be transformed into an absolutum in eternal being without becoming, eternal rest without movement, eternal mirroring without acting. But when an actual disposition is to rest exclusively on matter and its forces, it ceases to be a disposition; for a disposition in the stricter sense presupposes a plan, a plan in turn presupposes determining reflection, and determining reflection ideal energy: the functions are material only by being ideal.

\opage{398} Now if every actual disposition begins with a system of real possibilities determined by the principle, possibilities of ideal-material functions, then the constituting, the original function of the ground-cause, must be an activity of ideal and material energy, an activity of the principle, whose plan and purpose is reflected in the idea, whose carrying-out in the real system is conditioned by matter and its forces. From this it then follows that every totality closed in itself, whether it be otherwise abstract-mechanical or mechanical-dynamic or organic, must, as regards knowledge, reflecting and ideal energy, stand in an equally original relation to the idea and its functions of subjective knowledge: here there is to that extent no difference in the origin of the disposition. It is quite certainly an error if anyone assumes that the disposition of a mechanical totality may very well come about through the forces of matter alone, hence blindly, whereas the disposition of an organism, especially of an ensouled organism and then above all of the richly gifted human organism, should require a special conscious preparation, an extraordinary contribution of divine wisdom.

What holds of the disposition of the totalities holds also of the actual execution of the disposition, hence of articulation. For every totality, whatever stage of existence it may otherwise belong to, articulation is determined by a reciprocal action doubled in itself, the reciprocal action of ideal and material energy. It is only a speculative idée fixe that the idea in a system of mechanical functions should not have the power over matter, hence not the capacity to articulate the mechanical whole, that the idea in a system of organic functions has to articulate its organic whole; that, in other words, an organism should be more thoroughly pervaded by the eternal fundamental idea of knowledge and wisdom than, e.g., a solar system. It is two wholly different standpoints that have been allowed to run together into one, in that the relation between the relative totalities and \opage{399} the absolute principle of the totalities has been confounded with the relation of the relative totalities among themselves.

Seen from the absolute standpoint the idea has the same power over the solar system as over the human body; if the formation, articulation and maintenance of the solar system are to be explained by finite, material causes alone, then the same must hold of the human body. But if the organic articulation of the human body cannot be explained without a principle of totality for \emph{organic functions}, then neither can the mechanical articulation of the solar system be explained without a principle of totality for \emph{mechanical functions}:%
\footnote{``Because a merely mechanical reciprocal action between the central force and the tangential tendency suffices to explain the tautening of the cord in which the stone is swung, people imagine that a similar, merely finite mechanical reciprocal action suffices to explain how the earth originally came to move about the sun. From this first fundamental error a second follows; for once one has fixed oneself in the fancy that the forces in the great elementary totalities act as independent material causes, without there being any thought here of an idea-cause or cause of purpose, what concept is one then to form of the ideal causality that reveals itself in organisms? If there is no idea-cause active in the solar system, then neither is there any idea-cause in the formation of the fetus or in metamorphosis or in the manifestations of instinct and drive; conversely: if the manifestations of instinct and drive cannot be explained without a corresponding idea-cause, then neither can the formation of the geological whole, then neither can the formation of the solar system, the great mechanical whole, be explained without a corresponding idea-cause. To take the mechanical cause for an objective concept, the idea-cause for a merely subjective one, is an empty evasion; if it is anthropomorphic to think of idea-causes, then it is also anthropomorphic to think of mechanical causes, for it is indeed the human consciousness itself that must think the latter just as well as the former.'' Forelæsn.\ over ``Phil.\ Propæd.'' academic year 1861--62 p.~91.}
\opage{400} the mechanical and the organic functions are, each in its sphere and in its way, functions of the ground-cause.

It is otherwise when the question concerns the mutual relation of the relative totalities. From this standpoint a difference shows itself, and an essential difference at that, between a mechanical and an organic whole.%
\footnote{``By making clear to oneself the difference between natural totalities such as the solar system, the terrestrial globe and the organisms on the terrestrial globe, one will convince oneself that there must be a difference in principle between nature's organic and non-organic purposes \ldots The mechanical relation demands that the means step forth with such independence that the ideal power of the purpose over them allows itself to appear only in the system of laws according to which the motions take place in space; the organic relation, on the other hand, permits no finite separation of the means itself from the purpose \ldots Between the organic and the mechanical system of purposes the geological forms a middle of its own. The qualitative diversity of the constituents: the opposition between the solid, the liquid and the gaseous, between the strata of the earth among themselves, etc., has for the teleological development of the terrestrial whole an entirely different significance from that which the quality of the rotating masses has for the solar system; for the only quality on which the mechanical purpose lays decisive weight is properly specific gravity. Nonetheless the massive constituents in the body of the earth itself still relate to one another so externally that there can be talk here of organization and organism only in an improper sense; as elementary, i.e.\ inorganic organism, the earth is the maternal womb from which the organic organisms issue, but the mother of the living is not herself living.'' Forelæsn.\ over ``Phil.\ Propæd.'' 1861--62 pp.~91--93.}
Since the stress falls on the opposition of the heterogeneous totalities and thereby at the same time on the peculiar relation of the relative ideas to the fundamental idea, namely as these relations present themselves in the real system, hence in the sphere of objectivity, the task will be: to develop what is peculiar to the organic functions.
\opage{401} Each of the peculiar systems of functions must correspond to its peculiar ground-cause, and is then by the same token determined by a peculiar self-activity, a form of selfhood of its own: to the system of the mechanical functions corresponds a mechanical selfhood, to the system of the mechanical-dynamic functions an elementary selfhood, to the system of the organic functions an organic selfhood; the task will thus be: to demonstrate the peculiarity of organic selfhood. Finally, since in the real sphere the one system forms the basis for the other (the mechanical for the mechanical-dynamic, and both for the organic), it will be the task of logic in this connection: to set forth the ontological relation of the organic system to its particular bases.

Now, as regards first the difference between the abstract-mechanical and the mechanical-dynamic or \emph{elementary} function, the former is characterized by letting the cause coincide immediately with the formal ground, whereas the latter, by contrast, has as its distinguishing mark that it unites the cause with the real and mixed ground. The mechanical functions are abstract functions of the cause; the abstract causal relation goes together into one with the formal relation of ground. The mechanical functions are homogeneous in their existential essentiality; the universal and homogeneous lies at the ground of the punctual diversity; that is why rational mechanics is an exact science; that is why the mechanical lends itself so conveniently to mathematical analyses. The elementary functions must no doubt, since they rest on reciprocal determinations of the mechanical and the dynamic, necessarily retain the mechanical-punctual presuppositions; but since the mechanical form here takes up a heterogeneous dynamic content (thus the functions of sound and light, the functions of heat, the magnetic, electric functions, etc., are indeed dynamically different), the elementary functions are everywhere recognizable by the abstractions and concretions corresponding to the mixture. Each elementary \opage{402}function is in its concrete determinateness as it were a complex of functions; how many heterogeneous functions must not co-operate in the formation, e.g., of sea water, of greywacke, of mould, etc. In contrast to this the peculiarity of the organic function is now conspicuous. Between the elementary and the organic there is here the essential likeness that both kinds of function are concrete, whereas the mechanical is abstract; but to this there corresponds in turn an equally essential unlikeness. For the concretion in the elementary function belongs to otherness, whereas the concretion in the organic function belongs to selfhood. In transforming the heterogeneous constituents of the concretion into moments, the self-cause regains in the subjective the stamp of simplicity that the mechanical function has in the abstractly objective; the organic function has its peculiarity in the way in which it transforms and assimilates to itself the mechanical and elementary functions.

The mutual diversity of the functions must naturally reflect itself in the diversity of the corresponding totalities, and conversely. In this two-sided reflecting the peculiarity of the relative principles comes to view. As a principle, every real principle is an energetic selfhood; the energy at once both material and ideal.

The mechanical principle, as the material mirroring of power, is nevertheless still only an abstract and formal selfhood. This is a plain and necessary consequence of the essence of power. Power makes the members of the relation into active moments by making the active moments into existing members. Insofar as the existing members of the totality are essentially moments, the principle with its ideal supremacy is an all-determining self-cause, a powerful selfhood. But since the moments of the cause exist as thing-causes, as finitely separated members, actuality presents a system of functions whose punctual presuppositions belong to finitude to such a degree that in the phenomenon it looks as though the mechanical totality were a whole pieced together.
\opage{403} The formal unity of essence and phenomenon, namely the immediate unity of power, is an equally formal and yet powerful dualism; the real intermediate determinations by which the distance between the extremes is filled in first allow themselves to appear when the mechanical members of the relation are totalized anew, hence in the elementary totality. Insofar as articulation passes, though in a finite way, through the parts of the parts and thereby qualifies the connections, the articulating self-cause posits the moments of concretion as determinately existing elements and acts upon the singular mediately; ideality can indeed \emph{master}, but not \emph{assimilate}, the heterogeneous constituents; the energy loses itself in finite, if many-sided and multiform, jointings: the articulating principle is an elementary selfhood.

In the organic, on the other hand, the ideal energy of the principle is so taken up into the material that the articulating selfhood pervades its articulated stuff. In the organism selfhood is bodily, subjectivity objective; organization is the embodiment of selfhood, its immediate self-objectification.%
\footnote{``In the astronomical whole, the system of worlds that fills celestial space, the weight of matter prevails; here is the realm of masses. But the realm of masses is also the realm of light: sensible light is the symbol of spirit, a reflected gleam of the ideality active in the ground of reason. Already in the body of the earth the massive is forced under articulating forms; with the organism it disappears entirely; the cell, insignificant in mass, is here the significant point of departure. In the formation of the first cell the whole problem of life is compressed. What is peculiar about the problem of life is that two causes so dissimilar, so mutually opposed to one another, as self-cause and material cause act together in an immediately forceful unity; what is peculiar is that the principle pervades the material to such a degree that it is itself materialized, that selfhood is embodied: organism is embodied selfhood.'' Forelæsn.\ over ``Phil.\ Propæd.'' 1860--61 pp.~208--209.}
If one abstracts from the self-embodying selfhood, \opage{404}it is impossible to form any true concept either of organism or of organic functions. In the mechanical and the elementary totality the separate parts, members and elements come forward in so finite and external a way that the unity of connection conceals the self-unity, and mediacy prevails: all mechanical and elementary functions take on, owing to their partition and fragmentation, the stamp of functions of otherness. In decisive opposition to this, the self-unity in the organic totality is so pervasive that in all nutritive activities and typical formations it determines the unity of connection in an overarching way: the organic functions are, in virtue of their material energy, real functions of selfhood. In the system of real functions of otherness the principle, despite its omnipresence, is as it were something present nowhere; in the system of the real functions of selfhood the principle, although nowhere immediately present, is nonetheless recognizable as something essentially omnipresent.

Insofar as the organic functions reflect themselves in material objectivity, they branch out in an external way. In the comprehensive total system of the organism the one function, indeed the one subordinate system, may fall outside the other; but since, with all their mutual diversity, they nonetheless serve the overarching unity of selfhood, the single function reflects for its part the others. The more the organic totality doubles itself, so that the functional relations in the objective have separate functional reflections in the subjective, the higher the organism stands, the more perfect are the organic functions.

The functional doubling can, since it takes place on finite, material conditions, be carried through only on a basis of relative dualism. From the objective side the dualism is seen in this, that in the ensouled organism two spheres of organ systems form themselves, such that the one merely serves the vegetative, the other also the animal life. From the \opage{405} subjective side the dualism is seen in this, that one sphere of functions falls outside, another inside, the life of consciousness. That the organized unity in which such a dualism is sublated must nevertheless have its system of organs and functions determined by the overarching selfhood is comprehensible%
\footnote{``That each cell at every point is a part for itself only by being a moment for the whole, that its peculiar disposition is preformed in and with the disposition of the whole, its single life in and with the development of life in the whole, is then seen clearly enough in contemplating the series of planned changes, answering to the demands of the paradigm, that the germ undergoes \ldots `it separates into two layers, of which the upper is more transparent and is called the \emph{animal} or \emph{bestial} leaf, the lower is more milk-white and is called the \emph{vegetative} or \emph{mucous leaf} \ldots Soon, however, a third layer forms between these two. It is called the \emph{vascular leaf}, for in it the heart, the blood vessels and the blood itself develop, thus as it were a system common to both its classes of organs (the \emph{animal}: the nervous system, the organs of sense and motion; the \emph{vegetative}: the intestinal tube, the lungs the glandular ducts)''. Cf.\ Forelæsn.\ over ``Phil.\ Propæd.'' 1860--61 pp.~345--46.}.

From what has here been elucidated concerning the peculiarity of the organic function, the relation between philosophy and empirical physiology can then also be more closely determined. When the physiologists themselves admit that they do not know what life is, they must, from their standpoint, no doubt be right; only it should be added that if they do not know what life is, then neither do they know what an organic function is. However important and necessary the exact investigations of the special science may be, and however great the advances physiology must be able to make by its empirical procedure, this nevertheless stands firm: without insight into the essence of life the physiological facts can never be rightly interpreted; but without a logic of the fundamental ideas true scientific insight into the essence of life will be unattainable. Life is idea; in order to cognize what life is, one must in \opage{406} logical universality cognize what an idea is. Hegel has in this respect done right in assigning the doctrine of life as ``immediate idea'' a place in logic; as regards his logical treatment of the idea, on the other hand, it is especially the unclarity in the objective that we shall touch on once more.

``Die Bestimmungen des Gegensatzes'', says Hegel,%
\footnote{Die subjektive Logik p.~249.}
``sind die allgemeinen \emph{Bestimmungen} des \emph{Begriffs}, denn es ist der Begriff, dem die Entzweiung zukommt; aber die \emph{Erfüllung} derselben ist \emph{die Idee}. Das eine ist \emph{die Einheit} des Begriffs und der Realität, welche \emph{die Idee} ist, als die \emph{unmittelbare}, die sich früher als \emph{die Objektivität} gezeigt hat. Allein sie ist hier in anderer Bestimmung. Dort war sie die Einheit des Begriffs und der Realität, insofern der Begriff in sie übergegangen und nur in sie verloren ist; er stand ihr nicht gegenüber, oder weil er ihr nur \emph{Inneres} ist, ist er nur eine ihr \emph{äußerliche} Reflexion. Jene Objektivität ist daher das Unmittelbare selbst auf unmittelbare Weise. Hier hingegen ist sie nur das aus dem Begriffe Hervorgegangene, so daß ihr Wesen das Gesetztseyn, daß sie als \emph{Negatives} ist.''
In an abstract way all this can very well be said; but the root harm is once again the old one: the objective is treated in general, hence not objectively but subjectively; the specifications --- sensibility, irritability, reproduction ---, and what is taught about the ``life-process'' and the ``genus-process'' dissolve to such a degree into the universal and the merely synoptic that the singular, the exact and the punctual are lost sight of entirely.

That logic everywhere keeps to the universal development of concepts without mixing in stuff-like, irrelevant details is in order; but it is one thing to lay hands on particulars, another to demonstrate the significance of the punctual: by abstracting from the empirical moments of the phenomena of life \opage{407} the Hegelian logic has loosed the bonds that were to bind physiology to philosophy and make it possible for the doctrine of the idea to profit from the exact researches of the special science.

That the empirical and the a priori do not find themselves in a sluggish co-being but stand in living reciprocal action shines forth especially from the organic functions. Of all the functions belonging to the real system the organic is, as has been said, the most concrete of all. To its concretion belongs, besides the mechanical, also the dynamic in countless modifications and then in particular the chemical. As a function of selfhood the organic function is in its essence a priori; as the most concrete function in reality it is empirical. The a priori is that the heterogeneous constituents are real moments of self-activity; the empirical is that the real moments of self-activity step forth as heterogeneous constituents. The a priori is that the function of selfhood masters the various functions of the other; the empirical is that the functions of otherness, in their mutual diversity, have a particular subsistence outside the organism. If we let the a priori and the empirical in the ensouled organism separate as extremes in all immediacy, we run up at once against something absolutely incomprehensible. For is it not incomprehensible that bare willing can act upon a nerve, or that a nervous impression can awaken a feeling: what a difference between such empirical objects as ganglia and nerves and such a priori determinations of subjectivity as feelings and manifestations of will! If on the other hand we consider the organic function in its ideal-material energy, then the intermediate determinations show themselves, and when the intermediate determinations show themselves, the absolutely incomprehensible must vanish.

The consideration of the organic function cannot be immediate. The mechanical function can be presented by itself immediately, the chemical in a certain way likewise, for they are indeed both functions of otherness; but the organic function, as a real function of selfhood, is empirical when one considers it \opage{408} a priori, a priori when one considers it empirically. For the infinite reciprocal determination of the a priori and the empirical, punctual in its inner infinity, there is no single scientific method; the scientific method must then itself divide into an empirical and an a priori method. For the empirical method, then, the mechanical, mechanical-dynamic, chemical functions, etc., which make up the lateral determinations and elementary requisites of the organic function, are to be treated by themselves: here there are tasks enough for empirical physiology. But in the assurance that the functions of otherness, separated off in a finite way, have, together with the empirical phenomenal determinations, their true significance in this, that point by point they furnish presuppositions and conditions for the activity by which the organic function articulates the material energy of selfhood, philosophy must know how to lay out its a priori method in such a way that the development of concepts can enrich itself by the empirical researches.

Anatomy, organic chemistry, morphology, etc., mark out various spheres of empirical investigations, all of which are indispensable for physiology. The functions belonging under these are then also treated in a plainly physical way as functions of otherness, exact expressions of the fulfillment of the laws of otherness, the physical laws. From this standpoint the various bases of the organism and of organic life must come into consideration in particular. Thus the mechanical system, when we disregard the first origin, indeed furnishes the most universal and most comprehensive basis. But matters do not rest with the abstract basis; the mechanical functions enter into the organism itself and show themselves co-determining for the organic activity. The nearer and richer basis is supplied by the terrestrial totality, the inorganic organism; but here it is seen properly in all particulars what it means that the organic functions appropriate the elementary contributions as moments. Now in order to investigate the organic functions, \opage{409} empirical physiology must, so to speak, decompose them; it must thus separate the function of selfhood into lower, one-sided, elementary functions of otherness; by such a separation it comes precisely to treat the moments as relatively independent factors. But to know the moments as relatively independent factors without cognizing how it can come about that such factors become moments is insufficient in science. In relation to the organism and the organic functions the researches of the special science are all afflicted with a one-sidedness which, if it is not corrected point by point, must necessarily lead to misunderstanding and misrecognition of the essence of life.

``For \emph{chemical analysis} the organic compound is a product of affinity just as much as the inorganic; the product is selfless, therefore organic chemistry cannot possibly give us any direct elucidation of what life properly is''.\footnote{Forelæsn.\ over ``Phil.\ Propæd.'' 1860--61 p.~209.} What is objective in organic objectivity is one-sidedly brought out by these analyses; the manifold, the heterogeneous, in short, all the marks of the objective are present here. But the organically objective has in itself an indwelling subjectivity; if empirical inquiry has transformed the moments into independent factors, then philosophy must in turn transform the factors into moments. Instead of abstracting from the analyses of empirical inquiry, philosophy ought on the contrary to penetrate into their essence and reshape them into proofs of the pliancy with which matter, without giving up its materiality, the relative independence of its atoms, is able to be moved by the ground-cause and to satisfy the ideal demands of selfhood. ``Under the dominion of self-unity the two-sided nature of the atoms comes to view: they are points of life and yet points of stuff; they are called to life, their elementary function is raised in potency to a function of life. But it is only for a few moments that the single point of matter \opage{410} can thus endure the ideal tension; its stuff-nature is contrary to ideality, so it grows faint, so it is pushed out of the sphere of the points of life and returns again to the elementary.''\footnote{Anfr.\ Skr. p.~219.}

What has here been said of the chemical holds in a corresponding way of the morphological. Morphology lays decisive weight on the difference between an organic and a merely chemical product, and then at the same time logical weight on the difference between an organic objectivity and a merely chemical object. ``For chemical analysis the producing and the produced, the process and the product, must fall apart: insofar as the process is, the product is not yet; insofar as the product is, the process is over; in organic activity, on the other hand, the process is, essentially seen, reflected in its product, the produced simultaneous with its producing: the vessels with which the organism operates are formed, shaped and maintained by the operations themselves.''\footnote{Anfr.\ Skr. p.~220.} The organic productions are then, from the standpoint of morphology, to be apprehended as phenomena of selfhood, images of the essence of life, objective utterances of an indwelling subjectivity. But the subjectivity in the objective producing morphology cannot see through; it sees only the producing in the produced; it has indeed an inkling of the essence, but must nevertheless stop at the phenomenon. Now if the essence, logically seen, is reflected in its phenomenon, and if morphology understands precisely how to bring out the peculiar sides of the phenomena and to demonstrate the phenomenal transitions, then philosophy for its part must in turn know how to see the essence in the transparent phenomena and to let the inner articulation reflect itself in the outer. So many morphological phenomena, so many manifestations of selfhood, so many mirrorings of the \opage{411} idea of life, so many reflexes of organizing subjectivity in the organically objective: such an apprehension cannot remain standing at the universal; it has an original drive to go into the singular.

In the doctrine of the organism, the principial totalizing of the organic functions, the dualism of objectivity thus finds an at least provisional resolution. ``The organic opposition of inner and outer is taken up into organic unity\ldots\ For since all the parts of the organism are moments, the whole totality is thereby posited in the form of inwardness, of unity and of selfhood; the whole system of functions shows itself from the ideal side, here in this inner there is a whole, complete, ideal organism. But since the activity of form has its strength precisely in this, that the mastered real moments at the same time appear as existing members, since it is just the material immediacy of these members that furnishes the resistance on which the forming principle uninterruptedly tries its energy, the otherness through which selfhood incessantly confirms its ideality, the whole system of functions also has a material character, and here in this outer there is a whole complete, material organism; thus: \emph{two organisms in one organism}.''\footnote{Anfr.\ Skr. pp.~231--32.}

Since, however, the resolution of the dualism rests on this, that organic objectivity is itself subjective, because organic subjectivity is in an immediate way objective, the dualism thus resolved must repeat itself anew in a still higher form. The objectivity immediately pervaded by subjectivity, organized bodiliness, will then as the one extreme furnish objective existence, i.e.\ the actual body with its organs, while subjectivity, developed by objective means and set free on organic conditions, will articulate for itself an essential content of its own, an ideal objectivity, \opage{412} in which a system of thoughts and ideas can be reflected. This doubling leads us to what follows.

\begin{center}{\bfseries β)\quad Functions of the Mechanics of Spirit.}\end{center}
\addcontentsline{toc}{paragraph}{β) Functions of the Mechanics of Spirit}

Objectivity rests on objectifying, objectifying on subjective acting; it is then a matter of course that the dualism of objectivity cannot reach its culminating point without the dualism in the objective necessarily going together into one with a corresponding dualism in the subjective. At the psychological standpoint the widest chasm forms between the subjective and the objective, and that in the following way. At the subjective pole the functions are indwelling reflections, ideal effects of light without finite articulation; at the objective pole the functions are bound in finite relations to presuppositions that lose themselves in material immediacy; at the subjective pole all is clarity, ideality, selfhood: the functions vanish in pure light; at the objective pole all is obscurity, reality, otherness: the functions lose themselves in pure darkness. The intermediate sphere lying between the two poles now exhibits doublings of function of various kinds, doublings that are determined in every particular by the constitution of the extremes to whose uniting middle they belong.

On the one side we have, as regards the reciprocal action between world-consciousness and self-consciousness in the human being, the outwardly actual with the functions by which sensible, material objectivity is articulated; on the other side the inwardly actual with the functions by which the sphere of representations with its ideal objectivity is formed and clarified. Now insofar as the functional unity-in-difference of the outer and the inner, the material and the ideal, otherness and selfhood, is conditioned by the separate organs of the organism, the self-conscious functions of objectification, as regards both the material and the ideal objectivity, are to be apprehended and determined as \emph{functions of the mechanics of spirit}.
\opage{413} For a closer elucidation of these functions it is of decisive importance that both extremes, the subjective as well as the objective, each from its own side present starting points that fall within immediate observation: what goes on in consciousness can just as well be observed immediately as what happens in the external world. What goes on during the acts of objectification in the organism itself (in the brain and the nerves), on the other hand, cannot become an object of immediate observation; but since the organism with its organs constitutes precisely the living middle that connects the external bodily with the internal psychical, the material with the spiritual, it follows from this that the organs serving objectification must have the bodily in a psychical way, the psychical in a bodily way, the material in a spiritual way, the spiritual in a material way; must, in other words, express the active unity-in-difference of bodily and psychical, material and spiritual.

If observation now chooses its starting point in the external bodily, then the inquiry begins with what Fechner has called \emph{outer psychophysics}: the elucidation of the psychophysical functions then forms a natural introduction to what Herbart treats as the \emph{mechanics of spirit} in the narrower sense. What yield can a logic of fundamental ideas now expect from psychophysics? To have the relation between the activity in the organism and the activity in the inorganic worked out in an exact manner is of the highest importance, if the relation between spirit and matter is to be determinable at all.

``That it must necessarily be outer psychophysics, not inner, that forms the foundation,\ldots\ is evident. There is indeed a system of internal experiences just as there is of external ones; we can indeed observe our own states of soul, our sensations, feelings and representations, just as immediately as we can observe things and the causes of things in the external world; but --- and this is here a decisive point \opage{414} --- whereas we can immediately and yet exactly determine the difference between stimulation and stimulation, we cannot in the same way immediately and yet exactly state the difference between sensation and sensation, feeling and feeling. Certainly sensation must be measured by sensation, just as stimulation by stimulation; but --- and this is here a decisive point --- whereas, as regards stimulation, one can state the standard of measure directly, the unit of measure of the sensations can, on the other hand, only be stated indirectly: one thing is the means by which the standard of measure is found, another the standard of measure itself.''\footnote{Forelæsn.\ over ``Phil.\ Propæd.'' 1860--61 p.~446.}

When it is said in ``the speculative logic'' that the mechanical is a sublated moment in the dynamic, the dynamic sublated in the teleological, the teleological actualized in the organic, which in the concept presents ``the immediate idea'': then one easily hits upon the thought that every trace of mechanical activity must in the organism be, if not entirely vanished, then at least of quite subordinate significance for the organic manifestations of life. ``The whole category of the object,'' it says, ``stands in closest connection with that of causality. Mechanical activity presents the infinite series in which each member is equally immediately cause and effect. Dynamic activity presents the reflective relation of causality, namely the sublation of the cause in the effect. Finally, teleological activity presents reciprocal action, or the unity in which the cause is effect, and thereby its own cause.''\footnote{Heiberg. Specul.\ Logik. (Pros.\ Skr. 1st vol. p.~338.)} Since, now, life is idea, and everything in the idea proceeds by ``teleological syllogisms,'' one is further tempted to assume that the mechanical functions must be of little or no significance for the life of consciousness; but that this assumption is false, psychophysics shows.

However much we must count among the problematic in the first \opage{415} investigations of a beginning science, psychophysics can nevertheless, already at its present stage, place beyond all doubt that the mechanical, even as a moment in the organism, continues to correspond with the mechanical outside the organism; this is just the remarkable thing, that the mechanical function, after it has here been taken up into organic activity, still preserves its exact character: to such a degree is objectivity in earnest in the organic.

``Weber's law'' may indeed state: that the increase of sensation is proportional to the relative increase of stimulus; that the relation between the stimulus $\beta$ and the sensation $\gamma$, where $\gamma = f(\beta)$, is precisely determined by the differential equation
\[ d\gamma = \frac{k\,d\beta}{\beta} \]
as its ``\emph{fundamental formula}'' and, after integration, determination and choice of the appropriate constants, by the equation:
\[ \gamma = k\,(\log\beta - \log b) \]
as its ``\emph{formula of measurement},'' and that finally, since the equation can of course also be written:
\[ \gamma = k\,\log\frac{\beta}{b}, \]
the ratio $\frac{\beta}{b}$ states precisely ``\emph{the fundamental value of the stimulus}'', since b denotes precisely the limiting value of the stimulus $\beta$ at which the feeling $\gamma$ begins and vanishes;\footnote{Cf.\ Forelæsn.\ over ``Phil.\ Propæd.'' 1860--61 pp.~444--453.} --- all this, we say, may indeed contain much that is problematic when the question is one of application to the particular cases; but so much is evident, that the organic functions which condition the functions of consciousness are everywhere articulated in an exact manner by the mechanical. For the mechanical articulation is naturally not confined to feeling; it holds for every sensation, for sight, for hearing, in all sensory \apage{416} changes and aesthetic transitions. What the natural science of sound and of light makes clear from the physical-objective side, psychophysics thus makes clear from the organic-objective side: the unity-in-difference of the organic and the inorganic activities becomes in this way an object of the ever further-reaching investigations of philosophy as well as of the special sciences.

In opposition to psychophysics, which takes its starting point in the external bodily, mathematical psychology takes its starting point in the psychical internal; this difference is one of principle. Insofar as measure and standard of measure are originally sought in something external, the exact is reached by means of the empirical; insofar as measure and standard of measure are originally sought in the internal, the exact form must be determined by the abstractly rational. Certainly mathematical psychology is referred to experiences, namely internal experiences; but what is experienced in the internal cannot be held fast in immediate delimitation in the way that what is experienced in the external can. If the objects of inner experience, representations e.g., are nevertheless to be apprehended in exact form, the presuppositions of the analysis must rest on pure determinations of concepts. How this is to be understood more precisely can be seen from Herbart's method.

With his ``Mechanik des Geistes'' Herbart stands in modern philosophy as Hegel's most significant, most pronounced antipode. Panlogism, which could not appreciate the essential significance of the mechanics of light, because in all mechanics there are and must be elements of the antilogical, could naturally still less appreciate the significance of a mechanics which must seem to it to stand in open conflict with the essence of spirit. Laid out on discreteness and punctuality, the philosophy of the ``reals'' has gone to the opposite extreme; we have already, in another connection, stated the outlines of the Herbartian theory: ``The soul is only one, but the representations \opage{417} are many. The many and different representations cannot possibly be reconciled in one consciousness unless the difference is equalized. In the same degree as the representations are opposed to one another, they must, in the equalizing of the opposition, obscure and repress one another. Each single representation maintains itself with its degree of clarity in consciousness until it is displaced by other representations. The displaced representation does not flee away, but lingers at the threshold of consciousness (die Schwelle des Bewußtseyns), waiting for the opportune moment when the resistance is lifted, so that it can return again. The representations, more or less opposed to one another, are transformed by mutual contact into forces; one can compare them with steel springs or elastic bodies: they act and are acted upon, they are obscured and clarified, inhibited and furthered, in relations of tension which correspond precisely both to their relative strength and to the different degrees of the oppositions.''\footnote{Forelæsninger over ``Phil.\ Propæd.'' 1860--61. pp.~454--55.}

The metaphysical foundation of the theory is the doctrine of the ``reals''; the sphere of consciousness, with its manifoldly alternating representations, is formed, namely, during the struggle of the soul-real with the ``reals'' in its environment: ``The representations are testimonies of the self-preservation of the soul-real.'' That the metaphysical foundation is at the same time mechanical is seen at once from its atomistic character. But atomism here has the awkwardness about it that, instead of being material-ideal, or, determined still more closely, organic-psychical, it is on the contrary abstract-intellectual: the soul-real is, like every other real, something that is, in all immediacy, logical-antilogical, hence an enormous contradiction! This awkwardness in the foundation already makes itself known in the first propositions of the system.

To determine the relation in which the opposed representations obscure or inhibit one another mutually, the so-called \apage{418} ``relation of inhibition'', Herbart sets up the following three propositions: ``1) Every representation acts in proportion to its strength and its degree of opposition; 2) the representation acts in the proportion in which it suffers; 3) the representation suffers in inverse proportion to its strength.'' That these propositions, regarded from the formal side, are suited to be presuppositions for a mechanical analysis shines forth immediately from the manner in which they are brought into application. Let a and b be two representations, the degree of opposition m; what a loses is then proportional to mb, inversely proportional to b and to a; this then gives mb $\frac{1}{b}\,\frac{1}{a}$, hence the ratio $\frac{m}{a}$.
% Errata: ``Förholdet'' (``the ratio'') printed with an umlaut; standard Danish ``Forholdet''.
In the same way, what b loses is seen to be $\frac{m}{b}$. Hence these ratios can also be expressed by $\frac{1}{a}$ and $\frac{1}{b}$, since
\[ \frac{m}{a} : \frac{m}{b} = \frac{1}{a} : \frac{1}{b}. \]

If we now only had certainty that these very first principles were reliable, then much would already be gained; but that they neither are nor can be so can be proved at once. The representations on whose peculiar character the mechanics of spirit is here to be built are themselves functions of the mechanics of spirit. To want to make mechanics dependent on representations which themselves depend on mechanics is to go in a circle. It is, then, the basic conditions of the mechanics of spirit that must first of all be cleared up. The necessary presupposition is that to every activity of spirit there must correspond a material activity of organs. According to this presupposition every representation is then a product of two mutually heterogeneous factors, the spiritual factor of selfhood and the bodily factor of otherness; it is the unity-in-difference of these factors that determines the organopsychical function, in that it determines the representation. Through this two-sided determining, every representation then \opage{419} also has a double motive, a spiritual and a bodily one; and the question becomes how the two motives are related to each other mutually. The degree of opposition between two representations must thus here be determined from the spiritual as well as from the bodily side, from the conscious psychical as well as from the immediately organic side. The same holds of the representation's strength, its acting, suffering, etc. Since, now, the spiritual motive, in the punctual as well as in the general, is different from the bodily, this difference must of course each time have an influence on the determination of the functional unity: the basic view corresponding to this would then have to yield a still more intricate calculation than the Herbartian.

The mechanics of spirit is not to be sought in the spiritual activity of spirit, in the pure light of subjectivity, in the transparent medium of thoughts, reflections and concepts; for all thinking, reflecting and comprehending is indeed an ideal acting; but we have seen in the foregoing that the functions identical with pure thinking and reflecting cannot possibly assume an exact form, hence cannot possibly be mechanical. Nevertheless Herbart would seek the mechanics of spirit in spirit itself. It is this misunderstanding that from the beginning has misled him into abstracting from the very first mechanical condition of mechanics, matter with inertia. ``In the mechanics of bodies,'' it says, ``it is the differential of the velocity, but in the mechanics of spirit it is the velocity itself, that is determined by the force, for in the mechanics of spirit no inertia holds; here there is motion only insofar as the force acts.''

If one now holds fast to the assumption that the only thing that can give the attempt at a mechanics of spirit scientific value is the endeavor to open for oneself an insight into the possibility of the articulated reciprocal action of the spiritual and the organically bodily, then it undeniably makes a depressing impression when one learns that in a mechanics of spirit no \opage{420} account is to be taken of the bodily, since no account at all is to be taken of the material.

It must be granted that Herbart, with his foundation and his peculiar procedure, has arrived at results of calculation which in an ingenious manner shed light on psychological phenomena. Among these results of calculation we will in this connection recall those whose general equation is of the form: $y = f(x) = e^{-x}$. Since Herbart's basic view of the representations and their reciprocal action has, through exact consequences of exact presuppositions, led him to the application of this equation, he has found expressions for transitions of consciousness which are in a high degree characteristic. The transition from motion to rest is hereby presented in such a way that rest in the finite sense is suddenly attained, while in the absolute sense it is nevertheless unattainable.\footnote{Since the examples belonging here, with the corresponding calculations, have already been treated in my lectures on ``Philosophisk Propædeutik'' from the academic year 1860--61, it will be sufficient here to recall the results.

\emph{Attention.} Let the force of attention be $\varphi$; after the lapse of the time $t$ the strength of the representation has grown to $z$. Of the force of attention there then remains $\varphi - z$. If now $\beta$ is a constant denoting the intensity of the apprehension, we obtain: $z = \varphi(1 - e^{-\beta t})$, and
\[ \frac{dz}{dt} = \beta\varphi e^{-\beta t} \]
From which it is seen that the force with which attention apprehends and grasps its object decreases \emph{suddenly}, but \emph{never} becomes nothing. Forelsngr.\ over ``Phil.\ Propæd.'' p.~376 note.

\emph{The motion of representations.} If one or more representations suddenly press into consciousness, the equilibrium among the representations already there is disturbed. The sum of inhibition, the obscuring which the new representations cause in the others, now acts as a pressure that is to be distributed over the single representations, and will only vanish when equilibrium has again been brought about. If the sum of inhibition is now called $S$, and $\sigma$ the part of $S$ expended at a certain time $t$, one gets $S - \sigma = Se^{-t}$. From this it is seen that the mind, after a disturbance has entered the sphere of representations, \emph{suddenly} reaches a rest which it nevertheless in the absolute sense can \emph{never} reach. Forelæsngr.\ over ``Phil.\ Propæd.'' pp.~462--63.

\emph{The reawakening of representations.} A representation $P$ which under certain conditions has sunk below the horizon of consciousness will under favorable conditions rise again, and that according to a law which, when $p$ is what has entered consciousness by the time $t$, is stated by the equation: $p = P(1 - e^{-t})$. Here we again have a \emph{sudden attainment} of something nevertheless \emph{unattainable}. If now a remainder $r$ of $P$ is ``fused'' or ``complicated'' with $\varrho$, the remainder of another representation $\Pi$, then $\Pi$, as $P$ rises, receives an aid $\frac{r\varrho}{\Pi}$ corresponding to the fusion or complication. The aid $\frac{r\varrho}{\Pi}$ will then act on $\Pi$ in such a way that, when the part of $\Pi$ which has already been raised is set $= \omega$, and the part, therefore, which is still to be raised $= \varrho - \omega$, the equations:
\[ \frac{\varrho r}{\Pi} \cdot \frac{\varrho - \omega}{\varrho}\,dt = d\omega \]
\[ \text{and}\quad \omega = \varrho\left(1 - e^{-\frac{rt}{\Pi}}\right) \]
likewise point to a sudden attainment of something in itself unattainable. Forelæsngr.\ over ``Phil.\ Propæd.'' pp.~466--67.}

\opage{421} The attainment of something unattainable has an extraordinary significance in all forms of the life of spirit, and not least in the mechanics of spirit. We forget something, e.g.\ a name. If we do not at the same time forget that we have forgotten it, but try to recall it, then what was forgotten can, after a shorter or longer lapse of time, suddenly come forward to consciousness: how is this to be understood? In order to grasp the mechanics of memory we must seek to understand the mechanics of forgetting. What is forgotten is something vanished for consciousness; but if its trace in the sphere of representations and in the organic conditions of that sphere \opage{422} were actually $= 0$, then there would be no memory. Memory is possible only on the presupposition that what is forgotten is nevertheless not forgotten. If one says: what is remembered has not been forgotten; it has only been pushed back into the background of consciousness, then what is ``the background of consciousness'' other than a figurative expression for the thought that here is something vanished which has nevertheless not vanished? Mechanically, the matter stands thus: the organopsychical indication of a representation can be weakened ad infinitum, and yet the representation can be recalled as long as the trace is not downright effaced; thus a finite magnitude can shrink to a differential and be integrated again, but the pure $0$ cannot be integrated.

The profound basic thought, decisive for the mechanics of spirit, which the Herbartian theory has thus introduced, has now, as was said, the repellent feature about it that it excludes all regard to inertia, and lets the first differential coefficient express the force. Were it, as Herbart assumes, impossible to arrive at similar results with other mechanical presuppositions, then every attempt at a mechanics of spirit would certainly be fruitless. But the matter does not stand so. By building on the general principles of analytical mechanics, acknowledging inertia and letting the 2nd differential coefficient be the measure of force, one will attain results which express the acknowledged main thought in an equally satisfactory manner.\footnote{Thus, to mention only one example, the expression cited in the previous note, $S - \sigma = Se^{-t}$, is certainly obtained on the presupposition that $\frac{d\sigma}{dt} = S - \sigma$, that it is thus the first differential coefficient that denotes the force; but by letting the second differential coefficient denote the force and proceeding from the equation
\[ \frac{d^{2}\sigma}{dt^{2}} = h(S - \sigma) - k\frac{d\sigma}{dt} \]
one arrives, naturally by another mode of integration, at the result:
\[ \sigma = S\left[1 - e^{-\frac{kt}{2}}\left(\cos qt + \frac{k}{2q}\sin qt\right)\right], \quad\text{where}\quad \frac{k^{2}}{4} - h = -q^{2}, \]
a result which can satisfy the basic requirement just as well as the Herbartian one. That, moreover, the fundamental equations cannot be altered in particulars without these alterations in turn reacting on the whole is self-evident.}

\opage{423}That the functions of the mechanics of spirit have psychological significance is, then, certain; but now the question arises: does the mechanics of spirit not also have its ontological significance? In answering such questions it is chiefly a matter of following the law-determined transformations without letting oneself be dazzled by analogies. Were the relation between the psychological and the ontological --- whether it is otherwise a question here of the determination of the objective or of the subjective --- to be found out and fixed by means of analogies, then one could never carry it further than to conjectures and loose hypotheses. Thus in the present inquiry one would arrive at highly different results according as one made an intellectual, a mystical or an anthropomorphistic use of analogy.

From an abstractly intellectual standpoint the matter stands thus: ``Psychological subjectivity has the organism, the ensouled organism, as its real presupposition; the organopsychical functions from which the life of consciousness proceeds must necessarily have their organic-mechanical side, and the mechanics of spirit must maintain its validity in the psychological. It is otherwise with the ontological: ontological subjectivity is without organism, without organs, without soul, hence without organopsychical functions; to a subjectivity which, in the form of infinite self-power, posits and determines matter and everything material, the mechanics of spirit cannot possibly find any application.'' By such a negative use of analogy the critique of the \apage{424} understanding would guard itself against confusions of concepts, and does not notice that the negation must always be as the position is. If the position has only an analogy to offer, and the negation only an analogy to reject, then the negative can be just as misleading as the positive. If there is no analogy between psychological selfhood and the ontological unity of reason, then it is of course altogether improper to call the infinite unity of reason an ontological selfhood, still more improper to call it subjectivity, to speak of ontological comprehending, ontological acts of objectification, ontological sense-perceptions, etc.

If one has an immediate eye for the logically desperate in rejecting all analogy between psychological and ontological subjectivity, and at the same time an eye for the overreaching power of infinity over the finite, then one is led, if critique is lacking, to use analogy mystically. The mystical is then seen in this, that identity rules with a floating difference. From this standpoint the all-pervading unity of reason is apprehended as a living selfhood, a universal world-life, a world-soul, a world-spirit; and now it is quite floating whether one is to regard the reason living in everything as a thinking, knowing and conscious unity, or to let it work in the great whole of nature, the world-organism, the ensouled world-body, as a blind master-builder, until the dawn begins in the animal world and the light rises in man. How groping, unclear and floating the whole mode of view is can indeed also be discerned from this, that it is even left undetermined whether psychological subjectivity is an actual self-unity for itself in opposition to the universal unity of reason, or a concentration of the unity of reason itself under finite, transient conditions.

If one finally attempts, in order to get out of the unclarity, to set what is alike in the analogy against what is unlike, and thus to state the likeness between psychological and ontological subjectivity by itself, and to set the unlikeness up beside the likeness, then \opage{425} one uses analogy anthropomorphistically. According to the anthropomorphistic mode of view the following explanation might well be the most reasonable. ``Psychological subjectivity and eternal reason, or, in other words, the human and the divine spirit, resemble each other in that they are both knowing, thinking, willing; but the unlikeness is that divine thinking is infinitely perfect, human thinking imperfect, etc. The divine and the human spirit resemble each other in that they both act on matter; but the unlikeness is that the finite spirit is powerless and in its acting on matter dependent on organs, while the divine spirit, almighty in its divinity, can by mere will act upon all matter, and that still more immediately than we by our mere will can move hand or foot, etc.'' But who does not now see that comparisons of this kind must always dissolve into the contradiction that divine and human things resemble each other in being unlike each other?

And yet analogy is of essential significance, not only in empirical but also in a priori knowledge, not only in natural history but also in logic. But the point is that analogy must be used differently in a priori than in empirical investigations. In a priori knowledge analogy has only reflected, not immediate, validity; the reflected validity of analogy shows itself in this, that it serves as the basis for law-determined transformations of concepts. Thus in the present case the task is this: to determine what transformations the concept of the mechanics of spirit must undergo when this concept, after being elucidated from a psychological standpoint, is to be seen in its ontological significance?

If the organopsychical function were in truth the highest expression of the reciprocal action between spirit and matter, then there could undeniably be talk of a \opage{426} mechanics of spirit only as regards psychological subjectivity, human consciousness. But now it has of course been shown that the ideas life, soul and relative spirit have reality only insofar as they rest and move in the eternal, absolute idea of subjectivity; further, it has been shown that the relative systems, the mechanical, the elemental, the organic system of functions, have objective validity only insofar as they are objectifications in which a powerful knowledge, a knowing power, realizes the concepts of existence.

In human existence the soul forms the uniting middle of body and spirit. The soul is, as organic life, bodily in its unconscious functions, and, as a self-thinking being, spiritual in its conscious functions: the soul is in a dualistic state; it is as two different beings, and yet one being. With the dark, nocturnal side of the soul as its ground of life, the spirit of man cannot become immediately conscious of its relation to its bodily organs; the organs by which it observes everything else and itself, it cannot observe. The functions of the mechanics of spirit are, at the stage of human life, organopsychical; and this is again a proof that the relation between spirit and matter is not an original but a derived relation.

In the divine life, in the idea itself, the relation of spirit to matter is an original relation of power. Identical with power, the knowing subjectivity, the almighty spirit, apprehends matter and the material without organs: the organopsychical functions fall away; here there are ontological sense-perceptions, but no psychical middle term between matter and spirit. In the identity of selfhood with the ontological sense-perceptions the idea is life; but since this identity is in an infinite way reflected in the eternal light of subjectivity, in the unity of power with knowledge, life, with its infinite wealth of immediate determinations, is not a life of the soul but a life of the spirit: the spirit absolute in the idea has an eternal life, but no soul.

It is the relation of subjectivity to the objective, of knowledge \opage{427} to power, on whose articulating reciprocal determinations everything here depends. If knowledge in its extreme is to be characterized by the pure logical function of selfhood, which sublates all finitude and transforms the determinations in the material into negative reflexes in the ideal, then power in its extreme is to be characterized by the antilogical function of otherness, which posits the finite and thereby converts the ideal forms into material productions. The logical function of selfhood expresses the transition of power into knowledge, and denotes the ideal system, in that it denotes spirit; the antilogical function of otherness expresses the transition of knowledge into power, and denotes the real system, in that it denotes the material world. As the middle of the extremes, the mechanics of spirit is thus a categorical expression for the equilibrium with which the subjective freedom of the ideal system and the objective necessity of the real system hold themselves in teleological reciprocal action; for power is the end of knowledge just as knowledge is the end of power. It is through this reciprocal action that the subjective content of the realm of ideas unfolds itself in the multiplicity of functions, in the constant and yet varying forms of objectivity.

\begin{center}{\bfseries γ)\quad The Energy of the Functions.}\end{center}
\addcontentsline{toc}{paragraph}{γ) The Energy of the Functions}

The functions are known by dependence, dependence by change, change by reflecting and acting. Reflecting and acting are related as selfhood and otherness, yet with the express determination that the doubling is seen from both sides; there is acting in all reflecting and reflecting in all acting; there is a tinge of otherness in every selfhood, there is a semblance of selfhood in every otherness. The functions now belong to the intermediate region where the light of selfhood is refracted in the medium of otherness, and so they also have a double energy, a material one, which belongs to otherness, and an ideal one, which belongs to selfhood.

In the articulating reciprocal determination of selfhood and \opage{428} otherness, hence in the sphere of the functions, falls totality. There is no true totality without selfhood, because there is no essential whole without unity; there is no true totality without otherness, because there is no actual whole without parts. Just as the material energy of the functions shows itself especially from the standpoint of the parts, so the ideal energy shows itself especially from the standpoint of unity, of self-unity. In essence and ground totality is subjective, in existence and actuality objective. The subjective essence of totality is seen in its design, and the ideal energy of the functions determines the design; but actuality is seen in the execution of the design, and in the execution the functions especially prove their material energy. As little as a design can be without all execution, just as little can the ideal energy be without material conditions; but since totality can only be actualized by its design being executed, the material energy must at all points presuppose the ideal.

The actual opposition of the design and the execution of totality reflects itself in an originally active selfhood; the active selfhood, the cause of totality, is ground-cause; every ground-cause is self-cause. The actual opposition between the parts of totality reflects itself in active otherness as a multiplicity of divided causes, a division of causes which makes finitude actual and conditions the execution. The two-sided essentiality of the functions is now recognizable by this, that with their ideal energy they serve the ground-cause (so that each function, however particular it may seem to appear, is essentially a function of totality), while nevertheless, with their material energy, they are all, each in its own way, disposed to enter into finitude, into dividedness, into parceling-out, into the multiplicity of the objective.

As the totality is, so must its functions be: every totality is in actuality itself a function of its actual functions. Conversely, the functions belonging to a totality, \opage{429} however much they may seem, from a standpoint of finitude, to proceed from the parts and to produce the whole in question as a finite result of their finitely conditioned cooperation, are nevertheless, when the matter is taken from the ground up, originally determined by the totality which through them determines its development: the functions of totality develop only the totality by which they themselves are developed.

The transition from the design of totality to its execution, determined by the \emph{energy of the functions}, is cognized only when it is seen in a \emph{system of syllogisms}, indeed of \emph{teleological syllogisms}. ``Just as the rational in everything that exists consists in its being a syllogism, so'' --- says ``the speculative logic'' --- ``it shows itself, more definitely expressed, in this, that everything is a system of three syllogisms''; and in this we must grant ``the speculative logic'' that it is right, just as we must also grant it that it is right in saying that ``teleological activity can first be seen in its infinity in the completed system of syllogisms''. But the ambiguity and formalism from which the Hegelian presentation of the system of syllogisms evidently suffers is in equal degree misleading, whether we regard the matter from the objective or from the subjective standpoint.

The ambiguity which arises from this, that the movement of the concept, according to the immanent method, is lumped together with the actual movement of the thing, and syllogisms of reason, which according to their ideal nature rest on ontological acts of thought and belong to the ideal system, are lumped together with the acts of power which, in the existential transitions, with their now partial, now total mediations and middle terms, belong to the real system and express the logical content of the syllogisms in an antilogical way, has already been elucidated in the logic of subjectivity. In the present connection, where it is chiefly a matter of an understanding of the objective, insofar as developed objectivity in turn casts light back on the subjective, the following main points are to be especially emphasized: 1) the opposition between the universal, the \opage{430} particular and the singular is, however often it is repeated that each of these sides of the concept is itself the whole concept, too abstract and formal to bear any more concrete development of the system of syllogisms; 2) the separation and union of extremes and middle terms by which each single syllogism is characterized is, in the Hegelian logic, apprehended and presented in so formal a generality that schematism, barren schematism, shows itself immediately every time the validity of the syllogism is to be cognized in concreto. That both awkwardnesses evidently derive from this, that the functions of totality find no development in and with totality itself, can be shown without difficulty.

1) \emph{The opposition between the universal, the particular and the singular is to be determined in totality by the energy of the functions}. It is a one-sidedness, closely connected with the Hegelian holding fast to the formal unity of thinking and being, that the copula ``is'' is to denote the return of reflection to the immediacy of being reborn in the concept. A consequence of this is that the three main sides of the concept, which at one and the same time both coincide \emph{as being} and yet stand \emph{as being} against one another, must lose the ideal as well as the material energy and fall to schematism. The schematism from which the system of syllogisms suffers is so superficial that the typical terminology even comes into conflict with itself.

According to the schema E --- S --- A, E, the singular, should of course always be the subject, A, the universal, always the object, and S the particular, which united the subject with the object. But if we look more closely, the singular, according to the development of the concept, can represent the object exactly as well.\footnote{``Die sich auf sich selbst beziehende Bestimmtheit\ldots ist die \emph{Einzelnheit}. So unmittelbar die Allgemeinheit schon an und für sich selbst Besonderheit ist, so unmittelbar an und für sich ist die Besonderheit auch \emph{Einzelnheit}\ldots'' Hegel, Subj.\ Logik p.~51. What is said here of the singular fits exactly what is said later of the object.} ``For, in that
\opage{431} every one of the three terms, E, S and A alike, has shown itself to be the connecting middle\ldots, then the judgment E — A\ldots has returned to the immediacy which has mediation as its presupposition. But in this way it is no longer a judgment, for E is no longer subsumed under A; or A is no longer determined as the other of E, or as predicate, but as the sublated other (it is sublated, namely, because the mediation S, whereby it was determined as other, is itself sublated), consequently as the subject itself. Thus the subject is determined as \emph{object}; and the judgment E — A has the meaning that the subject is immediately object, or that the subject without mediation closes itself together with itself, and in this form is object.''\footnote{Heiberg. Specul.\ Logik.\ Pros.\ Skr. vol.~1, p.~321--22.} Should one not now think that the singular, after having run through this formal mediation, would have to stand as near to the object as to the subject?

``The object,'' it is said, ``is causality returned. In causality there are only two terms: cause and effect. In the object there is in addition the \emph{ground}. The cause is E, the ground S, the effect A.''\footnote{Heiberg. Anfr.\ Skr. p.~327.} But, we ask, how does it come about that the effect here acquires a greater extension than the cause, that the effect can be the universal although the cause is the singular? Or do cause and effect perhaps have the same extension? But then is it not preposterous to call the effect the universal and the cause the singular? If the syllogism E — S — A is teleological, then in ``the speculative logic'' E is ``the end as subjective,'' A ``the end as objective''; but on what ground? In actuality the relation is, after all, exactly the reverse. If E is to represent the end at a certain stage, \opage{432} then it must surely be the objective, not the subjective end; for a realized plan is not a plan in general, but this plan, in this actual execution.

With what arbitrariness the Hegelian logic, in its systems of syllogisms, applies the opposition of the universal and the singular is shown, moreover, by the philosophy of nature. In the solar system, e.g., it is the planet that is designated by the singular, while the sun represents the universal. Now that the sun, as the central body existing for itself, cannot be the universal is evident; but even if one went along to a certain degree with Hegel's conception of the sun as the universal source of light of the system, it would nevertheless, even from an unconditionally Hegelian standpoint, surely be difficult to hit upon a ground with which one could defend the assertion that the sun, ``the body of abstract centrality,'' should be the \emph{object} of the system, and the planet, ``concrete individuality,'' on the other hand the \emph{subject} of the system. But when the sun, although it is designated by A, is not the object, and the planet, though designated by E, is not the subject, then it does not hold either, after all, that E in the schema E — S — A always designates the subject, and A the object.

In order to be rid of the formalism, we let each of the sides, A, E and S, designate by itself the whole totality, since the concept as concept must necessarily be totality both in form and in existence. Since the one and the same totality in this way becomes three fundamentally different totalities, the unity-in-difference totalized as syllogism must be determined by corresponding functions. In the form A the totality is subjective, in the form E objective, in the form S subjective-objective. The subjective totality is the content of subjectivity; its unity is self-unity, its functions functions of selfhood: \emph{its idea is knowledge}. The objective totality consists of the terms of objectivity; its unity is connection, its functions functions of otherness: its \emph{idea is power}. The subjective-objective totality, the totality in concreto, consists in \opage{433} reciprocal determinations; these reciprocal determinations are functions of the overarching selfhood that comprehends the totality: \emph{its idea is truth}. That the three totalities are in truth one and the same totality is explicable only from the energy of the functions. By virtue of the ideal energy the functions of the totality are essentially functions of selfhood, hence functions of knowledge; by virtue of the material energy the functions are altogether functions of otherness, hence functions of power; by virtue of the totalized reciprocal determination of ideal and material energy the functions of knowledge and of power are, in articulated identity of inner infinity, functions of truth.

In the totality as totality there must absolutely be two extremes: that of knowledge or of subjectivity, and that of power or of objectivity; these two extremes have, in the concretion corresponding to truth, their uniting middle; but this middle of the totalities is again the totality itself. The totality itself is thus at once its own extremes, its own premises and its own conclusion. Each time the syllogism is at an end, it begins over again: the true totality has two extremes, whose syllogism it is; but each of the two extremes is a true totality, a syllogism of two extremes, etc. In signs this can be expressed thus: E — S — A is a syllogism; the extreme E is a syllogism e — s — a; e again a syllogism $\varepsilon$ — $\sigma$ — $\alpha$, and so on to infinity.\footnote{What lies hidden herein must be developed from another point of view in the doctrine of \emph{the logical functions}.} In the system of the three syllogisms there is a swarm of systems of syllogisms; and the question is how the relation between the unity of the system, the fundamental system, and the multiplicity of the system, the multiplicity of the relative systems, is to be more precisely determined.

``The system of the three syllogisms,'' it is said,\footnote{Heiberg. Anfr.\ Skr. p.\ (Pros.\ Skr. vol.~1, p.~315.)} ``can partly \opage{434} be regarded as a \emph{progression}, inasmuch as in each syllogism one appeals only to the preceding ones, partly as a \emph{regression} or \emph{circle}, inasmuch as each of them nonetheless presupposes the other two.'' But when the relation between progression and regression is to be more precisely determined, the relation between the fundamental system, the original totality, and the relative systems, the totalities derived from the fundamental totality, must be determined by the peculiar character of the corresponding functions. The emphasis can then no longer fall on the opposition between the singular and the universal; for the singular, which according to the material energy of the functions is the objective, is according to the ideal energy the subjective; conversely, the universal, which according to the energy of selfhood, the ideal energy, is the subjective, can according to the material energy, the energy of otherness, nevertheless also be the objective (the system of laws and of relations of power).

Since, then, the emphasis falls on the opposition between selfhood and otherness, between subjectivity with its inner and objectivity with its outer actuality, the extremes of the syllogism ought no longer to be designated by E and A, but by S and O. The middle term which now forms itself between S, the absolute selfhood, the original subjectivity, and O, objectivity, the actual totality having determinate being in the energy of otherness, certainly continues to be the particular; but since this particular nevertheless, in its determinateness, conforms to the more precisely determined extremes, we shall for the sake of clarity designate it in different ways, according as the context in each case requires.

2) \emph{The system of the totalities can be determined only by systematic changes in the energy of the functions}. Which and how many fundamental ideas are there? This question must naturally admit of being answered in different ways when one chooses different points of departure; but systematically viewed the answer will always depend on the fundamental division with which the systematics of the totalities \opage{435} originally begins. As unity of subjectivity and the objective, the idea is what is absolute in every totality: as many peculiar fundamental totalities, so many fundamental ideas. How the absolute idea lays out its content in a system of relative ideas can be seen only insofar as it is seen how the original subjectivity, which embraces all and encloses all within itself, objectifies its content in systems of lower and higher fundamental totalities.

\begin{center}{\bfseries αα)\quad The first fundamental totality of objectification is the abstract-mechanical.}\end{center}
\addcontentsline{toc}{subparagraph}{αα) The first fundamental totality of objectification is the abstract-mechanical}

The opposition of selfhood and otherness, of knowledge and power, of subjectivity and objectivity, strikes through immediately. The extremes $S$ — $O$ stand so sharply against each other that a word of power must dictate their union. If the emphasis is laid decisively on $S$, then subjectivity is cognized as the substantial selfhood, the sole-mighty self-cause; every term in the totality $O$ is then at once reduced to a moment for the overarching unity of power. If, on the other hand, the decisive emphasis falls on $O$, then the substantial moments are at one stroke transformed into relative substances; otherness rules: self-unity has given place to the unity of connection.

On whichever side the word of power throws itself, it yields only an extreme incomprehensible in itself; if word of power is set against word of power, it yields two extremes; but a syllogism, a rationally grounded cognition, it does not yield. If the gaze fastens on the unity, it meets universal gravity; before this unity of power the independence of the terms must vanish; gravity is the substantiality of the masses: how should the heavy masses with their atoms have independence over against gravity itself? As universal gravity the unity of power falls under subjectivity; and yet gravity as power is an objective unity which has its reality in a substantial multiplicity, in the immediate independence of the terms. It is, then, not gravity itself but the heavy bodies that \opage{436} are here the actual representatives of power. Instead of being a power for itself, gravity, seen in the extreme $O$, is a property of the material parts; this common property of the parts can manifest itself only in their mechanical reciprocal action: self-unity is transformed into unity of connection.

The unrest of the unity of power, an abrupt conversion of selfhood and self-unity into otherness and unity of connection, can be sublated only in the system of mechanical motions. The system is thus itself the middle of the extremes; but this middle is, be it noted, not an objective actuality (for $O$, the objective actuality, is itself an extreme), but the disposition of actuality, and as disposition the identity of the subjective and the objective. The question of the origin of the edifice of the world (the solar systems) is in the logical sense a question of the disposition of the mechanical totality. If we designate the disposition by $A$ and, in order to indicate its peculiar character as cosmic disposition, add $\varkappa$ as index, then $O$ must receive the same index, and then
\[ S - A_{\varkappa} - O_{\varkappa} \]
is the general expression of totality of the mechanical syllogism.

The more eagerly reflection now labors to resolve $A_{\varkappa}$, the greater will be the mechanical strife between self-unity and unity of connection. This strife arises precisely from the fact that the mechanical functions, according to what we have shown in the foregoing, are abstract functions of the cause: the self-cause devours the cause of otherness, and this in turn the former; both must be posited and both must be sublated in one and the same now. If one appeals one-sidedly to $S$, then $A_{\varkappa}$ represents the unity of connection and thereby otherness in the external causes and the finite conditions; one is thus referred to $O_{\varkappa}$: \emph{the functions of the totality are functions of otherness and act with material energy}. If one appeals one-sidedly to $O_{\varkappa}$ and wishes to remove $S$, or at least to let it fade away, then $A_{\varkappa}$ represents self-unity, the mighty selfhood; it is then not possible to piece together the system of finite causes with finite terms in finite reciprocal actions; for \opage{437} the disposition requires that the finite presuppositions must all have their ground in one presupposition: \emph{the functions of the totality are functions of selfhood and act with ideal energy}.

From whichever extreme the mechanical functions are regarded, the abstract relation of cause goes into one with the formal relation of ground; the abrupt conversion of selfhood into otherness, and conversely, reflects itself in an abiding, totalized unity: in the totalized unity the mechanical unrest is essentially rest, and the mechanical rest an essential unrest.

\begin{center}{\bfseries ββ)\quad The second fundamental totality of objectification is the mechanical-dynamic.}\end{center}
\addcontentsline{toc}{subparagraph}{ββ) The second fundamental totality of objectification is the mechanical-dynamic}

The difference between ``the mechanical'' and ``the chemical object'' is presented in Hegel's Logic in the following manner: ``Das chemische Objekt unterscheidet sich von dem mechanischen dadurch, daß das letztere eine Totalität ist, welche gegen die Bestimmtheit gleichgültig ist; bei dem chemischen dagegen gehört die \emph{Bestimmtheit}, somit die \emph{Beziehung auf Anderes}, und die Art und Weise dieser Beziehung, seiner Natur an.''\footnote{Hegel. Subj.\ Logik p.~200--201.} From this it is then explained that \emph{difference} as a whole is for chemical objects what \emph{indifference} is for mechanical ones, etc. ``Chemism'' and ``the chemical process'' thus stand higher in the self-development of the concept than ``mechanism'' and ``the mechanical process.''

But although the concept of chemism in Hegel's Logic is expressly designed for a higher development of the objective essence of the totality, it nevertheless happens that what according to ``the immanent method'' is thought as an advance becomes in actuality a step backward. For it is evident that chemism, in spite of its cycle of different processes, when it is conceived in the way Hegel has presented it, cannot possibly, as regards the syllogism of totality, keep itself at the same height as ``free mechanism.'' In systems of the \opage{438} mechanical totality the idea of the edifice of the world is expressed; but what idea, pray, is expressed in the system of processes which, on Hegel's view, is to unfold the chemical totality? ``Um der Unmittelbarkeit und Aeußerlichkeit willen\ldots, in deren Bestimmung die chemische Objektivität steht, \emph{fallen diese Schlüsse}\footnote{For in what immediately precedes there is talk of three chemical syllogisms, whose peculiar character is presented in the following manner: ``Die drei Schlüsse, welche sich ergeben haben, machen seine (the chemism's) Totalität aus; der erste hat zur Mitte die formale Neutralität und zu den Extremen die gespannten Objekte, der zweite hat das Produkt des ersten, die reelle Neutralität zur Mitte und die dirimirende Thätigkeit, und ihr Produkt, das gleichgültige Element, zu den Extremen; der dritte aber sich der sich realisirende Begriff, der sich die Voraussetzung setzt, durch welche der Proceß seiner Realisirung bedingt ist, — ein Schluß, der das Allgemeine zu seinem Wesen hat.'' One sees from this how necessary it is to treat the doctrine of syllogisms as many-sidedly as possible. In the example cited, the determinations which essentially belong to the objective coincide without further ado with those which rest exclusively on the subjective: the logic of functions is confounded with the logical functions.} \emph{noch auseinander}. Der erste Proceß, dessen Produkt die Neutralität der gespannten Objekte ist, erlischt in seinem Produkte, und es ist eine äußerlich hinzukommende Differentiirung, welche ihn wieder anfacht\ldots Ebenso muß die Ausscheidung der differenten Extreme aus dem Neutralen, ingleichen ihre Zerlegung in ihre abstrakten Elemente, von \emph{äußerlich hinzukommenden Bedingungen} und Erregungen der Thätigkeit ausgehen. Insofern aber auch die beiden wesentlichen Momente des Processes, einer Seits die Neutralisirung, anderer Seits die Scheidung und Reduktion, in einem und demselben Processe verbunden sind, und \emph{Vereinigung} und Abstumpfung der gespannten Extreme auch eine \emph{Trennung} in solche ist, so machen sie \opage{439} um der noch zu Grunde liegenden Aeußerlichkeit willen \emph{zwei verschiedene} Seiten aus\ldots''\footnote{Hegel. Anfr.\ Skr. 207.}

Here one sees at once what is amiss in the logic of objectivity and the logic of objects being lumped together into one without further ado. If, in all finitude, we set the mechanical object against the chemical, then we must surely concede to Hegel that the former stands lower in the concept than the latter. But the object in its finite determinate being is one thing; the sphere of objectivity to which the finite object must be reckoned is another. Each globe in the solar system can, when taken by itself, be regarded as object: to weigh the globes is precisely to determine them as mechanical objects; but the solar system itself is, in its total objectivity, more than an object; it is a realized idea, and a realized idea is more than an object. So too, then, the chemical products in the earth's crust can rightly be regarded as chemical objects; but the terrestrial globe itself, in its essential totality, is no object; it is a realized idea, and a realized idea is no object. If, now, the chemical totality is to indicate a higher stage in the sphere of objectivity than the mechanical, then the \emph{difference} of objects must not be restricted to designating the opposition between two substances which conditions their \emph{affinity}, but must be extended to a designation for the multiplicity of qualitative differences which originally belongs to the determination of an elementary totality closed in itself. The mechanical-dynamic function is now, according to its character, since it is, as we have seen, so bound up with the real and mixed grounds, precisely suited to embrace the qualitative moments and elementary differences in rich multiplicity, and just thereby suited to form an elementary totality. The syllogism corresponding to this,
\[ S - A_\gamma - O_\gamma, \]
\opage{440} whose index suggests the geological,\footnote{That by a geological totality here one is not to think of our terrestrial globe alone is presumably understood of itself.} has, like the mechanical, its peculiar essence in its peculiar middle.

What is peculiar to the middle is in turn determined by the objective extreme $O_\gamma$. This extreme differs from the extreme $O_{\varkappa}$ in that its functions are not abstract-mechanical, but mechanical-dynamic. In the dynamic there is already here a more concrete ideality and to that extent an inwardness; the extreme $O_\gamma$, then, does not form so abrupt an opposition to $S$ as the extreme $O_{\varkappa}$ does. Nonetheless, the opposition between the extremes is absolutely thoroughgoing. So long as the dynamic function merely supports itself on the mechanical, its ideality is material and its inwardness elementary. An elementary inwardness is essentially externality. The dynamic processes, the elementary formations, lose themselves in the external: the different objects are external, separations and combinations are conditioned by something external; and yet each phenomenon in this external essentially points toward something internal, indeed is even converted term by term into something internal; here lies the ambiguity in $O_\gamma$.

If, now, $O_\gamma$ is posited immediately in its objective extreme, then the dynamic becomes a moment for the mechanical. Thereby externality and material energy gain such a preponderance that the elementary functions sink down into mere functions of otherness; finitude, parceling out and piecing together, rule to such a degree that $O_\gamma$ has ceased to be totality. If $O_\gamma$ is reduced outright to its subjective extreme $S$, then the mechanical vanishes as a moment in the dynamic, and the dynamic as a moment in free inwardness: the elementary functions lose their material energy and transform themselves into mere functions of selfhood. In this way, then, $O_\gamma$ cannot become totality either. To an actually objective totality belongs a material otherness; but the extreme $S$ is, in \opage{441} immediate contact with every immediate otherness, every immediate objective thing, like a consuming fire that does not leave behind even ashes of objectivity. The reality of $O_\gamma$ thus depends on the ideality of $A_\gamma$, but the ideality of $A_\gamma$ depends on $S$. In its unity-in-opposition with $S$ the elementary $A_\gamma$ is essentially selfhood; in its unity-in-difference with $O_\gamma$ the selfhood in $A_\gamma$ is essentially elementary.

\begin{center}{\bfseries γγ)\quad The third fundamental totality of objectification is the organic-living.}\end{center}
\addcontentsline{toc}{subparagraph}{γγ) The third fundamental totality of objectification is the organic-living}

From the above it is seen that $S$ in the different syllogisms remains unchanged, while the changes take place in $O$ and thereby in $A$. This is, moreover, a direct consequence of the fact that subjectivity, the overarching selfhood, essentially expresses the absolute, while otherness, objectivity, expresses only the relative. Certainly objectivity, each time a peculiar whole is formed, takes on the stamp of something closed in itself and to that extent absolute; but since the relatively absolute is and remains a moment totalized in the absolutely absolute, the fundamental relation between $S$, the absolutely absolute, and $O$, the relatively absolute, must to that extent be constant. Between $S$ and the objective totalities of different order, $O_{\varkappa}$, $O_\gamma$ and now, as regards the organism, $O_o$, the fundamental relation is on the other hand variable, since the change in $O$ must necessarily presuppose a change, not in $S$, for $S$ is unchangeable in its absolute absoluteness, but in the relation in which $S$ has originally placed itself to $O$. The change in the fundamental relation between $S$ and $O$ is determined by the middle term $A$ in such a way that $O$ in its material wholeness can be said to be a function of $A$, and $A$ in its ideal wholeness a function of $S$. $O$ is thus, seen from this point of view, not directly a function of $S$, but, owing to the intervening $A$, a function of a function of $S$. If, now, the peculiar relation in which $S$ originally places itself to $A$ were determined immediately by formal freedom and to that extent posited with divine arbitrariness, then the relation between $S$ and $A$ would be a groundless \opage{442} relation, hence a groundless fundamental relation; but a groundless fundamental relation is a contradiction. The relation of freedom in which $S$ places itself to $A$ is in turn determined by the relation of necessity in which $A$ stands to $O$; $A$ is thus, owing to the intervening $O$, not directly a function of $S$ either, but a function of a function of $S$.

The relation of dependence in which $O$ originally stands to $S$ is determined by the ideal energy of finite functions, for the determination is mediated by $A$; the relation of dependence in which $A$ stands to $S$ is, on the other hand, determined by the material energy of the functions, insofar as the determination is mediated by $O$. If, now, the relation between the ideal and the material energy of the functions were to be conceived as an original fact, an eternal factum for which there was no ground, then $O$ would without further ado present itself as a function of $A$, and so here we would get a groundless fundamental relation. But in this context it is easy to see that if one of the fundamental relations is groundless, then they all are.

To make $O$ an immediate function of $A$ is, in other words, to make the consequence a formal paraphrase of the presupposition. Thus the arrangement of the solar system becomes a direct consequence of that presupposition of presuppositions which had its immediate determinate being in ``the primeval nebula''; thus the present stage of development of the terrestrial globe is a direct consequence of that totality of presuppositions whose first existence was a rolling vapor, torn loose from the original mass, etc. But in thus making the consequence a formal paraphrase of the presupposition, one at the same time makes the presupposition a formal paraphrase of the consequence. What motions, what strivings can there have been in ``the primeval nebula''? what basic substances, what elements in the earth's ``rotating vapor''? When these questions are to be answered, one must proceed from what is given; but what is given, which now becomes presupposition, \opage{443} is, after all, precisely a consequence of the presupposition that is to be found out as the consequence of the given presupposition.

That with the syllogisms one goes in a circle is in and for itself no fault; the fault lies in the formal: that one constantly grasps after a first immediate fact instead of comprehending the ground. If $O$ cannot be an immediate function of $A$, then $S$ must absolutely step in as mediator between them; $O$ then becomes, owing to the intervening $S$, a function of a function of $A$. Herewith the grounding of the totality is formally carried through, and by its formal carrying-through real. For the organism, $O_o$, the syllogism will thus be:
\[ S - A_o - O_o. \]
Already by the fact that $O_o$, in comparison with $O_\gamma$ and $O_{\varkappa}$, is so limited and in its phenomenon so surveyable, it makes itself in a special sense the representative of the revelation of the idea in material shape. Although $A_o$ unites selfhood with otherness in a far more inward and concrete way than either $A_{\varkappa}$ or $A_\gamma$, the opposition of the extremes is by no means weaker on that account; it is, on the contrary, owing to the individual limitation in $O_o$, carried through still further. In the first place, $A_o$ has in $O_o$ a more decisive struggle with the material than $A_{\varkappa}$ or $A_\gamma$ can have. As regards $O_\gamma$, it is already remarkable that matter, which with few properties satisfies the abstract-mechanical function in $O_{\varkappa}$, must have a multiplicity of heterogeneous properties in order to fulfill the many-sided demands of the elementary function. But the elementary qualifications by which the forces of matter are multiplied and articulated in $O_\gamma$ all stand in a direct relation to $A_\gamma$. It is otherwise with $O_o$ and $A_o$; the material must be taken up and appropriated in $O_o$ just as, quite just as, it is delivered by $O_\gamma$.\footnote{``On the coming together of the elements into a chemical compound the vital force has not the slightest influence; no element by itself alone is capable of serving for the nourishment and development of a plant or of the animal organism. All the substances which take part in the vital process are lower groups of single atoms, which through the influence of the vital force come together into atoms of a higher order.'' Liebig. Cf.\ my lectures on ``Philosophisk Propædeutik'' from the academic year 1860--61, p.~189 ff.} The organic \opage{444} principle works up something foreign: the idea of life borrows its matter from the geological idea.

In the second place, $A_o$ has in $S$ an ideal archetype, which nevertheless in an essential sense is its absolute counter-image. In organic selfhood the derived subjectivity acquires a quite different inwardness and ideal independence than in the mechanical or the elementary: $A_o$ stands in this respect nearer to $S$ than either $A_{\varkappa}$ or $A_\gamma$. If the principle of life is a principle of soul, and the principle of soul at the same time a principle of spirit, then the likeness between $A_o$ and $S$ is all the more essential. But precisely in the essential likeness the opposition has its very deepest ground. Only the principle of spirit in $A_o$ is able to develop a disposition of selfhood, a subjectivity that has ideal independence enough to place itself over against $S$ and, through the cognition of this its opposition, to cognize itself.

From the principle in $A_o$ the origin of life is now understood: from the unity-in-opposition with $S$ comes selfhood, from the unity-in-opposition with $O_o$ otherness; selfhood is at the beginning hidden by otherness, and yet otherness comes to be determined by selfhood. ``Life does not come as a foreign being into the body, but springs from the nature of the substances, flares up through the self-ignition of the atoms, a flame of ideal combustion. The beginning thus takes place not from an unconditioned unity, but from a multiplicity conditioned in itself\ldots Yet hardly has the unity of life gathered itself out of this manifold before it again concentrates itself into self-unity to such a degree that it raises itself above the multiplicity and subjects to itself the activity of the material points. But that the unity thus results in the \apage{445}\emph{consequent is just a proof that as principle it was already present in the presupposition}.''\footnote{Lect.\ on ``Phil.\ Propæd.'' 1860--61 p.~219.}

From the principle in $A_o$ it is further understood, what would otherwise, if $A_o$ were not conceived as middle between $S$ and $O_o$, be altogether incomprehensible, that the activity of life, although it proceeds unconsciously, can nevertheless be as teleological as if it were guided by a divine knowledge. From the idea-motion in $S$ follows the self-motion in $A_o$; by the real motion in $O_o$ the individual actuality and psychological independence of this self-motion are conditioned. The idea-motion and the real motion mutually presuppose each other; but since the principle, in reflection of the all-moving $S$, itself conditions the conditions by which it is conditioned in $O_o$, the idea-motion is overarching.

From the principle in $A_o$ the otherwise inexplicable unity of the perishable and the imperishable in which life is seen must, then, also admit of explanation. The matter in $O_o$ is subject to the law of otherness; and the law of otherness, the fundamental law of the elements, means perishability. But the law of selfhood in $S$ is the law of ideality; and the law of ideality, the fundamental law of life, means imperishability. ``It is the destination of otherness to serve selfhood; it is essential to selfhood to need otherness. That the plant withers, that the animal dies, shows clearly enough what it means that the living too is subject to the law of otherness and therewith to perishability. The idea of life needs the external other in order through it to develop its disposition, its internal other. But the external other is only a material condition; selfhood, which is conditioned only by the fact that with inner necessity it conditions itself, is what is unconditioned in the system of conditions: the development of life is essentially self-development, and self-development, teleologically seen, idea-development. On the law of self-development rests the imperishability of life\ldots To make the law of otherness the fundamental law of life is nonsense; to let \opage{446} the law of selfhood hold, and yet to assume that the ideal result of development could vanish and life lose itself in nothing because the material organism, subject to the law of otherness, falls prey to perishability, is a senseless contradiction.''\footnote{Lect.\ on ``Phil.\ Propæd.'' 1860--61 p.~281--82.}

3) \emph{In the energy of the functions the energy of the fundamental ideas is reflected}. It is a peculiar matter with making the absolute comprehensible. Immediately seen, the absolute is absolute, and more cannot be said about it. If the absolute is to be determined immediately, it presents itself involuntarily as the one, that which is, the resting. But here two alternating difficulties at once present themselves: either the manifold, what becomes and is moved, must transform itself into semblance and nothing for the absolute itself, or it must retain a finite reality falling outside the absolute. That the latter assumption is unreasonable everyone must concede when he considers that any reality, however relative, which does not have its ground in the absolute must, after all, set itself up as a limit against the absolute and thereby make the absolute itself relative. As regards the former assumption, on the other hand, the matter is infinitely intricate, since the categories semblance and nothing are so infinitely dialectical.

If the relative: the manifold, what becomes and is moved, is semblance and nothing for the absolute in the sense that it is without any significance, then the absolute becomes a unity without semblance of multiplicity, a being without semblance of becoming, a resting without semblance of motion, hence an empty abstraction. If, on the other hand, the categories semblance and nothing are to be understood in such a way that relativity, with its multiplicity of becoming things and moved elements, has indeed no immediate determinate being over against the absolute, but nevertheless belongs as a moment in its essence, then the question of what a moment is, and what validity the moments of the manifold, of \apage{447} what becomes and is moved, must be able to have in, by and for the absolute itself, will surely be of such a scope that its answer goes into one with a thoroughgoing development of concepts and presupposes a logic of the fundamental ideas.

The question of the absolute must in the last instance coincide with the question of the actual, the actual that is perfect in all respects. The first great contribution to a determination of the problem we must then seek in Aristotle. Aristotle rightly makes the demand that \emph{energy} shall be the highest expression of the idea, the absolute the perfect actuality (θεοῦ ἐνέργεια ἀθανασία, τοῦτό δ'ἐστι ζωὴ ἀΐδιος). But since, in accordance with his conception of matter and the finite, he had to lay an altogether one-sided emphasis on the opposition between the possible (τὸ δυνάμει ὄν) and the actual (τὸ ἐνεργείᾳ ὄν), he was thereby led to reduce motion to a finite intermediate determination between potential and actual being (καὶ γὰρ ἔστιν ἡ κίνησις ἐνέργειά τις ἀτελὴς μέντοι). The perfect energy was thus to be an actuality without possibility, an acting without motion. The contradiction that thereby arises is so insistent, so supple and mobile, that neither dogmatism nor criticism is able to master it. Only with the Hegelian logic has a prospect been opened toward the resolution of the contradiction.

When Hegel has been reproached for having, in his Logic, put becoming, multiplicity, motion and possibility into the eternal essence, reproaches of this kind obviously betray a rather gross misunderstanding. In order to grasp how crude the misunderstanding is, one need only notice that all the determinations of finitude—becoming, change, motion, etc.—which belong to the development of the concept are absolutely sublated in the absolute idea. Certainly Hegel's Logic speaks of the becoming of the idea, but it is shown at the same time that the idea, which itself sublates this its becoming, \emph{is in eternal} \emph{being}.
\opage{448}Instead of reproaching Hegel for there not being enough being, rest and immobility in his moving Idea, we must on the contrary emphasize that, at the close of the logical circuit, he ends with an all too immediate being, a being that lacks the absolute energy.

When Hegel describes his method, he is essentially describing the energy. ``Die Bereicherung geht an der \emph{Nothwendigkeit} des Begriffes fort, sie ist von ihm gehalten, und jede Bestimmung ist eine Reflexion in sich. Jede neue Stufe \emph{des Außersichgehens}, das heißt der \emph{weitern Bestimmung}, ist auch ein In-sich-gehen, und die größere \emph{Ausdehnung} ebenso sehr \emph{höhere Intensität}. Das Reichste ist daher das Konkreteste und \emph{Subjektivste}, und das sich in die einfachste Tiefe Zurücknehmende das Mächtigste und Uebergreifendste. Die höchste zugeschärfteste Spitze ist die \emph{reine Persönlichkeit}, die allein durch die absolute Dialektik, die ihre Natur ist, ebenso sehr \emph{Alles in sich befaßt} und \emph{hält}, weil sie sich zum Freisten macht, --- zur Einfachheit, welche die erste Unmittelbarkeit und Allgemeinheit ist.''\footnote{Subjkt.\ Logik p.~349. Werke vol.~5.} Herewith the highest is expressed that any philosophy will ever be able to present. But when the result of the circular movements is to be stated more precisely, the tension slackens, and the immediate, albeit infinitely rich in content, being then shows itself more as being than as acting. ``So ist denn auch die Logik in der absoluten Idee zu dieser einfachen Einheit zurückgegangen, welche ihr Anfang ist; die reine Unmittelbarkeit des Seyns, in dem zuerst alle Bestimmung als ausgelöscht oder durch die Abstraktion weggelassen erscheint, ist die durch die Vermittelung, nämlich die Aufhebung der Vermittelung zu ihrer entsprechenden Gleichheit mit sich gekommene Idee.''\footnote{Subjkt.\ Logik p.~351--52. Werke vol.~5.}

If the relation between the absolute and the relative in \opage{449} being and becoming, in rest and motion, acting and suffering, is to be developed in all its aspects, then it must first and last be emphasized that here there are two kinds of being and two kinds of becoming, two kinds of rest and two kinds of motion, an ideal and a real, and that this doubling consists precisely in and through the opposition, already illuminated from several sides, between the ideal system and the real system. To the opposition between the two systems there corresponds an essential opposition between two realms of ideas. What such an opposition has to signify can be shown at this stage, albeit only provisionally, when we set the real fundamental ideas---the cosmological, the geological, and the biological idea or the idea of life---over against the ideal ideas: knowledge, power, and truth. To show how this opposition presupposes unity and resolves itself into unity is to show how the energy of the ideas reflects itself in the energy of the functions. Such a reflection is more than a reflection totalized in abstract essentiality; it has in its totalities all the moments of the concept, the antilogical as well as the logical; as comprehending totality-reflection it is a comprehending of concrete syllogisms.

The cosmological idea has, as already shown, its reality in the syllogism:
\[ S - A_{\varkappa} - O_{\varkappa}, \]
the geological in the syllogism:
\[ S - A_\gamma - O_\gamma, \]
the biological in the syllogism:
\[ S - A_o - O_o. \]
Now if there is here a system of syllogisms, then not only must each of these syllogisms by itself be a system, but for the sake of unity they must together constitute a system of systems. This comes about, then, in that their middle terms $A_{\varkappa}$, $A_\gamma$, $A_o$ in turn form extremes and middle term, and their extremes $O_{\varkappa}$, $O_\gamma$, $O_o$ middle term and extremes, with one another. If, accordingly, $A_{\varkappa}$ and $A_o$ are posited as extremes with $A_\gamma$ as middle term, then this signifies that the organism's \opage{450} disposition is, by means of the earth's disposition, already anticipated in the disposition of the cosmic structure. If $O_{\varkappa}$ and $O_o$ are posited as extremes with $O_\gamma$ as middle term, then this signifies that the existence of the organism on the earth must be absolutely conditioned by the cosmic structure to which the earth belongs. If $S$ and $O_o$ are posited as extremes, not with $A_o$ but with $A_\gamma$ as middle term, then this expresses that life in the existing organism must have a selfhood reflected in the absolute selfhood but, in accordance with the peculiar disposition of the earth, peculiarly individualized, etc. But however many transpositions we may make, it is and remains only a survey, and we would unavoidably fall back into a Hegelian schematism if we could not expressly get each syllogism in each of the systems motivated by a special development of the corresponding energy of the functions.

That the systems of syllogisms must absolutely constitute a unity means, in other words, that the ideal system itself must be the system of systems. The unity of the ideal system is, according to the first phase of comprehension, a unity of knowledge by means of power; the unity of knowledge, infinitely rich in content, is, as ideal unity of power, a unity in infinite energy: this infinite energy is \emph{the Idea of Spirit}. But the energetic unity of the systems of syllogisms can only consist in an actual multiplicity of particular systems, since each system in turn consists of particular syllogisms, and each syllogism of particular totalities. That the unity, the absolute self-unity, is energetic, and the functions in the systems of syllogisms essentially selfhood-functions, is then expressed mediately in that the multiplicity in the real diversity of the syllogisms acts with such energy that the functions transform themselves into otherness-functions; this multiplicity of the syllogisms belongs to the real system. The real system is, according to the second phase of comprehension, a realm of power by means of knowledge: the energy in this realm of power is \emph{the Idea of Nature}.

\opage{451} Just as the Idea of Spirit has its infinite energy in a system of systems of syllogisms in which the ideality of power serves knowledge, so the Idea of Nature has its energy in a system of systems of syllogisms in which knowledge is a moment for the reality of power; here there is a difference between being and being, motion and motion, energy and energy: in the realm of spirit every real motion is a motion of the idea; in the realm of nature every motion of the idea is a real motion. The two different kinds of series of motion by no means run parallel; but just as each of the series by itself advances through syllogisms and is bent round into a circuit through systems of syllogisms, so too the connections of the series advance through syllogisms and are bent round through systems of syllogisms into a circuit. That the plant sprouts, grows, blossoms, sets fruit, etc., expresses the motion of the idea in the form of real motion. The real motion is not a merely external motion in space; it goes through syllogisms and systems of syllogisms: in this motion of totality the idea of nature has its energy. But the sprouting, growing, blossoming plant is a specimen; the specimen presupposes the individual idea, which in turn presupposes the idea of the species, which in turn presupposes the idea of the genus, etc. The motion of totality corresponding to this likewise goes through syllogisms and systems of syllogisms; but the motion is a motion of reflection: the infinite energy with which the motion of reflection is reflected in eternal knowledge is the energy of the idea of spirit. That both motions, each with its own energy, engage with one another and---through eternal recollecting and temporal producing---pass over into one another, rests then finally also on syllogisms and systems of syllogisms; for the species can, if it is to happen in all immediacy, just as little produce the specimen as the specimen can produce the species.

Only in syllogistic union with the idea of nature does the Idea of Spirit have its truth; only in syllogistic union with the Idea of Spirit does the idea of nature have truth: the Idea of Truth is therefore \opage{452} the unity of the Idea of Nature and of Spirit; but this unity is, be it noted, a totalized unity, a syllogistic unity. As totalized unity the syllogistic unity is grounded in a totalized opposition: the Idea of Truth is therefore the opposition of the Idea of Nature and of Spirit; but the totalized opposition is a syllogistic opposition and just thereby a totalized unity.

So long as the idea of truth, the most concrete of all the fundamental ideas of logic, has not yet reached its full development, the relation between spirit and nature cannot possibly be cleared up. Idealism carried through one-sidedly is essentially spiritualism; realism carried through one-sidedly is essentially naturalism. Since each of these theories has the idea, and in the idea the totality of the whole of existence, in its own one-sided way, each of them, insofar as it wants to be consistent, must move in its own circle; it is altogether impossible to make a transition from the one to the other. ``The Idea of Truth is the unity of the Idea of Nature and of Spirit; \emph{nature and spirit are in truth one}'': with this both spiritualism and naturalism begin. For the principle, true in and for itself, is distorted by two false explanations. ``If nature and spirit are in truth one, then nature must absolutely have its truth in spirit; if nature has its truth in spirit, then spirit must necessarily be the original, nature the derived; if nature is the derived, then it has reality only in the temporal, thus a subordinate and vanishing reality.'' This is the explanation of spiritualism. --- ``If spirit and nature are in truth one, then spirit, which after all has nature as its presupposition, must absolutely have its truth in nature; if spirit has its truth in nature, then nature is the original, spirit the derived; if spirit is the derived, then it has indeed an eternal ideality in eternal nature, but its existence is temporal like the soul, perishable like the organic individual.'' This is the explanation of naturalism.

\opage{453} Once a thinker has fixed himself in one-sided spiritualism, it is impossible for him to recognize the originality of nature; if he has become entangled in the presuppositions of naturalism, then it will be just as impossible for him to acknowledge the originality of spirit. Sound common sense, which has a natural eye for all glaring one-sidednesses but is not developed enough to get to the bottom of the investigations, will hold fast to nature and protest against spiritualism; but it will naturally also hold fast to spirit and protest against naturalism.\footnote{``The fundamental question thus becomes this: which form, that of necessity or that of freedom, is the original form of the idea? In the form of necessity the idea is nature, in the form of freedom the idea is spirit; the question can therefore be put thus: is it the realm of nature that is the original, spirit the derived, or is it conversely the realm of spirit that is the original, nature the derived? Or if one conceives the relation in such a way that nature must just as necessarily have spirit as spirit in turn has nature as its presupposition, and consequently, by ascribing the same originality to both realms, presupposes a ground-unity of nature and spirit: then it is asked again whether, according to this ground-unity, there is not just as well a real opposition between an imperfect and a perfect nature as between relative imperfect spirits and an absolute or perfect spirit?'' Forelæsn.\ over ``Phil.\ Propæd.'' 1861--62 p.~138--39.}

But when spiritualism is consistent from its standpoint and naturalism from its own, while the unclear semi-dualism is superficial and inconsistent, how then is one to be able to demonstrate that the consistent one-sidedness really is one-sided and in its one-sidedness untenable? Answer: by consistency itself! The energy of the ideas must be tested against the energy of the functions.

For consistent spiritualism cannot avoid weakening the material otherness that furnishes the necessary presuppositions and conditions of the exact \opage{454} and physical functions. For the spiritualistic reflection, forsaken by nature, all functions must show themselves as abstract selfhood-functions; but abstract selfhood-functions are intellectual functions in one-sidedly logical form. One therefore need only sharpen the demand that spiritualism explain actuality with its facts, and one will at once notice the slackness, the hovering-above, the insipid generality, which is most subjective of all just when it aims most of all at the objective. It is this slackness that we took note of in the foregoing when the dynamic was separated from its mechanical, the ideal from its material foundation. If spiritualism is in this way transformed into intellectualism, then the Idea of Spirit loses its energy, and with that the theory falls.

Consistent naturalism, which, be it noted, begins with the real fundamental idea, namely with the idea of nature, is called upon from its standpoint to recognize and assert the reality of the idea in all actuality. Now if naturalism is driven by its own irresistible consistency to an extreme in which it is compelled to let go of the idea and seize the finite fact, then it has undermined itself. But against such an extreme it cannot possibly defend itself. The idea of nature is, in the immediate sense, an idea of actuality; its concrete content is in the actual facts. But just as the content of the idea is immediately resolved into the facts, and the actual facts belong to finitude and thereby to otherness, the functions are transformed into otherness-functions. Naturalism now has the choice between one of two things: either it must, in order to assert its idea, keep to metaphysical universal forms, or seize actuality in its finite conditions. If it chooses the first alternative, then it lapses into mysticism: all abstract metaphysics of nature is essentially mysticism. The real idea, whose whole tension is conditioned by the mystical fermentation, will then be just as \opage{455} foggy, doubtful and empty as the fermentation-functions in which it has its obscure stirrings. If naturalism chooses the second alternative, then it is forced under observations and calculations. The calculation-functions are in reality mechanical otherness-functions: naturalism has then lost its idea and has become a plain and simple materialism.

Naturalism and spiritualism are both recognized by this: that they conceive the unity of nature and spirit as a unity of essence, without conceiving it as a syllogistic unity. When the unity of nature and spirit is seen in the Idea of Truth as a syllogistic unity, only then is it seen as a true concrete unity of beginning. But in recognizing the concrete unity of beginning as a syllogistic unity, this syllogistic unity must at the same time be recognized as a totalized opposition; the one pole of the opposition falls in the ideal system, the other in the real system. Just as the ideal energy of the functions is in all punctuality conditioned by the material, and conversely, so now the ideal energy of the ideas in the ideal system is conditioned by the material energy of the ideas in the real system. Now since the energy of the ideal ideas is spirit, the energy of the real ideas nature, nature and spirit at bottom constitute one realm in that they in actuality constitute two realms. In the realm of nature the energy of the idea everywhere presupposes a finite separation of being and becoming, rest and motion, possibility and actuality. But in the realm of spirit, by contrast, the energy of the idea presupposes an infinitely reflected unity of being in becoming, of rest in motion, of possibility in actuality.

The two realms are, by means of the overarching power of spirit, one realm in the Idea of Truth. Here, then, in this double realm there is an actuality that in its difference from possibility is one with possibility, because here there is an actuality that in its unity with possibility has possibility as its finite opposition; here in the realm of spirit there is a unity that in its self-movement through the manifold is immovable, because here in the realm of nature \apage{456} there is a multiplicity in which motion sublates rest, and rest conversely motion. Here there is, in other words, an energy of the ideas in external acting and temporal becoming, because here there is an energy of the ideas in inward acting and eternal being. Indeed, so thoroughgoing is the doubling that it runs through the whole realm of ideas; for the ideas in the pure logical form, the universal ideas of the ideal system, have their ideal energy in a becoming transformed into being, in a motion transformed into rest, while the ideas in the logical-antilogical form, the concrete ideas of the real system, have their material energy in a being transformed into becoming, in a rest transformed into motion. Now since the opposition of ideal and real ideas everywhere has its syllogistic unity, and thereby its concrete unity of beginning, in the fundamental ideas, the energy of the fundamental ideas, in the realm of spirit as well as in the realm of nature, is an energy infinitely doubled in itself; before the energy of the fundamental ideas spiritualism and naturalism must in the end succumb.

The true insight into the energy of the fundamental ideas is conditioned by a corresponding insight into the energy of the functions; but the energy of the functions has hitherto served only a development of objectification and of general objectivity; it is not the objectivity in the existing objects, but a knowledge of the objectivity determined by the objects, to which the development has attained; such a knowledge is in its generality ontological-subjective; the functions of general objectivity are all functions of subjective knowledge.

\end{document}
